% !TeX program = xelatex
% Standalone English volume. Use XeLaTeX until the references stabilize.
\documentclass[a4paper,11pt]{article}
\usepackage[T1]{fontenc}
\usepackage{lmodern}
\usepackage[english]{babel}
\usepackage[left=26mm,right=26mm,top=25mm,bottom=25mm,headheight=15pt,headsep=7mm,footskip=11mm]{geometry}
\usepackage{amsmath,amssymb,mathtools,amsthm}
\newtheorem{proposition}{Proposition}[section]
\newtheorem{lemma}[proposition]{Lemma}
\newtheorem{theorem}[proposition]{Theorem}
\newtheorem{corollary}[proposition]{Corollary}
\newtheorem{maintheorem}[proposition]{Main Theorem}
\theoremstyle{definition}
\newtheorem{definition}[proposition]{Definition}
\theoremstyle{remark}
\newtheorem{remark}[proposition]{Remark}
\makeatletter
\renewcommand*\l@section{\@dottedtocline{1}{0em}{2.4em}}
\renewcommand*\l@subsection{\@dottedtocline{2}{1.45em}{3.45em}}
\renewcommand*\l@subsubsection{\@dottedtocline{3}{3.8em}{4.55em}}
\makeatother
\usepackage{xcolor,fancyhdr,titlesec,needspace,xurl}
\usepackage{longtable,multicol}
\usepackage{hyperref}
\definecolor{ink}{HTML}{1A3542}
\definecolor{muted}{HTML}{4F606A}
\hypersetup{unicode=true,colorlinks=true,linkcolor=ink,filecolor=ink,urlcolor=ink,citecolor=ink,bookmarksnumbered=true,
 pdfdisplaydoctitle=true,pdflang={en},
 pdftitle={The generalized Riemann hypothesis for Dirichlet L-functions. Edition 930, Volume XII},
 pdfsubject={A natural-language reconstruction of the proposed proof},
 pdfauthor={A mathematical reconstruction}}
\titleformat{\section}{\Large\bfseries\color{ink}}{\thesection}{.65em}{}
\titlespacing*{\section}{0pt}{.8em}{.65em}
\setlength{\parindent}{1.2em}
\linespread{1.045}
\setlength{\parskip}{.22em plus .05em minus .03em}
\setlength{\emergencystretch}{2em}
\setcounter{topnumber}{2}
\setcounter{tocdepth}{3}\setcounter{secnumdepth}{3}
\widowpenalty=4500
\clubpenalty=4500
\pagestyle{fancy}\fancyhf{}
\fancyhead[L]{\small\color{muted}Dirichlet GRH \textperiodcentered\ Volume I}
\fancyhead[R]{\small\color{muted}Dirichlet zeros and power intervals}
\fancyfoot[C]{\small\thepage}
\renewcommand{\headrulewidth}{.3pt}
\newcommand{\lead}[1]{\par\smallskip\noindent{\bfseries\color{ink}#1.}\enspace}
\newcommand{\sour}[2]{\unskip\penalty150\hspace{.35em}{\small\hyperlink{ref#1}{[#2]}}}
\titleformat{\subsection}{\large\bfseries}{\thesubsection}{.6em}{}
\titlespacing*{\subsection}{0pt}{1.15em plus .15em}{.4em}
\setlength{\abovedisplayskip}{9pt plus 2pt minus 3pt}
\setlength{\belowdisplayskip}{9pt plus 2pt minus 3pt}
\newcommand{\RR}{\mathbb R}\newcommand{\CC}{\mathbb C}\newcommand{\QQ}{\mathbb Q}
\newcommand{\ZZ}{\mathbb Z}\newcommand{\ord}{\operatorname{ord}}
\newcommand{\Frob}{\operatorname{Frob}}\newcommand{\Gal}{\operatorname{Gal}}
\newcommand{\len}{\operatorname{length}}\newcommand{\Res}{\operatorname*{Res}}
\newcommand{\wmean}{\operatorname{av}}\newcommand{\Ker}{\operatorname{ker}}
\newcommand{\Lh}{\mathcal L}\newcommand{\Gg}{\mathcal G}
\newcommand{\real}{\operatorname{Re}}\newcommand{\imag}{\operatorname{Im}}
\newcommand{\abs}[1]{\left|#1\right|}\newcommand{\eps}{\varepsilon}
\newcommand{\anchoritem}[2]{\par\Needspace{2.8\baselineskip}\hypertarget{ref#1}{}\noindent{\bfseries\color{ink}#2.}\enspace}
\newcommand{\srcloc}[2]{{\small\nolinkurl{#1.lean}}, line~#2}
\numberwithin{equation}{section}
\raggedbottom

\newcommand{\Sel}{\mathcal S}
\newcommand{\Cl}{\operatorname{Cl}}
\newcommand{\Aut}{\operatorname{Aut}}
\newcommand{\GL}{\operatorname{GL}}
\newcommand{\OO}{\mathcal O}
\hypersetup{bookmarksdepth=3}


\addto\captionsenglish{\renewcommand{\partname}{Volume}}


\makeatletter
\expandafter\def\csname r@cor:coc145\endcsname{{25.15}{36}{}{proposition.25.15}{Dirichlet-GRH-v930-Volume-V.pdf}}
\expandafter\def\csname r@cor:finmark557\endcsname{{48.19}{121}{}{proposition.48.19}{Dirichlet-GRH-v930-Volume-IX.pdf}}
\expandafter\def\csname r@cor:pointalternative550\endcsname{{48.15}{117}{}{proposition.48.15}{Dirichlet-GRH-v930-Volume-IX.pdf}}
\expandafter\def\csname r@cor:tracesplit548\endcsname{{48.11}{113}{}{proposition.48.11}{Dirichlet-GRH-v930-Volume-IX.pdf}}
\expandafter\def\csname r@lem:bernoulliall436\endcsname{{55.5}{22}{}{proposition.55.5}{Dirichlet-GRH-v930-Volume-XI.pdf}}
\expandafter\def\csname r@lem:cuspidaltrace433\endcsname{{46.15}{72}{}{proposition.46.15}{Dirichlet-GRH-v930-Volume-IX.pdf}}
\expandafter\def\csname r@lem:uniformoddroots430\endcsname{{46.6}{66}{}{proposition.46.6}{Dirichlet-GRH-v930-Volume-IX.pdf}}
\expandafter\def\csname r@prop:actualcomponent557\endcsname{{48.18}{120}{}{proposition.48.18}{Dirichlet-GRH-v930-Volume-IX.pdf}}
\expandafter\def\csname r@prop:auxiliarycusp436\endcsname{{55.6}{23}{}{proposition.55.6}{Dirichlet-GRH-v930-Volume-XI.pdf}}
\expandafter\def\csname r@prop:finiteordinary558\endcsname{{48.21}{122}{}{proposition.48.21}{Dirichlet-GRH-v930-Volume-IX.pdf}}
\expandafter\def\csname r@prop:strongall559\endcsname{{48.23}{123}{}{proposition.48.23}{Dirichlet-GRH-v930-Volume-IX.pdf}}
\expandafter\def\csname r@sec:generaleven547\endcsname{{48.3}{109}{}{subsection.48.3}{Dirichlet-GRH-v930-Volume-IX.pdf}}
\expandafter\def\csname r@sec:missing\endcsname{{6.6}{32}{}{subsection.6.6}{Dirichlet-GRH-v930-Volume-I.pdf}}
\expandafter\def\csname r@sec:symplecticfrobenius433\endcsname{{46.4}{71}{}{subsection.46.4}{Dirichlet-GRH-v930-Volume-IX.pdf}}
\expandafter\def\csname r@sec:symplecticstrict429\endcsname{{46.1}{63}{}{subsection.46.1}{Dirichlet-GRH-v930-Volume-IX.pdf}}
\expandafter\def\csname r@sec:twomarks557\endcsname{{48.7}{118}{}{subsection.48.7}{Dirichlet-GRH-v930-Volume-IX.pdf}}
\expandafter\def\csname r@seriesguide\endcsname{{1.1}{4}{}{subsection.1.1}{Dirichlet-GRH-v930-Volume-I.pdf}}
\expandafter\def\csname r@thm:allglobalroots547\endcsname{{48.6}{109}{}{proposition.48.6}{Dirichlet-GRH-v930-Volume-IX.pdf}}
\expandafter\def\csname r@thm:allroottrace548\endcsname{{48.9}{112}{}{proposition.48.9}{Dirichlet-GRH-v930-Volume-IX.pdf}}
\expandafter\def\csname r@thm:automult559\endcsname{{48.25}{125}{}{proposition.48.25}{Dirichlet-GRH-v930-Volume-IX.pdf}}
\expandafter\def\csname r@thm:commonsptrace433\endcsname{{46.16}{72}{}{proposition.46.16}{Dirichlet-GRH-v930-Volume-IX.pdf}}
\expandafter\def\csname r@thm:conductor232\endcsname{{19.15}{12}{}{proposition.19.15}{Dirichlet-GRH-v930-Volume-IV.pdf}}
\expandafter\def\csname r@thm:criticalall310\endcsname{{28.6}{90}{}{proposition.28.6}{Dirichlet-GRH-v930-Volume-V.pdf}}
\expandafter\def\csname r@thm:derivation367\endcsname{{37.15}{50}{}{proposition.37.15}{Dirichlet-GRH-v930-Volume-VII.pdf}}
\expandafter\def\csname r@thm:fourier145\endcsname{{25.14}{35}{}{proposition.25.14}{Dirichlet-GRH-v930-Volume-V.pdf}}
\expandafter\def\csname r@thm:integralroots547\endcsname{{48.7}{110}{}{proposition.48.7}{Dirichlet-GRH-v930-Volume-IX.pdf}}
\expandafter\def\csname r@thm:lengthmultiplicity550\endcsname{{48.14}{116}{}{proposition.48.14}{Dirichlet-GRH-v930-Volume-IX.pdf}}
\expandafter\def\csname r@thm:levelstrict345\endcsname{{22.17}{75}{}{proposition.22.17}{Dirichlet-GRH-v930-Volume-IV.pdf}}
\expandafter\def\csname r@thm:oddrootlocalization430\endcsname{{46.8}{67}{}{proposition.46.8}{Dirichlet-GRH-v930-Volume-IX.pdf}}
\expandafter\def\csname r@thm:purityblocks549\endcsname{{48.13}{115}{}{proposition.48.13}{Dirichlet-GRH-v930-Volume-IX.pdf}}
\expandafter\def\csname r@thm:ramification549\endcsname{{48.12}{114}{}{proposition.48.12}{Dirichlet-GRH-v930-Volume-IX.pdf}}
\expandafter\def\csname r@thm:twoprimeadjoints436\endcsname{{55.7}{24}{}{proposition.55.7}{Dirichlet-GRH-v930-Volume-XI.pdf}}
\makeatother
\begin{document}
\setcounter{part}{11}
\fancyhead[L]{\small\color{muted}Dirichlet GRH \textperiodcentered\ Volume XII}
\fancyhead[R]{\small\color{muted}Cuspidal companions and Hecke fibers}
\thispagestyle{plain}\begin{center}
{\LARGE\bfseries The generalized Riemann hypothesis\\[2pt]for Dirichlet \(L\)-functions}\\[6pt]
{\small 27 November 2026 \quad\textperiodcentered\quad Edition 930}\end{center}
\part{Cuspidal companions and ordinary Hecke fibers}\label{part:XII}
\begin{abstract}\noindent We construct cuspidal points on auxiliary towers from quadratic induction and symmetric cubes. At the prescribed weights we study their traces, their Hecke fibers and the dual Selmer conditions on these families.\end{abstract}
\setcounter{section}{58}
\noindent The \hyperref[seriesguide]{guide in Volume~I} lists all twenty-seven volumes.
\begingroup\small\setlength{\parskip}{0pt}\tableofcontents\endgroup
\clearpage

\section{Cuspidal companions and the two even planes}\label{nav:chapter-companions-even-planes}\label{nav:volume-XII}

\subsection{Cuspidal companions over one imaginary quadratic field}
\label{sec:onerealization560}

Fix \(K=\QQ(i)\) throughout the interval of Section~\ref{sec:missing}.
Let \(V\) be one of the characteristic-zero representations of
Proposition~\ref{prop:strongall559}, of dimension \(n=4\) or \(2\).
Fix \(\iota:\overline{\QQ}_\ell\simeq\CC\). We say that a
cuspidal representation \(\pi\) of \(\mathrm{GL}_n(\mathbb A_K)\)
matches \(V\) if \(\pi\) is regular algebraic and the following
equality of multisets holds at all but finitely many places \(w\).
Write \(q_w\) for the residue cardinality, \(S_w\) for the
unramified Satake parameters of \(\pi_w\), and \(g_w\) for
geometric Frobenius. The equality is
\[
       \iota\big(\operatorname{Eig}(V(g_w))\big)
                             =q_w^{(n-1)/2}S_w .
\]
This is the Frobenius convention of the source in
\srcloc{Development/AutomorphicArithmetic/AutomorphicVarieties/AutomorphicForms/BlochKatoSelmerVanishes}{66}.
Below we specify when we take the inverse, arithmetic Frobenius.

\begin{proposition}[Polarization from a single cuspidal match]
\leavevmode\label{prop:onecompanion560}
A matching \(\pi\) admits a polarization whose Galois multiplier
on \(G_K\) equals \(\lambda\). Here \(\lambda=\mu|_{G_K}\)
in dimension four and \(\lambda=(\det V)|_{G_K}\) in dimension
two. The Galois representation attached to the polarized pair
is isomorphic to \(V|_{G_K}\). In dimension four, when \(T\)
is finite, it follows that \(n_\sigma=1\).
\end{proposition}
\begin{proof}
Denote the same two multipliers on \(G_\QQ\) by \(\lambda_0\).
They are products of a finite Dirichlet character and a power of
\(\kappa_\ell\), by the definition of \(\mu\) and the two
determinant identities in Theorem~\ref{thm:purityblocks549}.
In either dimension there is an alternating pairing showing
\(V^\vee\otimes\lambda_0\simeq V\). Conjugation by the
nontrivial automorphism of \(K/\QQ\) leaves the isomorphism
class of the restriction unchanged, since it extends to
\(G_\QQ\).

Take the algebraic Hecke character \(\chi\) of \(\QQ\)
corresponding to \(\lambda_0\kappa_\ell^{\,n-1}\) by class
field theory in geometric Frobenius normalization. It is a
finite character times a norm power. Its base change
\(\chi_K=\chi\circ N_{K/\QQ}\) has uniformizer value
\(\iota((\lambda_0\kappa_\ell^{\,n-1})(g_w))\) at every good
place. Since \(\kappa_\ell(g_w)=q_w^{-1}\), dualizing the
display defining a match gives
\[
 S_{cw}=\iota(\lambda_0(g_w))q_w^{1-n}S_w^{-1}
       =\chi_K(\varpi_w)S_w^{-1}.
\]
Here \(c\) acts on places by complex conjugation.
Strong multiplicity one for cuspidal \(\mathrm{GL}_n\) gives
\(\pi^c\simeq\pi^\vee\otimes(\chi_K\circ\det)\).
If necessary, multiply \(\chi\) by the quadratic character of
\(K/\QQ\) to give the required real sign in the definition of
a polarized pair. This quadratic character is trivial after
base change, so the character on \(G_K\) in the pairing is
still
\(\kappa_\ell^{1-n}r_{\ell,\iota}(\chi)|_{G_K}=\lambda\).
The existence and unramified compatibility theorem for regular
algebraic cuspidal polarized pairs now gives the associated
semisimple Galois representation. Its good Frobenius
polynomials give the same traces as \(V|_{G_K}\). Chebotarev,
continuity and the characteristic-zero character criterion give
the isomorphism, since the latter is irreducible by
Proposition~\ref{prop:strongall559}. Apply
Theorem~\ref{thm:automult559} in dimension four.
~\cite{XII-IX-mult560,XII-IX-blggt560}
\end{proof}

\begin{theorem}[The sign obstruction for either even plane]
\leavevmode\label{thm:evenplanecompanion560}
Neither of the two characteristic-zero planes in
Theorem~\ref{thm:purityblocks549} has a matching regular
algebraic cuspidal representation over \(K\).
\end{theorem}
\begin{proof}
If such a match existed,
Proposition~\ref{prop:onecompanion560} would identify the plane
over \(G_K\) with the irreducible representation attached to a
polarized cuspidal pair. Its conjugate-dual pairing with
multiplier \((\det V)|_{G_K}\) is symmetric by the sign theorem
of Bella\"iche and Chenevier, in the polarized formulation of
Barnet-Lamb, Gee, Geraghty and Taylor, Theorem~2.1.1(1).
Here the lift \(c\) of complex conjugation has order two.
Choose an alternating determinant matrix \(J_2\) on the plane.
The matrix \(J_2 V(c)\) is also a conjugate-dual pairing with
this very multiplier, since
\[
 V(g)^{\mathsf T}J_2V(c)\,V(cgc^{-1})
           =(\det V(g))J_2 V(c)\qquad(g\in G_K).
\]
But \(V(c)\) is either \(1\) or \(-1\) on each plane by
Theorem~\ref{thm:purityblocks549}. This second pairing is
alternating. Irreducibility over \(K\) and Schur's lemma make
the space of such pairings one-dimensional. A nonzero symmetric
form and a nonzero alternating form cannot be proportional in
characteristic zero. This proves the assertion.
~\cite{XII-IX-sign560,XII-IX-blggt560}
\end{proof}

\begin{theorem}[The integral Frobenius eigenvalues at a matching point]
\leavevmode\label{thm:integerpurity560}
Suppose a rank-four point has a cuspidal match over \(K\). For
every rational prime \(p\ne\ell\) at which \(V\) is unramified,
each eigenvalue of its arithmetic Frobenius is an algebraic
integer. Under every complex embedding every such eigenvalue
has absolute value \(p^{(r+1)/2}\). In particular the
Frobenius trace \(u\) appearing in
Theorem~\ref{thm:commonsptrace433} is an algebraic integer,
and every conjugate of \(u\) has absolute value at most
\(4p^{(r+1)/2}\).
\end{theorem}
\begin{proof}
Choose a place \(w\) of \(K\) above \(p\) of residue degree
\(f\). The eigenvalues of arithmetic Frobenius there are the
\(f\)-th powers of the rational ones, since \(V\) is unramified
at \(p\). By Proposition~\ref{prop:onecompanion560} the
restricted representation belongs to the compatible system of
a polarized cuspidal pair. At \(w\nmid\ell\) the Weil--Deligne
parameter of this representation is unramified with zero
monodromy. Local-global compatibility identifies its Frobenius
semisimplification with the normalized local Langlands parameter
of \(\pi_w\). In particular \(\pi_w\) is unramified.
Strict purity for this compatible system implies that the
Frobenius eigenvalues are algebraic numbers and all eigenvalues
at \(w\) have the same complex absolute value under every
embedding. This holds also for their inverses in the arithmetic
convention. The determinant on \(V\) is \(\mu^2\), whose value
at this arithmetic Frobenius has absolute value
\(p^{2f(r+1)}\). Taking first fourth roots and then \(f\)-th
roots gives the stated absolute value at \(p\).
~\cite{XII-IX-blggt560}

At a coefficient place above any prime \(p'\ne p\), the system
is unramified at \(w\) and the same polynomial is realized on
a stable lattice by local-global compatibility. Its arithmetic
eigenvalues are therefore integral there. It remains to check
coefficient places above \(p\). The \(p\)-adic member is
crystalline at \(w\) since \(\pi_w\) is unramified. Its
crystalline Frobenius, with the residue degree included,
realizes the geometric Frobenius polynomial. These assertions
are Theorem~2.1.1(4) of the same reference. Its Hodge--Tate
numbers at all embeddings, with the cyclotomic character of
weight \(-1\), are
\[
                         -(r+1),\quad-r,\quad-1,\quad0 .
\]
Indeed the ordinary flag at \(\ell\) gives these numbers for
the original representation by Proposition~\ref{prop:finiteordinary558}.
They are the same at both embeddings of \(K\) since the
representation comes from \(\QQ\). The algebraic weight formula
of Theorem~2.1.1(3) gives the same multiset for the \(p\)-adic
member, including after permutation of the embeddings.

Use the dual \(p\)-adic member, whose labeled Hodge--Tate
numbers are all nonnegative. Its filtered isocrystal is weakly
admissible. The induced Hodge filtration on its smallest-slope
subisocrystal has only nonnegative jumps, so the Hodge number
of this subobject is nonnegative. Weak admissibility puts its
Newton number at least as high. Hence its smallest slope is
nonnegative, which gives nonnegative valuation for every
eigenvalue of its residue-degree Frobenius. These are the
arithmetic eigenvalues at \(w\) for the original member. One
may forget the coefficient action in the slope argument, since
this retains all conjugate eigenvalues. Integrality now holds
at every finite place. The \(f\)-th powers of the rational
Frobenius eigenvalues are thus algebraic integers. Each
eigenvalue itself is integral by transitivity of integrality.
Taking the trace proves the last assertion.
\end{proof}

The source's condition in
\srcloc{Development/AutomorphicArithmetic/AutomorphicVarieties/ModularityArgument/BaseChangeUnits/IsAutQ}{13}
asks for a cuspidal match over every imaginary quadratic field
on which the restriction is irreducible. For the conclusions
above the single field \(K\) suffices throughout the interval.
The sign obstruction concerns the existence of a cuspidal
match for an individual plane. Eliminating the plane from the
deformation space itself requires a further argument.


\subsection{A common level for the cuspidal contribution to strict length}
\label{sec:commonlevel561}

Keep \(K=\QQ(i)\) and the common family \(E_x\) of
Section~\ref{sec:missing}. Use \(V\) for the lower endpoint of
the interval \([V,V^{3/2}]\), \(M\) for the dyadic conductor
of Lemma~\ref{lem:uniformoddroots430}, and \(N_*,n_*,t_*,L_*,Y\)
as in Section~\ref{sec:symplecticstrict429}. Write \(\rho_\sigma\)
for a rank-four point admitting a match in the sense of
Section~\ref{sec:onerealization560}. For a coefficient prime
put \(d_\ell=\ell M N_*\). Principal level \(b\) on
\(\mathrm{GL}_4(\mathbb A_{K,f})\) means the kernel of reduction
modulo \(b\) in \(\mathrm{GL}_4(\widehat{\mathcal O}_K)\).

\begin{proposition}[The level and weight of any matching system]
\leavevmode\label{prop:levelofmatch561}
Every match \(\pi\) of \(\rho_\sigma\) over \(K\) has a nonzero
vector of principal level \(d_\ell^4\). It occurs in cuspidal
cohomology with coefficients of highest weights bounded in
absolute value by \(r+4\) at both embeddings of \(K\).
\end{proposition}
\begin{proof}
Fix a finite place \(w\) of \(K\) over \(q\). First take
\(q\ne\ell\). The inertia of the global lift at \(q\) factors
through the finite abelian group given by its residual inertia
image. This group has order prime to \(\ell\), and the conjugation
action of Frobenius on it is trivial since the residual global
image is diagonal. The multiplicities of the inertia characters
of the lift equal those of \(d\). Indeed the character idempotents
of this finite group act on the integral lattice, and their
ranks can be read after reduction. Thus the semisimplified Weil
representation restricts on inertia to characters of conductor
dividing \(M N_*\). The monodromy operator vanishes since the
inertia image is finite. Under restriction to \(K_w\) each local
character is composed with the norm. This preserves the bound
by the extended ideal. In fact the norm of a unit congruent to
one modulo \(q^k\mathcal O_{K_w}\) is congruent to one modulo
\(q^k\ZZ_q\), by taking the determinant of multiplication on
the integral lattice.

For \(q=\ell\), the ordered local representation is potentially
crystalline by Proposition~\ref{prop:finiteordinary558}. The
Weil--Deligne functor on potentially semistable representations
is exact. On each of the four graded characters the unramified
factor and the cyclotomic powers have unramified Weil parameter.
The only possible ramified factor is
\(\omega_\ell^{-1}\), of conductor dividing \(\ell\).
After restricting to \(K_w\) the same norm argument applies.
The semisimplified Weil parameter of \(\pi_w\), up to the
unramified norm twist, agrees with this parameter by local-global
compatibility up to semisimplification at the coefficient prime.
Away from \(\ell\) use Theorem~2.1.1(2) of Barnet-Lamb, Gee,
Geraghty and Taylor. Thus at every \(w\) the semisimplified
inertia parameter has four character factors, each of conductor
exponent at most \(b_w=v_w(d_\ell)\). When \(b_w=0\), the
argument at \(q\ne\ell\) also gives zero monodromy on the
automorphic parameter and \(\pi_w\) is unramified.
~\cite{XII-IX-lgp361,XII-IX-blggt560}

When \(b_w\ge1\), write \(a_j\) for the four character conductor
exponents. The Artin conductor formula for a Weil--Deligne
parameter \((R,N)\) adds
\(\dim R^{I_w}-\dim(\ker N)^{I_w}\) to the Weil conductor.
The first dimension is the number of unramified character
factors. Consequently the local conductor exponent of \(\pi_w\)
is at most
\[
       \sum_{a_j>0}a_j+\#\{j:a_j=0\}\ \le\ 4b_w .
\]
Here we use compatibility of local Langlands with conductors.
The local components of a cuspidal \(\mathrm{GL}_4\) representation
are generic. Their conductor newvectors are fixed by the subgroup
whose last row is that of the identity modulo the local
conductor. It contains the principal congruence subgroup at
exponent \(4b_w\). The tensor product of these local vectors
proves the level assertion.~\cite{XII-IX-newvector561}

The labeled Hodge--Tate multiset at both embeddings is the one
computed in Theorem~\ref{thm:integerpurity560}. In the algebraic
weight formula the four highest-weight entries are obtained by
subtracting \(3,2,1,0\) from the ordered Hodge--Tate numbers.
Taking the dual coefficient system reverses the entries and
changes their signs. In either convention the absolute bound
is \(r+4\). The cohomological characterization of regular
algebraic cuspidal representations gives nonzero relative Lie
algebra cohomology with this algebraic coefficient system.
Cuspidal cohomology injects into the cohomology of the locally
symmetric space, with the action of the finite Hecke operators.
~\cite{XII-IX-clozel561,XII-IX-franke561}
\end{proof}

In Section~\ref{sec:symplecticfrobenius433} choose the common
prime \(p\) in \((Y+32,2(Y+32))\) to split in \(K\). For all
sufficiently large \(V\) such a choice is possible by the prime
number theorem in the fixed class \(1\bmod 4\). It is prime to
\(3d_\ell\) for every \(\ell\) in the indicated range. Fix
one of its two places \(w\) in \(K\).

\begin{theorem}[One finite list for every matching trace on the interval]
\leavevmode\label{thm:commoncusplist561}
For each \(2\le r\le t_*\) there is a rationally Galois-stable
multiset \(\mathcal C_r\) of algebraic integers, all of whose
complex absolute values are at most \(4p^{(r+1)/2}\), with
\[
    |\mathcal C_r|\ \le\
        C(r+5)^{20}\sum_{\substack{5\le\ell\le L_*\\\ell\ {\rm prime}}}
                                      (3d_\ell^4)^{32}
\]
for an absolute constant \(C\). Fix any triple in
Section~\ref{sec:generaleven547} with \(T\) finite and one
coefficient embedding \(\iota\). The arithmetic Frobenius traces
at \(p\) of all the distinct rank-four points having a match
under \(\iota\), counted once per point, form a submultiset of
\(\mathcal C_r\) under that embedding.
\end{theorem}
\begin{proof}
Use the locally symmetric spaces for \(\mathrm{GL}_4/K\)
modulo the positive real split center at infinity, at principal
levels \(3d_\ell^4\). Take all dominant pairs of coefficient
weights, one for each embedding of \(K\), with entries bounded
by \(r+4\). The complex algebraic coefficient representation
with such a pair of weights descends to \(K\). Indeed the
restriction of scalars of \(\mathrm{GL}_4\) splits into two
copies over \(K\), and each highest-weight representation is
defined over the split field. Restrict scalars of this
coefficient space to \(\QQ\) to obtain rational local systems.
Their complexifications include every indicated coefficient
space and its conjugate. The finite Hecke correspondences act
over \(\QQ\) on the resulting cohomology.

The Borel--Serre compactification gives a finite cell model
at the fixed level three. A stabilizer in the symmetric space
modulo the positive real center preserves an integral lattice
of rank eight. Its determinant is a Gaussian unit, of complex
absolute value one, so it belongs to a compact stabilizer
before taking that quotient. Its discreteness makes it finite.
A finite-order integral matrix congruent to the identity modulo
three is the identity. Thus the level-three arithmetic groups
act freely. The level-\(3d_\ell^4\) spaces are covers of degree
at most \((3d_\ell^4)^{32}\), by reduction of sixteen entries
in \(\mathcal O_K\). There are \(O((r+5)^8)\) highest-weight
pairs. The Weyl dimension formula on the two factors gives
rank \(O((r+5)^{12})\) for each coefficient system, including
restriction of scalars. Multiplying by the number of cells
bounds the sum of all cohomology dimensions by the displayed
estimate up to an absolute constant.~\cite{XII-IX-V-bs}

Give the maximal compact subgroup at \(w\) volume one. Let
\(U_w\) be the Hecke operator of the double coset of
\(\operatorname{diag}(1,1,1,\varpi_w^{-1})\). Use
\(p^{-3}U_w\) on rational cohomology.
In the convention in which the double coset of
\(\operatorname{diag}(\varpi_w,1,1,1)\) has eigenvalue
\(p^{3/2}\sum_j\beta_j\) on Satake parameters \(\{\beta_j\}\),
the inverse coset has eigenvalue \(p^{3/2}\sum_j\beta_j^{-1}\).
These are the two minuscule formulas of the normalized Satake
isomorphism. After division its eigenvalue is
\(\sum_j(p^{3/2}\beta_j)^{-1}\), the arithmetic Frobenius
trace for a match. The choice of a split prime identifies the
Frobenius at \(w\) with the rational one on this representation.

Take the product of the characteristic polynomials on all
the cohomology groups, levels and weights above. It is a
rational polynomial. Keep each root which is an algebraic
integer and whose every conjugate has absolute value at most
\(4p^{(r+1)/2}\), with twice its multiplicity. This defines
\(\mathcal C_r\) and proves its size and stability assertions.
Proposition~\ref{prop:levelofmatch561} and
Theorem~\ref{thm:integerpurity560} put each matching trace
among the retained roots.

It remains to account for repeated values. Distinct cuspidal
systems supply linearly independent simultaneous Hecke
eigenvectors in the direct sum, because strong multiplicity one
distinguishes them away from the union of the levels. Thus a
value shared by \(j\) such systems has algebraic multiplicity
at least \(j\) before doubling. Two rank-four points with the
same matching system have isomorphic restrictions to \(G_K\)
by Proposition~\ref{prop:onecompanion560}. Fix an isomorphism
on this subgroup. For each \(g\in G_\QQ\) its conjugate by
the two actions is another such intertwiner, since \(G_K\) is
normal. Their ratio is scalar by
Proposition~\ref{prop:strongall559} and Schur's lemma. These
scalars form a character of \(G_\QQ/G_K\), which has only two
characters. There are therefore at most two global isomorphism
classes for the given system. The traces of a global class
determine its point of \(T\). The factor two supplies the
required upper bound for repeated values.
\end{proof}

\begin{corollary}[The interval bound under finiteness and cuspidal matching]
\leavevmode\label{cor:uniformstrict561}
On each \(E_x\) sum over any set of triples from
Section~\ref{sec:generaleven547} with the following property.
For each triple \(T\) is finite, and for some common embedding
for that triple, each nonsplit geometric point is realized by
a rank-four representation with a cuspidal match over \(K\).
There is one positive integer \(D_{\rm Sp}\) for the whole
interval, with \(\log D_{\rm Sp}\ll V^{C'}\) for a fixed
\(C'>0\), such that
\[
 \sum_{(\alpha,r,\ell)}
     \operatorname{length}_O S(\alpha,r)\log\ell
                          \ \le\ 2n_*t_* \log D_{\rm Sp}.
\]
The same conclusion holds if the matching assumption is imposed
instead on every irreducible constituent of the semisimple
realizations in Theorem~\ref{thm:purityblocks549}.
\end{corollary}
\begin{proof}
Theorem~\ref{thm:commonsptrace433} applies to the lists just
constructed. All its size factors, including \(p,M,N_*,L_*,n_*\)
and \(t_*\), are bounded by fixed powers of \(V\), so its
integer has the asserted logarithmic bound. At every nonsplit
point the multiplicity is one by
Proposition~\ref{prop:onecompanion560}. The trace difference
at the same \(p\) is nonzero by Lemma~\ref{lem:cuspidaltrace433}
or by its triangle inequality using the stronger magnitude
bound just proved. The strict-length inequality of
Theorem~\ref{thm:lengthmultiplicity550} now gives a bound of
\(2v_\ell(D_{\rm Sp})\) for the individual length. At a
fixed \(\ell\) there are at most \(n_*\) geometric characters
and at most \(t_*\) weights. Summing and using the prime
factorization of \(D_{\rm Sp}\) gives the formula. Under the
alternative hypothesis, Theorem~\ref{thm:evenplanecompanion560}
excludes the two-plane realization at each point, giving the
first hypothesis there.
\end{proof}

The corollary constructs the level, spectral list and common
integer from field and interval data. Finiteness and existence
of the indicated cuspidal matches are still needed to apply it
to all cases in Section~\ref{sec:missing}.


\subsection{The prescribed auxiliary conductors on the period interval}
\label{sec:sourcelevels563}

Use the polynomial-conductor intervals of Theorem~\ref{thm:conductor232}.
Write \(F\) for the conductor of the original character \(\psi\),
and \(o\) for its order. At \(x\in[V,V^{3/2}]\) let
\[
 E_x=K_\psi C_J(\sqrt{-q},\sqrt{-d_0}),\qquad
 N=\operatorname{lcm}(2^{J+2},F)q d_0 .
\]
Let \(N^\natural\) use the largest \(J\) on the interval in this
formula, and let \(n_*=[E_*:\QQ]\) be the degree at that layer.
Fix a nontrivial character \(\alpha\) of \(E_x\) and a weight
\(2\le r<t(x)\) with \(\alpha(-1)=(-1)^r\). Throughout this
subsection impose
\[
 \ell\ge5,\qquad \ell\notin\mathcal B_\psi,\qquad
 \ell\nmid n_*N^\natural,\qquad \ell-1\mid 2r,\qquad
 (\bar\alpha\bar\kappa_\ell^{\,r})|_{D_\ell}\ne1 .
\]
Here \(\mathcal B_\psi\) consists of \(3\) and the odd primes
dividing \(oFp_K h_*p_1\). The fixed integers \(p_K,h_*\) are
the auxiliary prime and quadratic class number of the source.
The integer \(p_1\) is its least prime congruent to one modulo
the odd part of \(o\) and prime to \(2Fp_K\). Dirichlet's
theorem gives such a prime. Keep
\(\Psi=\alpha\kappa_\ell^r\), and let \(\omega_\ell\) be the
Teichm\"uller lift of \(\bar\kappa_\ell\).

The chosen modulus gives exactly the required field. Indeed,
the subgroup of \((\ZZ/N)^\times\) acting trivially on the
four factors is cut out by \(\ker\psi\) modulo \(F\), by
\(\{\pm1\}\) modulo \(2^{J+2}\), and by the squares modulo
each of \(q,d_0\). The last two descriptions use
\(q,d_0\equiv3\bmod4\). The fixed field of the intersection
is the compositum. The auxiliary primes also satisfy the size
condition of the source on this very interval. In the proof of
Theorem~\ref{thm:conductor232} both exceed \(H^{401W}\), while
\(V^2\le H^{K_0}\) and \(W>K_0+2\). Since
\(t_*=\lfloor V^{1203/2}\rfloor\), they exceed \(2t_*+1\)
and every fixed integer \(F,o,p_K,h_*,p_1\) at large \(V\).

Write \(\alpha=\alpha_2\alpha'\) for its factorization into
the characters at two and at the odd primes. Put
\(2^{c_2}=\operatorname{cond}(\alpha_2)\), allowing \(c_2=0\).
Then \(\operatorname{ord}(\alpha')\mid 2o\).
In fact \(\operatorname{Gal}(E_x/\QQ)\) embeds in the product
of the four Galois groups. A character of a subgroup of a
finite abelian group extends to the group with values in the
roots of unity. To see this, adjoin one generator at a time,
choosing a root of the prescribed value of its first power in
the subgroup. Hence \(\alpha\) is a product of a power of
\(\psi\), a dyadic character and the two quadratic characters.
Taking the odd-prime components proves the divisibility claim.
An \(\alpha\) of odd order factors through \(K_\psi\), since
\([E_x:K_\psi]\) is a power of two. It has conductor dividing
\(F\), trivial dyadic part and order dividing the odd part of
\(o\). In general \(\operatorname{ord}(\alpha)\mid 2^{c_2+1}o\).

Define
\[
 f_0=\operatorname{ord}_{\operatorname{lcm}(4,\operatorname{ord}\alpha')}
                  (\ell),\qquad n_0=v_2(\ell^{f_0}-1).
\]
Here the order is defined since \(\operatorname{ord}\alpha'\)
divides \(\operatorname{ord}\alpha\), which divides \(n_*\).
On the indicated triples the source calls the even-order case
II and the odd-order case III. Its component away from two
is \(\alpha'\), since \(\alpha\) has order prime to \(\ell\).
Its component away from \(\ell\) is all of \(\alpha\), since
\(\ell\nmid N\). Thus its exact dyadic conductor \(M_\alpha\)
(which also depends on \(\ell,x\)) is the least power of two
at least the respective quantity
\[
\begin{array}{ll}
 \max\{\ell N4^{n_0},\ 2^{c_2+6}o\ell^3 F^2\},
       &\operatorname{ord}\alpha\ \text{even},\\[2pt]
 \max\{\ell Fp_1 4^{n_0},\ 8\ell^2\},
       &\operatorname{ord}\alpha\ \text{odd}.
\end{array}
\]
These are the definitions in
\srcloc{Development/AutomorphicArithmetic/AutomorphicVarieties/HeightAndVanishing/DirichletPrimeToLPartMulLPowerPart}{211}
and the character prescription in
\srcloc{Development/AutomorphicArithmetic/AutomorphicVarieties/AutomorphicCases/SymplecticGSp4}{13}.

\begin{lemma}[Local separation at the prescribed conductor]
\leavevmode\label{lem:prescribedaux563}
There is an odd character \(\xi\) unramified at \(\ell\)
whose dyadic component has conductor exactly \(M_\alpha\).
In case II it is purely dyadic. In case III its only other
component has conductor \(p_1\) and order \(\operatorname{ord}\alpha\).
For \(A=\kappa_\ell\omega_\ell^{-1}\xi\) the following holds.
At an arithmetic Frobenius \(s\in D_\ell\) with \(\kappa_\ell(s)=1\),
\(\bar A(s)\) is outside the subgroup \(\bar\Psi(D_\ell)\).
There is an arithmetic Frobenius \(h\) separating all four
residual entries of \(d=(A\Psi,\Psi,A,1)\).
\end{lemma}
\begin{proof}
Write \(M_\alpha=2^d\) in this proof. Send \(-1\) to \(-1\)
and \(5\) to a primitive \(2^{d-2}\)-th root to define
\(\xi_2\) on the units modulo \(2^d\). Its restriction to the
kernel of reduction modulo \(2^{d-1}\) is nontrivial, so it
is primitive. In II use \(\xi=\xi_2\). In III the cyclic
group of units modulo \(p_1\) admits a character \(\xi_0\)
of order \(\operatorname{ord}\alpha\), since this nontrivial
odd order divides \(p_1-1\). It is even and primitive.
Set \(\xi=\xi_2\xi_0\). All finite character values used here
have order prime to \(\ell\) and are available after a finite
unramified coefficient extension. Under such an extension a
finite length over the integer ring is preserved. For the
strict cohomology groups this base change follows by writing
the new ring as a finite free module over the old one, which
commutes with continuous cochains and their restriction maps.

Inertia at \(\ell\) surjects onto the cyclotomic image, so we
may choose \(s\) with the asserted normalization. Put
\(a=\pm\ell\equiv1\bmod4\) and \(b=v_2(a-1)\), so that
\(2^b\le\ell+1\). The inequality \(2^d>\ell+1\) gives
\(b<d\). Each squaring of \(a\) adds one to the valuation of
its difference from one. Since \(\xi_2\) is faithful on the
cyclic subgroup congruent to one modulo four, its value at
\(a\) has order \(2^{d-b}\ge M_\alpha/(\ell+1)>2\).
Multiplication by a sign preserves this order. The odd-order
factor in III also preserves its two-part, as does reduction
in odd characteristic. Thus the two-part of the order of
\(\bar A(s)=\bar\xi(\ell)\) is at least \(M_\alpha/(\ell+1)\).

On \(s\), \(\bar\Psi\) has the value of \(\bar\alpha\); its
inertia image has order at most two since \(\ell-1\mid 2r\).
Consequently the two-part of the order of \(\bar\Psi(D_\ell)\)
is at most \(2^{c_2+1}o\) in II and at most two in III.
The second term in the defining maximum for \(M_\alpha\)
makes \(M_\alpha/(\ell+1)\) larger than these bounds.
If \(\bar\Psi(s)\ne1\) take \(h=s\). Otherwise the local
nontriviality assumption gives an inertia element by which
one can multiply \(s\) to obtain \(\bar\Psi(h)\ne1\).
The residual \(\bar A\) is unramified there. Therefore
\(\bar A(h)\bar\Psi(h),\bar\Psi(h),\bar A(h),1\)
are pairwise distinct.
\end{proof}

\begin{proposition}[The six odd strict groups with the source characters]
\leavevmode\label{prop:gooddata563}
Set \(\mu=A\Psi\). For each of
\[
       \Gamma=A^{\pm1},\quad(\Psi/A)^{\pm1},\quad(A\Psi)^{\pm1}
\]
the strict group for \((O/\ell^j)(\Gamma)\) vanishes at every
\(j\ge1\). The strict conditions are zero on \(D_\ell\)
and zero on inertia at every other finite prime. The prescribed
characters and marks satisfy the local and residual hypotheses of
\srcloc{Development/AutomorphicArithmetic/AutomorphicVarieties/SymplecticGSp4/DecompSigmaOrd}{90}.
\end{proposition}
\begin{proof}
The Teichm\"uller character lifting \(\bar\Gamma\) is odd.
Factor its primitive Dirichlet character as \(\phi\eta\)
with odd and dyadic conductors \(f\) and \(M_\alpha\).
For the latter assertion observe that its dyadic factor is
\(\xi_2^{\pm1}\) times a character of conductor at most
\(2^{c_2}\), and \(d=\log_2 M_\alpha>c_2\) in both cases.
It is nontrivial on \(1+2^{d-1}\). At the odd primes the
factors come from \(\alpha'\), \(\omega_\ell^r\) and, in III,
\(\xi_0\), or their inverses. Each has order dividing
\(\ell^{f_0}-1\), since \(\operatorname{ord}\xi_0=
\operatorname{ord}\alpha'\) in III and \(\ell-1\mid\ell^{f_0}-1\).
The conductor \(f\) divides \(\ell N\) in II and \(\ell Fp_1\)
in III. Hence \(M_\alpha\ge f4^{n_0}\), and
\(\ell^{f_0}\equiv1\bmod4\). Theorem~\ref{thm:fourier145}
makes both \(B_{1,\phi\eta}\) and its inverse-character
counterpart units above \(\ell\).

The residual symplectic diagonal has prime-to-\(\ell\) inertia
image away from \(\ell\). Consequently a cocycle unramified on
the kernel of that diagonal on inertia gives zero as a class on
inertia, by averaging on the finite quotient. A strict residual
cocycle thus vanishes on \(I_v\cap\ker\bar\Gamma\) at every
finite \(v\), since a coboundary vanishes on that subgroup.
Corollary~\ref{cor:coc145} applies. Its class-group argument
uses the unit just obtained in the integral Stickelberger
annihilator. To lift the vanishing, subtract a lift of a global
coboundary from a cocycle at depth \(j>1\) to make its values
divisible by \(\ell\). At \(D_\ell\) each \(\Gamma(s)-1\) is a
unit by Lemma~\ref{lem:prescribedaux563}. The element exhibiting
the local coboundary is therefore also divisible by \(\ell\).
Elsewhere inertia has finite prime-to-\(\ell\) image. A nontrivial
inertia action supplies the same unit difference, and at a
trivial one the cocycle already vanishes on inertia. Division
by \(\ell\) gives a strict cocycle at depth \(j-1\), so induction
finishes the proof of vanishing.

For the remaining data take all prime divisors of \(2\ell N\)
in II and of \(2\ell Np_1\) in III. Outside these primes the
finite characters are unramified. The auxiliary character is
odd, \(\Psi\) is even, and the only prime in the distinguished
local set is \(\ell\), since \(\ell\nmid\operatorname{ord}\alpha\).
At \(s\) one has \(\bar\xi(s)\bar\Psi(s)^j\ne1\) for all
integers \(j\). Both each odd root and its Tate dual have
nontrivial reduction there, using \(\bar\kappa_\ell(s)=1\).
Choose \(g_0\) outside the union of the two proper kernels of
\(\bar A,\bar\Psi\). In the procyclic closure of \(s\) the
sequence \(s^{\ell^{j!}}\) converges to its prime-to-\(\ell\)
part, since this holds in each finite quotient. Finally the
intersection of inertia with the kernels of the finite
characters has finite index in inertia. Its cyclotomic image
contains an open subgroup of \(1+\ell\ZZ_\ell\), giving the
required marked element of infinite cyclotomic order. This
constructs the two marks of Section~\ref{sec:twomarks557}
together with the separating \(h\) of Lemma~\ref{lem:prescribedaux563}.
\end{proof}

\begin{theorem}[A common strict bound from the prescribed arithmetic data]
\leavevmode\label{thm:sourcecommon563}
For each \(x\) select any of the above triples for which the
following three assertions hold for the characters just constructed.
Use the global multiplier \(\mu\), the two marks, the ordinary
local component through \(d\) at \(\ell\), and the kernel of
the residual diagonal on inertia away from \(\ell\).
\begin{enumerate}
\item In every such integral lift over a complete Noetherian
local coefficient algebra, the subalgebra generated by
products of matrix entries of total torus weight zero is
finite over \(O\).
\item For each characteristic-zero specialization preserving
both coordinate planes \(02,13\), the matrices on plane \(02\)
span a proper subspace of its endomorphisms.
\item There is a single coefficient embedding for the triple
under which every absolutely irreducible characteristic-zero
specialization has a regular algebraic cuspidal match over
\(\QQ(i)\), with the good Frobenius polynomials of
Section~\ref{sec:onerealization560}.
\end{enumerate}
Then a single positive integer \(D_{\rm Sp}^{\rm src}\)
for the interval satisfies, for every \(x\) on it,
\[
 \sum_{\text{selected }(\alpha,r,\ell)\text{ on }E_x}
       \operatorname{length}_O S(\alpha,r)\log\ell
     \ \le\ 2n_*t_* \log D_{\rm Sp}^{\rm src}\ \ll\ V^C
\]
for a constant \(C\) depending only on the fixed interval
construction. The automorphic requirement of the source over
imaginary quadratic fields implies assertion~3 here.
\end{theorem}
\begin{proof}
First apply Corollary~\ref{cor:finmark557} with these data.
Its argument uses the two disjoint residual Frobenius pairs,
the two distinct residual eigenvalues of the compressed
self-adjoint mark, and the separated \(h\), all verified above.
It identifies the finiteness in assertion~1 with finiteness of
the trace ring \(T\). Proposition~\ref{prop:actualcomponent557}
identifies its local component with the ordered ordinary flag,
also over nonreduced quotients. In the proof of
Theorem~\ref{thm:allglobalroots547}, the odd-root inputs are
strict vanishing and the local dimensions. The former is
Proposition~\ref{prop:gooddata563}. At \(\ell\) the vanishing
of the invariants of a residual odd root and its Tate dual
gives these dimensions by local duality and the Euler
characteristic formula. The localization calculation at each
finite depth now follows from the same finite duality formula
as in Theorem~\ref{thm:oddrootlocalization430}. The even root
is still \(\Psi\), so the integral negative-root module measures
exactly \(S(\alpha,r)\). The trace presentation and the
strict-length inequality of Theorem~\ref{thm:lengthmultiplicity550}
therefore apply with the prescribed \(\xi\).

At a nonsplit point the two-plane alternative in
Theorem~\ref{thm:purityblocks549} has an integral block
realization. Putting it in the two-mark frame preserves those
blocks. Indeed the projector from the prime-to-\(\ell\) part
of \(s\) is the identity on the first block and zero on the
second, as follows by reduction and its prime-to-\(\ell\)
eigenvalues. The second change of frame diagonalizes a
compression within the first block and uses the dual basis on
the second. The first plane is absolutely irreducible. Its
image spans an algebra since group elements are closed under
multiplication, and this algebra is the full matrix algebra
by Burnside's theorem. Assertion~2 excludes this case.
For a rank-four point Proposition~\ref{prop:strongall559}
gives irreducibility on \(G_{\QQ(i)}\), using the ordinary
line, the separated \(h\), and the simple spectrum of every
positive power of the marked inertia element. This proves
the last claim about the source's matching requirement.
Under assertion~3, Proposition~\ref{prop:onecompanion560}
gives multiplicity one at each such point.

It remains to bound the changing conductors in a common
integer. Put \(L_*\) as in Section~\ref{sec:symplecticstrict429}.
Since \(f_0\le4n_*\), repeated squaring gives
\(2^{n_0}\le(\ell^2-1)f_0\). Indeed for \(f_0\) odd the
valuation is \(v_2(\ell-1)\), and for \(f_0\) even it is
\(v_2(\ell-1)+v_2(\ell+1)+v_2(f_0)-1\). Also
\(2^{c_2}\le N^\natural\), \(N\mid N^\natural\),
and \(\ell\le L_*\). Taking the least power of two above
a positive real number costs a factor of at most two.
The maximum \(M_{\max}\) of the prescribed dyadic conductors
on the interval is therefore polynomial in \(V\) and a multiple
of each; take it to be one for an empty family.
The conductor of \(\xi\) divides \(p_1 M_{\max}\) and its
order divides \(n_* M_{\max}\).

Use \(d'_\ell=\ell N^\natural p_1M_{\max}\) in
Proposition~\ref{prop:levelofmatch561}. It bounds all four
local character conductors, including the tame character
\(\omega_\ell^{-1}\) at \(\ell\). Thus its fourth power is a
common principal level at that \(\ell\). Choose one rational
prime \(p\equiv1\bmod4\) between \(Y'\) and \(2Y'\) with
\(Y'=32+L_*+N^\natural+p_1\), using the prime number theorem
in this fixed progression. It is prime to every \(3d'_\ell\).
Theorem~\ref{thm:commoncusplist561} now supplies one multiset
of integral cuspidal Hecke traces for each weight, with
polynomial cardinality and conjugate magnitudes at most
\(4p^{(r+1)/2}\). Use the polynomial in the proof of
Theorem~\ref{thm:commonsptrace433} with the first root of
unity of order dividing
\(\operatorname{lcm}(\ell-1,n_*M_{\max})\) and the second
of order dividing \(n_*\). It includes the split trace for
every one of the present choices of \(\xi\). Its values on
the Hecke multiset have nonzero integral product, by the
same archimedean separation and Galois stability. Multiplying
their absolute products over the weights gives the common
integer. The polynomial degrees, the multiset sizes, the
weights and \(\log p\) are all polynomial in \(V\).
The strict-length inequality at multiplicity one bounds
each length by \(2v_\ell(D_{\rm Sp}^{\rm src})\). There are
at most \(n_*t_*\) pairs \((\alpha,r)\) on a given field at
each \(\ell\). Prime factorization proves the stated bound.
\end{proof}

The conductor choices and the odd strict vanishings in this
subsection are unconditional. The three global arithmetic
assertions in Theorem~\ref{thm:sourcecommon563} are hypotheses.
They concern the actual deformation spaces for the fields
used by the integral projected theta period.


\subsection{Two character sums for the even planes}
\label{sec:mirrors565}

Use the fields \(E_x\), the character \(\alpha\), the colliding
weight \(r\) and the prime \(\ell\) of Section~\ref{sec:sourcelevels563}.
In the first two results we also allow \(\bar\Psi|_{D_\ell}=1\).
Put \(k=O/\ell O\), \(\tau=\bar\Psi=\bar\alpha\bar\kappa_\ell^r\)
and let \(\omega=\bar\kappa_\ell\). The global \(\tau\) is even
and nontrivial. Indeed at the half residue it is nontrivial on
inertia and at the zero residue it is the nontrivial \(\bar\alpha\).
For \(\epsilon\in\{+1,-1\}\) set
\[
          M_\epsilon=k(\tau^\epsilon),\qquad
          \delta_\epsilon=\omega\tau^{-\epsilon}.
\]
Let \(H_{\rm rel}(M)\) denote the classes in \(H^1(\QQ,M)\)
whose restriction to inertia is a coboundary away from \(\ell\).
Let \(S(M)\) impose in addition zero restriction as a class at
\(D_\ell\).

\begin{lemma}[The two ordinary class groups for relaxed extensions]
\leavevmode\label{lem:reflectedrel565}
Let \(L_\epsilon\) be the cyclic field cut out by \(\delta_\epsilon\)
and \(\Delta_\epsilon=\operatorname{Gal}(L_\epsilon/\QQ)\).
Then
\[
 \dim_k H_{\rm rel}(M_\epsilon)
   =\dim_k \operatorname{Hom}_{\Delta_\epsilon}
          \bigl(\Cl(L_\epsilon),k(\delta_\epsilon)\bigr).
\]
The homomorphisms on the right are additive.
Write \(f_\epsilon\) for the conductor of the primitive Dirichlet
character lifting \(\delta_\epsilon\) by Teichm\"uller representatives.
It divides \(\ell N\) and is divisible by \(\ell\).
Take \(1\le b_\epsilon\le f_\epsilon\) prime to \(f_\epsilon\)
with \(b_\epsilon\bmod\ell\ne\delta_\epsilon(b_\epsilon)\).
Such a choice exists. If the finite sum
\[
 C_\epsilon=\sum_{a\in(\ZZ/f_\epsilon)^\times}
        \left\lfloor\frac{b_\epsilon[a]}{f_\epsilon}\right\rfloor
                         \delta_\epsilon(a)^{-1}\quad\in k
\]
is nonzero, then \(H_{\rm rel}(M_\epsilon)=0\). Here \([a]\)
is the representative in \([1,f_\epsilon]\). Whether the sum
is nonzero is independent of the allowed \(b_\epsilon\).
\end{lemma}
\begin{proof}
The Tate dual character of \(M_\epsilon\) is \(\delta_\epsilon\).
It is odd and its inertia character at \(\ell\) is
\(\omega^{1-\epsilon r}\). The collision condition gives
\(r\equiv0\) or \((\ell-1)/2\bmod\ell-1\); for \(\ell\ge5\)
the indicated inertia character is nontrivial. This proves the
assertion about \(f_\epsilon\), since \(\alpha\) has conductor
dividing \(N\). Global invariants of both modules vanish.

Use the finite Poitou--Tate order formula as in
Theorem~\ref{thm:oddrootlocalization430}. At each finite
\(q\ne\ell\) the unramified local conditions are orthogonal
complements. Indeed their cup product factors through the
unramified quotient of cohomological dimension one. On either
module the unramified dimension equals the local invariant
dimension, by the kernel and cokernel formula for Frobenius
minus one on inertia invariants. These dimensions add to the
full local \(H^1\) dimension by local duality and the local
Euler characteristic. Thus the order-formula factor at
\(q\ne\ell\) is one. At \(\ell\) the relaxed group takes all
of \(H^1\). Its dimension exceeds the invariant dimension
by one, since the dual has no local invariants. This holds
also when \(M_\epsilon\) is locally trivial. At the real place
the factor is \(1/|k|\), since \(M_\epsilon\) is even and
\(\ell\) is odd. These two factors cancel. The orthogonal
condition at \(\ell\) is zero, so
\(\dim_k H_{\rm rel}(M_\epsilon)=\dim_k S(k(\delta_\epsilon))\).

The group \(\Delta_\epsilon\) has order prime to \(\ell\).
Inflation--restriction and averaging identify classes on the
right with equivariant homomorphisms on \(G_{L_\epsilon}\)
unramified everywhere and zero on decomposition at primes
over \(\ell\). The same argument on inertia and decomposition
groups verifies the local converse. By class reciprocity these
homomorphisms factor through the ordinary class group, with
the prime classes above \(\ell\) killed. In fact every equivariant
homomorphism on the ordinary class group already kills those
classes. An inertia subgroup at \(\ell\) fixes each such prime
in the abelian field, whereas it acts nontrivially on the target
line. This gives the dimension formula.

For the finite test abbreviate \(\delta=\delta_\epsilon\),
\(f=f_\epsilon\), and use the full cyclotomic field
\(\QQ(\zeta_f)\). A strict cocycle restricts there to an
unramified additive homomorphism on the class group, equivariant
for \(\delta\). Set
\(\Theta_f=\sum_{a}([a]/f)\sigma_a^{-1}\), where
\(\sigma_a(\zeta_f)=\zeta_f^a\). The integral Stickelberger
element \((b-\sigma_b)\Theta_f\) annihilates that class group.
Its coefficient at \(\sigma_a^{-1}\) is
\((b[a]-[ba])/f\). Hence its action on the target is precisely
the scalar \(C_\epsilon\). If this scalar is nonzero the
restricted homomorphism is zero. The original class then
inflates from the abelian cyclotomic Galois group. Its
\(H^1\) on this line is zero: commuting complex conjugation
with each group element gives a coboundary since \(\delta\)
is odd and two is a unit. This proof uses no restriction on
the index \([\QQ(\zeta_f):L_\epsilon]\).

Since \(\ell\mid f\), the character \(\omega\) also factors
modulo \(f\). Its ratio to \(\delta\) is the nontrivial
\(\tau^\epsilon\). An allowable \(b\) therefore exists.
Evaluating the same group-ring element at the Teichm\"uller
lift \(\widetilde\delta\) gives
\[
 (b-\widetilde\delta(b))\,B_{1,\widetilde\delta^{-1}},\qquad
 B_{1,\widetilde\delta^{-1}}
     =f^{-1}\sum_a [a]\,\widetilde\delta(a)^{-1}.
\]
The first factor is a unit in \(O\), and the product is integral
by the floor formula. Reduction proves the asserted independence
of the choice of \(b\).~\cite{XII-IX-V-stick}
\end{proof}

\begin{theorem}[Triangularity and a vanishing test at the colliding weights]
\leavevmode\label{thm:planeunits565}
If either \(C_+\ne0\) or \(C_-\ne0\), a continuous two-dimensional
lift with diagonal residual characters of ratio \(\tau\) and
prime-to-\(\ell\) inertia image away from \(\ell\) is triangular
after a change of frame. This holds over a complete Noetherian
local coefficient ring with finite residue field \(k\), including
in positive characteristic. In particular there is no absolutely
irreducible plane of this kind. If \(C_-\ne0\), then also
\[
                         S(\alpha,r)=0
\]
for the full strict group with divisible coefficients, even when
\(\tau|_{D_\ell}=1\).
\end{theorem}
\begin{proof}
Let \(R\) be the coefficient ring and \(\mathfrak m\) its maximal
ideal. Choose \(h\in G_\QQ\) with \(\tau(h)\ne1\). The matrix
at \(h\) has two distinct residual roots. Hensel factorization
and the two spectral idempotents give a basis of its free
eigenlines lifting the original coordinates. In that frame write
\[
      \rho(g)=
        \begin{pmatrix}a(g)&b(g)\\ c(g)&d(g)\end{pmatrix},
      \qquad \bar a/\bar d=\tau .
\]
Let \(J_b\) be the ideal generated by the \(b(g)\). If it is
nonzero, Nakayama's lemma supplies a \(k\)-linear functional
\(\lambda\) on \(J_b/\mathfrak m J_b\) nonzero on some entry.
Reduction of
\(b(gg')=a(g)b(g')+b(g)d(g')\) then makes
\(z(g)=\lambda(b(g))/\bar d(g)\) a nonzero cocycle on
\(k(\tau)\). It is continuous. Indeed Artin--Rees gives
\(J_b\cap\mathfrak m^N\subseteq\mathfrak m J_b\) for large
\(N\), so the finite linear quotient is continuous in the
subspace topology. The cocycle is zero at \(h\). A coboundary
with this value is zero everywhere, because \(\tau(h)-1\)
is a unit. On inertia away from \(\ell\) the cocycle factors
through its finite prime-to-\(\ell\) matrix image and is a
coboundary there by averaging. Thus \(J_b\ne0\) implies
\(H_{\rm rel}(M_+)\ne0\). Interchanging the two lines proves
that a nonzero lower-entry ideal implies \(H_{\rm rel}(M_-)\ne0\).
Lemma~\ref{lem:reflectedrel565} now kills the corresponding
ideal under the asserted test. This proves triangularity over \(R\).

For the strict group put \(D=(\operatorname{Frac}(O)/O)(\Psi^{-1})\).
Its global invariants vanish because \(\tau^{-1}\ne1\).
The coefficient sequence for multiplication by \(\ell\) gives
a unique finite-coefficient class in \(H^1(\QQ,D[\ell])\) above
every element of \(S(\alpha,r)[\ell]\). At any finite
\(q\ne\ell\) the inertia action on \(D\) is finite of order
prime to \(\ell\). Its invariant group is either zero or the
entire divisible module. Thus \(D^{I_q}/\ell D^{I_q}=0\).
The finite-coefficient class is unramified there by the same
coefficient sequence on inertia. At \(\ell\) impose the relaxed
condition. We have obtained an injection of \(S(\alpha,r)[\ell]\)
into \(H_{\rm rel}(M_-)\). When \(C_-\ne0\) the latter vanishes.
Every continuous cocycle from a compact group to the discrete
\(D\) has finite image, so its class is killed by a power of
\(\ell\). A nonzero such group contains an element of order
\(\ell\). This gives the full vanishing.
\end{proof}

\begin{corollary}[Removing the plane hypothesis on the original interval]
\leavevmode\label{cor:planeinterval565}
In Theorem~\ref{thm:sourcecommon563}, on triples with
\(C_+\ne0\) one may omit assertion~2. For triples with
\(C_-\ne0\) the strict length is zero independently of all
three assertions. The two tests use a polynomial number of
finite-field summands on the whole interval.
\end{corollary}
\begin{proof}
The two-plane alternative of Theorem~\ref{thm:purityblocks549}
provides an integral lift of the plane \(13\) with residual
characters \(\tau,1\) and prime-to-\(\ell\) inertia image away
from \(\ell\). Its absolute irreducibility contradicts
Theorem~\ref{thm:planeunits565} when \(C_+\ne0\).
The proof of Theorem~\ref{thm:sourcecommon563} uses assertion~2
only to remove this alternative. The remaining length claim
follows directly from Theorem~\ref{thm:planeunits565}.
On the interval all characters come from \(E_*\), so there
are at most \(n_*\) geometric choices. At each of the at most
\(L_*\) coefficient primes, \(\omega^r\) has only the two
indicated possibilities. Each conductor divides \(\ell N^\natural\).
The two sums require together at most \(4n_* L_*^2 N^\natural\)
terms. All these factors are polynomial in \(V\); the calculation
uses ordinary residue arithmetic even at moduli divisible by
the coefficient prime.
\end{proof}

For example take \(\psi=1\), \(\ell=13\), and the even
quartic character modulo \(16\) specified by
\(\bar\alpha(5)=5\in\mathbb F_{13}\). Its field
\(\QQ(\zeta_{16})^+\) occurs in each sufficiently large \(E_x\).
For \(12\mid r\) in the stated range one has \(\tau=\bar\alpha\)
and \(\tau(13)=8\). Both primitive conductors are
\(16\cdot13=208\), and \(b=3\) is allowed for both.
The rows below give the data for grouping the sum by \(c\bmod16\).
In the last row sum over \(a=c+16j\), \(0\le j\le12\),
omitting \(13\mid a\); each term is \(\lfloor3a/208\rfloor/a\)
in \(\mathbb F_{13}\).
\[
 \begin{array}{c|rrrrrrrr}
 c                         &1&3&5&7&9&11&13&15\\ \hline
 \tau(c)                   &1&8&5&12&12&5&8&1\\
 \text{sum in the class }c &9&0&3&7&7&3&0&9
 \end{array}
\]
Multiplying the last row by \(\tau(c)\) and summing gives
\(C_+=8\). Using \(\tau(c)^{-1}\) gives \(C_-=0\).
Thus one test alone excludes the plane on every such weight.
A zero value of the other sum leaves open the strict-length
question for that test. More generally, simultaneous zero sums
are necessary for either even plane to occur. Global finiteness
and cuspidal realization of rank-four points still enter the
length bound when the negative-root test gives zero.


\subsection{Critical values and the trace congruence at an even plane}
\label{sec:planecongruence566}

Work on the fields and triples of Section~\ref{sec:sourcelevels563},
with the separating \(h\in D_\ell\) and the prescribed \(\xi\).
For a nonsplit characteristic-zero point with the two-plane
realization of Theorem~\ref{thm:purityblocks549}, write \(P\)
for the plane on coordinates \(13\) in its integral \(h\)-eigenbasis.
Its determinant is \(\Psi=\alpha\kappa_\ell^r\), and its ordered
residual lines have characters \(\bar\Psi,1\). Take a finite
coefficient field \(J'\) of realization, with integer ring \(O'\),
uniformizer \(\pi\), and \(e'=v_\pi(\ell)\). We use
\(v_\ell=v_\pi/e'\) on \(J'\). Put
\[
 P(g)=\begin{pmatrix}a(g)&b(g)\\c(g)&d(g)\end{pmatrix},
 \qquad
 \mathfrak b=(b(g):g\in G_\QQ),\qquad
 \mathfrak c=(c(g):g\in G_\QQ)
\]
as ideals in \(O'\). Both ideals are nonzero by irreducibility
of the plane and both lie in the maximal ideal by reduction.
Denote the entire four-dimensional realization by \(\rho\)
and the split representation by \(d_{\rm sp}\).

\begin{theorem}[The exact trace ideal and its critical-value bound]
\leavevmode\label{thm:planedepth566}
The ideal generated by
\(\operatorname{tr}\rho(\operatorname{Frob}_p)-
\operatorname{tr}d_{\rm sp}(\operatorname{Frob}_p)\)
at the unramified rational primes away from \(\ell\) and the
prescribed level is exactly \(\mathfrak b\mathfrak c\).
Writing \(m=v_\pi(\mathfrak b\mathfrak c)\), one has
\[
                  0<\frac{m}{e'}\le v_\ell(L(1-r,\alpha)).
\]
Here \(L\) is the primitive Dirichlet \(L\)-function in the
specified coefficient embedding. The result does not require
automorphic realization of the plane.
\end{theorem}
\begin{proof}
Modulo \(\mathfrak b\mathfrak c\), the two diagonal functions
on \(P\) are multiplicative by the matrix product law. On
inertia at \(\ell\) they are \(\Psi,1\) because the local
representation is upper triangular in the ordered \(h\)-frame.
At every other finite prime the inertia matrix image has finite
prime-to-\(\ell\) order. A character on it congruent to one has
trivial image in the principal units of \(O'/\pi^m\).
Indeed those units form an \(\ell\)-group. Apply this to the
ratios of the two diagonal characters to \(\Psi,1\) on inertia,
using also that the finite character \(\alpha\) has order
prime to \(\ell\). Both ratios are everywhere unramified at
finite primes. Abelian class reciprocity over \(\QQ\) gives
triviality, since the narrow class number of \(\QQ\) is one.
Consequently \(a\equiv\Psi\), \(d\equiv1\bmod\pi^m\) globally.

The alternating pairing identifies the other plane with
\(\mu P^\vee\simeq A\otimes P\), since \(\det P=\Psi\)
and \(\mu=A\Psi\). Thus all four-dimensional trace differences
are in \(\mathfrak b\mathfrak c\). For the converse take
\(\mathfrak t\) generated by these differences on all of
\(G_\QQ\). Modulo \(\mathfrak t\), the four eigenvalues of
\(\rho(h)\) are the corresponding entries of \(d_{\rm sp}(h)\).
In fact their characteristic polynomials agree by the traces at
\(h,h^2,h^3,h^4\) and Newton's identities, with \(1,2,3,4\)
invertible. Each actual eigenvalue is congruent to its labeled
split value: evaluation in the split polynomial makes all the
other factors units, by residual separation. For any \(g\), the
four trace identities at \(h^j g\), \(0\le j\le3\), now say that
the four diagonal differences of the matrices at \(g\) are killed
by the Vandermonde matrix of those split eigenvalues. Its
determinant is a unit. In particular \(a(g)\equiv\Psi(g)\)
and \(d(g)\equiv1\bmod\mathfrak t\). The matrix product then
gives \(b(g)c(g')=a(gg')-a(g)a(g')\in\mathfrak t\) for all
\(g,g'\). This proves equality with \(\mathfrak b\mathfrak c\).
Some element attains its least valuation \(m\). Reduction of
the two continuous representations modulo \(\pi^{m+1}\) has
finite image. Chebotarev supplies a good rational prime in the
same conjugacy class in that finite quotient, so the asserted
good-prime ideal also has valuation \(m\).

Choose \(j=v_\pi(\mathfrak b)\). The map
\(z(g)=\pi^{-j}b(g)\bmod\pi^m\) is a cocycle for
\((O'/\pi^m)(\Psi)\). This follows from
\(b(gg')=a(g)b(g')+b(g)d(g')\) and the diagonal congruences
just proved. Some \(z(g)\) is a unit. Since \(z(h)=0\) and
\(\Psi(h)-1\) is a unit, a scalar multiple of \(z\) that is
a coboundary must vanish identically. Thus its class generates
a copy of \(O'/\pi^m\). It is unramified away from \(\ell\)
by inertia averaging. The map from these finite coefficients
into \((J'/O')(\Psi)\) by division by \(\pi^m\) is injective
on global \(H^1\), by the multiplication sequence and the
vanishing of the global invariants of \((J'/O')(\Psi)\).
Its image belongs to the relaxed positive-twist group.
Theorem~\ref{thm:criticalall310} computes the length of that
group over \(O\) before this coefficient extension. Its
hypotheses hold because \(r\ge2\), the sign of \(\alpha\)
is \((-1)^r\), and \(\bar\Psi\ne1\). Tensoring with \(O'\)
commutes with continuous cochains and the inertia conditions,
since \(O'\) is finite free over \(O\); each \(O\)-length is
multiplied by \(e'\) as an \(O'\)-length. The length of the
cyclic submodule is \(m\). This proves the inequality.
\end{proof}

\begin{lemma}[One integral trace test for finitely many points]
\leavevmode\label{lem:averagevaluation566}
Let \(O\) be unramified over \(\ZZ_\ell\), with residue field
of cardinality \(q_0\). Let \(t_1,\ldots,t_j\) be integral
trace differences evaluated at a finite set of geometric points
\(\sigma\), with nonnegative integer weights \(n_\sigma\).
Suppose that at each point some difference is nonzero and put
\(w_\sigma=\min_i v_\ell(\sigma(t_i))\). Some \(O\)-linear
combination \(t\) is nonzero at every point and satisfies
\[
   \sum_\sigma n_\sigma v_\ell(\sigma(t))
       \ \le\ \sum_\sigma n_\sigma
                       \left(w_\sigma+\frac{1}{q_0-1}\right).
\]
\end{lemma}
\begin{proof}
Choose the coefficients independently with normalized additive
Haar measure on \(O\). At a point fix an index attaining the
minimum and condition on all other coefficients. After dividing
by the value of that entry, the linear combination at the point
is \(c-z\) with \(c\) uniform in \(O\) and \(z\) a fixed element
of the valuation ring of a finite extension. For \(n\le t<n+1\),
\(n\ge0\), the condition \(v_\ell(c-z)>t\) confines \(c\) to
at most one class modulo \(\ell^{n+1}O\). In fact any two
solutions differ by an element of integral valuation greater
than \(t\). The probability is at most \(q_0^{-n-1}\).
Integration over \(t\ge0\) bounds the expectation of the
excess over \(w_\sigma\) by \(1/(q_0-1)\). An exact zero has
probability zero. Summing the expectations over the finitely
many points proves the existence claim.
\end{proof}

\begin{corollary}[Counting the plane contribution by its multiplicity]
\leavevmode\label{cor:planeweight566}
Suppose \(T\) is finite for each selected triple and that all
its absolutely irreducible rank-four points admit the cuspidal
matches of Theorem~\ref{thm:sourcecommon563}, with one
embedding per triple. Let \(P_{\rm pl}=\sum n_\sigma\) on
the two-plane points of that triple, and let \(H_{\rm irr}\)
be its number of irreducible points in the nonsplit quotient.
There exist two positive integers \(D_{\rm cusp},D_{\rm crit}\)
for the entire interval, with logarithms polynomial in \(V\),
such that on every selected triple
\[
\begin{split}
 \operatorname{length}_O S(\alpha,r)
 \ \le\ 2\bigg(&v_\ell(D_{\rm cusp})
         +P_{\rm pl}\, v_\ell(D_{\rm crit})\\[-1pt]
         &+\frac{H_{\rm irr}+P_{\rm pl}}{|k|-1}\bigg).
\end{split}
\]
In particular a polynomial bound on \(P_{\rm pl}\) suffices
for a polynomial bound on the weighted sum of strict lengths,
even if some even planes occur.
\end{corollary}
\begin{proof}
For each \(2\le r\le t_*\) and each nontrivial geometric
character \(\gamma\) of \(E_*\) of sign \((-1)^r\), let
\(f_\gamma\) be its primitive conductor. Define \(D_{\rm crit}\)
to be the absolute value of the product, over these pairs, of
\[
                  r f_\gamma(r+1)!\,L(1-r,\gamma).
\]
The primitive functional equation and the Euler product at
\(r>1\) give nonvanishing with this parity. Each factor is an
algebraic integer. Indeed the generalized Bernoulli formula is
\[
       L(1-r,\gamma)=
       -\frac{f_\gamma^{\,r-1}}{r}
               \sum_{a=1}^{f_\gamma}\gamma(a)B_r(a/f_\gamma).
\]
In the polynomial expansion the Bernoulli denominators through
order \(r\) divide \((r+1)!\). This follows by induction from
the Bernoulli recurrence after multiplication by \((j+1)!\)
at order \(j\). Multiplication also by \(f_\gamma\) clears
the one possible remaining power of the conductor. The product
is Galois invariant, hence rational integral. Every conjugate
factor has magnitude at most \(\exp(Cr\log(2rf_\gamma))\).
For example the expansion of \(z/(e^z-1)\) on \(|z|=1\)
gives \(|B_j|\le C_0 j!\) by Cauchy's estimate; insertion in
the displayed finite sum gives the claimed bound. Therefore
\[
          \log D_{\rm crit}\ll n_*t_*^2\log(2t_*N^\natural).
\]
At the chosen embedding all the other factors are integral.
Since the clearing multiplier is an ordinary integer,
Theorem~\ref{thm:planedepth566} bounds the minimum valuation
at each plane by \(v_\ell(D_{\rm crit})\).

Use the common prime and the cuspidal multiset in the proof of
Theorem~\ref{thm:sourcecommon563}. Their construction uses
only the matching points and gives \(D_{\rm cusp}\) even
when additional points have two-plane realizations.
The matching points have multiplicity one by
Proposition~\ref{prop:onecompanion560}. The sum of their
valuations at this prime is at most \(v_\ell(D_{\rm cusp})\)
by Theorem~\ref{thm:commonsptrace433}. Enlarge this single
test by a good prime attaining the minimum of
Theorem~\ref{thm:planedepth566} for each plane point.
Finiteness of \(T\) makes this a finite set. At the matching
points the minimum can only decrease; at each of the other
points it has the ideal valuation already computed. Apply
Lemma~\ref{lem:averagevaluation566} with the multiplicities
of Theorem~\ref{thm:lengthmultiplicity550}. The resulting
integral linear combination belongs to its trace-difference
ideal and is nonzero at all nonsplit points. That theorem gives
the displayed bound. With an empty nonsplit fiber the same
theorem gives zero directly.

Finally \(H_{\rm irr}\) is bounded by the size of the common
Hecke multiset, which is polynomial on the interval. If also
\(P_{\rm pl}\le V^C\), sum the displayed estimate with weight
\(\log\ell\). At each \(\ell\) there are at most \(n_*t_*\)
pairs \((\alpha,r)\). For each integer \(D\) the sum of
\(v_\ell(D)\log\ell\) is \(\log D\), and
\(\sum_{5\le\ell\le L_*}\log\ell/(|k|-1)
\ll(1+\log L_*)^2\) by comparison with the sum over all
integers. These bounds give the claimed polynomial estimate.
\end{proof}

The critical value bounds the depth at each even plane without
an assumption on its origin. The total multiplicity of the plane
points, finiteness of the trace algebra, and the cuspidal matches
of the rank-four points remain to be controlled in general.


\subsection{The total congruence depth on the one-variable plane locus}
\label{sec:totalplane567}

For this subsection put
\[
 a_\pm=\dim_k S(k(\bar\Psi^{\pm1})),\qquad
                 a_+=0,\quad a_-=1,
\]
on the fields and triples of Section~\ref{sec:sourcelevels563}.
These are conditions on the two strict residual groups. The
positive relaxed group can still be nonzero. We bound all the
two-plane trace congruences together when each point is counted
with the multiplicity coming from the plane equations.

\begin{theorem}[The sum over the plane equations]
\leavevmode\label{thm:intrinsicplane567}
Let \(T_{\rm pl}\) be the quotient of the trace ring obtained
by setting the two rank-four trace coordinates \(V,W\) to
zero, retaining the block coordinate \(U\) defined below.
Its quotient by \(\ell\)-power torsion is finite over \(O\).
For a geometric characteristic-zero point with \(U_\sigma\ne0\)
put
\[
             m_\sigma^{\rm pl}
              =\dim_{\overline J}
                     (T_{\rm pl}\otimes_O\overline J)_\sigma .
\]
Then
\[
          \sum_{U_\sigma\ne0}
                m_\sigma^{\rm pl} v_\ell(U_\sigma)
                       \ \le\ v_\ell(L(1-r,\alpha)).
\]
At each such point \(v_\ell(U_\sigma)\) is the minimum
valuation of the good-prime trace differences of
Theorem~\ref{thm:planedepth566}.
\end{theorem}
\begin{proof}
In Theorem~\ref{thm:allglobalroots547} the three positive
coordinates after taking the first torus invariants are
\(x,v,w\) of weights \(1,1,2\); the single negative coordinate
is \(y\) of weight \(-1\). Here \(x\) is the upper even-root
coordinate on the local ordinary flag. There is one relation
\(E\) of weight \(1\) and two \(F_1,F_2\) of weight \(-1\).
The trace coordinates of Theorem~\ref{thm:allroottrace548} are
\(U=xy,\ V=vy,\ W=wy^2\). On the block locus set \(v=w=0\).
Equivalently one may first set the three off-block entries
to zero in the full two-torus chart. Their invariant products
are \(v,w\); every neutral monomial involving either of these
contains \(V\) or \(W\). Projecting the equations onto their
weights therefore gives the same quotient on trace coordinates.
For some \(g,f_1,f_2\in O[[U]]\), one has
\[
 E(x,0,0,y)=xg(xy),\qquad F_i(x,0,0,y)=yf_i(xy).
\]
Multiplication by the opposite-weight monomials gives precisely
the neutral ideal
\[
               T_{\rm pl}=O[[U]]/(Ug(U),\,Uf_1(U),\,Uf_2(U)).
\]

The scalar \(g(0)\) has valuation exactly \(v_\ell(L(1-r,\alpha))\).
To prove this, over \((O/\ell^j)[\varepsilon]/(\varepsilon^2)\)
use only the positive root and put \(x=\varepsilon z\).
The sole equation is \(g(0)z=0\) in \(O/\ell^j\).
In matrix entries these lifts have the form
\((1+\varepsilon b(g')X_+)d_{\rm sp}(g')\) with
\(X_+=E_{02}+E_{13}\) as in Theorem~\ref{thm:integralroots547}.
The group law gives cocycles for \((O/\ell^j)(\Psi)\).
They are normalized by \(b(h)=0\). Conversely every relaxed
class has a unique representative with this normalization
because \(\Psi(h)-1\) is a unit. On \(D_\ell\) an upper-root
cocycle preserves the ordinary flag and fixes its graded
characters. Away from \(\ell\) the inertia averaging used in
Theorem~\ref{thm:integralroots547} gives precisely the prescribed
kernel condition. Thus these lifts give all the relaxed classes,
with linear correspondence since \(\varepsilon^2=0\). The
other coordinates vanish on these lifts: their first-order
entries have other weights and there are no additional
relative positive even-root coordinates since \(a_+=0\).

Multiplication by \(\ell^j\) on
\((\operatorname{Frac}(O)/O)(\Psi)\) identifies this finite
group with the \(\ell^j\)-torsion of the divisible-coefficient
relaxed group. The global invariants vanish and the invariants
on inertia away from \(\ell\) are divisible, giving the
identification also with those local conditions.
Theorem~\ref{thm:criticalall310} gives a finite relaxed group
of length \(v_\ell(L(1-r,\alpha))\). For all sufficiently
large \(j\) its full length is attained, whereas the
length of the kernel of multiplication by \(g(0)\) is
\(\min(j,v_\ell(g(0)))\). This proves the claim, in particular
\(g(0)\ne0\).

Let \(\mu_0\) be the least valuation of a coefficient of \(g\).
Weierstrass preparation gives
\(g=\ell^{\mu_0}u_0 H\) with \(u_0\) a power series unit
and \(H\) a monic distinguished polynomial, possibly one.
Its constant term is nonzero. In the torsion-free quotient
of \(T_{\rm pl}\) one has \(U H(U)=0\). This quotient is
therefore a quotient of \(O[[U]]/(UH)\), finite free as a
module before taking that further quotient by Weierstrass
division. At a nonzero geometric root of \(H\), its local
dimension after imposing the remaining equations is at most
the root multiplicity of \(H\). All roots of a nonconstant
distinguished polynomial have positive valuation. Counting
with the latter multiplicities, the sum of their valuations
equals \(v_\ell(H(0))=v_\ell(g(0))-\mu_0\). This proves
the inequality with \(m_\sigma^{\rm pl}\) and does not impose
finiteness on the full trace ring.

At a point with \(U_\sigma\ne0\) the ideal of specialized
trace differences is \(U_\sigma O'\) in any integral field
of realization. Indeed the augmentation ideal in the block
trace ring is generated by \(U\) and by the trace differences,
using Theorem~\ref{thm:allroottrace548}. Equivalently, to
pass from closed trace generation to this identity after
specialization, approximate \(U\) by polynomials in the
traces and subtract their split values. The resulting ideals
in the integer ring \(O'\) are closed. Chebotarev as in
Theorem~\ref{thm:planedepth566} detects the valuation at a
good prime.
\end{proof}

The local multiplicity of the full characteristic-zero trace
algebra at such a point can also include the directions \(V,W\).
We first record an exact form of the simultaneous choice of a
trace test that works over the original normalized valuation.

\begin{lemma}[Simultaneous minimum valuations after extending coefficients]
\leavevmode\label{lem:exacttests567}
For any triple of Section~\ref{sec:sourcelevels563} with finite
trace ring, let \(\Sigma'\) be the geometric points of the
nonsplit quotient in Theorem~\ref{thm:lengthmultiplicity550}.
Take a finite list of integral trace differences \(t_i\) with
some nonzero evaluation at each point. Then
\[
  \operatorname{length}_O S(\alpha,r)
        \ \le\ 2\sum_{\sigma\in\Sigma'} n_\sigma
                          \min_i v_\ell(\sigma(t_i)).
\]
There is no restriction on \(a_+,a_-\) in this lemma.
\end{lemma}
\begin{proof}
Choose a finite unramified extension \(O_1/O\) whose residue
field \(k_1\) has cardinality greater than \(|\Sigma'|\),
inside the fixed algebraic closure of \(J\). The presentations
and the strict-length inequality of Theorems~\ref{thm:allglobalroots547}
and~\ref{thm:lengthmultiplicity550} may be tensored with \(O_1\).
In detail, finite free scalar extension commutes with the power
series rings, their homogeneous weight spaces, the finite lists
of relations and the minors in the negative-root module.
It preserves lengths over the respective unramified integer
rings. The continuous cochains and all the local restriction
maps computing the strict group also commute with it, by a
finite free basis. The resulting generic algebra tensored with
\(\overline J\) is the original geometric algebra, with the
same multiplicities. These are exactly the algebraic data
used in the strict-length inequality.

At each \(\sigma\) divide the list of evaluations by one
attaining the displayed minimum. The residues of the resulting
integral entries give a nonzero linear form on the list of
coefficient variables over \(\overline{k}\). On \(k_1^j\),
where \(j\) is the list length, its zero set has at most
\(|k_1|^{j-1}\) elements: fix all coordinates except one
with nonzero coefficient. The union of the bad sets misses
a tuple in \(k_1^j\). Lift that tuple to \(O_1^j\). Its
linear combination belongs to the extended trace-difference
ideal and attains the minimum valuation at every point.
The strict-length inequality over \(O_1\) proves the result.
For an empty point set the same inequality gives zero directly.
\end{proof}

\begin{corollary}[The contribution of the two transverse coordinates]
\leavevmode\label{cor:transverseplane567}
Suppose \(a_+=0,a_-=1\), \(T\) is finite, and each rank-four
point has the cuspidal match in
Corollary~\ref{cor:planeweight566}. At a two-plane point let
\(\mathfrak j_\sigma=(V,W)\) in the local geometric Artin
algebra of the nonsplit quotient. Take \(d\ge1\) with
\(\mathfrak j_\sigma^d=0\) at every such point. Then
\[
 \operatorname{length}_O S(\alpha,r)
   \ \le\ 2v_\ell(D_{\rm cusp})+
                  d(d+1)v_\ell(L(1-r,\alpha)).
\]
The same fields admit a polynomial bound on the weighted sum
over any such triples for which \(d\) is polynomial in the
interval parameter \(V_{\!0}\), where \(V_{\!0}\) here denotes
the left endpoint of the interval. If all the local geometric
algebras at plane points are reduced, one can take \(d=1\).
\end{corollary}
\begin{proof}
At \(U_\sigma\ne0\) the trace ring and its nonsplit quotient
have the same local geometric algebra. This follows directly
from the neutral relations: the equations \(U f_i=0\) give
\(f_i=0\) when \(U\) is a unit. Equivalently the split summand
has disappeared after localization. Quotienting the local
algebra by \(V,W\) therefore gives precisely the corresponding
local algebra of \(T_{\rm pl}\). In the filtration by
\(\mathfrak j_\sigma^i\) the \(i\)-th quotient is generated
over this block algebra by the \(i+1\) monomials of degree
\(i\) in \(V,W\). Summing through \(i=d-1\) gives
\[
                   n_\sigma\le \tfrac12d(d+1)m_\sigma^{\rm pl}.
\]
Nilpotence holds at each point since \(V,W\) are zero in its
residue field. In a reduced local Artin algebra their ideal
is zero.

Use the common prime at the rank-four points and one prime
attaining each plane minimum in Lemma~\ref{lem:exacttests567}.
The matching points have multiplicity one and their total
valuation is bounded by \(v_\ell(D_{\rm cusp})\) as in
Corollary~\ref{cor:planeweight566}. For the others combine
the multiplicity inequality with
Theorem~\ref{thm:intrinsicplane567}. This proves the
individual bound without any factor counting distinct
plane points. To sum over each field and all the indicated
weights, bound the critical valuations by those of the
common integer \(D_{\rm crit}\) of
Corollary~\ref{cor:planeweight566} and apply prime
factorization to the two common integers. There are at most
\(n_*t_*\) terms at each prime.
\end{proof}

The one-variable block calculation gives unconditional control
of the weighted sum internal to that block. The two directions
in the full algebra still require a uniform nilpotence bound
when nonreduced plane points occur.


\subsection{Ordinary extensions between the two even planes}
\label{sec:ordinarycross568}

Fix a triple and its characters on the fields of
Section~\ref{sec:sourcelevels563}. In the order \(0,1,2,3\) used
for the local flag, put \(P=\langle e_1,e_3\rangle\) and
\(P'=\langle e_0,e_2\rangle\). In this subsection the given
representation is block diagonal on \(P'\oplus P\). The two
summands are Lagrangians, and the pairing identifies
\(P'=\mu P^\vee\). We allow a finite local Artin coefficient ring
\(R\) over \(O\) with residue field \(k'\) extending \(k\).
Assume the block representation reduces to the specified diagonal,
is ordinary in the specified frame at \(\ell\), is diagonal at
\(h\), and kills \(N_q=I_q\cap\ker\bar d\) at the finite
level primes \(q\ne\ell\). It is unramified outside the level
and \(\ell\). Reductions of the integral realizations of
Theorem~\ref{thm:planedepth566} satisfy these hypotheses.

An infinitesimal symplectic map with only its \(P\)-to-\(P'\)
block nonzero is a symmetric form on \(P\) with values in the
multiplier line. Reversing the two spaces gives the other
off-diagonal block. Denote the resulting \(G_\QQ\)-modules by
\[
      M_+=\mu\,\operatorname{Sym}^2(P^\vee),\qquad
      M_-=\mu^{-1}\operatorname{Sym}^2(P).
\]
Both have rank three. The signs in \(M_\pm\) specify the
direction of a map and are independent of the notation \(S^\pm\)
for the even character groups. Let \(B_\pm\subset M_\pm\) be
the strictly upper triangular part in the original order
\(0,1,2,3\). Conjugation by an upper triangular block
representation preserves it on \(D_\ell\). In \(B_+\) one
may freely specify entries \(01\) and \(03\), with \(23\)
determined by \(01\). In \(B_-\) the free entry is \(12\).
Thus these are direct summands as \(R\)-modules of ranks two
and one. Write \(H_{\rm rel}(N)\) for the global \(H^1\)
classes unramified at every finite place away from \(\ell\),
and \(H_{\rm str}(N)\) for the kernel of their restriction to
\(D_\ell\).

\begin{lemma}[Localization for the entire off-diagonal block]
\leavevmode\label{lem:crosslocal568}
For either \(M=M_\pm\) and every finite \(R\)-module \(C\),
restriction gives a natural isomorphism
\[
       H_{\rm rel}(M\otimes_R C)
           \ \xrightarrow{\ \sim\ }\ H^1(D_\ell,M\otimes_R C).
\]
The right side is \(H^1(D_\ell,M)\otimes_R C\), where
\(H^1(D_\ell,M)\) is free of rank three. For \(B=B_\pm\),
the map from \(H^1(D_\ell,B\otimes_R C)\) to this group is
injective. Its source and cokernel are obtained by tensoring
free \(R\)-modules of ranks \(\operatorname{rank} B\) and
\(3-\operatorname{rank}B\), respectively, with \(C\).
\end{lemma}
\begin{proof}
The residual lines of \(M_+\) have characters \(A/\Psi,A,A\Psi\)
after reduction. Those for \(M_-\) have their inverses. Each
line is odd and has zero strict group by
Proposition~\ref{prop:gooddata563}. Scalar extension to \(k'\)
preserves this calculation by additivity of cochains. The
residual eigenvalues of the element \(s\) of that proposition
on these six lines and on their finite Tate duals all differ
from one. Hence the local invariants at \(\ell\) vanish on
every module built by extensions of these lines and likewise
on its Tate dual. This applies also to the successive
constituents of \(B\) and \(M/B\), since the residual
representation is diagonal in the specified frame.

Take a composition series of \(C\) and argue by induction.
For a short exact sequence \(0\to C_1\to C\to C_2\to0\)
in such a series, tensoring with \(M\) is exact. The image
of a strict class in the last module vanishes by induction.
Subtract a global coboundary lifting a witness in the
quotient. The resulting cocycle takes values in
\(M\otimes C_1\). At \(\ell\) it is a coboundary, and a
witness for that coboundary belongs to \(M\otimes C_1\)
because its image in \(M\otimes C_2\) is invariant there.
Away from \(\ell\), being unramified is equivalent to the
cocycle vanishing on the kernel of the residual inertia
action: the kernel acts trivially here, and its quotient in
inertia has order prime to \(\ell\), on which averaging kills
\(H^1\). Vanishing on this kernel survives passage to the
submodule. The class in the submodule is therefore strict,
and induction kills it. The split residual lines give the
starting case. This proves \(H_{\rm str}(M\otimes C)=0\).

Apply the finite Poitou--Tate order formula and the exact
sequence for enlarging local conditions as in
Theorem~\ref{thm:oddrootlocalization430}
\cite{XII-IX-V-milne98}. For the strict conditions both global
invariant groups vanish by the local invariant calculation.
The factor at \(\ell\) is one. Away from \(\ell\), Frobenius
minus one on inertia invariants has kernel and cokernel of
equal orders. At the real place all residual characters are
odd, so the invariants vanish there as well, since two is
invertible. The order formula gives zero dual Selmer group.
The exact sequence from the strict to the relaxed conditions
now gives the asserted localization isomorphism.

For each of \(D=M,B,M/B\), local duality gives
\(H^2(D_\ell,D\otimes C)=0\), and \(H^0\) vanishes as above.
Consequently \(C\mapsto H^1(D_\ell,D\otimes C)\) is exact.
Using a finite presentation of \(C\) and additivity on free
modules identifies this functor with
\(H^1(D_\ell,D)\otimes_R C\). The local Euler characteristic
at the residue field gives dimension \(\operatorname{rank}_R D\).
Exactness up a composition series then gives
\(|H^1(D_\ell,D)|=|R|^{\operatorname{rank}_R D}\).
Nakayama's lemma gives a surjection onto this group from a
free module of that rank, and equality of orders makes it an
isomorphism. Applying the local cohomology sequence to
\(0\to B\to M\to M/B\to0\) proves the last assertion.
\end{proof}

\begin{theorem}[One-sided ordinary extensions with fixed blocks]
\leavevmode\label{thm:ordinarycross568}
For every finite \(R\)-module \(C\), the global off-diagonal
classes with prescribed ordinary flag at \(\ell\) in the two
directions are respectively
\[
         H^1(D_\ell,B_+)\otimes_R C,\qquad
         H^1(D_\ell,B_-)\otimes_R C .
\]
They have unique global cocycle representatives whose value
at \(h\) is zero. Each such representative defines an actual
one-sided symplectic extension with the original diagonal
blocks. More precisely, for a finite local \(R\)-algebra
\(R'\) and a nilpotent ideal \(I\subset R'\), use \(C=I\)
and pull the given block representation to \(R'\).
These representatives parametrize all the indicated
extensions in that frame which reduce to the given
representation modulo \(I\). Every one of them has the
original trace at every group element.
\end{theorem}
\begin{proof}
On \(M\), conjugation by the diagonal matrix at \(h\) minus
the identity is invertible: every eigenvalue is a ratio of
distinct residual roots. The same holds on \(B\) and \(M/B\).
Subtracting a unique coboundary therefore normalizes any
cocycle to have zero value at \(h\). A normalized cocycle
on \(D_\ell\) belongs to \(B\otimes C\) if and only if its
class comes from \(H^1(D_\ell,B\otimes C)\). Indeed a
representative in \(B\otimes C\) can be normalized there,
and normalized representatives in \(M\otimes C\) are unique.
The original flag is the coordinate flag in this normalized
frame by its four distinct residual roots. Thus this is
precisely the ordinary condition on an off-diagonal
extension. Lemma~\ref{lem:crosslocal568} supplies the unique
global class with each prescribed local value.

Write \(\rho_0\) for the block representation. For a cocycle
\(z\) with values in \(M_+\otimes I\) put
\(\rho_z(g)=(1+z(g))\rho_0(g)\). Two maps supported in
this same off-diagonal block have zero product, even if
\(I^2\ne0\). The multiplication law for \(\rho_z\) is
therefore exactly the cocycle law. Its multiplier is the
prescribed one because the linear symplectic condition is
the symmetry defining \(M_+\), and the quadratic pairing
of two images is zero on the Lagrangian \(P'\).
The restriction at \(\ell\) preserves the flag, the
normalization at \(h\), and the graded characters. At
every other finite prime the cocycle kills the required
inertia kernel by the averaging argument in the lemma.
Conversely, multiply any extension of this form by
\(\rho_0(g)^{-1}\) on the right and subtract one. The
result is the required normalized cocycle, by the same
matrix identities. Interchanging \(P,P'\) proves the
assertion for \(M_-\). A matrix supported in just one
off-diagonal block changes neither diagonal block in this
formula, proving the trace assertion.
\end{proof}

\begin{corollary}[The three ordinary directions at a plane point]
\leavevmode\label{cor:threecross568}
At an integral two-plane point over a finite extension \(J'\)
of the coefficient field, the normalized one-sided ordinary
classes over \(O'\) are free of ranks two and one. Over \(J'\)
their dimensions are two and one. Both spaces are invisible
to all trace functions on one-sided extensions.
\end{corollary}
\begin{proof}
Reduce a fixed integral block realization through
\(R=O'/\pi^j\). The local cohomology groups for \(B_\pm\)
commute with these reductions by Lemma~\ref{lem:crosslocal568}.
A basis at residue level lifts compatibly to bases at every
level by freeness and Nakayama's lemma. The normalized
global representatives also reduce compatibly by uniqueness.
Taking their inverse limits gives continuous cocycles with
values in the integral lattice, and every normalized integral
cocycle gives such a system. Over \(J'\), a continuous
cocycle on the compact global group has bounded image.
Multiplying it by a power of \(\pi\) puts it in the
lattice. The normalized local flag and inertia conditions
are preserved by this operation and by division in \(J'\).
This identifies the two rational spaces with the integral
ones tensored with \(J'\). The last assertion follows from
the matrix formula in Theorem~\ref{thm:ordinarycross568}.
\end{proof}


\subsection{The mixed obstruction at an even plane}
\label{sec:mixedplane569}

Let \(\rho_0\) be an integral two-plane realization as in
Section~\ref{sec:ordinarycross568}. Work over a finite coefficient
field \(L\) containing its entries. Let \(X_+\) and \(X_-\) be the
spaces of normalized global one-sided ordinary cocycles over
\(L\). They have dimensions two and one by
Corollary~\ref{cor:threecross568}. An element \(g\) acts on
matrix coefficients by conjugation by \(\rho_0(g)\). We take
cochains on \(G_\Sigma=\Gal(\QQ_\Sigma/\QQ)\), where
\(\QQ_\Sigma\) is the maximal extension unramified outside the
finite level primes, \(\ell\), and the real place.

On the neutral block of the fixed-multiplier Lie algebra we
identify \(N=\operatorname{End}_L(P)\) with the endomorphisms
acting by \(w\) on \(P\) and by \(-w^\vee\) on \(P'\).
Write \(N_0=\{w:\operatorname{tr}_P w=0\}\). For \(D=N,N_0\)
let \(\mathcal C^1(D)\) be the space of continuous \(D\)-valued
one-cochains on \(G_\Sigma\) with the following constraints.
They vanish on all \(N_q\) at level primes \(q\ne\ell\).
On \(D_\ell\) they are upper triangular on \(P\), at \(h\)
they are diagonal, and on \(I_\ell\) both diagonal entries are
zero. For a cochain set
\[
     \partial w(g,g')=w(gg')-w(g)-g w(g').
\]
Given \(u\in X_+\) and \(v\in X_-\), form the continuous
two-cochain on the neutral block
\[
     H(u,v)(g,g')=\tfrac12\bigl([u(g),g v(g')]
                                   +[v(g),g u(g')]\bigr).
\]
The bracket is the matrix commutator. Project onto \(N_0\)
by subtracting half the trace on \(P\) times the identity on
\(P\). Write \(K_v(u)\) for the class of this projection in
\[
    C^2_{\rm cts}(G_\Sigma,N_0)\big/\partial\mathcal C^1(N_0).
\]
This quotient allows the obstruction to be stated directly
in continuous cochains. The normalization at \(h\) makes
\(K_v(u)\) bilinear without any choices of representatives.

\begin{proposition}[The traceless obstruction to mixing two extensions]
\leavevmode\label{prop:mixedobstruction569}
Put \(D_2=L[\varepsilon,\eta]/(\varepsilon^2,\eta^2)\).
The two one-sided deformations \(u,v\) on the two axes extend
to one ordinary fixed-multiplier representation over \(D_2\),
with the prescribed inertia kernels and diagonal frame at
\(h\), if and only if \(K_v(u)=0\).
\end{proposition}
\begin{proof}
Every matrix congruent to the identity modulo
\((\varepsilon,\eta)\) admits a logarithm by the degree-two
power series. In fixed multiplier this logarithm lies in the
symplectic Lie algebra: taking the symplectic adjoint
intertwines logarithm, and the logarithm of the inverse has
opposite sign. Write a candidate lift as
\[
    \exp\bigl(\varepsilon u(g)+\eta v(g)+
                        \varepsilon\eta w(g)\bigr)\rho_0(g).
\]
Multiplying the two truncated exponentials gives the group
law exactly when \(\partial w=H(u,v)\). This also shows that
\(H(u,v)\) is a two-cocycle, either by associativity of
three such products or by the cocycle identities for \(u,v\).
The right side is neutral. If any \(w\) works, its neutral
component works, since the block decomposition is stable,
and projection to its diagonal blocks preserves the stated
local constraints. For neutral \(w\), the constraints are
precisely \(w\in\mathcal C^1(N)\). Indeed \(u,v\) already
vanish on the forbidden inertia kernels and at \(h\).
Locally at \(\ell\) they are strictly upper triangular,
so every quadratic product there has zero diagonal and
respects the ordered flag.

It remains to solve for the scalar part of \(w\). The scalar
coefficient on \(P\) has trivial Galois action. We claim
\(H^2(G_\Sigma,L)=0\). At every finite place in \(\Sigma\)
the dual invariant group for \(L(1)\) is zero, hence so is
the local \(H^2\). Global duality reduces the claim to the
injectivity of localization on \(H^1(G_\Sigma,L(1))\)
\cite{XII-IX-V-milne98}. By Kummer theory and the trivial ideal
class group of \(\ZZ[1/\prod_{q\in\Sigma_f}q]\), the latter
space is spanned freely over \(L\) by the Kummer classes of
the rational primes in the finite subset \(\Sigma_f\).
Their valuations at those same primes form the identity
matrix. The claim follows.

There is consequently a continuous scalar primitive \(w_s\)
for the scalar component of \(H(u,v)\). At \(q\ne\ell\)
the two cocycles and the action there factor through
\(D_q/N_q\). This local quotient has finite inertia of
order prime to \(\ell\) and procyclic unramified quotient,
so its \(H^2\) for \(L\) is zero. Choose a scalar primitive
on that quotient. Its difference from \(w_s|_{D_q}\) is
an additive character. By local reciprocity every such
\(L\)-valued character of \(\QQ_q^\times\) kills the units,
since their open pro-\(q\) subgroup and their finite
quotient both have zero image. Thus \(w_s\) already
vanishes on \(N_q\). At \(\ell\) the scalar two-cochain
is zero because both products in the bracket are strictly
upper triangular. The restriction of \(w_s\) is now an
additive character. Its restriction to inertia is a scalar
multiple of the restriction of \(\log_\ell\kappa_\ell\),
by local reciprocity and the rank one of the principal
units of \(\QQ_\ell\). Subtract that multiple of the global
cyclotomic logarithm. The resulting scalar primitive has
all the required local conditions. Adding it to any
traceless primitive proves sufficiency. Conversely a
permitted neutral primitive projects to a permitted
traceless one. This proves necessity.
\end{proof}

\begin{theorem}[The transverse trace tangent space]
\leavevmode\label{thm:mixedtrace569}
Suppose \(a_+=0,\ a_-=1\) and the trace algebra \(T\) is finite.
At a nonzero two-plane point \(\sigma\) let \(A_\sigma\) be
its local characteristic-zero trace algebra after a finite
extension \(L\) splitting the point. Put
\[
    I=(V,W),\qquad
    m=\dim_L(A_\sigma/I),\qquad n=\dim_L A_\sigma .
\]
For any \(0\ne v\in X_-\), the image of the tangent space of
\(A_\sigma\) in the two trace coordinates \(V,W\) is
isomorphic to \(\ker K_v\). In particular
\[
   \dim_L \mathfrak m_\sigma/\mathfrak m_\sigma^2
           =\dim_L\ker K_v+\mathbf 1_{\{m>1\}},
\]
where \(\mathfrak m_\sigma\) is the maximal ideal of \(A_\sigma\).
\end{theorem}
\begin{proof}
Use the two-torus chart of Theorem~\ref{thm:allglobalroots547}
after eliminating the two neutral parameters. Its variables
are \(a,x,b,c,y\). Before taking invariants the equations
are \(E,F_1,F_2\). They depend only on \(x,ab,bc,y\)
by their zero first torus weight. As in
Theorem~\ref{thm:intrinsicplane567}, the trace variables are
\[
                U=xy,\qquad V=yab,\qquad W=y^2bc .
\]
There are series \(e,f_1,f_2\) in \(U,V,W\) such that
\(yE=e\) and \(F_i=yf_i\). Indeed \(yE\) is neutral and
every monomial of weight \(-1\) is \(y\) times a neutral
monomial. The integral block realization gives values
\(x_0,y_0\) in the maximal ideal of the coefficient integers
with \(x_0y_0=U_\sigma\ne0\); all of \(a,b,c\) are zero
there. Enlarge \(L\) to contain these values if needed.
At this point the trace equations are exactly \(e,f_1,f_2\),
because \(U\) is invertible in \(A_\sigma\). A tangent
vector \((dU,dV,dW)\) must annihilate the differentials
of these three series at \((U_\sigma,0,0)\), and that
condition suffices. The first-order Taylor evaluations
converge since \(U_\sigma\) has positive valuation.
Conversely every map from \(T\otimes_O L\) to dual numbers
at this point is continuous on the finite \(O\)-module \(T\);
approximating its series by polynomials gives the same
Taylor rule.

At the block realization the first torus acts trivially on
the point. Its weights on \(a,c\) are \(+1\), on \(b\)
are \(-1\), and on \(x,y\) are zero. Therefore on a
one-sided positive tangent only \(a,c\) can vary, say by
\(a_1,c_1\), and on a negative tangent only \(b\) can
vary, say by \(b_1\). These coordinates detect the
normalized cocycles: if their variations all vanish the
local six-parameter chart is unchanged, since the
eliminated neutral coordinates also have zero derivative
on a nonzero weight. The normalized local cocycle is then
zero, hence so is the global one by
Lemma~\ref{lem:crosslocal568}. By the dimensions in
Corollary~\ref{cor:threecross568}, the maps to
\((a_1,c_1)\in L^2\) and \(b_1\in L\) are isomorphisms.
In particular \(b_1\ne0\) for the chosen \(v\).

If a mixed lift exists, the trace is constant on each
axis. On \(D_2\) its trace coordinates consequently change
only in degree \(\varepsilon\eta\), with
\[
                  dV=y_0 a_1 b_1,\qquad
                  dW=y_0^2 c_1 b_1 .
\]
To apply the chart also to a continuous lift over \(D_2\),
choose a uniformizer \(\pi\) of \(L\). The coefficients of
its three nilpotent terms on the compact Galois group are
bounded. For large \(k\) all entries lie in the finite
free integer algebra on \(\pi^{-k}\varepsilon,\pi^{-k}\eta\)
inside \(D_2\), with the prescribed residual representation.
Reduction of this algebra through every power of \(\pi\)
and completeness give the desired chart evaluation.

Conversely, take a tangent vector with the displayed
\(dV,dW\) and some \(dU\). In the original chart set
\[
  y=y_0,\quad x=x_0+\varepsilon\eta\,dU/y_0,\quad
  a=\varepsilon a_1,\quad c=\varepsilon c_1,\quad b=\eta b_1 .
\]
All power series converge by Taylor expansion about the
integral point. Their traces are precisely the assigned
tangent vector multiplied by \(\varepsilon\eta\). Thus
\(e=f_1=f_2=0\) there, and invertibility of \(y_0\) over
\(L\) gives all three original equations. The specialization
of the universal matrices is continuous, as one also sees
by scaling the nilpotent generators as above and evaluating
in their complete integer algebra. Its two axis restrictions
have exactly the prescribed normalized cocycles by the
coordinate isomorphisms. Proposition~\ref{prop:mixedobstruction569}
gives the stated equivalence. Since \(y_0b_1\ne0\), the
displayed map from \((a_1,c_1)\) to \((dV,dW)\) is an
isomorphism of the two ambient planes.

Finally, tangents with zero projection to \(V,W\) factor
through \(A_\sigma/I\). This local block algebra is a
quotient of the one-variable completed local ring at
\(U_\sigma\), by Theorem~\ref{thm:intrinsicplane567}.
Its tangent dimension is zero if \(m=1\) and one otherwise.
Taking dimensions proves the formula.
\end{proof}

\begin{corollary}[Multiplicity when the mixed obstruction is injective]
\leavevmode\label{cor:mixedlength569}
If \(m=1\), the full trace algebra at \(\sigma\) is reduced
exactly when \(K_v\) is injective. For arbitrary \(m\), if
\(K_v\) is injective then \(n\le md\) whenever \(I^d=0\).
Thus if all the two-plane points in
Corollary~\ref{cor:transverseplane567} have injective \(K_v\),
its length estimate improves to
\[
  \operatorname{length}_O S(\alpha,r)
      \le 2v_\ell(D_{\rm cusp})+2d\,v_\ell(L(1-r,\alpha)).
\]
\end{corollary}
\begin{proof}
For \(m=1\), the tangent space vanishes exactly when the
maximal ideal is zero, by Nakayama's lemma. For \(m>1\)
and \(K_v\) injective, the cotangent space is a line
generated by the class of \(t=U-U_\sigma\): it maps onto
the cotangent line of the block quotient, where this class
generates. Nakayama's lemma and nilpotence then give
\(A_\sigma=L[t]/(t^n)\). Indeed \(t\) generates the
maximal ideal, successive powers span its nonzero graded
pieces, and the nilpotence order of \(t\) equals the
length of the algebra. Its ideal with quotient of length
\(m\) is \((t^m)\). Hence \(I^d=0\) implies \(n\le md\).
Use this inequality in the sum over two-plane points in
Theorem~\ref{thm:intrinsicplane567}; the rank-four term
is the one in Corollary~\ref{cor:transverseplane567}.
\end{proof}

The scalar part of the mixed equation always has a solution.
The action of the traceless obstruction on the two-dimensional
space \(X_+\) still has to be determined arithmetically.

\section{Auxiliary towers and cuspidal points}\label{nav:chapter-towers-cuspidal-points}

\subsection{An explicit tame tower for the odd-order characters}
\label{sec:tametower571}

Return to the fields \(E_x\) and triples of
Section~\ref{sec:sourcelevels563}. Consider its case III, where
\(m=\operatorname{ord}(\alpha)>1\) is odd. The character
\(\alpha\) has conductor dividing the fixed conductor \(F\),
is unramified at two, and \(r\) is even. Write
\[
   \tau=\bar\alpha\,\bar\omega_\ell^r,\qquad
   \chi=\bar A=\bar\xi_2\bar\xi_0,\qquad
   \Delta=\chi^2/\tau .
\]
The character \(\xi_2\) is odd with conductor \(M_\alpha\) a
power of two; \(\xi_0\) has order \(m\) and conductor \(p_1\).
Bars denote residual characters. All the finite character
values occurring here have order prime to \(\ell\). We
construct the auxiliary totally real and CM fields in
\srcloc{Development/AutomorphicArithmetic/AutomorphicVarieties/WeakShimuraTheory/NearSplit}{141}
for the ordered residual tuple \((\chi\tau,\tau,\chi,1)\).

\begin{lemma}[One auxiliary quadratic pair throughout the interval]
\leavevmode\label{lem:twoprimes571}
There are distinct primes \(p,p'\equiv7\bmod8\), the same
for all the triples on the interval, such that every prime
in the prescribed levels and every coefficient prime of those
triples splits in both \(K=\QQ(\sqrt{-p})\) and
\(K'=\QQ(\sqrt{-p'})\). The prime \(p\) also splits in \(K'\).
They can be chosen with
\[
       \log(pp')\ll_\psi (t_*+1)\log(2t_*N^\natural).
\]
\end{lemma}
\begin{proof}
Take the set \(\mathcal S\) of prime divisors of
\(2N^\natural Fp_1\) together with all primes at most
\(2t_*+1\). A coefficient prime is included since
\(\ell-1\mid 2r\) and \(r<t(x)\le t_*\).
Let \(B=8\prod_{q\in\mathcal S,\,q\ {\rm odd}}q\).
Choose primes \(p\equiv-1\bmod B\) and
\(p'\equiv-1\bmod Bp\). The discriminants are respectively
\(-p,-p'\), both one modulo eight. At every odd prime of
\(\mathcal S\) these discriminants are one modulo that prime,
and \(-p'\) is also one modulo \(p\). The quadratic splitting
criterion gives all the assertions about places. Linnik's
least-prime theorem gives \(p\ll B^{c}\) and
\(p'\ll(Bp)^c\) for an absolute \(c\). The logarithm of
\(B\) has the asserted bound by summing the logarithms of
the integers up to \(2t_*+1\) and adding
\(\log(8N^\natural Fp_1)\).
\end{proof}

\begin{theorem}[An abelian totally real base with separated residual characters]
\leavevmode\label{thm:tametower571}
Let \(H_0\) be the cyclic field cut out by \(\Delta\), and set
\[
          H=H_0\,\QQ(\sqrt{pp'}),\qquad H'=HK .
\]
Then \(H\) is totally real and abelian over \(\QQ\) of even
degree, and \(H'/H\) is CM quadratic and unramified at every
finite place. One has
\[
     [H':\QQ]=4\operatorname{ord}(\Delta)\le 4M_\alpha m,
        \qquad \ell\nmid[H':\QQ].
\]
The degree and logarithmic conductor of \(H'\) are bounded by
a polynomial in the interval parameter, uniformly in the triple.
The following properties hold for these fields.
\begin{enumerate}
\item On \(G_H\) one has \(\tau=\chi^2\). On \(G_{H'}\) there
is an element \(g_0\) with \(\bar\omega_\ell(g_0)=1\) at
which \(\chi\) has order \(2m\). The four ordered residual
characters are separated there.
\item At every place of \(H\) above \(\ell\) the residual
diagonal is trivial on inertia, while \(\bar\omega_\ell\)
is nontrivial on that inertia. The ramification index over
\(\QQ_\ell\) divides two, and the place splits in \(H'\).
\item On \(G_{H'}\) all eight symplectic roots are nontrivial,
and each differs from \(\bar\omega_\ell\). The only quadratic
character factoring through \(\bar\omega_\ell\) which permutes
the four residual characters by multiplication is the trivial
one.
\end{enumerate}
\end{theorem}
\begin{proof}
The residual character \(\Delta\) is even since \(\tau\)
is even and \(\chi\) is odd. Its field \(H_0\) is totally
real and unramified outside \(2Fp_1\ell\). It is disjoint
from \(\QQ(\sqrt{pp'})\) since the latter ramifies at \(p\).
Their compositum has degree \(2\operatorname{ord}(\Delta)\).
Adjoining \(K\) doubles this degree since \(H\) is totally
real. The resulting quadratic extension is totally imaginary,
with conjugation fixing \(H\). A finite prime different
from \(p\) is unramified in \(K/\QQ\), and remains unramified
or splits after base change. At \(p\) one has \(H'=HK'\)
since \(H\) already contains \(\sqrt{pp'}\). The splitting
of \(p\) in \(K'\) gives splitting there after base change
as well.

The order of \(\Delta\) divides \(M_\alpha m\): its dyadic
factor has order dividing \(M_\alpha\), its odd-order factors
have orders dividing \(m\), and \(\bar\omega_\ell^r\) has
order at most two. Also \(\Delta\) takes values in the
multiplicative group of a finite field of characteristic
\(\ell\). These facts give the degree claims. A cyclotomic
modulus for \(H'\) divides \(\ell M_\alpha Fp_1pp'\).
In the formula for \(M_\alpha\) of
Section~\ref{sec:sourcelevels563}, \(f_0\) is bounded in terms
of \(\psi\) since \(\operatorname{ord}\alpha\mid o\).
Thus \(4^{v_2(\ell^{f_0}-1)}\le\ell^{2f_0}\) and \(M_\alpha\)
is polynomial in \(1+\ell\). Lemma~\ref{lem:twoprimes571}
now gives the logarithmic conductor bound.

To construct \(g_0\), choose a residue at
\(\operatorname{cond}(\alpha)\) where \(\bar\alpha\) equals
a primitive \(m\)-th root \(z\). Squaring is an automorphism
of the group of \(m\)-th roots, so choose its unique solution
\(y^2=z\) in that group. Choose a residue at \(p_1\) where
\(\bar\xi_0=y\), the residue \(-1\) at \(M_\alpha\), and
the residue one at \(\ell,p,p'\). These moduli are pairwise
coprime: \(p_1\) is prime to \(2F\), \(\alpha\) has odd
conductor, and the other exclusions follow from
Section~\ref{sec:sourcelevels563} and the two progressions.
The Chinese remainder theorem and the Galois group of the
cyclotomic field give an element of \(G_\QQ\) with all these
values. It has \(\chi=-y\), \(\tau=z=\chi^2\), and fixes
both quadratic fields. It therefore belongs to \(G_{H'}\).
The order of \(-y\) is \(2m\ge6\). On \(G_H\), by the
definition of \(H_0\), the four residual characters are
\(\chi^3,\chi^2,\chi,1\); their values at this element
are distinct.

At \(\ell\) both \(\chi\) and \(\bar\alpha\) are unramified.
The restriction of \(\Delta\) to inertia is
\(\bar\omega_\ell^{-r}\), of order at most two, and the
two auxiliary quadratic fields split there. This proves the
ramification bound and, on inertia over \(H\), forces
\(\tau=\chi=1\). Since \(\bar\omega_\ell\) has order
\(\ell-1\) on rational inertia and this inertia subgroup
has index at most two, its restriction has order at least
\((\ell-1)/2>1\). Splitting at \(\ell\) allows this same
element of inertia to be used on \(G_{H'}\).

The roots are \(\chi^{\pm1},\tau^{\pm1},
(\tau/\chi)^{\pm1},(\chi\tau)^{\pm1}\), with their indicated
multiplicities as root spaces. On \(G_H\) their exponents
as powers of \(\chi\) belong to \(\{\pm1,\pm2,\pm3\}\).
They are all nontrivial at \(g_0\), while they are all
trivial at the inertia element just chosen where
\(\bar\omega_\ell\ne1\). Finally a character factoring
through \(\bar\omega_\ell\) has value one at \(g_0\).
Multiplying by it can fix the set of the four characters
only by fixing each character, because their values there
are distinct. Cancellation gives the trivial character.
\end{proof}

These calculations also check the remaining elementary
conditions on the tower in
\srcloc{Development/AutomorphicArithmetic/AutomorphicVarieties/WeakShimuraTheory/NearSplit}{211}.
Use the original order of the four lines. Each rational level
prime splits in \(K\). For every place of \(H\) over \(\ell\)
choose an extension to the algebraic closure inducing the
prescribed one of the two places of \(K\); this is possible
since \(K\) is split there. The root \(\chi\) has order
greater than four at \(g_0\). At a prime away from \(\ell\),
triviality of the residual diagonal on an inertia element
implies \(\bar\alpha=1\) there since \(\bar\omega_\ell\)
is unramified. It implies \(\alpha=1\) since its order is
prime to \(\ell\). The same order \(m\) kills \(\alpha\)
even on the entire group. Thus there is no distinguished
ramification prime besides \(\ell\). The reduction of the
second root equals the square of the first root on \(G_H\).
This verifies the field and character requirements in case III;
the construction of the automorphic family over this base is
a further step.

\begin{corollary}[Descent of the strict conditions to the period field]
\leavevmode\label{cor:towerdescent571}
Put \(B=E_x\) or \(B=\QQ\), and \(B'=BH'\). For a discrete
\(\ell\)-primary \(G_B\)-module \(W\), restriction identifies
the global \(H^1\) classes over \(B\) that are zero at the
places above \(\ell\) and unramified elsewhere with the
\(\Gal(B'/B)\)-invariants of the group defined by the same
conditions over \(B'\). The assertion also holds with no
condition at \(\ell\).
\end{corollary}
\begin{proof}
The extension \(B'/B\) is Galois of degree dividing
\([H':\QQ]\), hence prime to \(\ell\). Cohomology in positive
degree of this finite group with \(\ell\)-primary coefficients
is zero by averaging. Inflation--restriction identifies the
two global groups before imposing local conditions. At each
coefficient place restriction on \(H^1\) is injective by the
same argument for the local extension. At every other finite
place the inertia subgroups have quotient of order dividing
the global degree, and the same injectivity holds on their
\(H^1\). Consequently a global class has each of the specified
local restrictions zero exactly when its restriction upstairs
does. This gives the assertion on the original \(E_x\) for
each \(x\) in the interval.
\end{proof}


\subsection{The unramified quadratic character in the even-order tower}
\label{sec:eventower572}

For case II of Section~\ref{sec:sourcelevels563} put
\(m=\operatorname{ord}\alpha\). Thus \(m\) is even,
\(\ell\nmid m\), \(\tau=\bar\alpha\bar\omega_\ell^r\), and
\(\chi=\bar A=\bar\xi_2\) is the reduction of the odd
primitive dyadic character, of conductor \(M_\alpha=2^d\).
Write \(2^{c_2}\) for the conductor of the dyadic part
of \(\alpha\), and set \(b=\max(2,c_2)\).
Retain the primes \(p,p'\), modulus \(B\), and imaginary
quadratic fields \(K,K'\) of Lemma~\ref{lem:twoprimes571}.

\begin{lemma}[A real quadratic field common to all even-order triples]
\leavevmode\label{lem:realquadratic572}
One may choose a single prime \(p_3\equiv1\bmod Bpp'\)
throughout the interval. Write \(K_0=\QQ(\sqrt{p_3})\)
and \(\phi\) for its quadratic character. This field
splits at every prime of \(\mathcal S\), and \(p_3\)
splits in both \(K,K'\). Moreover
\(\log p_3\ll_\psi(t_*+1)\log(2t_*N^\natural)\).
Let \(H_0\) be the field of the character \(\tau/\phi\), and
put \(H=H_0\QQ(\sqrt{pp'})\), \(H'=HK\). Then
\[
   [H:\QQ]=2m,\qquad [H':\QQ]=4m,\qquad
                        \ell\nmid [H':\QQ] .
\]
Here \(H\) is totally real abelian, \(H'/H\) is CM
quadratic and unramified at every finite place, and every
place of \(H\) above \(\ell\) splits in \(H'\).
\end{lemma}
\begin{proof}
Use Linnik's theorem once more for the residue class one
modulo \(Bpp'\), using the bound on \(\log(pp')\) from
Lemma~\ref{lem:twoprimes571}. The discriminant \(p_3\)
is one modulo eight and one at each odd prime in
\(\mathcal S\), so these primes split in \(K_0\).
The Dirichlet characters of \(K,K'\) have conductors
\(p,p'\), and both take the value one on \(p_3\).

The character \(\tau/\phi\) is even and its order divides
\(m\), since \(\bar\omega_\ell^r,\phi\) have order at most
two. Choose a residue giving \(\bar\alpha\) order \(m\)
and prescribe one at \(\ell,p_3\). These three moduli
are pairwise coprime, so this gives a value of order \(m\)
for \(\tau/\phi\). The field \(H_0\) is therefore totally
real of degree \(m\) and unramified at \(p,p'\).
The real quadratic field \(\QQ(\sqrt{pp'})\) ramifies at
\(p\), so their intersection is trivial. Adjoining \(K\)
doubles the degree again. At \(p\) write \(H'=HK'\) and
use the splitting of \(p\) in \(K'\). At every other finite
place \(K/\QQ\) is unramified. The conclusions about \(H'/H\)
follow, including splitting at \(\ell\) since \(\ell\) splits
in \(K\).
\end{proof}

\begin{theorem}[The tame even-order tower with its two local witnesses]
\leavevmode\label{thm:eventower572}
On \(G_H\) the even root \(\tau\) equals the character of
\(HK_0/H\). This character is nontrivial, unramified at every
finite place, and trivial on each decomposition group above
\(\ell\). The following further properties hold.
\begin{enumerate}
\item There is a single \(g_0\in G_{H'}\) with
\(\bar\omega_\ell(g_0)=1\), \(\tau(g_0)=-1\), and
\(\chi(g_0)^4\ne1\). The residual entries
\((\chi\tau,\tau,\chi,1)\) there are distinct.
\item At every place above \(\ell\) the inertia index
over \(\QQ_\ell\) divides two. The residual diagonal is
trivial there on inertia, and \(\bar\omega_\ell\) is
nontrivial on that same inertia. Above \(\ell\) the
decomposition group contains an element on which
\(\chi^4\ne1\).
\item At every place above two the inertia group contains
an element on which \(\chi^4\ne1\).
\end{enumerate}
These assertions about local groups hold over \(H\); the
witnesses above two and \(\ell\) may also be taken over \(H'\).
\end{theorem}
\begin{proof}
The equality \(\tau=\phi\) on \(G_H\) is its defining
character equation. Away from \(p_3\) the character \(\phi\)
is unramified. At \(p_3\) the character \(\tau\) is unramified,
so their common restriction to inertia over \(H\) is trivial
there too. Since \(\ell\) splits in \(K_0\), \(\tau\) is
trivial on the entire decomposition group over \(\ell\).

We first record the available variation at two. The group
\(1+2^b\ZZ_2\) in the units modulo \(2^d\) has order
\(2^{d-b}\). It is generated by \(5^{2^{b-2}}\):
one has \(v_2(5^{2^j}-1)=j+2\) by induction on \(j\).
Primitivity of the dyadic character makes it faithful on
the cyclic subgroup \(1+4\ZZ_2\) modulo \(2^d\).
Indeed any nontrivial kernel there contains its subgroup
of order two, which is the kernel of reduction from
\(2^d\) to \(2^{d-1}\). Reduction in characteristic
\(\ell\ne2\) preserves the orders. Thus on \(1+2^b\ZZ_2\)
the order of the image of \(\chi\) is exactly \(2^{d-b}>4\),
by the lower bound for \(M_\alpha\) in case II.
This whole subgroup acts trivially on \(\alpha_2\).

Choose a residue at \(\operatorname{cond}\alpha\) giving
\(\bar\alpha=-1\), which is possible since its image is cyclic
of even order with faithful reduction. Lift its dyadic component
to the units modulo \(2^d\), and vary it by \(1+2^b\ZZ_2\).
The value of \(\bar\alpha\) stays fixed. The possible values
of \(\chi\) form a coset of size greater than four, so at
least one has fourth power different from one. Prescribe
one modulo \(\ell,p,p'\) and choose a nonsquare modulo
\(p_3\), so that \(\phi=-1\). The Chinese remainder theorem
applied to the separate prime powers and cyclotomic Galois
theory realize all these values on one \(g_0\). It fixes
\(H_0,K,K'\), satisfies \(\tau(g_0)=-1\), and belongs to
\(G_{H'}\). This proves nontriviality of \(\phi|_{G_H}\).
The indicated four entries at \(g_0\) are
\(-\chi(g_0),-1,\chi(g_0),1\), pairwise distinct since
\(\ell\) is odd and \(\chi(g_0)^2\ne1\).

At \(\ell\), the restriction of \(\tau/\phi\) to rational
inertia is \(\bar\omega_\ell^r\), of order at most two.
The quadratic fields \(K,K'\) split there. Hence the
ramification index divides two, \(\chi\) is unramified,
and \(\tau=\phi\) is trivial on inertia over \(H\).
The cyclotomic reduction still has order at least
\((\ell-1)/2>1\) there. For the decomposition group choose
the arithmetic Frobenius lift \(s\) at \(\ell\) with
\(\kappa_\ell(s)=1\) from Lemma~\ref{lem:prescribedaux563}.
It satisfies
\[
       \operatorname{ord}(\chi(s))\ge M_\alpha/(\ell+1),
       \qquad j=\operatorname{ord}(\tau(s))\le 2^{c_2+1}o .
\]
The second inequality uses \(\tau(s)=\bar\alpha(\ell)\).
Then \(s^j\) fixes \(H_0,K,K'\), since \(\phi,K,K'\)
are split at \(\ell\). The order of \(\chi(s^j)\) is
at least \(M_\alpha/((\ell+1)j)>4\) by the case II
conductor bound. Conjugating the chosen decomposition
group gives a witness at every place since the characters
are abelian.

Finally, rational inertia at two acts on the \(2^d\)-th
roots of unity through all the dyadic units, by total
ramification. Choose an element mapping to a generator
of the subgroup \(1+2^b\ZZ_2\) modulo \(2^d\). Its
\(\bar\alpha\)-value is one by the preceding computation
and unramifiedness of the odd-conductor part. The
cyclotomic character at \(\ell\) is unramified there.
Also \(K_0,K,K'\) split at two. The element belongs to
inertia over \(H'\) and has \(\chi^4\ne1\). Conjugation
again supplies the other places.
\end{proof}

\begin{corollary}[Uniform descent in both source cases]
\leavevmode\label{cor:bothaux572}
For all the triples of Section~\ref{sec:sourcelevels563}
the two original auxiliary field cases admit the ordered
tower, residual separation, and inertia conditions of
\srcloc{Development/AutomorphicArithmetic/AutomorphicVarieties/WeakShimuraTheory/NearSplit}{211}.
The finite characters used for \(K,K',K_0\) may be fixed
on the entire interval. Each tower has polynomial degree
and polynomial logarithmic conductor, and the strict
and relaxed groups on \(E_x\) descend from \(E_x H'\)
by the invariant formula of Corollary~\ref{cor:towerdescent571}.
\end{corollary}
\begin{proof}
The odd-order assertions are Theorem~\ref{thm:tametower571}.
In the even-order case a cyclotomic modulus for \(H'\)
divides \(\ell Np_3pp'\). Its logarithm has the required
bound by the two prime constructions, and its degree is
\(4m\le4n_*\) by Lemma~\ref{lem:realquadratic572}.
The extension is prime to \(\ell\), which proves descent
by the same inflation--restriction argument.

For the remaining character conditions use the original
ordered tuple. All symplectic roots are ratios of distinct
entries at \(g_0\) and hence nontrivial there. The inertia
witness above \(\ell\) has all residual entries one and
cyclotomic reduction different from one; in particular
every monomial in the residual roots differs from the
cyclotomic character. At \(g_0\) the latter is one, so a
quadratic character factoring through it can permute the
four characters only by fixing each, as in the odd-order
case. The root \(\chi\) has two-power order greater than
four on \(g_0\). The field \(K_0\) required for the even
root is split at two and at \(\ell\). Away from \(\ell\)
triviality of the residual diagonal on inertia forces
\(\alpha=1\) on that element, since \(\alpha\) has order
prime to \(\ell\); its order also kills it globally.
The prescriptions for the marked places above \(\ell\)
are possible because \(K\) splits there, and all level
primes split in \(K\). This proves the stated conditions.
\end{proof}


\subsection{One splitting field for the trace algebras on the interval}
\label{sec:intervaltrace573}

Keep the period fields \(E_x\), the upper field \(E_*\),
and both auxiliary field constructions of
Corollary~\ref{cor:bothaux572}. We compare integral trace
algebras before taking any characteristic-zero points.
For a finite Galois extension \(B/\QQ\), write \(T_B\)
for the closed \(O\)-subalgebra of the global trace algebra
\(T\) generated by the traces on \(G_B\). We use the
topology induced by the complete local representation ring.

\begin{proposition}[A common field of polynomial Galois exponent]
\leavevmode\label{prop:commonaux573}
There is a single finite abelian field \(\mathcal F/\QQ\)
for the interval with the following properties. It contains
all \(E_x\) and all \(H'\) of both source cases, and every
residual diagonal character of every indicated triple is
trivial on \(G_{\mathcal F}\). The degree is prime to each
coefficient prime of the triples. Its Galois exponent \(D\)
and its logarithmic conductor are polynomial in \(V\).
In particular the strict and relaxed groups at each \(E_x\)
can be recovered as the \(\Gal(\mathcal F/E_x)\)-invariants
of the corresponding groups over this single field.
\end{proposition}
\begin{proof}
Let \(M_*\) be the largest of the dyadic conductors
\(M_\alpha\) on the interval, using \(8\) if there are
none. Let \(o_{\rm odd}\) be the odd part of \(o\), and let
\(C_1\) be the unique subfield of degree \(o_{\rm odd}\)
in \(\QQ(\mu_{p_1})\), allowing \(C_1=\QQ\). This exists
since the Galois group is cyclic and
\(p_1\equiv1\bmod o_{\rm odd}\). Let \(\mathcal L\) consist
of the coefficient primes that occur. Take the compositum of
\[
 E_*,\quad \QQ(\mu_{M_*}),\quad C_1,\quad K,\quad K',
 \quad K_0,\quad
 \QQ\bigl(\sqrt{(-1)^{(\ell-1)/2}\ell}\bigr)
                                      \quad(\ell\in\mathcal L).
\]
The last fields are the quadratic subfields of the prime
cyclotomic fields. They kill \(\bar\omega_\ell^r\) since
\(\ell-1\mid2r\). The field \(E_*\) kills \(\bar\alpha\);
the second field kills all dyadic \(\bar\xi_2\); and \(C_1\)
kills every \(\bar\xi_0\) since its order divides
\(o_{\rm odd}\). This proves triviality of the residual
diagonal. The displayed fields kill both characters
\(\chi^2/\tau\) and \(\tau/\phi\) used in the respective
constructions and contain \(KK'K_0\). Thus they contain
every required \(H'\).

The Galois group of the compositum embeds in the product
of those of the constituent fields. All their degrees
are prime to each \(\ell\in\mathcal L\) by the exclusions
\(\ell\nmid n_*o\) and \(\ell\ge5\). Their exponents give
\[
        D\mid\operatorname{lcm}(n_*,M_*,o_{\rm odd},2).
\]
For both formulas of Section~\ref{sec:sourcelevels563},
\(f_0\) is bounded in terms of the fixed \(\psi\),
\(4^{n_0}\le\ell^{2f_0}\), and \(2^{c_2}\le N^\natural\).
Thus \(M_*\) is polynomial in \(V\), since \(N^\natural\)
and \(t_*\) are polynomial and \(\ell\le2t_*+1\).
A cyclotomic modulus containing the compositum is
\[
        8N^\natural M_*p_1pp'p_3
                          \prod_{\ell\in\mathcal L}\ell .
\]
Its logarithm is polynomial by Lemmas~\ref{lem:twoprimes571}
and~\ref{lem:realquadratic572}. Finally apply the
inflation--restriction argument of
Corollary~\ref{cor:towerdescent571} to \(\mathcal F/E_x\).
\end{proof}

\begin{theorem}[Integral traces from primes split on the entire interval]
\leavevmode\label{thm:splitintegral573}
For every indicated triple the global trace algebra \(T\)
is a finite module over \(T_{\mathcal F}\). This subalgebra
is generated topologically by good rational Frobenius traces
at primes split completely in \(\mathcal F\).
Every trace \(\operatorname{tr}\rho(g)\) satisfies a monic
equation of degree \(D^4\) with coefficients in
\(T_{\mathcal F}\), including on nonreduced coefficient
rings. In particular \(T\) is finite over \(O\) exactly
when \(T_{\mathcal F}\) is finite over \(O\).
\end{theorem}
\begin{proof}
For a positive integer \(d\) and indeterminates \(X_i\)
form the polynomial
\[
  \prod_{\zeta_1,\ldots,\zeta_4\in\mu_d}
       \left(Z-\sum_{i=1}^{4}\zeta_i X_i\right).
\]
Its coefficients as a polynomial in \(Z,X_i\) are
cyclotomic integers fixed by every automorphism of the
cyclotomic field, hence rational integers. Multiplying
any one \(X_i\) by a \(d\)-th root of unity preserves
the product. Every exponent of \(X_i\) in a nonzero
monomial is consequently divisible by \(d\). Permutation
of the four variables also preserves the product. The
fundamental theorem of symmetric polynomials now writes
the coefficients in \(Z\) as integral polynomials in
the elementary symmetric functions of the \(X_i^d\).
The polynomial is monic of degree \(d^4\) and vanishes at
\(Z=\sum_i X_i\).

For any four-by-four matrix \(A\) this gives an equation
for \(\operatorname{tr}A\) in the characteristic coefficients
of \(A^d\). To justify it over arbitrary coefficient rings,
first take the generic matrix over the polynomial ring
in sixteen independent entries over \(\ZZ\). The equation
holds over a splitting field of its characteristic polynomial
by the eigenvalue computation. Its two sides belong to
the integer polynomial ring, which injects into that
field. Specialization therefore gives the identity even
over rings with nilpotents.

Use \(d=D\) and \(A=\rho(g)\). Then \(g^D\in G_{\mathcal F}\).
Newton's identities express the characteristic coefficients
of \(A^D\) using the traces of its first four powers with
denominators invertible in \(O\), since \(\ell\ge5\).
These coefficients all lie in \(T_{\mathcal F}\).
By Corollary~\ref{cor:tracesplit548} the augmentation ideal
of \(T\) is generated by a finite list of good-prime trace
differences. Their traces generate \(T\) topologically
over \(O\): adjoining \(\ell\) to those differences
generates its maximal ideal, and successive approximation
in that ideal gives the assertion. If the list has length
\(b\), adjoining its traces to \(T_{\mathcal F}\) as an
ordinary algebra requires at most \(D^{4b}\) module
generators by the monic equations. This algebra is already
closed: the coefficient algebra \(T_{\mathcal F}\) is
compact, and the finite list of monomials gives a continuous
map from a finite power of it onto their module span.
Thus the span is all of \(T\). One direction of the
finiteness equivalence follows by transitivity. In the
other, \(T_{\mathcal F}\) is a submodule of a finite
module over the Noetherian ring \(O\).

For any \(g\in G_{\mathcal F}\) and any finite quotient of
the integral representation through a power of the maximal
ideal, take the compositum of its fixed field with
\(\mathcal F\). The conjugacy class of \(g\) in this finite
quotient restricts to the identity on \(\mathcal F\).
Chebotarev gives a good rational prime in the class, with
the same trace in the quotient. Varying the precision proves
the asserted generation by split primes.
\end{proof}

\begin{proposition}[Equality of the integral algebras over a separated base]
\leavevmode\label{prop:separatedalgebra573}
For the particular auxiliary field \(H'\) of each triple,
one has \(T=T_{H'}\). The augmentation ideal is generated
as a closed ideal by the good-prime trace differences at
primes split in \(H'\). In case II with
\(\ell-1\nmid r\), both equalities hold with \(E_xH'\) in
place of \(H'\).
\end{proposition}
\begin{proof}
Here is an integral group algebra argument. Let \(\mathcal R\)
be the complete Noetherian local coefficient ring of the
given lift with finite residue field, and \(G\) its compact
matrix image. Suppose a normal subgroup has closed image
\(G_1\) of index prime to \(\ell\), containing an element
whose residual diagonal has four distinct entries. The
residual image \(D_0\) of \(G\) is a finite group of diagonal
matrices of order prime to \(\ell\). Its kernel \(P\) is
pro-\(\ell\), since each successive principal congruence
quotient embeds in an additive group in residual
characteristic. Hence \(P\subset G_1\).

There exists a finite subgroup of \(G\) mapping
isomorphically to \(D_0\). In detail, lift the identity
section of \(D_0\) successively through the finite images
modulo powers of the maximal ideal. At each step the
inverse image of the previous section is an extension of
\(D_0\) by an abelian \(\ell\)-group. A set-theoretic
section has a multiplicative defect in its second
cohomology, which vanishes by averaging since \(|D_0|\)
is invertible on the kernel. Correcting the section
gives a homomorphism lifting the preceding one. Compactness
of \(G\) identifies the inverse limit of the compatible
sections with a section into \(G\). Denote its image by
\(D_1\), so \(G=P D_1\).

Take \(s\in D_1\) lifting the separated residual element.
It is in \(G_1\) because \(P\subset G_1\).
The eigenvalues of \(s\) are the Teichm\"uller lifts
of its four residual entries in \(O\): its finite order
is prime to \(\ell\), and the corresponding simple roots
of unity lift uniquely. More explicitly the roots of
the prime-to-\(\ell\) order polynomial split over \(O\)
for this element since its order is the order of its
diagonal reduction. Lagrange interpolation in \(s\) gives
four projectors in \(O[s]\) with rank-one free images over
\(\mathcal R\). Every element of the abelian group \(D_1\)
acts by a scalar on each such line. Its scalar is again
the Teichm\"uller lift of its residual character, by
uniqueness of the prime-to-\(\ell\) root of unity in that
residue class. The coefficient ring \(O\) contains these
character values by the choice of coefficients in
Section~\ref{sec:sourcelevels563}. It follows that
\(D_1\subset O[s]\), whence
\[
                        O[G]=O[G_1]
                \quad\hbox{inside } M_4(\mathcal R).
\]
Taking traces shows that every global trace is already an
\(O\)-linear combination of traces on the subgroup. Applying
the augmentation of the universal lift to that same equality
and subtracting gives the assertion for trace differences.
Chebotarev at finite precision replaces the subgroup elements
by split rational primes, as above.

For \(G_{H'}\) the two requirements are the prime-to-\(\ell\)
index and the element \(g_0\) of
Theorem~\ref{thm:tametower571} or~\ref{thm:eventower572}.
In the stated subcase of case II prescribe one modulo the
odd part of \(N\) and one modulo \(2^{v_2(N)}\). This fixes
\(E_x\). At the dyadic modulus \(M_\alpha\), vary the latter
residue by \(1+2^{\max(2,v_2(N))}\ZZ_2\). Its \(\chi\)-image
has order \(M_\alpha/2^{\max(2,v_2(N))}>4\) by the
conductor bound and the cyclic calculation in
Theorem~\ref{thm:eventower572}. Choose \(\chi^2\ne1\).
Independently choose \(\bar\omega_\ell^r=-1\), possible
because \(\ell-1\mid 2r\) and \(\ell-1\nmid r\).
Require \(\phi=-1\) at \(p_3\) and one at \(p,p'\).
These prescriptions give \(\bar\alpha=1,\ \tau=-1=\phi\),
so the element fixes \(E_xH'\) and its four entries
\(-\chi,-1,\chi,1\) are distinct. The field \(E_xH'\)
has degree prime to \(\ell\). The same proof applies.
\end{proof}

The finite-module result uses all the primes split in one
explicit field on the interval. Estimating the resulting
restricted trace algebra is still an arithmetic requirement.


\subsection{An integral denominator for traces at common split primes}
\label{sec:denominator574}

Keep the single field \(\mathcal F\) of
Proposition~\ref{prop:commonaux573}. Every residual diagonal
character is trivial on its Galois group. We use the known
ordinary inertia characters in the integral lift to recover
the four diagonal projectors with a bounded denominator.
Write \(G\) for the compact image of \(G_\QQ\) in the
universal normalized representation over \(\mathcal R\), and
\(P\) for its kernel of reduction. Write \(O[G]\) for the
subalgebra of \(M_4(\mathcal R)\) algebraically generated by
the group matrices over \(O\), and similarly for subgroups.

\begin{lemma}[The diagonal finite subgroup in the ordinary frame]
\leavevmode\label{lem:ordcomplement574}
One has \(P=\rho(G_{\mathcal F})\). The group \(G\) contains
the finite diagonal lifts of its residual image in the
frame normalized at \(h\): for each residual matrix it
contains the diagonal matrix of the Teichm\"uller lifts
of its four entries on the diagonal.
\end{lemma}
\begin{proof}
Triviality of the residual diagonal gives
\(\rho(G_{\mathcal F})\subset P\). The quotient
\(G/\rho(G_{\mathcal F})\) has order dividing
\([\mathcal F:\QQ]\), which is prime to \(\ell\). The
kernel \(P\) is pro-\(\ell\), as in
Proposition~\ref{prop:separatedalgebra573}. Its image in
this quotient is therefore trivial, proving equality.

Let \(D_0\) be the residual image and \(C_0\subset D_0\)
the cyclic subgroup generated by \(\bar\rho(h)\). Its
lift by the diagonal Teichm\"uller matrix \(s\) is already
in \(G\). Indeed \(\rho(h)\) is diagonal and
\(\rho(h)^{\ell^{j!}}\) tends to \(s\): in each finite
quotient the principal-unit part is killed for large \(j\),
and \(\ell^{j!}\) is one modulo the prime-to-\(\ell\)
order for all sufficiently large \(j\).

We can choose a section \(D_0\to G\) agreeing with this
specified lift on \(C_0\). Use the successive finite
congruence quotients from Proposition~\ref{prop:separatedalgebra573}.
Suppose the agreement holds at one level. After lifting the
section to the next level using the vanishing of \(H^2\),
its restriction and the specified section of \(C_0\) differ
by a one-cocycle in the abelian congruence kernel. Averaging
over \(C_0\) makes that cocycle a coboundary. Conjugating the
new section by the corresponding kernel element restores
agreement and leaves its previous reduction fixed. Induction
and compactness give the section into \(G\).
Its image commutes with \(s\) since \(D_0\) is abelian.
The diagonal entries of \(s\) differ pairwise by units;
commuting with \(s\) thus forces all off-diagonal entries
to be zero. The diagonal entries have order prime to \(\ell\)
and the prescribed residues, so they are their unique
Teichm\"uller lifts. This proves the assertion in the
specified frame.
\end{proof}

\begin{theorem}[The common-split matrix denominator]
\leavevmode\label{thm:matrixden574}
For \(c=3+v_\ell(r(r-1)(r+1))\) one has
\[
        \ell^c O[G]\ \subseteq\
             O[\rho(G_{\mathcal F})]\ \subseteq\ O[G] .
\]
The same inclusions hold after every specialization of the
global representation.
\end{theorem}
\begin{proof}
Rational wild inertia at \(\ell\) acts trivially on
\(\mathcal F\), whose degree is prime to \(\ell\).
Its cyclotomic image is \(1+\ell\ZZ_\ell\): full inertia
surjects onto \(\ZZ_\ell^\times\) since the local cyclotomic
extensions are totally ramified, and the quotient by wild
inertia has order prime to \(\ell\) in every finite quotient.
Choose an element \(u\) of wild inertia with
\(\kappa_\ell(u)=1+\ell\). In the ordered ordinary flag
the diagonal of \(\rho(u)\) consists of
\[
     \lambda_i=(1+\ell)^{a_i},\qquad
                   (a_0,a_1,a_2,a_3)=(r+1,r,1,0).
\]
The finite inertia factors are trivial on this element.
If \(N\) is the order of the diagonal matrix \(s\) in
Lemma~\ref{lem:ordcomplement574}, then
\[
         Z=\frac1N\sum_{j=0}^{N-1}s^j\rho(u)s^{-j}
            =\operatorname{diag}(\lambda_0,\lambda_1,
                                             \lambda_2,\lambda_3)
              \quad\hbox{belongs to }O[\rho(G_{\mathcal F})].
\]
Indeed \(N\) is a unit and every conjugate belongs to the
normal subgroup \(P\). Each off-diagonal entry is multiplied
in the sum by the powers of a nontrivial \(N\)-th root of
unity whose difference from one is a unit. Their sum is zero.

Set \(d_i=\prod_{j\ne i}(\lambda_i-\lambda_j)\). For
distinct nonnegative integers \(a,b\),
\[
  v_\ell\big((1+\ell)^a-(1+\ell)^b\big)=1+v_\ell(a-b).
\]
In fact raising a unit \(1+z\), \(v_\ell(z)\ge1\), to
a power prime to \(\ell\) preserves the valuation of its
difference from one, by the leading term of the binomial
expansion. Raising it to \(\ell\) increases that valuation
by one, since the other terms have strictly larger
valuation for odd \(\ell\). Applying these two facts to
\(|a-b|\) proves the identity. The absolute differences
at each of our four exponents are either \(1,r,r+1\) or
\(1,r,r-1\). At most one of \(r-1,r,r+1\) is divisible
by \(\ell\ge5\). Thus \(\max_i v_\ell(d_i)=c\).
For the coordinate projector \(e_i\) this gives the identity
of integral matrices
\[
        \ell^c e_i=\frac{\ell^c}{d_i}
                            \prod_{j\ne i}(Z-\lambda_j).
\]
The scalar coefficient on the right lies in \(O\).
The identity is checked on the four diagonal entries,
so also holds in quotients with \(\ell\)-torsion.

Every \(g\in G\) has the form \(t p\), with \(p\in P\)
and \(t\) the finite diagonal lift of its reduction by
Lemma~\ref{lem:ordcomplement574}. The product \(\ell^c t\)
is an \(O\)-linear combination of the \(\ell^c e_i\).
It follows that \(\ell^c g\) belongs to \(O[P]\).
Taking \(O\)-linear combinations gives the first inclusion
since the algebra of a matrix group is its linear span.
All statements pass to images of the matrices under
coefficient maps.
\end{proof}

\begin{corollary}[Trace ideals and preservation of characteristic-zero multiplicities]
\leavevmode\label{cor:splitloss574}
Let \(\mathfrak a\) be the ideal in \(T\) generated by all
differences from the split trace and \(\mathfrak a_{\mathcal F}\)
the ideal in \(T\) using only good rational primes split in
\(\mathcal F\). Then
\[
     \ell^c\mathfrak a\subseteq\mathfrak a_{\mathcal F}
           \subseteq\mathfrak a,\qquad
     T_{\mathcal F}[1/\ell]=T[1/\ell].
\]
At any integral point whose least normalized valuation of a
good-prime trace difference is finite and equals \(v\), the
least valuation using these common split primes lies between
\(v\) and \(v+c\). Uniformly on the interval,
\[
          c\log\ell\ \le\
             3\log(2t_*+1)+3\log(t_*+1).
\]
The equality of generic algebras includes their nilpotents.
\end{corollary}
\begin{proof}
For each group element the matrix theorem writes
\(\ell^c\rho(g)\) as a finite \(O\)-linear combination of
matrices of elements of \(G_{\mathcal F}\). Take traces,
then apply the split augmentation to the same expression
and subtract the resulting equality. This puts the scaled
trace difference in the ideal generated by the restricted
differences. At each finite precision any restricted trace
difference is approximated by one at a good split rational
prime as in Theorem~\ref{thm:splitintegral573}, taking
both the universal and split matrices in the finite
quotient if necessary. Ideals in the complete Noetherian
local ring \(T\) are closed. Its intrinsic topology agrees
with that in the representation ring by compactness and
Hausdorffness. This gives the asserted inclusions.

For the equality after inversion, let \(b\) be the size of
the finite trace generating list in
Theorem~\ref{thm:splitintegral573}. Each generator multiplied
by \(\ell^c\) belongs to \(T_{\mathcal F}\) by the matrix
identity. That theorem realizes \(T\) as the ordinary
\(T_{\mathcal F}\)-module spanned by monomials in the
generators of exponents less than \(D^4\) in each variable.
In particular
\[
                \ell^{c b(D^4-1)}\,T\ \subseteq\ T_{\mathcal F},
\]
which proves equality after inversion. Thus all
characteristic-zero local algebras, including their
multiplicities when finite, coincide. Evaluating the two
ideal inclusions in a coefficient DVR gives both valuation
bounds with \(v_\ell(\ell)=1\). Finally
\(\ell\le2t_*+1\) and
\(v_\ell(r(r-1)(r+1))\log\ell\le3\log(t_*+1)\).
\end{proof}


\subsection{Nonsplit block witnesses on the auxiliary towers}
\label{sec:blockwitness575}

On a triple of Section~\ref{sec:sourcelevels563}, let \(H,H'\)
be its fields in Corollary~\ref{cor:bothaux572}. Use \(P\)
for the plane on coordinates \(13\), and \(P'\) for the plane
on coordinates \(02\) in a block realization, as in
Section~\ref{sec:ordinarycross568}. The following construction
gives the orientation of the nonsplit extension appearing in
\srcloc{Development/AutomorphicArithmetic/AutomorphicVarieties/WeakShimuraTheory/RedIdeal}{65}.

\begin{lemma}[Invariant subspaces on the auxiliary cyclotomic kernel]
\leavevmode\label{lem:cyclospaces575}
Let an integral realization over a complete DVR with finite
residue field reduce to the specified diagonal. The coefficient
field may have characteristic zero or \(\ell\). Every invariant
subspace over the algebraic closure of its fraction field for
\(G_{H'(\mu_{\ell^\infty})}\) is already invariant for \(G_\QQ\).
\end{lemma}
\begin{proof}
The elements \(g_0\) in Theorems~\ref{thm:tametower571}
and~\ref{thm:eventower572} were chosen using prescribed
cyclotomic residues, with the residue one at \(\ell\).
In their construction one can prescribe the entire
\(\ell\)-adic cyclotomic component to be one. Indeed all
other moduli used in specifying this element are prime
to \(\ell\), and
\(\Gal(\QQ(\mu_\infty)/\QQ)=\prod_p\ZZ_p^\times\).
Lifting this element to \(G_\QQ\) gives
\(g_0\in G_{H'(\mu_{\ell^\infty})}\) with the same four
distinct residual diagonal entries.

Hensel's lemma lifts these four simple roots for \(\rho(g_0)\)
to the coefficient DVR. Its four integral eigenlines give
a basis by their unit-separated spectral projectors, and
their reductions are the original coordinate lines. Every
subspace invariant under the indicated subgroup is a sum
of a subset of these eigenlines, since the projectors are
polynomials in \(\rho(g_0)\). The subgroup is normal in
\(G_\QQ\) because the underlying extension is abelian.
Hence any \(g\in G_\QQ\) carries such a sum to another sum
of the same eigenlines. The reductions of the two sums
are identical: the original residual matrix at \(g\) is
diagonal, and an invertible integral matrix transports
the saturated lattice of a subspace to that of its image.
Distinct subsets have distinct reductions. Thus \(g\)
preserves the original subspace. This argument also covers
all subspaces over the algebraic closure, since a sum of
the specified eigenlines is already defined over the
coefficient field.
\end{proof}

\begin{theorem}[An ordinary nonsplit extension with exactly the plane traces]
\leavevmode\label{thm:blockwitness575}
In either characteristic, every two-plane DVR realization
\(\rho_0=P'\oplus P\) of Theorem~\ref{thm:purityblocks549}
admits another integral realization \(\rho_1\) with the
same multiplier, residual representation, ordered local
conditions and all Frobenius polynomials. Over the
fraction field it is an extension
\[
                     0\longrightarrow P\longrightarrow
                      \rho_1\longrightarrow P'\longrightarrow0 .
\]
Its only nonzero proper invariant subspace over the algebraic
closure is \(P\). The same holds on \(G_{H'(\mu_{\ell^\infty})}\).
\end{theorem}
\begin{proof}
Write \(R\) for the DVR and \(\pi\) for a uniformizer.
At each \(R/\pi^j\), Theorem~\ref{thm:ordinarycross568}
gives a free rank-one space of ordinary normalized
cocycles in the direction \(P'\to P\). Reduction maps
these spaces onto one another by Lemma~\ref{lem:crosslocal568}.
Lift a nonzero residual generator compatibly through all
levels. Uniqueness of representatives normalized at \(h\)
gives an actual continuous inverse limit cocycle \(z\)
with values in \(M_-\). It has a unit entry at some group
element since its reduction is nonzero. No part of this
inverse-limit argument depends on the characteristic
of the DVR.

Put \(\rho_1(g)=(1+\pi z(g))\rho_0(g)\). Products of two
of these off-diagonal corrections vanish even in \(R\):
their images are in \(P\), on which they are zero. Thus
the cocycle law gives an actual representation. Symplecticity,
the ordered local conditions and the inertia kernels follow
from the same matrix verification as in
Theorem~\ref{thm:ordinarycross568}. The factor \(\pi\)
retains the original residual diagonal. The characteristic
polynomial of a block triangular matrix is the product of
those of its diagonal blocks, so all characteristic
polynomials agree with those of \(\rho_0\).

An invariant complement of \(P\) would be a graph from
\(P'\) to \(P\). Since \(z(h)=0\) and the two blocks at
\(h\) have disjoint spectra, a graph invariant at \(h\)
must have zero graph map. But \(P'\) itself is moved
by an element where \(z\ne0\). Thus there is no
complement, even after scalar extension. The two blocks
are absolutely irreducible by
Theorem~\ref{thm:purityblocks549}. The intersection of a
stable subspace with \(P\) is therefore zero or all of
\(P\), and its image in \(P'\) is zero or all of \(P'\).
The only additional nontrivial possibility would be the
excluded complement. Lemma~\ref{lem:cyclospaces575}
preserves this conclusion on the cyclotomic kernel.
\end{proof}

\begin{corollary}[Detecting failure at a common split prime]
\leavevmode\label{cor:commonwitness575}
If \(T\) is infinite over \(O\), there is an ordinary lift
over a complete finite-residue-field DVR of characteristic
\(\ell\) with a transcendental trace at a good rational
prime split completely in the single field \(\mathcal F\).
It is either absolutely irreducible or a nonsplit extension
with exactly the proper subspace in
Theorem~\ref{thm:blockwitness575}. In characteristic zero
every nonzero two-plane trace point has a witness of the
latter kind for which, at every unramified place of every
finite extension over a good rational prime \(q\), the
ratio of the largest to the smallest complex Frobenius
absolute value is at least \(q^{f_q}\), where \(f_q\) is
the residue degree.
\end{corollary}
\begin{proof}
In case of nonfiniteness take the one-dimensional
special-fibre quotient and its integral realization in
the proof of Corollary~\ref{cor:pointalternative550}.
At least one rational Frobenius trace is transcendental
over the finite residual coefficient field. If every
common split-prime trace were algebraic over that field,
each would lie in the constant field of the target DVR.
Indeed reduction is injective on its elements algebraic
over a finite field: the nonzero ones are roots of unity
of order prime to \(\ell\). They thus form a finite field
in the target, since its residue field is finite. It is
closed, so the whole image of \(T_{\mathcal F}\) would
lie in this finite field by Theorem~\ref{thm:splitintegral573}.
The monic equation in that theorem would make every trace
algebraic, a contradiction. Thus one can take a good
prime split in \(\mathcal F\).

Use the rank-four realization when it is irreducible,
and replace the block realization by the nonsplit one
of Theorem~\ref{thm:blockwitness575} otherwise. This
preserves the transcendental trace. In characteristic
zero the same replacement preserves all Frobenius
polynomials, so the determinant comparison on the two
planes in Theorem~\ref{thm:purityblocks549} gives the
asserted ratio for any coefficient embedding into
\(\CC\).
\end{proof}

We record precisely how the finite-weight condition on a
family would exclude the equal-characteristic witness.
Let \(Y\) be a local algebra over a coefficient DVR \(O_1\)
of characteristic zero and residue characteristic \(\ell\),
with finite residue field. Suppose a specialization \(y\)
to a field of characteristic \(\ell\) supplies the traces
of the witness on \(G_B\) for a finite extension \(B/\QQ\)
as elements of \(y(Y)\). Suppose further that an ideal
\(J\subseteq\ker y\) and a collection of weight elements
\(\theta_{v,i}(s)\in Y\) satisfy
\[
   \mathfrak m_Y^N\subseteq
        J+\mathfrak m_{O_1}Y+
                   \bigl(\theta_{v,i}(s)-1\bigr)_{v,i,s},
          \qquad y(\theta_{v,i}(s))=1
\]
for some \(N\ge1\), with \(s\) ranging over inertia above
\(\ell\). These are the parts of the finite-weight condition
in
\srcloc{Development/AutomorphicArithmetic/AutomorphicVarieties/ModularityArgument/ThresholdArguments/DenseThresholdCases/PotentialModularityTower}{38}
used here. When the weights specialize to the prescribed
diagonal inertia on \(H\), their reductions there are one
by Theorems~\ref{thm:tametower571}
and~\ref{thm:eventower572}.

Indeed \(y\) kills \(\mathfrak m_{O_1}\), since a power
of its uniformizer is a unit times \(\ell\), and a field
has zero nilradical. The inclusion gives \(y(\mathfrak m_Y)=0\)
and hence a finite image \(y(Y)\). For any \(g\in G_\QQ\)
some positive power \(g^j\) lies in \(G_B\), by its action
on the finite coset space. The traces on \(g^{ij}\) for
\(1\le i\le4\) are now algebraic over the finite field.
Newton's identities and \(\ell\ge5\) show that all
eigenvalues of \(\rho(g^j)\) are algebraic. Those of
\(\rho(g)\) are \(j\)-th roots of these eigenvalues and
are therefore algebraic as well. This contradicts the
common split-prime witness. Thus such family realizations
for all the indicated equal-characteristic witnesses would
prove finiteness of \(T\). In characteristic zero the
corollary gives the separate contradiction for a family
giving Frobenius purity at even one good place of a
nonsplit block witness. Establishing the required family
realizations and the rank-four cuspidal matches remains
necessary for this argument.


\subsection{A cuspidal symmetric cube on the odd-order towers}
\label{sec:cubicseed577}

Use the odd-order case of Section~\ref{sec:tametower571}, the
real field \(H\), its CM extension \(HK\), and the element
\(g_0\) of Lemma~\ref{lem:cyclospaces575}. On \(G_H\) the
ordered residual characters are \((\chi^3,\chi^2,\chi,1)\),
where \(\chi=\bar\xi\), and \(\chi(g_0)\) has order \(2m\ge6\).
Use the cyclotomic character \(\kappa_\ell\) with arithmetic
Frobenius convention. The required multiplier is
\(\mu=A\Psi=\kappa_\ell^{r+1}\omega_\ell^{-1}\xi\alpha\).
We construct an initial automorphic point for these towers
with its integral representation and ordered flags. Its weight
is allowed to vary in the definition of the initial point.

\begin{lemma}[One level prime and an ordinary modular plane]
\leavevmode\label{lem:cuspidalplane577}
There is one rational prime \(w\) split completely in
\(\mathcal F\) of Proposition~\ref{prop:commonaux573}, common
to the triples on the interval, with \(\log w\) polynomial in
\(V\). For each odd-order triple and each positive multiple
\(j\) of \(\ell-1\), choose a character \(\theta\) of order
\(\ell\) and conductor \(w\). There is an ordinary cuspidal
newform of weight \(j+1\), exact level \(\operatorname{cond}(\xi)w\)
and character \(\xi\theta\). Its Galois plane \(W\) has
determinant \(\kappa_\ell^j\xi\theta\), residual
semisimplification \(\chi\oplus1\), and Zariski-dense
projective image in \(\mathrm{PGL}_2\).
\end{lemma}
\begin{proof}
Choose \(w\equiv1\bmod B\), where \(B\) is a common multiple of
the conductor of \(\mathcal F\), \(2N^\natural M_*p_1\), and
all the primes up to \(2t_*+1\); the least such multiple will
do. Proposition~\ref{prop:commonaux573} and
Lemma~\ref{lem:twoprimes571} bound \(\log B\) by a polynomial
in \(V\). Linnik's theorem supplies \(w\ll B^{C}\) with an
absolute \(C\). In particular \(w\equiv1\bmod D\ell\) for
\(D=\operatorname{cond}(\xi)\), so \((\ZZ/w)^\times\) admits
the required character. Construct forms for any fixed
identification of the algebraic closures used for their
complex and \(\ell\)-adic coefficients.

Here \(j\) is even and \(\xi\) is primitive and odd.
At weight \(k=j+1>2\) the normalized \(E_k(1,\xi)\)
has integral constant term at \(\ell\) by
Lemma~\ref{lem:bernoulliall436} and integral nonconstant
coefficients. On the prime-to-\(\ell\) full cover the
expansion principle gives integrality at the cusps, as in
Proposition~\ref{prop:auxiliarycusp436}. At level \(Dw\) reduce
\(E(z)-E(wz)\). For a cusp represented by a primitive column
\((b,c)^{\mathsf T}\), set \(d=\gcd(w,c)\). The second constant
term is the first-level constant at
\((wb/d,c/d)^{\mathsf T}\) times \((d/w)^k\). The two cusp
columns agree modulo \(D\) since \(w\equiv1\bmod D\), and the
factor reduces to one since \(w\equiv1\bmod\ell\). The
difference is consequently cuspidal in the residue field.
Its first coefficient is one. It is an eigenvector for the
Hecke operators at primes \(p\nmid Dw\) with eigenvalues the
reductions of \(1+\xi(p)p^j\), including the unit eigenvalue
at \(\ell\).

The character lifting argument of
Proposition~\ref{prop:auxiliarycusp436} applies with the single
prime \(w\). More explicitly, the \(\ell\)-Sylow diamond group
acts freely on \(X_1(Dw)\), since smooth-point stabilizers
have order divisible only by two and three and cusp stabilizer
orders divide \(2Dw\). On the quotient, descend the weight-\(k\)
cusp line with the prescribed \(\ell\)-power character. Its
residual twist is trivial; the other diamond characters are
separated by integral projectors. Serre duality on the special
fibre gives vanishing of the first cohomology of the descended
cusp line, since its dual sections pull back to weight \(2-k<0\).
Thus the residual eigensystem occurs in the reduction of the
torsion-free cusp module with character \(\xi\theta\).
A minimal prime of its finite Hecke algebra below this
eigencharacter has characteristic zero: a uniformizer is a
nonzero divisor on that algebra by its faithful action on the
torsion-free module. This yields a characteristic-zero
cuspidal eigensystem. The character is primitive of conductor
\(Dw\), so its newform has this exact level. Deligne's
representation and the ordinary local description give the
determinant, the stated residual semisimplification by
Chebotarev, and an unramified local quotient with unit
Frobenius value reducing to one.~\cite{XII-IX-V-deligne,XII-IX-ordinary361}

The plane of the cuspidal newform is absolutely irreducible.
The two Hodge--Tate weights \(j,0\) exclude finite projective
image. If the connected projective Zariski closure were a
proper positive-dimensional subgroup, it would be solvable.
A nontrivial unipotent radical has a unique fixed line
preserved by the whole group, contradicting irreducibility.
Otherwise the connected group is a torus, whose normalizer
permutes two lines. A transitive permutation supplies a
quadratic self-twist by conjugating with the diagonal matrix
with entries \(1,-1\) on these lines. Reduction would give
a nontrivial quadratic character permuting \(\{1,\chi\}\)
by multiplication. This forces \(\chi^2=1\), contrary to
its value at \(g_0\). This proves the density.
\end{proof}

\begin{proposition}[The nonsplit residual point and its ordinary flags]
\leavevmode\label{prop:residualcube577}
There is a lattice in \(W\) in which \(g_0\) is diagonal and
the residual plane is a nonsplit extension
\[
             0\longrightarrow {\bf1}\longrightarrow\bar W
                           \longrightarrow\chi\longrightarrow0 .
\]
On \(G_H\) its symmetric cube is lower triangular with
diagonal \((\chi^3,\chi^2,\chi,1)\) and has both simple
negative root entries nonzero as functions on the group.
At every place above \(\ell\) its ordinary flag is represented
by a lower unitriangular symplectic matrix in this frame.
These assertions persist on every finite totally real abelian
extension \(H_1/\QQ\) containing \(H\) with
\([H_1:H]\) an \(\ell\)-power, using a power
\(g_0^{\ell^a}\in G_{H_1}\) for the normalization.
\end{proposition}
\begin{proof}
Start with a stable lattice over a coefficient DVR \(O\)
large enough for \(W\) and the characters, with uniformizer
\(\pi\). Since \(g_0\) has distinct residual eigenvalues,
Hensel's lemma and its spectral projectors give an integral
eigenbasis ordered by \(\chi(g_0),1\). The traces at \(g\)
and \(g_0g\) determine the two diagonal entries modulo \(\pi\)
as \(\chi(g),1\) by a two-by-two system of unit determinant.
Write \(b_g,c_g\) for the upper and lower off-diagonal entries.
The first diagonal entry at \(gg'\) gives
\(b_gc_{g'}\equiv0\bmod\pi\) for all \(g,g'\).
By irreducibility both sets of entries have finite minimum
valuation. If \(c=\min_g v_\pi(c_g)\), replace the second
basis vector by \(\pi^c\) times itself. Now all the upper
entries vanish on reduction, some lower entry is a unit,
and \(g_0\) remains diagonal. A stable complement of the
second residual line would be the first eigenline of \(g_0\);
the nonzero lower entry moves it.

The extension class belongs to \(H^1(G_\QQ,k_O(\chi^{-1}))\),
where \(k_O\) is the residue field. Restriction to \(G_{H_1}\)
is injective on this group by inflation--restriction, since
\(\chi(g_0^{\ell^a})\ne1\) for every \(a\), and some such
power belongs to \(G_{H_1}\). This applies also to \(H\).
Thus some lower entry on the restricted group remains
nonzero. In the monomial basis of the symmetric cube the
first two lower subdiagonal entries of a residual matrix
with plane first column \((\chi(g),c_g)^{\mathsf T}\) are
\(3\chi(g)^2c_g\) and \(2\chi(g)c_g\). Both are nonzero
when \(c_g\ne0\), since \(\ell\ge5\).

For the symplectic frame use the monomials
\(X^3,X^2Y,-3XY^2,Y^3\). Indeed the alternating form induced
on the symmetric tensors by the cube of the determinant
pairing has values \(1,-1/3\) on the two ordered complementary
pairs in the unscaled monomial basis. This follows by
averaging the three tensor factors: among the pairings of
\(X^2Y\) with \(XY^2\) one third are nonzero, with negative
sign. The rescaling gives the prescribed antidiagonal form;
it changes the two subdiagonal entries only by units. The
multiplier of the symmetric cube is the cube of the determinant.

At any rational place above \(\ell\) the ordinary subline
of \(W\) has residual character \(\chi\), since the unramified
quotient reduces to one. This character differs from one
on the rational decomposition group by
Lemma~\ref{lem:prescribedaux563}, which gives
\(\chi(\ell)\ne1\). The reduction of the saturated subline
therefore differs from the global invariant residual line.
It has an integral generator \(X+a_vY\) in the ordered plane
basis. The lower unitriangular matrix sending \(X\) to this
vector and fixing \(Y\) gives the ordinary flag, and its
symmetric cube gives the required symplectic flag. At a
place upstairs use the restriction of this same rational
ordinary line. Conjugating the place transports the line
because the positive-weight ordinary subcharacter is
distinguished from its weight-zero quotient. The flags
therefore restrict compatibly through all the indicated
towers, including at places where the two local residual
characters become equal.
\end{proof}

\begin{theorem}[A cuspidal initial point for every odd-order tower]
\leavevmode\label{thm:cuspidalinitial577}
Let \(C_{w,\ell}\) be the degree-\(\ell\) subfield of
\(\QQ(\mu_w)\). For every \(H_1\) of
Proposition~\ref{prop:residualcube577} containing \(C_{w,\ell}\)
there is an ordinary symplectic lift of that same residual
point with multiplier \(\mu\). It is absolutely irreducible on
\(G_{H_1K(\mu_{\ell^\infty})}\), and over \(H_1K\) an explicit
scalar twist of it is attached to a regular algebraic cuspidal
representation of \(\mathrm{GL}_4\). Its ordered local characters,
weights, and ramification satisfy the initial-point conditions
of \srcloc{Development/AutomorphicArithmetic/AutomorphicVarieties/WeakShimuraTheory/NearSplit}{312}
on the towers of Section~\ref{sec:tametower571}. One can take
\(H_1=HC_{w,\ell}\).
\end{theorem}
\begin{proof}
The character \(\mu/(\det W)^3\) on \(G_H\) is one on
reduction because \(\tau=\chi^2\) there. Let \(b\) be its
unique square root with values in \(1+\mathfrak m_O\).
Existence and uniqueness follow from Hensel's lemma since
\(\ell\) is odd. Uniqueness proves multiplicativity, and
applying the same lifting modulo all powers proves continuity.
Use the integral representation
\[
               R=b\,\operatorname{Sym}^3 W
\]
on \(G_H\) and its restrictions, in the basis of the
proposition. Its multiplier is \(\mu\) and its residual
point is exactly the constructed one. Every prescribed
normalizing power of \(g_0\) is diagonal. Before adjoining
\(C_{w,\ell}\) allow the single extra ramification prime \(w\).

At every prime \(p\mid Dw\) the conductor exponent of the
plane equals that of its determinant by primitivity of the
character and local-global compatibility for newforms.
For a ramification group the codimension of fixed vectors
is at least that for the determinant. Equality of the Artin
conductors forces equality for inertia, including the
nonnegative monodromy contribution. More explicitly the
inertial fixed space has codimension at least one since the
determinant is ramified, hence exactly one. Finite inertia
then has characters \(1,\xi\theta\), distinct on inertia,
and a commuting nilpotent monodromy must be zero. At the
other primes away from \(\ell\) the plane is unramified.
These are the conductor and inertia arguments of
Theorem~\ref{thm:twoprimeadjoints436}.~\cite{XII-IX-VI-saito40}
On \(G_{H_1}\) the character \(\theta\) is trivial. At an
inertia element away from \(\ell\) in the kernel of the
residual four characters, \(\xi\) and \(\alpha\) are both
one: their reductions are one and their orders are prime
to \(\ell\). The plane, the multiplier and its scalar
correction are therefore trivial at this element. On \(H\)
the same calculation works outside \(w,\ell\).

Above \(\ell\) use the integral flags of the proposition.
The four residual local characters in their given order are
\(\chi^3,\chi^2,\chi,1\). The two ordinary characters of
the plane are an unramified character and
\((\det W)\) times its inverse. On elements whose image in
the local abelianization has \(\ell\)-power order the
unramified characters are one, since the unramified quotient
is torsion-free. The cyclotomic character also takes value
one there as \(\ZZ_\ell^\times\) has no \(\ell\)-power torsion,
and the finite factors in \(\mu,\det W\) are unramified
there or of order prime to \(\ell\). Thus all four ordinary
characters and \(b\) take value one on these elements. On
an open subgroup of inertia the squared ordered characters
of \(R\) are the cyclotomic powers with exponents
\[
      h_i=2(3-i)j+r+1-3j,\quad 0\le i\le3 .
\]
Their consecutive differences are \(2j\ge4\). The inverse
twist \(b^{-1}\) has square on the same subgroup equal to
\(\kappa_\ell^{a}\), \(a=3j-r-1\), and \(a+h_0=6j\) is even.

The projective image of \(W\) on the absolute commutator
subgroup is Zariski dense in \(\mathrm{PGL}_2\). Indeed the
closure contains all commutators of the original Zariski
closure by density and the commutator morphism, and
\(\mathrm{PGL}_2\) is its own derived group. Since
\(H_1K(\mu_{\ell^\infty})/\QQ\) is abelian, its absolute
Galois group contains these commutators. In characteristic
zero the third symmetric power of the standard
\(\mathrm{SL}_2\)-module is irreducible, so no proper
subspace of \(R\) can be stable on this group.

The plane comes from a non-CM form by the density in
Lemma~\ref{lem:cuspidalplane577}. Symmetric power
functoriality in degree three gives a regular algebraic
cuspidal transfer over \(\QQ\). In the algebraic Frobenius
normalization its four parameters at a good prime are the
degree-three monomials in the two parameters of \(W\).
This is Theorem A of Newton and Thorne,
\href{https://pmihes.centre-mersenne.org/articles/10.1007/s10240-021-00126-4/}{\emph{Symmetric
power functoriality for holomorphic modular forms, II}}.
Use cyclic base change successively through the finite abelian
extension \(H_1K/\QQ\). Cuspidality could be lost at a cyclic
step of prime degree only by a nontrivial self-twist of the
representation before that step. The equality of good
polynomials and Chebotarev would then give a self-twist of
\(\operatorname{Sym}^3 W\) on that field. An intertwiner
would commute with its restriction to the kernel of the
twisting character. This kernel has finite index and the
restriction is absolutely irreducible by the same density
argument (or by connected projective monodromy). The
intertwiner would be scalar, which forces the twisting
character to be trivial. Thus the base changes stay
cuspidal. Base change preserves the distinct algebraic
weights and raises the good Frobenius eigenvalues to the
residue degrees. The resulting representation over \(H_1K\)
is attached to \(b^{-1}R\) there.~\cite{XII-IX-I-ac123}

Finally, the characters of \(G_\QQ\) of order dividing
\(\ell\) unramified away from \(w\) all factor through
\(\Gal(C_{w,\ell}/\QQ)\). By rational class field theory
their conductors divide \(w\), since the wild units at \(w\)
are pro-\(w\), \(\ell\ne w\), and \(\ell\) is odd. The tame
quotient \((\ZZ/w)^\times\) is cyclic. Hence killing these
characters is exactly the condition \(C_{w,\ell}\subset H_1\).
This subfield is totally real and \(HC_{w,\ell}/H\) is cyclic
of degree \(\ell\), since \([H:\QQ]\) is prime to \(\ell\).
Together with the restriction of the flags, the two nonzero
negative root entries and the inertia calculation above,
this proves the claimed initial-point conditions.
\end{proof}

The prime \(w\) is auxiliary to this construction; the
cohomology and signed periods of Section~\ref{sec:missing}
are still taken over \(E_x\). The theorem supplies one
cuspidal initial point for each odd-order triple, together
with its restrictions to the required towers. Proving that
the components through the other ordinary deformations
occur in an automorphic family remains a separate step.


\subsection{Opposite ordinary quotients over the real quadratic field}
\label{sec:oppositequotients578}

Consider an even-order triple of Section~\ref{sec:sourcelevels563}.
Write \(D_0=K_0\) for the real quadratic field of
Lemma~\ref{lem:realquadratic572}, \(\phi\) for its quadratic
character and \(HK\) for the auxiliary CM field. The rational
prime \(\ell\) splits in \(D_0\) as \(v_0v_1\), as does two.
Let \(\lambda\) be the finite Hecke character of \(D_0\)
obtained by restricting the Galois character \(\xi\), or
equivalently by composing its idele character with the norm.
It is totally odd, of conductor \(\mathfrak n\) supported
at the two places over two, and unramified at \(\ell\).
Its values at \(v_0,v_1\) are both \(\xi(\ell)\). Their common
reduction differs from one by Lemma~\ref{lem:prescribedaux563}.
Use the arithmetic convention for the eigenvalues of Frobenius.

\begin{lemma}[Raising the weight with two different unit roots]
\leavevmode\label{lem:weightpair578}
For arbitrarily large positive \(j\) divisible by
\(2(\ell-1)\) there is a parallel-weight \(j+1\) Hilbert
Hecke eigensystem over \(D_0\) of character \(\lambda\),
tame level \(\mathfrak n\), and Iwahori level at both places
over \(\ell\). Its residual \(T_\mathfrak q\)-eigenvalues
away from \(\mathfrak n\ell\) are \(1+\bar\lambda(\mathfrak q)\).
Its two \(U\)-eigenvalues reduce to
\(\bar\lambda(v_0),1\) at \(v_0,v_1\), respectively.
\end{lemma}
\begin{proof}
Start with the parallel weight-one Eisenstein series
\(E_1(1,\lambda)\), with nonconstant ideal coefficients
\(\sum_{\mathfrak d\mid\mathfrak m}\lambda(\mathfrak d)\);
extend \(\lambda\) by zero at ideals meeting its conductor.
It exists for this totally odd primitive character of nontrivial
finite conductor. At each \(v\mid\ell\) its Hecke roots are
\(1,\lambda(v)\). The oldform \(V_v E\) has coefficient at
\(\mathfrak m\) equal to that of \(E\) at \(\mathfrak m/v\),
interpreted as zero when the latter ideal is nonintegral.
On the span of \(E,V_v E\), the formula for \(U_v\) gives
\(U_vV_v E=E\) and
\(U_v E=(1+\lambda(v))E-\lambda(v)V_v E\).
Thus apply \(1-V_{v_0}\) and \(1-\lambda(v_1)V_{v_1}\).
The resulting form \(E^*\) has the two required roots,
integral nonconstant coefficients and first coefficient one.
The operators at the two distinct places commute. Its
constants at the standard cusps are algebraic. We use the
ideal-coefficient conventions and Eisenstein expansions in
Dasgupta, Darmon and Pollack,
\href{https://services.math.duke.edu/~dasgupta/papers/Dasgupta-Darmon-Pollack.pdf}
{\emph{Hilbert modular forms and the Gross--Stark conjecture}, Sections 2.1--2.2}.

Take \(j=2(\ell-1)\ell^a\) for \(a\) tending to infinity.
The level-one parallel form \(E_j(1,1)\), for \(j\ge4\),
has normalized constant term \(2^{-2}\zeta_{D_0}(1-j)\)
on each standard component. Its remaining coefficients
are the integral divisor sums with exponent \(j-1\).
The valuation of this constant tends to minus infinity.
Indeed \(\zeta_{D_0}(1-j)=\zeta(1-j)L(1-j,\phi)\).
Von Staudt--Clausen gives
\(v_\ell(\zeta(1-j))=-1-v_\ell(j)\).
On this residue class of weights the Kubota--Leopoldt
interpolation formula identifies
\((1-\phi(\ell)\ell^{j-1})L(1-j,\phi)\) with
\(L_\ell(1-j,\phi)\). The latter tends to
\(L_\ell(1,\phi)\ne0\) by the \(\ell\)-adic class number
formula. For this real quadratic field the regulator is
nonzero: the fundamental unit in either completion
\(\QQ_\ell\) has nonzero logarithm, since a unit of
\(\QQ_\ell\) with zero logarithm is a root of unity. The
Euler factor tends to one. This proves the valuation claim.
The interpolation and the abelian \(\ell\)-adic class number
formula are the classical formulas of Chapter 5 of
\href{https://link.springer.com/book/10.1007/978-1-4612-1934-7}
{Washington, \emph{Introduction to Cyclotomic Fields}}.

Normalize \(E_j(1,1)\) by its constant term. Its expansion on
every standard component is congruent to one to arbitrarily
large \(\ell\)-adic precision, uniformly in the Fourier index.
To see this also in the unnormalized power series, take the
finitely many narrow class representatives prime to \(\ell\).
The changes between ideal coefficients and these power series
multiply by powers of their rational ideal norms, which are
\(\ell\)-adic units. Changing the weight by \(j\) changes
these factors by powers \(N(\mathfrak t)^{\pm j/2}\). These
are one on reduction by the choice of \(j\).
Multiply \(E^*\) by the normalized form. In the product all
nonconstant coefficients reduce to those of \(E^*\): the
precision can exceed the denominators of the finitely many
constants of \(E^*\), and the other coefficients are integral.
The coefficient convolutions at fixed totally positive index
are finite. This gives a genuine weight-\(j+1\) form with
the indicated residual coefficients.

At a good prime the second term in the Hecke coefficient
formula now acquires the factor \(N(\mathfrak q)^j\),
one on reduction. At \(v\mid\ell\) the \(U_v\) coefficient
is the coefficient at \(v\mathfrak m\) in either weight.
Consequently the product has the required eigencharacter on
reduction. For clarity, work in its Hecke span over a finite
coefficient field completed at \(\ell\), and let \(\mathcal L\)
be the lattice of forms in that span with integral nonconstant
ideal coefficients. Finitely many coefficient functionals
embed the span in a finite-dimensional coordinate space:
vanishing of all such coefficients would give a constant
function on each classical component, impossible in positive
weight unless it is zero by the modular transformation law.
Thus \(\mathcal L\) is a finite module over the coefficient
DVR. It spans the Hecke span since the coefficient formulas
preserve integrality and the span is generated by translates
of the product. Divisibility of every nonconstant coefficient
by a uniformizer is exactly divisibility in \(\mathcal L\).
The unit first coefficient therefore gives a nonzero
eigenvector on reduction. The finite Hecke algebra acts
faithfully on \(\mathcal L\), so a uniformizer is a nonzero
divisor on it. A minimal prime below the resulting maximal
ideal lies in characteristic zero. In its fraction field the
corresponding system occurs in the Hecke span by the faithful
action and simultaneous triangularization. This proves the
lemma at arbitrarily large weights.
\end{proof}

\begin{theorem}[A cuspidal plane with opposite ordinary quotients]
\leavevmode\label{thm:mixedplane578}
Every eigensystem just obtained is cuspidal. Its plane \(W\)
has determinant \(\kappa_\ell^j\lambda\), semisimple
reduction \(1\oplus\bar\lambda\), projectively dense image,
and an unramified ordinary quotient at each of \(v_0,v_1\).
These quotients reduce to \(\bar\lambda,1\), respectively.
Its exact finite conductor away from \(\ell\) is
\(\mathfrak n\); the underlying automorphic newform is
unramified above \(\ell\).
\end{theorem}
\begin{proof}
In parallel weight \(j+1>2\) the good eigenvalues of an
Eisenstein constituent are
\(\eta(\mathfrak q)+\nu(\mathfrak q)N(\mathfrak q)^j\)
for finite-order characters \(\eta,\nu\), with
\(\eta\nu=\lambda\). This is the classical Eisenstein
decomposition for holomorphic Hilbert forms of parallel
weight. The unweighted character is one global character:
the absolute-value exponents of a Hecke character at the
real places must be parallel by the unit theorem. Oldforms
retain the same pair. At a place over \(\ell\) the Iwahori
invariants require the two characters to be unramified.
Indeed ramification of either forces ramification of both
by the unramified central character, giving principal-series
conductor exponent at least two. Its two oldform roots at
that place are therefore \(\eta(v),\nu(v)N(v)^j\).
Exactly the first is a unit.

By Chebotarev the good residual traces and determinants
identify \(\{\bar\eta,\bar\nu\}=\{1,\bar\lambda\}\) as
characters, since \(\bar\kappa_\ell^j=1\). For example one
may compare the characteristic polynomials in a common finite
quotient and apply the semisimple character criterion.
It follows that the reduction of the unit root at both
places comes from \(1\) at both places or from
\(\bar\lambda\) at both places. This contradicts
Lemma~\ref{lem:weightpair578} since
\(\bar\lambda(v_0)=\bar\lambda(v_1)\ne1\).
An eigencharacter in the quotient by cusp forms would be
Eisenstein, so the eigensystem is cuspidal.

At a tame prime the conductor is bounded by the given level
and bounded below by the conductor of the central character.
These exponents agree. Above \(\ell\) the only ramified
Iwahori-spherical cuspidal local possibility is an unramified
twist of Steinberg. In the arithmetic normalization its
\(U_v\)-eigenvalue has square
\(\lambda(v)N(v)^{j-1}\), of positive valuation. Both
prescribed roots are units, so the local newforms are
unramified. The ordinary Galois representation theorem for
Hilbert newforms gives the determinant and the unramified
quotient with Frobenius eigenvalue the chosen unit root.
The good polynomials identify the residual semisimplification
by Chebotarev. These are the ordinary local descriptions
of Wiles and the local-global compatibility theorem for
Hilbert newforms.~\cite{XII-IX-VI-saito40}

The plane is absolutely irreducible by the Galois
representation theorem for regular cuspidal Hilbert newforms.
Its distinct weights \(j,0\) exclude finite projective image.
A nontrivial quadratic self-twist would permute
\(\{1,\bar\lambda\}\) on reduction, forcing
\(\bar\lambda^2=1\). At either place over two its inertia
has the full order of the original primitive dyadic
character on inertia since two splits in \(D_0\); that
order exceeds four by Section~\ref{sec:sourcelevels563}.
The algebraic subgroup argument in
Lemma~\ref{lem:cuspidalplane577} now gives projective density.
\end{proof}

\begin{corollary}[Cuspidal quadratic induction at parallel weight]
\leavevmode\label{cor:parallelind578}
Put \(U=\operatorname{Ind}_{G_{D_0}}^{G_\QQ}W\).
Its restriction to \(G_{HK(\mu_{\ell^\infty})}\) is absolutely
irreducible. Over \(HK\) it has the good Frobenius polynomials
of a cuspidal quadratic automorphic induction in the arithmetic
normalization. On \(G_H\) an integral scalar twist of \(U\)
has multiplier \(\mu\) and residual semisimplification
\((\chi\tau,\tau,\chi,1)\).
\end{corollary}
\begin{proof}
For the nontrivial conjugation \(\sigma\) of \(D_0/\QQ\),
\(W^\sigma\) is inequivalent to \(W\). Otherwise the
unramified ordinary quotients at the two places over
\(\ell\) would coincide, which is excluded on reduction
by Theorem~\ref{thm:mixedplane578}. The quotient at each
place is intrinsic, since the other character of the
ordinary flag has strictly positive weight.

Both planes are absolutely irreducible on \(G_{\QQ^{\rm ab}}\).
In fact this group contains the commutators of \(G_{D_0}\),
whose projective images are Zariski dense by the proof
of Theorem~\ref{thm:cuspidalinitial577}. If they became
isomorphic there, the line of intertwiners would be stable
under \(G_{D_0}\), since this subgroup is normal. Thus
\(W^\sigma\simeq W\otimes\delta\) for a character \(\delta\)
on \(G_{D_0}\). Their determinants agree because \(\lambda\)
is the restriction of \(\xi\). Hence \(\delta^2=1\); any
nontrivial choice again contradicts the residual pair
\(\{1,\bar\lambda\}\). The trivial choice contradicts
the preceding paragraph.

Let \(G'=G_{HK(\mu_{\ell^\infty})}\). The group
\(G'\cap G_{D_0}\) contains \(G_{\QQ^{\rm ab}}\), so the
restriction of \(U\) there splits into the inequivalent
irreducibles \(W,W^\sigma\). The element \(g_0\) of
Lemma~\ref{lem:cyclospaces575} in the even-order construction
belongs to \(G'\) and has \(\tau(g_0)=\phi(g_0)=-1\).
It interchanges these summands. Hence no nonzero proper
subspace is \(G'\)-invariant. For the automorphic assertion,
base change the Hilbert cusp form to \(HKD_0\). All the
cyclic steps remain cuspidal, since a self-twist on an
intermediate field would give a scalar intertwiner on the
kernel of its character by projective density on open
subgroups. The cuspidal base change is inequivalent to its
conjugate over \(HK\), as is seen from good polynomials
and the proved inequivalence on \(G_{\QQ^{\rm ab}}\).
Cyclic automorphic induction along \(HKD_0/HK\) is therefore
cuspidal. At split primes its parameters are the union of
the two pairs, and at inert primes its Frobenius polynomial
is obtained from the plane polynomial at the square by
substituting \(X^2\). These are the polynomials of \(U\),
in the normalization in which the plane has weight \(j\).
~\cite{XII-IX-I-ac123}

Write \(d=\kappa_\ell^j\xi\) on \(G_\QQ\) and take any
\(t\notin G_{D_0}\). In the two plane summands one can
write \(t(x,y)=(W(t^2)y,x)\). Give the two planes the
alternating determinant forms with coefficients \(1,-d(t)\),
respectively. On \(G_{D_0}\) both determinants are \(d\).
At \(t\), evaluation of the two pairings after applying
this matrix gives multiplier \(-d(t)\), since
\(\det W(t^2)=d(t)^2\). Thus the multiplier of \(U\) is
\(d\phi\). Its residual semisimplification is
\(1\oplus\phi\oplus\chi\oplus\chi\phi\) by induction of
the two residual characters. On \(G_H\) one has
\(\tau=\phi\), so \(d\phi\) and \(\mu\) have the same
reduction. The unique square root of \(\mu/(d\phi)\)
congruent to one on this group supplies the integral twist.
\end{proof}

The two labelled weights of the plane at each real embedding
are still \(j,0\). Thus the four weights of the induction
are \(j,j,0,0\), and a common scalar twist preserves both
repetitions. The regular initial point in the even-order
case requires different weight gaps at the two embeddings.
The result above supplies the cuspidal system with the
opposite ordinary quotients on the particular real quadratic
field where that change of weights is needed.


\subsection{A regular cuspidal point on the even-order tower}
\label{sec:regularind579}

Retain \(D_0,v_0,v_1,\lambda\) and one parallel exponent \(j\)
from Theorem~\ref{thm:mixedplane578}. Choose a positive multiple
\(t\) of \(\ell-1\). We raise the weight gap at \(v_1\).
At \(v_0\) the old ordinary quotient reduces to
\(\bar\lambda\); at \(v_1\) it reduces to one. Fix the two
embeddings of \(D_0\) indexed by \(v_0,v_1\) under the chosen
identification of coefficient fields.

\begin{proposition}[A change of one weight at the exact tame level]
\leavevmode\label{prop:weightchange579}
There is a cuspidal Hilbert plane \(W\) with determinant
\(d=\kappa_\ell^j\xi|_{G_{D_0}}\), residual semisimplification
\(1\oplus\bar\lambda\), and tame conductor exactly
\(\mathfrak n\). At the two places it is crystalline and
ordinary. Its ordered weights, with the subcharacter first,
are \((j,0)\) at \(v_0\) and \((j+t,-t)\) at \(v_1\).
The two ordered residual pairs there are \((1,\bar\lambda)\)
and \((\bar\lambda,1)\), respectively.
\end{proposition}
\begin{proof}
Use the quaternion algebra over \(D_0\) ramified at its two
real places and split at every finite place. The
Jacquet--Langlands correspondence realizes the cuspidal
system of Theorem~\ref{thm:mixedplane578} on this algebra,
with the same finite Hecke action and both chosen Iwahori
refinements. We use the algebraic modular forms realization
of this correspondence as in Hida, \emph{On \(p\)-adic Hecke
algebras for \(\mathrm{GL}(2)\) over totally real fields},
\href{https://web.iwr.uni-heidelberg.de/fileadmin/groups/arithgeo/templates/data/Hauptseminare/Literature-WS13/hida-quat_AF_and_JL.pdf}
{Theorem 2.1}. We give the finite-module argument for changing
the coefficient at \(v_1\).

Over a sufficiently large coefficient DVR, the original
algebraic coefficient at each split place is
\(\operatorname{Sym}^{j-1}\). At \(v_1\) replace it by
\(\operatorname{Sym}^{j+2t-1}\otimes\det^{-t}\), leaving
the \(v_0\)-factor as it is. For each factor
\(\operatorname{Sym}^{n}\otimes\det^m\), the holomorphic
weight is \(n+2\), the ordered arithmetic weights are
\((m+n+1,m)\), and the algebraic central exponent is
\(n+2m\). Thus both coefficient systems have central
exponent \(j-1\) at both places. At finite places fix the
same level, finite central character and tame nebentypus
as for the initial form.

Here are details of the integral comparison. In the
quaternionic model integral forms are functions of the
finite adelic argument satisfying left rational invariance,
right level equivariance through the inverse coefficient
action, and the prescribed central action. These conventions
result by sending a form equivariant on the rational left
through the algebraic coefficient \(\rho\) to its product
by \(\rho(g_\ell)^{-1}\) at the argument \(g\). The central
action in this integral model is the same for the two
coefficients, since their central exponents coincide and
their finite characters are fixed. It has unit values:
the valuation at \(\ell\) of the algebraic central factor
cancels the corresponding valuation of the finite adelic
norm to exponent \(j-1\). Away from \(\ell\) impose the
original character on level units.

Values at a finite set of double-coset representatives
identify this module with a direct sum of invariant
submodules of the coefficient lattice. The relevant
stabilizers modulo the center are finite by definiteness
and discreteness. Their orders are prime to \(\ell\).
Indeed a rational projective quaternion element of order
\(\ell\) would have two eigenvalues of ratio a primitive
\(\ell\)-th root \(\zeta\). Its squared reduced trace
divided by its reduced norm, minus two, equals
\(\zeta+\zeta^{-1}\) in \(D_0\). This has degree more
than two for \(\ell\ge7\); for \(\ell=5\) it generates
\(\QQ(\sqrt5)\), whereas \(D_0=\QQ(\sqrt{p_3})\) with
\(p_3\ne5\). Cauchy's theorem excludes divisibility by
\(\ell\). The compatible central character has already
been imposed in taking this quotient. Averaging over each
stabilizer now shows that the reduction of the integral
function module is the full module of residual functions
with the same equivariance, for either weight.

At a place above \(\ell\) use the upper Iwahori \(I\)
and the double coset \(I\operatorname{diag}(1,\ell)I\).
Its elements reduce to matrices of rank one with image
the first coordinate line. With the coefficient action
just described the operator sends a function \(f\) at \(g\)
to \(\sum_h\rho(h)f(gh)\) with \(h\) running through
right-coset representatives. For \(\det^m\) normalize
this operator by \(\ell^{-m}\). This multiplies the
determinant factor by the inverse of its uniformizer
part, leaving \((\det(h)/\ell)^m\) a unit. For \(m=0\)
this is the usual arithmetic \(U\) with eigenvalues the
two roots of the good Hecke polynomial for a spherical
representation. One may use either Iwahori direction for
that calculation: the element interchanging the two ends
of the Iwahori edge conjugates the two double cosets, and
the roots are those of the same spherical Hecke polynomial.
For general \(m\) the operator before normalization is
the same unweighted adelic double coset under the displayed
change of model. The scalar \(\ell^{-m}\) consequently
gives the normalized roots \(\ell^{-m}\) times the
arithmetic ones. The spherical Hecke normalization is
common to these weights since the central exponent \(j-1\)
is fixed.

On residual coefficients the image of the normalized
operator in its factor lies in the line spanned by \(X^n\),
where \(X\) is the first basis vector. Reduction of
\(\operatorname{Sym}^n h\) has this image or is zero on
a particular input. The Iwahori preserves the line, and
on it a level matrix acts by \(h_{11}^n(\det h)^m\).
For a double-coset summand the scalar is
\(h_{11}^n(\det h/\ell)^m\) on reduction. Both \(h_{11}\)
and the indicated determinant unit have residues in
\(\mathbb F_\ell^\times\), since the places split.
Thus the line with its level and normalized operator actions
is identical for the two choices of \((n,m)\) at \(v_1\);
their exponent differences \(2t,-t\) are divisible by
\(\ell-1\). The tensor of the two residual lines is also
the same for the actions at all other places and at the
center.

Choose a primitive integral eigenvector of the old system.
Its reduction takes values in both lines because its two
normalized \(U\)-eigenvalues are units. Identify the lines
in the new coefficient and obtain a nonzero eigenfunction
there with the same residual eigenvalues. The reduction
surjectivity proved above lifts it into the new integral
module. The finite algebra generated by the integral good
Hecke operators and the two normalized \(U\)'s acts faithfully
on that free module. The characteristic-zero minimal-prime
argument in Lemma~\ref{lem:weightpair578} lifts the
eigencharacter. One-dimensional quaternionic automorphic
representations have trivial type on the norm-one group at
infinity. Neither coefficient here has trivial type. Thus
the recovered system transfers to a cuspidal Hilbert form.
Its tame conductor equals the lower bound from the same
primitive central character at every tame prime. The determinant
in the arithmetic normalization is \(\kappa_\ell^j\lambda\).
The good Hecke polynomials give the same residual pair.

At a place of exponent \(m=0\) or \(-t\), the new normalized
root is a unit. Steinberg would give valuation
\((j-1)/2-m>0\) by the fixed central character, so the
underlying representation is spherical. The Hilbert
local-global compatibility theorem makes its Galois plane
crystalline. After the local twist by \(\kappa_\ell^{-m}\)
the two weights are \(n+1,0\) and there is a unit Frobenius
root with the specified reduction. The unit-root criterion
gives an unramified quotient with this root: in the dual
crystalline convention the slope-zero eigenline has Hodge
weight zero by weak admissibility and hence corresponds to
a crystalline line. Dualizing gives the stated quotient in
the arithmetic convention. Twisting back contributes
\(\bar\kappa_\ell^m=1\) on reduction. Taking the determinant
determines the two complementary subcharacters, giving
exactly the orders stated in the proposition. These
compatibility statements apply to both motivic weight
systems, including when the totally real degree is even.
\href{https://ems.press/content/serial-article-files/26055}
{C. Skinner, \emph{A note on the \(p\)-adic Galois representations
attached to Hilbert modular forms}}.
\end{proof}

\begin{theorem}[A regular cuspidal initial point for every even-order triple]
\leavevmode\label{thm:regularind579}
In the fields of Theorem~\ref{thm:eventower572} there is an
integral ordinary symplectic lift of multiplier \(\mu\)
whose residual matrices and local flags have the prescribed
even-order initial-point conditions. It is absolutely
irreducible on \(G_{HK(\mu_{\ell^\infty})}\); a scalar
twist over \(HK\) is attached to a regular algebraic
cuspidal representation of \(\mathrm{GL}_4\). The quadratic
field for the induction is the single \(D_0\) already fixed
on the interval.
\end{theorem}
\begin{proof}
Induce the plane in Proposition~\ref{prop:weightchange579}
from \(G_{D_0}\) to \(G_\QQ\). Its projective image on
\(G_{D_0}\) is dense by the argument in
Theorem~\ref{thm:mixedplane578}; residual orders at two
and distinct Hodge--Tate weights again exclude a torus
normalizer and finite image. The two conjugate planes
are inequivalent even after finite-order twisting, since
their weight gaps at \(v_0\) are \(j\) and \(j+2t\).
Their determinants agree. The intertwiner argument of
Corollary~\ref{cor:parallelind578} gives inequivalence on
\(G_{\QQ^{\rm ab}}\), irreducibility of the induction on
the required cyclotomic kernel, and cuspidal automorphic
induction over \(HK\). Its four weights are now the strictly
decreasing list
\[
                m_0=j+t,\qquad m_1=j,\qquad m_2=0,\qquad m_3=-t .
\]
This list at every complex embedding also gives regular
algebraicity in the arithmetic Satake normalization of
that corollary.

We check the lattice and flags. Write \(s=g_0\) in the
even-order tower; \(\phi(s)=-1\) and \(\chi(s)^4\ne1\).
In the plane choose an integral eigenbasis \(e,f\) of
\(W(s^2)\) with residual eigenvalues \(1,\chi(s)^2\).
Take the stable-lattice argument of
Proposition~\ref{prop:residualcube577} with the opposite
entry normalized: rescale until the minimum valuation of
the upper entries is zero. The products of upper and lower
entries lie in the maximal ideal by the two diagonal trace
equations at \(s^2\). Thus all lower entries reduce to zero
and the residual plane is a nonsplit upper extension of
\(\chi\) by one. This remains nonsplit on \(G_{HD_0}\):
inflation--restriction is injective since \(s^2\in G_{HD_0}\)
and \(\chi(s)^2\ne1\).

On the two summands of the induction \(s(x,y)=(W(s^2)y,x)\).
Choose square roots \(\alpha_0,\beta_0\) of the two
eigenvalues, with \(\alpha_0\equiv1\) and
\(\alpha_0\beta_0=-d(s)\), writing \(d=\kappa_\ell^j\xi\)
also on \(G_\QQ\). This is possible by Hensel's lemma and
\(\det W(s^2)=d(s)^2\). Use in order the four vectors
\[
  (\beta_0 f,f),\quad(-\alpha_0 e,e),\quad
  (\beta_0 f,-f),\quad(\alpha_0 e,e).
\]
They form an integral basis since \(2,\alpha_0,\beta_0\)
are units. For the alternating determinant forms on the two
summands with coefficients \(1,-d(s)\) as in
Corollary~\ref{cor:parallelind578}, both ordered antidiagonal
pairings are \(-2\alpha_0\beta_0\). Divide the form by this
unit. The matrix at \(s\) is diagonal with entries
\((\beta_0,-\alpha_0,-\beta_0,\alpha_0)\).
The unique square root \(b\equiv1\) of \(\mu/(d\phi)\)
on \(G_H\) gives the desired multiplier after twisting by
\(b\).

On \(G_{HD_0}\) let \(c,c'\) denote the two residual upper
entries of the plane and its conjugate, evaluated at the
same group element. Their diagonal pairs are both
\((1,\chi)\). Reading the induced action in the four
vectors above, its only possible off-diagonal positions
are \(10,12,30,32\), with indices starting at zero.
The entries at \(10,30,12\), in that order, are
\[
 \tfrac12(c'-\gamma c),\qquad \tfrac12(c'+\gamma c),
              \qquad -\tfrac12(c'+\gamma c),
       \qquad \gamma=\overline{\beta_0/\alpha_0}.
\]
The support assertion extends to the other coset by
multiplying by the diagonal at \(s\). The full residual
diagonal is \((\chi\tau,\tau,\chi,1)\) since \(\phi=\tau\)
on \(G_H\). Either the \(10\)-entry occurs or
\(c'=\gamma c\) identically, in which case both \(12,30\)
occur, since some \(c\ne0\). These are exactly the two
alternatives for the residual entries in the even-order
initial point.

Above \(\ell\) distinguish the summand with weight gap
\(j+2t\) and the other summand. The top line in the former
reduces to a generator \(f+ae\), since it has residual
character \(\chi\) different from one on the decomposition
group at the split rational quadratic-field place. In the
other summand the top line reduces to \(e\): it has character
one and this eigenline is unique at an element separating
\(\chi\) and one. The complete flag in descending weight
order consists successively of the first of these lines,
both top lines, and the first line together with the whole
other summand. These are isotropic in dimension two and
self-dual as a full flag for the displayed form.

Normalize its first column to have coordinate zero equal
to one. Its coordinate two is then \(1\) or \(-1\),
according to the summand containing the first line: the
two \(f\)-coordinates of a vector in either single summand
have this ratio in the displayed basis. In the other
summand the second top vector has, modulo the maximal ideal,
zero in coordinates zero and two and a unit in coordinate
one. Subtract a multiple of the first column to clear
coordinate zero, and divide by the remaining unit pivot.
Its coordinate two is still zero on reduction. The third
subspace is the perpendicular of the first line. In it the
vector with first two coordinates zero and coordinate two
one has uniquely determined last coordinate, since the
first pivot is one. Use it as third column and the last
coordinate vector as fourth. The resulting lower
unitriangular matrix is symplectic: the two prescribed
pairings are one by the pivots and all other pairings vanish
by isotropy, perpendicularity and the first two zero
coordinates. Thus its \(20\)-entry is a unit and its
\(21\)-entry reduces to zero. All other lower positions
are allowed. In this flag the ordered residual local
characters are \((\chi,1,\chi,1)\), as required since
\(\tau=1\) on the decomposition groups of \(H\) at \(\ell\).

Away from \(\ell\) the tame conductor of \(W\) equals that
of \(\lambda\); its inertia has the two characters
\(1,\lambda\) with zero monodromy by the conductor argument
of Theorem~\ref{thm:cuspidalinitial577}. At every finite
place \(HD_0/H\) is unramified by
Theorem~\ref{thm:eventower572}. Inertia over \(H\) therefore
acts through the two summands; it acts trivially where the
residual diagonal is trivial, since \(\xi\) has order prime
to \(\ell\), and the same holds for the scalar correction
there. At \(\ell\) the crystalline graded characters on the
split quadratic-field places are cyclotomic powers times
unramified characters. They and the correction kill the
\(\ell\)-power torsion of local abelianization, as in
Theorem~\ref{thm:cuspidalinitial577}. On open inertia the
squared characters after correction have exponents
\(h_i=2m_i+r+1-j\). The inverse twist \(b^{-1}\), which
recovers the automorphic induction, has squared exponent
\(a=j-r-1\). We have \(a+h_0=2m_0\) and the consecutive
gaps in the \(m_i\) are at least two. This verifies the
weight and twist conditions with the original multiplier.
The field, matrix and flag conditions just proved are
those in \srcloc{Development/AutomorphicArithmetic/AutomorphicVarieties/WeakShimuraTheory/NearSplit}{305}.
\end{proof}

\begin{corollary}[Initial points in both residual character cases]
\leavevmode\label{cor:bothinitial579}
For every triple of Section~\ref{sec:sourcelevels563}, the
required auxiliary character case admits a regular cuspidal
initial point with its indicated integral residual shape on
the towers of Corollary~\ref{cor:bothaux572}. The fields
\(K,K',D_0\) and the single auxiliary prime needed in
the odd-order construction can be fixed throughout the
period interval.
\end{corollary}
\begin{proof}
Apply Theorem~\ref{thm:regularind579} for even order and
Theorem~\ref{thm:cuspidalinitial577} for odd order. The
common base choices are those of Lemmas~\ref{lem:twoprimes571},
\ref{lem:realquadratic572} and~\ref{lem:cuspidalplane577}.
\end{proof}

The initial points supply cuspidal points with particular
residual extensions and ordinary flags. To bound the
strict lengths at every triple by this route one must
still propagate automorphic support to the deformations
that witness a failure of the bound. The signed integral
period comparison remains a second requirement on the
original fields.


\subsection{Ordinary deformations at the two cuspidal points}
\label{sec:seedweights580}

Write \(F=H\) in Theorem~\ref{thm:regularind579} and \(F=H_1\)
in Theorem~\ref{thm:cuspidalinitial577}, with \(C_{w,\ell}\subset H_1\).
Put \(L=FK\), a CM field with maximal totally real subfield \(F\).
Let \(R\) be the indicated four-dimensional representation of \(G_F\)
and put \(U=b^{-1}R\), removing the scalar correction of the
construction. Choose a finite coefficient extension \(E/\QQ_\ell\)
over which the matrices, the alternating form \(J\), and the ordinary
flags are defined. Fix the multiplier of \(R\). The corresponding
Lie algebra of infinitesimal changes is
\[
   V=\{X\in\operatorname{End}_E(E^4): X^{\mathsf T}J+JX=0\},
\]
with \(G_F\) acting by conjugation. \(R\) and \(U\) give the same
action. All cohomology below is continuous.

\begin{proposition}[The two finite Selmer groups on the auxiliary field]
\leavevmode\label{prop:seedadjoints580}
At each of these points both \(H^1_f(F,V)\) and
\(H^1_f(F,V^\vee(1))\) vanish. Here the finite condition is the
unramified condition away from \(\ell\) and the Bloch--Kato
crystalline-period condition above \(\ell\).
\end{proposition}
\begin{proof}
On \(G_L\) the representation \(U\) comes from a regular algebraic
cuspidal representation \(\Pi\), and its restriction to
\(G_{L(\mu_{\ell^\infty})}\) is absolutely irreducible by the two
theorems just cited. We specify the polarization to which we apply
adjoint vanishing. The multiplier \(\nu\) of \(U\) is
\((\det W)^3\) for the symmetric cube and \(d\phi\) for the quadratic
induction. In both cases it is an integral cyclotomic power times
a finite character of \(G_F\), and \(\nu(c)=-1\) at every real
conjugation \(c\). For one such \(c\) put \(C=U(c)\). The form
\(B=JC\) is symmetric since \(C^2=1\) and \(C^{\mathsf T}JC=-J\).
It satisfies
\[
 U(g)^{\mathsf T}B\,U(cgc^{-1})=\nu(g)B \qquad(g\in G_L).
\]
In particular the good parameters of \(\Pi^c\) are the dual
parameters of \(\Pi\) times a character from \(F\). Indeed the
display gives this identity in the Galois normalization. Restoring
the Satake normalization changes the twisting character by an
integral norm power, the square of the normalizing factor in
rank four. A finite character times a norm power is an algebraic
Hecke character, by class field theory. Strong multiplicity one
gives the corresponding global identity. One can normalize the
real sign of the polarizing character by the quadratic character
of \(L/F\), which is trivial on norms from \(L\).
This verifies polarizability with the specified restriction of
the multiplier. For that restriction \(B\) is unique up to
scalar by absolute irreducibility. These are the conventions
of Barnet-Lamb, Gee, Geraghty and Taylor, Section~2.1 of
\href{https://annals.math.princeton.edu/2014/179-2/p03}
{\emph{Potential automorphy and change of weight}}.

On the polarized adjoint over \(F\), \(G_L\) acts by conjugation
and \(c\) acts by \(X\mapsto-B^{-1}X^{\mathsf T}B\). If \(X\in V\)
the latter map equals \(CXC^{-1}\). The self-adjoint endomorphisms
for \(J\) give a stable complement, as the formula also shows.
Thus \(V\) with its original \(G_F\)-action is a direct summand.
Thorne's vanishing theorem over the CM field \(L\) applies using
the irreducibility on the cyclotomic kernel verified above. It
annihilates the finite Selmer group of the polarized adjoint over
\(F\), hence that of \(V\). This argument applies to the induced
representation even though it splits over the inducing field.
\href{https://arxiv.org/pdf/2207.04925}{J.~A. Thorne,
\emph{On the vanishing of adjoint Bloch--Kato Selmer groups of
irreducible automorphic Galois representations}, Theorem~6.1}.

For the dual assertion apply the global Selmer dimension formula
with the finite local conditions, which are mutual annihilators
under local duality. Both global invariant terms vanish.
For \(V\) this follows from Schur's lemma, since a scalar in
\(V\) is zero. The trace form identifies \(V^\vee\) with \(V\).
A nonzero invariant in \(V(1)\) would give an isomorphism
\(U\simeq U(1)\) by absolute irreducibility. Determinants would
then give \(\kappa_\ell^4=1\) on \(G_F\).
At a finite \(v\nmid\ell\) unramified cohomology has the same
dimension as the local invariants. At \(v\mid\ell\), \(V\) is
crystalline since \(U\) is crystalline. Its ten weights at each
embedding have four positive root weights, their negatives and
two zeros, using the convention that \(\kappa_\ell\) has weight
one. The local difference \(h^1_f-h^0\) is consequently
\(4[F_v:\QQ_\ell]\) by the Bloch--Kato dimension formula. At a
real place the two eigenspaces of \(C\) are complementary
Lagrangians since \(C\) reverses the form. An infinitesimal
symplectic transformation commuting with \(C\) acts on one by
an arbitrary two-by-two matrix and on the other by its negative
transpose in dual bases. The real term is therefore \(-4\).
Summing gives zero, since \(F\) is totally real. Thus the dual
finite Selmer group has the same dimension as the first.
\cite{XII-IX-VI-bk117,XII-IX-V-milne98}
\end{proof}

For \(v\mid\ell\) let \(\mathfrak b_v\subset V\) preserve the
chosen flag and let \(\mathfrak n_v\) be its strictly upper
triangular part. The quotient \(\mathfrak t=\mathfrak b_v/
\mathfrak n_v\) is the trivial rank-two module of infinitesimal
diagonal changes with fixed multiplier.

\begin{lemma}[Ordinary and fixed-weight local directions]
\leavevmode\label{lem:seedlocal580}
Write \(G_v=G_{F_v}\). The map \(H^1(G_v,\mathfrak b_v)\to
H^1(G_v,V)\) is injective. Restriction of its diagonal characters
to local units gives an exact sequence
\[
 0\longrightarrow H^1_f(F_v,V)\longrightarrow
 H^1(G_v,\mathfrak b_v)\longrightarrow
 \operatorname{Hom}_{\rm cont}(\mathcal O_{F_v}^{\times},E)^2
 \longrightarrow 0 .
\]
Over a small extension of local Artin \(E\)-algebras every ordinary
local deformation lifts with any prescribed lift of its two
diagonal inertia characters.
\end{lemma}
\begin{proof}
The ordinary characters of \(U\) are unramified characters times
cyclotomic powers. Their consecutive weight gaps are at least two.
Thus \(\mathfrak n_v\) has a filtration by four positive-root
characters, each of weight at least two, and \(V/\mathfrak b_v\)
has the four opposite weights. It follows that
\(H^0(G_v,V/\mathfrak b_v)=0\). Local duality gives
\(H^2(G_v,\mathfrak n_v)=H^2(G_v,\mathfrak b_v)=0\), since
the duals twisted by one have no constituent of weight zero.
The first vanishing of invariants gives the injection. The
vanishing of \(H^2(G_v,\mathfrak n_v)\) gives surjectivity
onto \(H^1(G_v,\mathfrak t)\). By local reciprocity, restriction
from \(F_v^\times\) to its units is onto on additive characters,
as one can choose their values at a uniformizer.

Put \(d_v=[F_v:\QQ_\ell]\). The kernel of the map to units
has dimension \(h^0(G_v,\mathfrak b_v)+6d_v-2d_v\), by the
local Euler characteristic formula. It contains
\(H^1_f(F_v,\mathfrak b_v)\) since the finite local condition
on the trivial module \(\mathfrak t\) is unramified. The
Bloch--Kato dimension formula for the crystalline submodule
\(\mathfrak b_v\) gives exactly the displayed kernel dimension.
Under the injection to \(H^1(G_v,V)\) this finite subspace
maps into \(H^1_f(F_v,V)\). The two spaces have equal dimension,
since \(h^0(G_v,V)=h^0(G_v,\mathfrak b_v)\) and both have four
positive weights. This proves the asserted identification.

For the lifting assertion choose a symplectic basis for the
lifted flag, possible by lifting a basis of each successive
summand and adjusting the dual vectors. First lift the two
diagonal characters by prescribing the given values on units
and lifting their values at a uniformizer. The fixed multiplier
then determines the other two. For a small extension with kernel
\(I\) annihilated by the maximal ideal, arbitrary upper triangular
symplectic-similitude matrix lifts with these diagonals have
multiplication defect a continuous two-cocycle in
\(\mathfrak n_v\otimes_E I\). Continuous choices exist using
coordinates on the upper unipotent group and an \(E\)-linear
section of the coefficient surjection. Its class vanishes by
the \(H^2\) computation above. Subtracting a one-cochain whose
coboundary is the defect gives the required homomorphism.
\end{proof}

\begin{theorem}[The formal ordinary family at either initial point]
\leavevmode\label{thm:seedweight580}
Consider deformations of \(R\) on \(G_F\) with fixed multiplier,
the indicated ordered ordinary flags at \(\ell\), and inertia
conjugate to the original finite inertia away from \(\ell\).
Take continuous deformations over local Artin \(E\)-algebras
with residue field \(E\), up to symplectic conjugation reducing
to the identity. For each such algebra the restriction to units
of the first two diagonal characters at every \(v\mid\ell\)
identifies this functor with the functor of arbitrary lifts of
these unit characters. In particular the unframed completed
deformation ring at the point is a power series ring over \(E\)
in \(2[F:\QQ]\) variables. Fixing all the inertial weights gives
the local ring \(E\).
\end{theorem}
\begin{proof}
Enlarge a finite set \(S\) to include the archimedean places,
those over \(\ell\) and all ramification of \(R\). Outside \(S\)
require unramifiedness. At finite \(v\in S\) away from \(\ell\)
the tangent condition for fixed inertia is unramified cohomology.
Indeed a class restricts to zero on inertia exactly when an
infinitesimal conjugation makes the two inertia actions agree.
This local condition lifts over small extensions. Having
identified the inertia action with the original one, a Frobenius
matrix is the original Frobenius times an element centralizing
finite inertia in the symplectic group. This centralizer is a
smooth algebraic group in characteristic zero and its elements
lift. The resulting matrix, together with finite inertia, defines
a continuous lift: the closure of the powers of the Frobenius
matrix is compact. In fact its eigenvalues are units and the
powers of each unipotent Jordan block are bounded by the
integrality of the binomial coefficients. This argument uses
the faithful \(E\)-linear realization of matrices over the
finite-dimensional Artin algebra. At a real place the group
of order two has zero higher cohomology.

Now let \(A'\to A\) be a small extension with kernel \(I\).
Start with a global deformation over \(A\) and a lift of its
unit characters. At places over \(\ell\) choose local lifts
by Lemma~\ref{lem:seedlocal580}. At the other places of \(S\)
use the preceding paragraph. The obstruction to a global
fixed-multiplier lift belongs to \(H^2(G_{F,S},V)\otimes_E I\);
it is the multiplication defect of arbitrary continuous matrix
lifts, computed as in the lemma with the full Lie algebra.
The chosen local lifts make its restrictions zero. By
Poitou--Tate duality the kernel of this localization in degree
two is dual to the kernel of localization in degree one for
\(V^\vee(1)\). The latter kernel is contained in its finite
Selmer group, which vanishes by Proposition~\ref{prop:seedadjoints580}.
Hence a global lift exists before imposing the local choices.

Compare this lift with the chosen ones at \(S\). Their
differences are local classes modulo \(H^1_f(F_v,V)\), tensored
with \(I\), by Lemma~\ref{lem:seedlocal580} and the fixed-inertia
calculation. The same duality sequence gives a surjection
\[
 H^1(G_{F,S},V)\longrightarrow
 \bigoplus_{v\in S,\ v\ {\rm finite}}
       H^1(F_v,V)/H^1_f(F_v,V),
\]
since the dual finite Selmer group is zero. Correct the global
lift by a preimage of the discrepancies. It now realizes the
local conditions with the chosen weights. Two such corrected
lifts of the same deformation differ up to strict conjugation
by a class in \(H^1_f(F,V)\otimes_E I\), which is zero. A lift
of the ordered flag for a fixed matrix representation is unique
at each small extension: the difference of two is invariant in
\((V/\mathfrak b_v)\otimes_E I\), whose invariants vanish by the
lemma. Induction on the length of the algebra proves the asserted
bijection, naturally in the algebra.

Finally the formal ratios of lifted unit characters to the
originals have values in \(1+\mathfrak m_A\). The finite logarithm
identifies this group with the additive \(\mathfrak m_A\).
Local units modulo their finite torsion have a pro-\(\ell\)
free part of rank \(d_v\); their open principal units have this
rank by the local logarithm. Two characters at every place
therefore give \(2\sum_{v\mid\ell}d_v=2[F:\QQ]\) formal
parameters. Sending all of them to zero gives the fixed-weight
assertion.
\end{proof}

The coefficient field in this theorem has characteristic zero.
The remaining support problem includes components in characteristic
\(\ell\), where the ordinary weights can collide. The theorem gives
no bound for the integral lengths of such components on the
original period fields.

\section{Quadratic induction and symmetric cubes}\label{nav:chapter-quadratic-induction-symmetric-cubes}

\subsection{Quadratic induction throughout the even-order ordinary family}
\label{sec:quadraticfamily583}

Take an even-order triple on the interval of
Corollary~\ref{cor:bothinitial579}. Write \(F=H\),
\(M=FD_0\), and let \(\phi\) be the nontrivial character of
\(\operatorname{Gal}(M/F)\). The extension is unramified at
every finite place and split above \(\ell\), by
Theorem~\ref{thm:eventower572}. Use a coefficient field \(E\)
as in Theorem~\ref{thm:seedweight580}. In its ordered ordinary
flags the initial representation has the description
\[
  R_0=b\,\operatorname{Ind}_{G_M}^{G_F}W_0,\qquad
  \det W_0=d|_{G_M},\qquad b^2 d\phi=\mu .
\]
Here \(W_0\) is the restriction to \(M\) of the Hilbert modular
plane of Proposition~\ref{prop:weightchange579}. Fix the first
summand of the induction over \(G_M\) to be the one labeled
\(W_0\). All Artin algebras below are local \(E\)-algebras
with residue field \(E\).

\begin{theorem}[Induction on the entire formal ordinary family]
\leavevmode\label{thm:quadraticfamily583}
Induction followed by multiplication by \(b\) gives a natural
bijection from the ordinary deformations of \(W_0\) over \(M\)
of fixed determinant \(d|_{G_M}\) and fixed inertia away from
\(\ell\) to the ordinary deformations of \(R_0\) of
Theorem~\ref{thm:seedweight580}. For the plane one retains
the indicated upper line at each place over \(\ell\) and uses
strict changes of basis. For the four-dimensional representation
one uses the original order of the four graded lines. On the
coset \(G_F\setminus G_M\) the trace vanishes in every coefficient
algebra, including those with nilpotents.
\end{theorem}
\begin{proof}
Write the induced module in the two summands specified above,
and let \(S\) act by \(1\) on the first and by \(-1\) on the
second. The alternating form in Theorem~\ref{thm:regularind579}
is an orthogonal sum of determinant forms on the two planes.
Thus \(S\) is symplectic and satisfies
\[
       S R_0(g)S^{-1}=\phi(g)R_0(g),\qquad S^2=1.
\]
It preserves the initial flags. In fact at each split place
the first line is in the plane with larger weight gap, the
second subspace is the sum of the two upper lines, and the
third subspace is the sum of the first line and the entire
other plane.

Let \(R_A\) be an ordinary deformation over \(A\), with flags
\(\mathcal F_v\). The matrices \(\phi(g)S R_A(g)S^{-1}\)
with flags \(S\mathcal F_v\) give another deformation of
the same point and multiplier. At a finite place away from
\(\ell\), \(\phi\) is trivial on inertia and \(S\) centralizes
the original inertia. Conjugating a strict inertia identification
by \(S\) gives one for the new matrices. At \(v\mid\ell\)
the character \(\phi\) is trivial on the whole local group.
The ordered graded characters, in particular their restrictions
to units, agree for the two deformations. The uniqueness in
Theorem~\ref{thm:seedweight580} supplies a strict symplectic
change of basis \(P\) from the new deformation to the given one.
It carries the indicated flags as well. Consequently \(J_A=PS\)
preserves those flags, reduces to \(S\), and satisfies
\[
        J_A R_A(g)J_A^{-1}=\phi(g)R_A(g).
\]

The commutant of \(R_A(G_F)\) consists of scalars. Indeed
absolute irreducibility of \(R_0\) implies that the \(E\)-span
of its group matrices is the entire matrix algebra. This follows
from Burnside's theorem after extending to an algebraic closure
and then by equality of dimensions over \(E\). Choose sixteen
group elements giving a basis. The determinant of their lifted
coordinates is a unit in \(A\), so the lifted group matrices span
\(\operatorname{Mat}_4(A)\). The square \(J_A^2\) is therefore a
symplectic scalar reducing to \(1\). Such a scalar \(u\) satisfies
\((u-1)(u+1)=0\) with \(u+1\) a unit. Thus \(J_A^2=1\).
Two symplectic intertwiners reducing to \(S\) likewise differ
by a strict symplectic scalar and coincide. This uniqueness
commutes with change of coefficients and strict conjugation.

The two idempotents \((1\pm J_A)/2\) give free summands of rank
two, since their images are projective over a local ring and
their reductions have dimension two. The summands are orthogonal
and the form on each is perfect. The group \(G_M\) preserves
both and the other coset interchanges them. On the positive
summand let \(W_A\) be the restriction of \(b^{-1}R_A\) to
\(G_M\). Its determinant is \(b^{-2}\mu|_{G_M}=d|_{G_M}\),
since on a plane the multiplier of a determinant form is the
determinant. The map from the induction which sends \(g\otimes x\)
to \(b(g)^{-1}R_A(g)x\) is well-defined and equivariant for
\(b^{-1}R_A\). On each of the two coset summands it is an
isomorphism onto the corresponding plane. The interchange of
the planes gives the asserted zero trace.

Each of the two idempotents preserves every local flag submodule
above \(\ell\). On each free rank-one graded quotient an idempotent
is either zero or one and has precisely its original value on
reduction. Applying the idempotents to the flag therefore recovers
the two plane filtrations with their prescribed upper lines.
For the two places of \(M\) above one place of \(F\), these are
the filtrations on the first and on the conjugate plane,
identified by a fixed element of \(G_F\setminus G_M\).
At a place away from \(\ell\), identify the four-dimensional
inertia with its constant lift by strict conjugation. In these
coordinates the positive projector \(p\) and its reduction
\(p_0=(1+S)/2\) both commute with the constant inertia.
The matrix \(p_0p+(1-p_0)(1-p)\) also commutes with it,
is the identity on reduction, and intertwines \(p\) with \(p_0\).
It identifies the inertia on the lifted plane with the original
one. Multiplication by the fixed scalar character \(b^{-1}\)
preserves this identification.

Conversely fix \(\sigma\in G_F\setminus G_M\) and start with
such a plane \(W_A\). In a basis lifting the original basis,
use the determinant forms with coefficients \(1,-d(\sigma)\)
on the two induced summands and the original overall scalar
normalization. On the subgroup both determinants are \(d\).
The matrix of \(\sigma\) sends
\((x,y)\) to \((W_A(\sigma^2)y,x)\); it multiplies this form
by \(-d(\sigma)\) because \(\det W_A(\sigma^2)=d(\sigma)^2\).
The full multiplier after twisting is \(\mu\). Interleave
the two lifted upper lines in the order described at the
beginning of the proof. This produces the required self-dual
flag above \(\ell\). Away from \(\ell\), inertia lies in
\(G_M\). The strict inertia identifications in the two planes
may each be given determinant one by multiplying by a scalar
congruent to one, since the square map is invertible on
\(1+\mathfrak m_A\) by the finite binomial series. Their
orthogonal sum is symplectic.

The same scalar adjustment makes strict plane changes of basis
compatible with symplectic equivalence. In the inverse construction
choose a basis of the positive summand lifting its original
basis with the prescribed determinant pairing, by rescaling
one vector by a strict unit. Transport it to the other summand
by \(b(\sigma)^{-1}R_A(\sigma)\). The two resulting forms have
the coefficient ratio \(-d(\sigma)\) just computed, so the
induction isomorphism is a strict symplectic one after identifying
the original bases. It respects the recovered flags. Finally, on
an induced lift the sign operator on its two summands already
has the defining properties of \(J_A\). Uniqueness recovers its
first plane. This proves both inverse identities and naturality.
\end{proof}

\begin{corollary}[All self-twisting characters at the two initial points]
\leavevmode\label{cor:allselftwists583}
For every deformation in Theorem~\ref{thm:seedweight580}, a
continuous \(A^\times\)-valued character giving an isomorphism
with its twist is trivial in the odd-order case and belongs
to \(\{1,\phi\}\) in the even-order case. Both characters in
the latter set give isomorphisms. In particular \(M\) is the
unique quadratic inducing field of an even-order lift.
\end{corollary}
\begin{proof}
Taking determinants shows that every value of the character
is a fourth root of unity. Reduction identifies these roots
with those in \(E\) because a root lifts as a constant and the
quotient of two lifts of the same root is a fourth root
congruent to one, hence one. Reducing the isomorphism now
gives a self-twist of \(R_0\). In the odd-order case the
projective image of the underlying plane stays dense in
the connected group \(\mathrm{PGL}_2\) on every finite-index
subgroup. Its third symmetric power remains absolutely
irreducible there, as in Theorem~\ref{thm:cuspidalinitial577}.
An intertwiner commutes with the restriction to the kernel
of the character, so it is scalar and the character is trivial.

In the even-order case restrict to the intersection of that
kernel with \(G_M\). Both plane constituents remain absolutely
irreducible there by connected projective density. They remain
inequivalent because at any extension of the embedding in
Proposition~\ref{prop:weightchange579} their Hodge--Tate weight
gaps are \(j\) and \(j+2t\), which remain distinct on every
finite field extension. A common scalar twist cancels in
this comparison. Hence the intertwiner over \(E\) is scalar
on each plane, and commutes with all of \(G_M\). The twisting
character is therefore trivial on \(G_M\).
Theorem~\ref{thm:quadraticfamily583} realizes the remaining possibility
\(\phi\) over \(A\). Any induction from a quadratic field has
the self-twist by its nontrivial quadratic character, as one
sees by the sign operator in the two coset summands. This gives
the uniqueness assertion.
\end{proof}

\begin{corollary}[Trace identities on components through an even-order point]
\leavevmode\label{cor:quadraticcomponent583}
Let a family with generic coefficient characteristic zero have a
point with residue field \(E\) and Noetherian local ring \(B\).
Suppose its reductions over all \(B/\mathfrak m_B^n\) are
deformations in Theorem~\ref{thm:quadraticfamily583} in compatible
coordinates. For every \(g\in G_F\setminus G_M\) its trace
vanishes in \(B\), hence on each irreducible component through
the point whenever the trace is defined on the ambient family.
\end{corollary}
\begin{proof}
The trace belongs to every power of \(\mathfrak m_B\) by the
theorem. Their intersection in a Noetherian local ring is zero.
In an affine neighborhood where the trace is defined it is
therefore killed by an element invertible at the point.
That element is outside every minimal prime contained in the
prime of the point, so the trace lies in each such minimal prime.
\end{proof}

This identifies the ordinary family at the specified
characteristic-zero point for every even-order triple on the
interval. Bounding the integral local cohomology still requires
control of the components carrying the torsion classes.


\subsection{The symmetric cube and the prescribed integral weights}
\label{sec:cubeweights584}

In the odd-order construction write \(F=H_1\), \(R_0=b\operatorname{Sym}^3 W_0\),
and \(\mu=b^2(\det W_0)^3\), as in Theorem~\ref{thm:cuspidalinitial577}.
Here \(W_0\) denotes the plane after restriction to \(G_F\).
At \(v\mid\ell\) number the four ordered diagonal characters from
zero to three, starting with the top of the ordinary flag.
Given an ordinary lift as in Theorem~\ref{thm:seedweight580},
let \(\epsilon_{v,i}\) be the ratio of its \(i\)-th character on
local units to the corresponding character of \(R_0\).

\begin{theorem}[The equations for the formal symmetric cubes]
\leavevmode\label{thm:cubeweights584}
Over every local Artin \(E\)-algebra with residue field \(E\),
the lifts of \(R_0\) obtained by applying \(b\operatorname{Sym}^3\)
to ordinary lifts of the plane of fixed determinant and fixed
inertia away from \(\ell\), with the induced ordered flags, are
precisely the points satisfying
\[
                  \epsilon_{v,0}=\epsilon_{v,1}^3
                         \qquad(v\mid\ell).
\]
They form a smooth closed formal subscheme of codimension
\([F:\QQ]\) in the ordinary deformation space. The same equations
hold if the plane determinant and scalar twist are both allowed
to vary subject to their original reductions and the fixed
four-dimensional multiplier.
\end{theorem}
\begin{proof}
For the plane put \(V_2=\mathfrak{sl}(W_0)\). The derivative
of its symmetric cube embeds \(V_2\) into the fixed-multiplier
adjoint \(V\) of Proposition~\ref{prop:seedadjoints580}.
The trace pairing on this image equals ten times
\((X,Y)\mapsto\operatorname{tr}_{W_0}(XY)\). Indeed both
pairings are invariant, and the invariant symmetric forms on
\(\mathfrak{sl}_2\) are multiples of its trace form. This last
fact follows from invariance under the diagonal torus and the
brackets of the upper and lower root vectors. For
\(X=Y=\operatorname{diag}(1,-1)\) the eigenvalues in the cube
are \(3,1,-1,-3\), giving \(20\) in place of \(2\). The
orthogonal complement therefore splits off over \(E\), even
when \(\ell=5\). The finite local conditions commute with
direct sums, as do their duals. The two vanishings of
Proposition~\ref{prop:seedadjoints580} imply
\[
            H^1_f(F,V_2)=H^1_f(F,V_2^\vee(1))=0 .
\]

We record why every choice of the single plane weight at each
place is attained. Above \(\ell\) the trace-zero flag-preserving
algebra \(\mathfrak b_2\) has its upper root \(\mathfrak n_2\)
and a trivial one-dimensional quotient. The upper root has
weight \(j\ge2\), and the quotient \(V_2/\mathfrak b_2\) has
weight \(-j\). Thus the latter has no invariants, while local
duality gives \(H^2(F_v,\mathfrak n_2)=H^2(F_v,\mathfrak b_2)=0\).
It follows that ordinary classes inject into \(H^1(F_v,V_2)\)
and map onto the local diagonal classes, hence onto additive
characters of local units by local reciprocity. The kernel of
the map to units has dimension
\(h^0(F_v,\mathfrak b_2)+[F_v:\QQ_\ell]\), by the local Euler
characteristic formula. It contains the finite subspace of
\(\mathfrak b_2\), since a finite class in the trivial diagonal
quotient is unramified. By the crystalline dimension formula
this finite subspace has the displayed dimension, as does
\(H^1_f(F_v,V_2)\), since the local invariants coincide and
there is one positive weight. Thus the kernel is precisely
the finite local condition. Over a small extension the lifting
obstruction with a given diagonal lift lies in the second
cohomology of the upper root and vanishes. Away from \(\ell\)
keeping finite inertia fixed gives the unramified tangent
condition and lifts over a small extension by lifting Frobenius
in its smooth centralizer, as in Theorem~\ref{thm:seedweight580}.

With these local lifts the global obstruction restricts to
zero at every place. By Poitou--Tate duality the kernel in degree two is dual
to the everywhere locally zero classes for \(V_2^\vee(1)\),
all of which are finite. It vanishes. The discrepancy between
an unrestricted global lift and the prescribed local lifts,
viewed modulo the local finite conditions, can be removed
by a global class because the dual finite Selmer group is
zero. With the weights fixed the ambiguity is the first
finite Selmer group, also zero. The local flag itself is unique
over a small extension since \(V_2/\mathfrak b_2\) has no
invariants. Induction on the length of the coefficient algebra
gives a unique plane up to strict equivalence for each choice
of its upper character on units.

Write \(z_v\) for the ratio of the lifted upper character to
the original one. The lower ratio is \(z_v^{-1}\). The first
two ratios for the cube are \(z_v^3,z_v\). Conversely take
\(z_v=\epsilon_{v,1}\) and use the plane just constructed.
Its twisted cube has the required multiplier, inertia and flags.
If the equations in the statement hold, its first two characters
agree with the given ones on units, so the two lifts agree
up to strict symplectic conjugation by
Theorem~\ref{thm:seedweight580}. In logarithmic weight
coordinates the equations subtract three times the second
coordinate from the first at each principal-unit generator.
Each has coefficient one on a distinct first coordinate.
There are \(\sum_{v\mid\ell}[F_v:\QQ_\ell]=[F:\QQ]\) such
generators. This proves the assertion on formal schemes.
Finally multiply a plane of variable determinant by the
inverse square root of its determinant ratio. This square
root is a character given by the finite binomial series on
strict units. Absorb its third power in the outside twist.
After this operation the determinant is fixed and the
remaining twist divided by \(b\) is a strict character of
square one, hence is identically one.
\end{proof}

We now use the prescribed integral characters, independently
of the existence of a lift of the initial point. Take any
triple of Section~\ref{sec:sourcelevels563}. Write
\[
  A=\kappa_\ell\omega_\ell^{-1}\xi,\quad
  \Psi=\alpha\kappa_\ell^r,\quad \mu=A\Psi,\qquad
  (\delta_0,\delta_1,\delta_2,\delta_3)=(A\Psi,\Psi,A,1).
\]
Let \(O\) be a coefficient DVR containing their values. For a
local inertia group \(I\), let \(\mathfrak c(I)\subset O\)
be generated by all
\(\delta_1(s)^3/(\mu(s)\delta_0(s))-1\) for \(s\in I\).
For an ordered scalar-twisted cube of a plane with diagonal
characters \(u,v\) and scalar \(b'\), its first two characters
are \(b'u^3,b'u^2v\) and its natural multiplier is
\((b')^2(uv)^3\). They give exactly zero in this equation
over any coefficient ring.

\begin{theorem}[Exact integral depth of the ordered weight equation]
\leavevmode\label{thm:cubedepth584}
At each \(v\mid\ell\) of the original period field \(E_x\),
throughout the common interval,
\[
 \mathfrak c(I_{E_{x,v}})=
 \begin{cases}
    O, &\ell-1\nmid r,\\
    \ell^{1+v_\ell(r-2)}O, &\ell-1\mid r .
 \end{cases}
\]
In the second case \(r>2\). In the odd-order case take the
auxiliary \(F=H C_{w,\ell}\) of Theorem~\ref{thm:cuspidalinitial577}.
At all places over \(\ell\) of either \(F\) or \(FE_x\) the ideal
is \(\ell^{1+v_\ell(r-2)}O\) for \(r>2\), and it is zero for
\(r=2\). These are equalities of ideals regardless of the
ramification index of \(O\).
\end{theorem}
\begin{proof}
The character in the definition of \(\mathfrak c\) before
subtracting one equals \(\Psi/A^2\).
Write \(\langle\kappa_\ell\rangle=\kappa_\ell/\omega_\ell\).
Since \(E_x\) is unramified at \(\ell\) and \(\alpha,\xi\)
are also unramified there, the character on this inertia is
\[
             \omega_\ell^r\langle\kappa_\ell\rangle^{r-2}.
\]
The extension generated by the \(\ell\)-power roots of unity over \(\QQ_\ell\) is totally
ramified, so this inertia realizes every value in
\(\ZZ_\ell^\times\) under \(\kappa_\ell\). A nontrivial tame
factor gives an element whose character value differs from one
in the residue field, hence the unit
ideal. If the tame factor is trivial, \(r\ge\ell-1\ge4\).
The pro-\(\ell\) part has topological generator \(1+\ell\).
For a nonzero integer \(n\) the ideal generated by \(u^n-1\)
as \(u\) runs over \(1+\ell\ZZ_\ell\) is
\(\ell^{1+v_\ell(n)}O\). All values have that divisibility
by the binomial expansion; equality occurs at \(u=1+\ell\),
since a prime-to-\(\ell\) power preserves the first nonzero
digit and an \(\ell\)-th power raises its valuation by one.
This proves the first formula.

On \(G_H\) the finite characters \(\xi^2\) and
\(\alpha\omega_\ell^r\) coincide. They have prime-to-\(\ell\)
order and have the same reduction by Theorem~\ref{thm:tametower571}.
Thus \(\Psi/A^2=\langle\kappa_\ell\rangle^{r-2}\) on \(G_H\).
The local ramification index of \(H\) divides two, and
\(C_{w,\ell}\) and \(E_x\) are unramified at \(\ell\).
Adjoining unramified extensions locally preserves this
ramification index. The inertia subgroups upstairs thus
have index prime to \(\ell\) in the original rational
inertia. Their images onto the pro-\(\ell\) cyclotomic
factor are still surjective. The same binomial calculation
gives the asserted ideal; at exponent zero the character is
identically one.
\end{proof}

\begin{corollary}[Comparison with the torsion cutoff on the period fields]
\leavevmode\label{cor:cubecutoff584}
On the triples of Theorem~\ref{thm:cubedepth584} the normalized
lengths of the evaluated weight quotients \(O/\mathfrak c(I)\)
on \(E_x\), summed with weight \(\log\ell\) over the odd-order
characters and \(2\le r<t=t(x)\), are
\(O_\psi(t\log^2(2t))\). The same bound holds over the indicated
odd-order auxiliary fields upon omitting \(r=2\).
For all sufficiently large \(x\), the image in \(U_{\ell,r}\)
of any \(O/\mathfrak c(I)\)-module under a \(\ZZ_\ell\)-linear
map is killed by the quotient
\(U_{\ell,r}\to R_{\ell,r}\) of Section~\ref{sec:missing}
in these finite-length cases. At \(r=2\) the odd-order
character summands are already removed by the averaging
operator of that section.
\end{corollary}
\begin{proof}
Normalized length means \(O\)-length divided by
\(v_O(\ell)\), with \(v_O\) equal to one at a uniformizer.
For \(r>2\) it is at most \(1+v_\ell(r-2)\). All odd-order
characters factor through the fixed cyclic field \(K_\psi\),
since \(E_x/K_\psi\) has two-power degree. Their number is
at most \([K_\psi:\QQ]\). A contributing prime satisfies
\(\ell-1\mid2r\). For a fixed prime the number of such \(r<t\)
is at most \(2t/(\ell-1)\). Consequently the contribution
of the constant one in the bound is at most a fixed multiple of
\[
        2t\sum_{2\le n\le2t+1}\frac{\log n}{n-1}
                         \ll t\log^2(2t).
\]
For each \(r>2\) the sum of \(v_\ell(r-2)\log\ell\)
over its contributing primes is at most \(\log(r-2)\).
At \(r=2\) the only possible prime is five and the
quotient on \(E_x\) has length zero.

Finally \(\ell\le2t+1\le P\) and \(1+\log_5 t<b_0\)
for all sufficiently large \(x\), by the definitions in
Section~\ref{sec:missing}. Every evaluated module in the
finite cases is annihilated by \(\ell^{b_0}\), so its image
is in the kernel \(U_{\ell,r}[\ell^{b_0}]\). At \(r=2\)
one has \(r\le x^\eta\) eventually and \(K_\psi\subset L_x\).
On the character summand the entire group
\(\operatorname{Gal}(E_x/L_x)\) therefore acts trivially,
so its average acts as the identity. The same holds for the
inverse character if one uses the dual indexing convention.
\end{proof}

The length here measures an equation on the ordered characters.
To deduce the full strict length bound one must still account
for the cohomology carried by representations outside the
symmetric cube.


\subsection{Quadratic induction at the rational fixed weights}
\label{sec:rationalinduction585}

Take an even-order triple of Section~\ref{sec:sourcelevels563}.
Use its rational fixed-weight trace ring \(T\) of
Theorem~\ref{thm:allroottrace548}, its coefficient DVR \(O\)
with residue field \(k\), and the universal trace \(t\).
Write \(F=H\), \(M=FD_0\), and let \(\phi\) denote the rational
character of \(D_0\). On \(G_F\) it is the nontrivial character
of \(M/F\). Define the ideal
\[
        J_\phi=\bigl(t(g):g\in G_F,\ \phi(g)=-1\bigr)\ \subset T .
\]
Theorem~\ref{thm:quadraticfamily583} gives these trace identities
at the varying-weight even-order initial point. Here we evaluate
them in the rational trace space with its prescribed weights.

\begin{proposition}[The exact intersection with the split trace]
\leavevmode\label{prop:splitinduction585}
Under the split evaluation \(T\to O\) the image of \(J_\phi\)
is \(\ell^{1+v_\ell(r)}O\). In \(T\) the ideal can also be
generated by good rational Frobenius traces at primes split
completely in \(F\) and inert in \(D_0\).
\end{proposition}
\begin{proof}
On \(G_F\) one has \(\alpha\omega_\ell^r=\phi\), by the
equality on reduction of finite prime-to-\(\ell\) characters
in Theorem~\ref{thm:eventower572}. Write
\(\langle\kappa_\ell\rangle=\kappa_\ell/\omega_\ell\).
For \(g\in G_F\setminus G_M\) the split trace is
\[
    (1+A(g))(1+\Psi(g))
       =(1+A(g))(1-\langle\kappa_\ell(g)\rangle^r).
\]
Every value has the stated divisibility. Choose \(g_0\) on
this coset with \(\kappa_\ell(g_0)=1\) and
\(\bar A(g_0)^4\ne1\), as in Lemma~\ref{lem:cyclospaces575}.
The ramification index of \(F\) above \(\ell\) divides two
and \(D_0\) splits there. Its inertia over \(F\) contains
an element \(s\in G_M\) with \(\kappa_\ell(s)=1+\ell\),
since the wild cyclotomic factor is pro-\(\ell\) and
surjective on rational wild inertia. At \(g_0s\) the first
factor is a unit and the second has valuation \(1+v_\ell(r)\)
in the normalization \(v_\ell(\ell)=1\), by binomial expansion.
For the generating assertion apply Chebotarev in every finite
Galois quotient which records membership in \(G_F,G_{D_0}\)
and the trace modulo an arbitrarily high power of the maximal
ideal. Such quotients exist by continuity on the compact group,
using an open normal subgroup on whose cosets this reduction
is constant. Conjugacy preserves the trace. Ideals of the
complete Noetherian local ring \(T\) are closed by the Krull
intersection theorem applied to the quotient. Hence the good
traces generate the ideal.
\end{proof}

\begin{theorem}[Which fixed-weight representations can have the quadratic trace identity]
\leavevmode\label{thm:rationalinduction585}
Let \(\sigma:T\to O'\) be a continuous local specialization to a
complete DVR with finite residue field, and suppose it kills
\(J_\phi\). Its coefficient field has characteristic \(\ell\).
If some trace coordinate of Theorem~\ref{thm:allroottrace548}
is nonzero at \(\sigma\), the realization is absolutely
irreducible of rank four, \(\operatorname{ord}\alpha=2\), and
after a finite coefficient extension it is induced over
\(\QQ\) from a plane over the real quadratic field cut out by
\(\tau=\bar\Psi\). This quadratic field composited with \(F\)
is \(M\). In particular, every nonzero two-plane specialization
in either characteristic has a nonzero trace at a good rational
prime split completely in \(F\) and inert in \(D_0\).
\end{theorem}
\begin{proof}
If all trace coordinates are zero, the evaluation factors
through the split trace. The proposition excludes characteristic
zero. Otherwise Theorems~\ref{thm:ramification549}
and~\ref{thm:purityblocks549} realize the trace, after a
finite coefficient extension, by an absolutely irreducible
four-dimensional representation or a sum of two absolutely
irreducible planes \(P'\oplus P\). We work over an algebraic
closure of the fraction field. All these irreducibles stay
irreducible on \(G_F\). Indeed a subspace invariant there
is invariant on \(G_{FK(\mu_{\ell^\infty})}\), and hence is
rationally Galois invariant by Lemma~\ref{lem:cyclospaces575}.

The restrictions of the semisimple realization and of its
\(\phi\)-twist have the same traces on \(G_F\). Newton's
identities applied to powers up to four give equality of
characteristic polynomials, since \(4!\) is invertible.
Brauer--Nesbitt gives an isomorphism there. In the plane
case real conjugation acts as opposite scalars on the two
planes, and both \(F\) and \(M\) are totally real. Twisting
therefore preserves each constituent separately. Let \(Z\)
denote the full irreducible representation in the first case
and the plane \(P\), with residual characters \(1,\tau\), in
the second. In both cases \(Z|_{G_F}\) is absolutely
irreducible and isomorphic to its \(\phi\)-twist.

An intertwining operator for this twist has scalar square by
Schur's lemma. Over a finite extension rescale it to an
involution. Its two eigenspaces are preserved on \(G_M\)
and exchanged on the other coset. Both occur, since otherwise
the operator would be scalar. They have equal dimension, and
inducing the first eigenspace from \(G_M\) recovers \(Z|_{G_F}\).
The first eigenspace is irreducible on \(G_M\), because a
proper nonzero subspace stable
there, together with its conjugate in the other eigenspace,
would give a proper \(G_F\)-submodule. The two constituents
on \(G_M\) are inequivalent. Indeed Frobenius reciprocity
identifies the endomorphisms of the induction with the maps
from the first constituent to the sum of the two on \(G_M\);
the endomorphism space has dimension one.

Since \(M/\QQ\) is abelian, \(G_M\) is normal in \(G_\QQ\).
The whole rational group permutes the two isotypic summands
of \(Z|_{G_M}\). It gives a quadratic character \(\eta\)
whose restriction to \(G_F\) is \(\phi\). This character is
continuous because it factors through \(\operatorname{Gal}(M/\QQ)\).
Writing \(D\) for its quadratic field, we have \(FD=M\)
since \(\eta|_{G_F}\) cuts out \(M/F\). The two summands
also express \(Z\) as an
induction over \(\QQ\) from \(D\) and yield the rational
self-twist by \(\eta\). These constructions descend to a
finite extension of the coefficient field, since the
intertwiner is found by linear equations over that field
and its eigenvalues require only a finite extension.

Reduction gives a rational residual self-twist. For the
plane its set of residual characters is \(\{1,\tau\}\);
for rank four it is \(\{1,\tau,\chi,\chi\tau\}\), with
\(\chi=\bar A\). The nontrivial quadratic character must
belong to that set since \(1\) does. In the rank-four set
neither \(\chi\) nor \(\chi\tau\) is quadratic, as an
inertia element over two in Theorem~\ref{thm:eventower572}
has \(\tau=1\) and \(\chi^4\ne1\). Thus \(\eta=\tau\)
in both cases, identifying the signs in odd characteristic.
On the rational group \(\tau\) has order \(\operatorname{ord}\alpha\)
because its two factors have coprime conductors and the Teichm\"uller
factor has order at most two, dividing the even order of
\(\alpha\). It follows that \(\operatorname{ord}\alpha=2\).

Suppose the coefficient characteristic is zero. The rational
decomposition group at \(\ell\) contains an element with
\(\eta=-1\), by the standing condition
\(\bar\Psi|_{D_\ell}\ne1\). This element exchanges the two
summands. Choose a wild inertia element with cyclotomic
value \(1+\ell\). It belongs to \(G_D\) since its image
in a group of order two is trivial. Its four eigenvalues
in the full ordinary realization are
\[
            (1+\ell)^{r+1},\quad (1+\ell)^r,\quad
                                 1+\ell,\quad 1 .
\]
They are distinct because \(r\ge2\). The same simplicity
holds on the invariant subrepresentation \(Z\). In either
case \(Z\) has a line stable under the rational decomposition
group, by intersecting the ordinary flag with \(Z\) and
taking the first nonzero step. This line is an eigenline
for the chosen inertia element. An eigenline for a simple
eigenvalue of an operator preserving both summands belongs
to one of them. The exchanging element cannot preserve it,
a contradiction.

Finally suppose the coefficient characteristic is \(\ell\)
and \(Z\) is the plane. It is induced from a character of
\(G_D\). That character has finite inertia at every place.
Above \(\ell\) its values on inertia are among those of
the fixed residual graded characters of the ordinary flag.
Away from \(\ell\) the fixed inertia kernel has finite
index, and the representation is unramified outside the
prescribed finite set. Some power of the character is
therefore everywhere unramified at finite places and factors
through the finite narrow class group by class field theory.
The original character has finite order. In characteristic
\(\ell\) its values are roots of unity algebraic over the
prime field. All traces of its induction are likewise
algebraic over the prime field. On such constants reduction
is injective, since a nonzero element has finite multiplicative
order and valuation zero. The residual trace is \(1+\tau\),
whose values are constants. Equality of traces on reduction
thus gives equality with \(1+\tau\) in the fraction field.
The degree-two Newton identities and Brauer--Nesbitt now
make \(Z\) reducible, a contradiction. This proves the
classification and, by the good-prime generating assertion,
the last conclusion.
\end{proof}

\begin{corollary}[The integral intersection in the rational trace ring]
\leavevmode\label{cor:integralinduction585}
The complete algebra \(C=T/J_\phi\) is killed by a power of
\(\ell\). It has finite \(O\)-length if
\(\operatorname{ord}\alpha>2\). For every even-order triple
its quotient obtained by imposing \(U_{pj}=W_{jk}=0\) for
all \(j,k\) has finite \(O\)-length. These are the coordinates
cutting out the two-plane locus in Theorem~\ref{thm:purityblocks549}.
\end{corollary}
\begin{proof}
Write \(\pi\) for a uniformizer of \(O\). If \(\pi\) were
not nilpotent in \(C\), a quotient by a prime avoiding it
would give a complete local domain in which \(\pi\) is
nonzero and belongs to the maximal ideal. While its dimension
is at least two, choose in its maximal ideal an element
outside all minimal primes of \((\pi)\). These primes have
height one. Mod out by a minimal prime of the chosen element.
It too has height one and cannot contain \(\pi\), as otherwise
it would be a minimal prime of \((\pi)\). The dimension
strictly decreases, and \(\pi\) survives. This yields a
one-dimensional complete local domain \(B\) with \(\pi\ne0\).
Its quotient modulo \(\pi\) is Artinian and finite over \(k\).
The \(\pi\)-adic and maximal-ideal topologies on \(B\) coincide,
since some power of the maximal ideal lies in \(\pi B\).
Lifting a basis modulo \(\pi\) and expanding successively in
powers of \(\pi\) shows that \(B\) is finite over \(O\).
The integral closure of \(O\) in its fraction field is a
complete DVR, and the resulting specialization is local
and continuous. This is a forbidden characteristic-zero
point of Theorem~\ref{thm:rationalinduction585}. Thus a power
of \(\pi\), and hence a power of \(\ell\), is zero in \(C\).

If the special fibre has positive dimension, quotient by a
prime and cut down to a one-dimensional complete local domain
of characteristic \(\ell\). One can choose a prime immediately
below the maximal ideal in a chain of maximum length. In that
domain choose a nonzero parameter \(z\) in its maximal ideal.
The same successive expansion gives finiteness over \(k[[z]]\),
which embeds because a nonzero series is a power of \(z\)
times a unit. Passing to the integral closure in the fraction
field gives an equal-characteristic complete DVR with finite
residue field. At least one trace coordinate has nonzero image
since these coordinates topologically generate the special
fibre. If \(\operatorname{ord}\alpha>2\) the theorem excludes
such a point. After imposing the displayed two-plane equations
it excludes such a point for every even order. In either
situation the quotient is zero-dimensional and Noetherian,
hence Artinian with finite residue field and finite \(O\)-length.
\end{proof}

The split evaluation determines an explicit exponent in the
first proposition. A uniform bound for the nilpotence exponent
in the full trace intersection, and control of the traces outside
that intersection, are still needed for the strict-length argument.


\subsection{Strict classes over the quadratic inducing field}
\label{sec:inducingstrict586}

In the last possible case of Theorem~\ref{thm:rationalinduction585}
let \(D\) be the rational quadratic field of \(\tau=\bar\Psi\).
It has a single place \(w\) over \(\ell\), since
\(\tau|_{D_\ell}\ne1\). Write \(\chi=\bar A\).
The character \(\chi\) is nontrivial on \(G_{D_w}\).
Indeed \(FD=M\) and Theorem~\ref{thm:eventower572} gives a
place over \(\ell\) at which \(M/F\) is split and a local
element of \(G_F\) with \(\chi^4\ne1\). With compatible
choices of places this element lies in \(G_{D_w}\).

\begin{lemma}[Vanishing of both strict groups over the inducing field]
\leavevmode\label{lem:inducingstrict586}
For \(\lambda=\chi|_{G_D}\) and for its inverse, the subgroup
of \(H^1(D,k(\lambda))\) restricting to zero on \(G_{D_w}\)
and on inertia at all places away from \(\ell\) is zero.
The assertion persists under finite residue field extension.
\end{lemma}
\begin{proof}
Let \(\lambda_0\) be the corresponding rational character
\(\chi\) or \(\chi^{-1}\). Induction from \(G_D\) gives
\(k(\lambda_0)\oplus k(\lambda_0\tau)\) as a rational
module. Indeed the permutation module of the two cosets
splits into its two sign eigenspaces because \(2\) is a unit.
Shapiro's isomorphism carries the strict conditions into the
rational strict conditions on both summands. For a rational
decomposition group the local isomorphism is a sum over its
orbits on the two cosets, with stabilizers the local groups
upstairs. The same coset description for inertia gives the
restrictions to the conjugate inertia subgroups upstairs.
Thus a class satisfying the stated vanishings has zero local
images in both rational summands. Their global strict groups
vanish by Proposition~\ref{prop:gooddata563} at residual
precision because the characters are among the reductions of
\(A^{\pm1},(\Psi/A)^{\pm1},(A\Psi)^{\pm1}\), since
\(\tau^2=1\). Extension of the finite coefficient field
commutes with these kernels by a basis expansion on cochains.
\end{proof}

\begin{proposition}[The local ordinary line is globally invariant]
\leavevmode\label{prop:inducingline586}
Let \(B\) be a complete Noetherian local coefficient ring
with finite residue field \(k'\) over \(k\). Let \(W\) be
a continuous representation of \(G_D\) on \(B^2\) whose
residual semisimplification has characters \(1,\chi|_{G_D}\).
Suppose it preserves a free direct summand of rank one at
\(w\) and its inertia away from \(\ell\) factors through
finite groups of order prime to \(\ell\). Then the local
line is preserved by \(G_D\).
\end{proposition}
\begin{proof}
Take \(h\in G_{D_w}\) with \(\chi(h)\ne1\). The characteristic
polynomial of \(W(h)\) has distinct residual roots and splits
into two linear factors over \(B\) by Hensel's lemma. Their
fiberwise distinctness gives two free rank-one eigensummands
by the polynomial projectors. The given local line is one of
these summands since the action of \(h\) on it is a root with the
specified reduction, and the complementary projector acts by
zero. Choose this summand first in the eigenbasis. Write
\(a_g,b_g,c_g,d_g\) for the entries in row order, so \(c_g=0\)
for \(g\in G_{D_w}\).

First the residual first line is globally invariant. Otherwise
the residual plane still has a global invariant line because
its composition factors have dimension one. Such a line must
be an eigenspace of \(h\), so it must be the second line.
The residual matrices are then lower triangular with distinct
diagonal characters \(\lambda_1,\lambda_2\) in the indicated
order. The lower entry divided by \(\lambda_1(g)\) is a
cocycle for \(\lambda_2/\lambda_1=\chi^{\pm1}|_{G_D}\).
It vanishes at \(w\). Its class on every other inertia group
is zero by averaging over the finite prime-to-\(\ell\) image
through which the residual matrices factor. The lemma makes
it a global coboundary. Its value at \(h\) is zero in our
diagonal basis, and \((\lambda_2/\lambda_1)(h)-1\ne0\).
The coboundary is therefore identically zero, contrary to the
first line being noninvariant.

Let \(\mathfrak m\) be the maximal ideal of \(B\), and let
\(L\) be the ideal generated by all \(c_g\).
The preceding paragraph gives \(L\subseteq\mathfrak m\).
Now the residual diagonal characters in the upper triangular
representation are again denoted by \(\lambda_1,\lambda_2\).
The identity \(c_{gg'}=c_g a_{g'}+d_g c_{g'}\) shows that
\[
          g\longmapsto \lambda_1(g)^{-1}
                         (c_g\bmod\mathfrak m L)
\]
is a cocycle for \(\lambda_2/\lambda_1\) with multiplicity
space \(L/\mathfrak m L\). This space is finite-dimensional
over \(k'\). The cocycle is continuous, since Artin--Rees
gives \(L\cap\mathfrak m^N\subseteq\mathfrak m L\) for
sufficiently large \(N\). It vanishes on the local group
at \(w\) and factors through the given finite images on
other inertia groups. Lemma~\ref{lem:inducingstrict586}
applied to each coordinate, followed again by evaluation
at \(h\), gives zero identically. This says \(L=\mathfrak m L\).
Nakayama's lemma applies to the finitely generated ideal
and gives \(L=0\). The chosen first line is globally invariant.
\end{proof}

\begin{theorem}[Finite length of the quadratic trace intersection for every even order]
\leavevmode\label{thm:allinduction586}
On every even-order triple on the interval, the rational
fixed-weight algebra \(T/J_\phi\) has finite \(O\)-length.
Every nonzero trace-coordinate specialization to a complete
DVR with finite residue field, in either characteristic,
has a nonzero good-prime trace at a rational prime split
completely in \(F\) and inert in \(D_0\).
\end{theorem}
\begin{proof}
The only outstanding DVR case in
Theorem~\ref{thm:rationalinduction585} is an absolutely
irreducible rank-four induction from a plane \(W\) over \(D\)
in equal characteristic. We check the lattice conditions in
the proposition. Work over a finite coefficient extension
on which the inducing plane is defined. Its intersection
with a stable full lattice gives a stable lattice in \(W\).
The direct sum of this lattice and its translate by an
element outside \(G_D\) is a stable induced lattice in the
four-dimensional space. Its residual characteristic
polynomials agree with those for the original lattice because
both are obtained by reducing the same integral characteristic
polynomials. On \(G_D\) the four residual constituents are
two copies each of \(1\) and \(\chi|_{G_D}\), since
\(\bar\Psi|_{G_D}=1\). The two conjugate plane lattices
have the same multiplicities for each of those characters,
as both characters extend to \(G_\QQ\). Each plane thus has
both constituents once. This uses Brauer--Nesbitt on the
residual semisimplifications.

Intersect the four-dimensional ordinary flag with \(W\)
over the fraction field. Its first nonzero intersection
is a line, since a single flag step increases the dimension
of an intersection by at most one. At the sole place \(w\)
it is locally stable. Its saturation in the plane lattice
is a direct summand over the DVR. All the other finite inertia
images of the rational representation have order prime to
\(\ell\), by the defining conditions of
Section~\ref{sec:generaleven547}. Their images on \(W\) retain
that property. Proposition~\ref{prop:inducingline586} now produces
a global invariant line in \(W\). Its induction is a proper
nonzero rational invariant subspace in rank four, contradicting
irreducibility.

It follows that \(T/J_\phi\) has no DVR specialization with a
nonzero trace coordinate. Corollary~\ref{cor:integralinduction585}
already killed a power of \(\ell\) there. A positive-dimensional
special fibre would give a DVR specialization with a nonzero
coordinate by the one-dimensional quotient argument in that
corollary. Thus the algebra is Artinian, and has finite
\(O\)-length because its residue field is finite.
Proposition~\ref{prop:splitinduction585} supplies the good
primes in the last assertion.
\end{proof}


\subsection{The exact exponent of the quadratic trace intersection}
\label{sec:quadraticexponent587}

Use the even-order data \(T,F,M,\phi,J_\phi\) of
Section~\ref{sec:rationalinduction585}. In particular \(T\)
is the ring of rational traces at the prescribed weights.
For a place \(v\) of \(F\) above \(\ell\), fix a compatible
rational decomposition group and denote the four characters
on the graded lines of its ordinary flag by
\(\theta_0,\theta_1,\theta_2,\theta_3\).

\begin{theorem}[Paired local characters and the coefficient ideal]
\leavevmode\label{thm:quadraticexponent587}
In \(C=T/J_\phi\) the four characters on the whole local group
satisfy
\[
      \theta_0|_{G_{F_v}}=\theta_2|_{G_{F_v}},\qquad
      \theta_1|_{G_{F_v}}=\theta_3|_{G_{F_v}} .
\]
The kernel of the coefficient map \(O\to C\) is exactly
\[
                       J_\phi\cap O=\ell^{1+v_\ell(r)}O .
\]
In particular the least power of \(\ell\) annihilating the
full ring, including its nilpotents, equals the power detected
by the split trace.
\end{theorem}
\begin{proof}
Let \(\rho\) denote the universal framed representation of
Section~\ref{sec:generaleven547}. By construction \(T\)
is the closed subalgebra of its coefficient ring generated
by the traces. Choose \(h\) in the rational decomposition
group at \(\ell\) separating all four residual characters,
as in Lemma~\ref{lem:prescribedaux563}. The characteristic
polynomial of \(\rho(h)\) has coefficients in \(T\), by the
Newton identities since \(\ell\ge5\). Its roots \(a_i\)
with the prescribed distinct residues are in \(T^\times\)
by Hensel's lemma. The operators
\[
         e_i=\prod_{j\ne i}
                   \frac{\rho(h)-a_j}{a_i-a_j}
\]
project onto its four eigenlines. In an eigenbasis in the
given order, the ordinary flag is the coordinate flag.
Indeed the \(e_i\) preserve it, and the summands of each
step are determined by its reduction. For every rational
Galois element \(g\) the diagonal entry
\[
                  z_i(g)=\operatorname{tr}(e_i\rho(g))
\]
belongs to \(T\), as the right side is a \(T\)-linear
combination of \(t(h^j g)\). It is a unit since the split
residual representation has nonzero diagonal characters.
On the rational decomposition group these entries are
the characters \(\theta_i\).

Fix any \(g\in G_F\setminus G_M\). Conjugating \(\rho(g)\)
by \(\rho(h)\) multiplies its entry \(ij\) by \(a_i/a_j\).
The polynomial
\[
      P(X)=\prod_{i\ne j}\frac{X-a_i/a_j}{1-a_i/a_j}
                      \ \in T[X]
\]
therefore projects this matrix onto its diagonal when evaluated
at the conjugation operator. Every denominator is a unit.
For any \(y\in G_{F_v}\) this gives
\[
              \sum_{i=0}^3 z_i(g)\theta_i(y)\ \in J_\phi .
\]
Indeed the sum is the trace of the diagonal matrix times
\(\rho(y)\). Expansion of \(P\) writes that trace as a
linear combination of \(t(h^a g h^{-a}y)\) for \(0\le a\le12\).
Both \(F\) and \(M\) are Galois over \(\QQ\), and \(\phi\)
is a rational character. Each conjugate of \(g\) thus belongs
to \(G_F\setminus G_M\), even if \(h\) does not fix \(F\).
The place \(v\) splits in \(M/F\), so \(y\in G_M\) and all
those arguments stay in the required coset.

Choose \(u\in G_{F_v}\) with \(\chi(u)^4\ne1\) by
Theorem~\ref{thm:eventower572}. The residual values of the
four \(\theta_i(u)\) are \(\chi(u),1,\chi(u),1\). Consider
the two monic polynomials
\[
       f_{02}(X)=(X-\theta_0(u))(X-\theta_2(u)),\qquad
       f_{13}(X)=(X-\theta_1(u))(X-\theta_3(u)).
\]
Their resultant is a unit because their reductions have disjoint
sets of roots. The Chinese remainder theorem gives a polynomial
\(q\in T[X]\), of degree at most three after division by
\(f_{02}f_{13}\), congruent to one modulo \(f_{02}\) and to zero modulo
\(f_{13}\). In the preceding sum replace \(y\) successively by
\(u^j y\), and combine the results with the coefficients
of \(q\). Since diagonal characters are multiplicative this
keeps precisely the terms with indices \(0,2\). Using \(1-q\)
keeps the terms with indices \(1,3\). Thus in \(C\) one has
\[
     z_0(g)\theta_0(y)+z_2(g)\theta_2(y)=0,\qquad
     z_1(g)\theta_1(y)+z_3(g)\theta_3(y)=0 .
\]
At \(y=1\) the two coefficients in each pair are negatives
of one another. They are units. Substitution yields the two
claimed character identities.

By Theorem~\ref{thm:eventower572} the local ramification index
of \(F/\QQ\) divides two. The rational wild inertia acts
trivially on this local extension, and its cyclotomic image
contains \(1+\ell\). Choose \(s\) in it with exactly this
cyclotomic value. At \(s\) the rational graded characters
are the powers of \(1+\ell\) with exponents \(r+1,r,1,0\),
respectively. The first character identity shows that
\((1+\ell)^r-1\) belongs to \(J_\phi\), giving the inclusion
\(\ell^{1+v_\ell(r)}O\subseteq J_\phi\cap O\).
Conversely split evaluation fixes every scalar and sends
\(J_\phi\) into this same ideal by
Proposition~\ref{prop:splitinduction585}. This proves equality.
\end{proof}

\begin{corollary}[The length supported by this exponent on the original fields]
\leavevmode\label{cor:quadraticcutoff587}
At \(x\) on the common interval write \(t=t(x)\), and use
the common upper bounds \(n_*,N_*\) for the degree and
conductor in Theorem~\ref{thm:levelstrict345}. For each prime
and weight in the even-order triples take the sum of any
collection of images in \(U_{\ell,r}\) of modules over
\(T/J_\phi\), via \(\ZZ_\ell\)-linear maps. Different
characters at the same \((\ell,r)\) may all occur in the
collection. Write \(B_{\ell,r}\) for this sum. Then on
the actual period fields \(E_x\),
\[
 \sum_{\ell>2}(\log\ell)\sum_{2\le r<t}
         \operatorname{length}_{\ZZ_\ell} B_{\ell,r}
       \ll_\psi n_* \log(2N_*)\,t\log^2(2t).
\]
For sufficiently large \(x\) every \(B_{\ell,r}\) has zero
image in \(R_{\ell,r}\) of Section~\ref{sec:missing}.
\end{corollary}
\begin{proof}
The theorem annihilates the entire sum at fixed \((\ell,r)\)
by \(\ell^{h_r}\), where \(h_r=1+v_\ell(r)\) is independent
of the character and its coefficient field. It is contained
in \(M_{\ell,r}(E_x)[\ell^{h_r}]\). The Kummer bound used in
Theorem~\ref{thm:levelstrict345} gives
\(\dim_{\mathbf F_\ell}M_{\ell,r}(E_x)[\ell]\ll_\psi
n_*\log(2N_*)\). Each cyclic summand of this finite module
contributes at most \(h_r\) to its truncated length.
For the constant one in \(h_r\), each prime occurs at most
\(2t/(\ell-1)\) times since \(\ell-1\mid2r\).
Summing \(\log n/(n-1)\) for \(2\le n\le2t+1\) gives
\(O(\log^2(2t))\). For each \(r\) the other terms contribute
at most \(\sum_{\ell\mid r}v_\ell(r)\log\ell=\log r\).
These estimates give the display without a multiplicity
factor for characters. Finally \(\ell\le2t+1\le P\) and
\(h_r\le1+\log_5 t<b_0\) for all sufficiently large \(x\).
The defining quotient by \(U_{\ell,r}[\ell^{b_0}]\) kills
every stated image.
\end{proof}

The corollary places every image from this quadratic trace
quotient within the existing cohomology allowance on the
same interval. The remaining strict classes would require
an argument controlling the other trace components or a
construction placing those classes in the stated images.


\subsection{Unordered symmetric-cube spectra on the period fields}
\label{sec:unorderedcube588}

The equations of Theorem~\ref{thm:cubedepth584} use the
order of the graded lines. A condition on characteristic
polynomials alone also gives a uniform integral bound. It
applies to the rational fixed-weight trace ring of each triple
of Section~\ref{sec:sourcelevels563} after restriction to
the original period field \(E_x\).

\begin{lemma}[An integral equation in the characteristic coefficients]
\leavevmode\label{lem:unorderedcube588}
For four roots define
\[
 {\cal P}(X_1,X_2,X_3,X_4)=
 \prod_{j=1}^4\ 
 \prod_{\substack{1\le i<k\le4\\i,k\ne j}}
                          (X_iX_k-X_j^2).
\]
It is an integer polynomial in their four elementary symmetric
functions. Its value on the characteristic coefficients of any
scalar multiple of a symmetric cube of a two-dimensional
matrix is zero over every commutative ring.
\end{lemma}
\begin{proof}
Permutations preserve both the distinguished index \(j\) and
the unordered pair of endpoints as types of indices, hence
permute the twelve factors. The fundamental theorem on
symmetric polynomials over \(\ZZ\) gives the first assertion.
For a generic two-by-two matrix with independent indeterminate
entries, work in a splitting field of its characteristic
polynomial over their rational function field. If the two
roots are \(a,b\), the four roots after the indicated
operations are \(z a^3,z a^2b,z ab^2,z b^3\), with \(z\)
also indeterminate. The product of the first and third is
the square of the second. Thus the value is zero in the
function field. The value in terms of characteristic
coefficients was a polynomial over \(\ZZ\) in the entries
and \(z\), so it was zero in the polynomial ring itself.
Specialization proves the identity in all characteristics
and in rings with nilpotents.
\end{proof}

For \(g\in G_{E_x}\) evaluate \({\cal P}\) on the coefficients
of the universal characteristic polynomial of \(g\), using
the elementary symmetric functions with their usual alternating
signs in that polynomial. Denote the value by \(p(g)\), and
let \(J_{\rm sc}(E_x)=(p(g):g\in G_{E_x})\subset T\).
These coefficients belong to \(T\) by the Newton identities.
In particular any specialization whose restricted characteristic
polynomials are those of a scalar-twisted symmetric cube kills
this ideal, regardless of any identification of flags. The ideal
also makes sense for four-dimensional determinants specified
only by their characteristic polynomials.

\begin{theorem}[An integral obstruction from one inertia element]
\leavevmode\label{thm:unorderedcube588}
For \(r>2\) put
\[
    D(r)=(r-2)(r-1)(r+1)(r+2)(2r-1)(2r+1),\qquad
    b_{\ell,r}=12+2v_\ell(D(r)).
\]
For \(r=2\) the only possible coefficient prime is five, and
put \(b_{5,2}=6\). At every \(v\mid\ell\) of \(E_x\) there is
a single \(s\in I_{E_{x,v}}\) for which \(p(s)\) is a scalar
generating \(\ell^{b_{\ell,r}}O\). Thus the quotient
\(T/J_{\rm sc}(E_x)\) is annihilated by that power, including
at every nonreduced specialization. If \(r>2\), the same
assertion holds on restricting further to any auxiliary field
whose local ramification index over \(E_{x,v}\) is prime to
\(\ell\).
\end{theorem}
\begin{proof}
The rational local ordinary flag gives the four prescribed
characters on inertia, hence determines the whole
characteristic polynomial there even when the representation
is a nonsplit extension. The period field is unramified
at \(\ell\), and both \(\alpha,\xi\) are unramified there.
Put \(u=1+\ell\). For \(r>2\) choose wild inertia with
cyclotomic value \(u\). To see this also after the indicated
local extension, observe that the image in the pro-\(\ell\)
cyclotomic factor of an inertia subgroup of index prime to
\(\ell\) is the whole factor. Its wild inertia still
surjects since the tame quotient has order prime to \(\ell\)
in every finite quotient. At this element the four roots are
\(u^e\) for
\[
                            e\in\{0,1,r,r+1\}.
\]
For a factor with middle exponent \(e_j\) the valuation is
\(1+v_\ell(|e_i+e_k-2e_j|)\), by the binomial valuation in
Theorem~\ref{thm:cubedepth584}. For middle exponent zero the
three differences are \(r+1,r+2,2r+1\). For middle exponent
one they are \(r-2,r-1,2r-1\). Reflection
\(e\mapsto r+1-e\) gives their negatives for the other two
middle exponents. All are nonzero, so adding the twelve
valuations gives the stated integer.

If \(r=2\), then \(\ell\ge5\) and \(\ell-1\mid4\) force
\(\ell=5\). The full cyclotomic image on inertia of \(E_x\)
is \(\ZZ_5^\times\). Choose \(s\) with
\(\langle\kappa_5(s)\rangle=6\) and \(\omega_5(s)\) of
order four. The four values, in increasing order of the
exponents of \(u=6\), are
\[
                              1,\quad u,\quad -u^2,\quad -u^3.
\]
For endpoints of opposite signs the product differs from
the middle square already modulo five. For the two positive
endpoints the absolute exponent differences with the two
possible middle terms are \(3,5\), giving valuations \(1,2\).
For the two negative endpoints the differences are \(5,3\),
giving \(2,1\). The sum is six. This computation uses the
tame part of rational inertia on the unramified period field.
\end{proof}

\begin{corollary}[Uniform comparison with the strict length cutoff]
\leavevmode\label{cor:unorderedcutoff588}
Take the spans of images in \(U_{\ell,r}\) on \(E_x\) of
any collection of \(\ZZ_\ell\)-linear maps from modules
over the quotients \(T/J_{\rm sc}(E_x)\) on the stated
triples, allowing both orders of \(\alpha\). Denote the
span at a given \((\ell,r)\) by \(B'_{\ell,r}\). Uniformly
on the common intervals,
\[
    \sum_{\ell>2}(\log\ell)\sum_{2\le r<t}
          \operatorname{length}_{\ZZ_\ell}B'_{\ell,r}
       \ll_\psi n_*\log(2N_*)\,t\log^2(2t).
\]
These spans map to zero in \(R_{\ell,r}\) for all sufficiently
large \(x\).
\end{corollary}
\begin{proof}
By the theorem the span is contained in
\(M_{\ell,r}(E_x)[\ell^{b_{\ell,r}}]\). Use the
\(\ell\)-torsion rank estimate of
Corollary~\ref{cor:quadraticcutoff587}. The constant twelve
in \(b_{\ell,r}\) contributes \(O(t\log^2(2t))\) to the
sum of exponents weighted by \(\log\ell\), since
\(\ell-1\mid2r\). For \(r>2\) prime factorization gives
\(\sum_{\ell\mid D(r)}v_\ell(D(r))\log\ell=\log D(r)\),
which is \(O(\log(2r))\) as this is a fixed polynomial of
degree six. The exceptional pair contributes at most
\(6\log5\). This proves the display. Also
\(b_{\ell,r}=O(\log(2t))\) uniformly, so
\(b_{\ell,r}<b_0\) and \(\ell\le2t+1\le P\) eventually.
The quotient defining \(R_{\ell,r}\) kills the entire span.
\end{proof}

In particular a characteristic-polynomial congruence with
a cube cannot account for higher precision on these fields.
The comparison with all of \(U_{\ell,r}\) still requires
control beyond the trace quotients treated here.


\subsection{Polarized ordinary deformations at the two cuspidal points}
\label{sec:polarweights593}

Use the even-order and odd-order initial points of
Section~\ref{sec:seedweights580}. Write \(L=FK\), \(d=[F:\QQ]\),
and \(U=b^{-1}R\) with multiplier \(\nu=b^{-2}\mu\).
The coefficient field \(E\) in this section is a finite
extension of \(\QQ_\ell\). Write \(J\) for the fixed alternating
form. All places of \(F\) above \(\ell\) split in \(L\).
Choose one place \(w\) above each such \(v\), so that
\(G_{L_w}=G_{F_v}\). On this group use the full ordinary flag
of \(U\).

We describe polarized deformations with fixed polarizing
character. Put \(\iota(A)=J^{-1}(A^{-1})^{\mathsf T}J\) on
\(\mathrm{GL}_4\), let \(e(g)\) be zero on \(G_L\) and one on its
other coset, and define an action of \(G_F\) by
\[
             g*A=U(g)\,\iota^{\,e(g)}(A)\,U(g)^{-1}.
\]
Indeed \(\iota^2=1\), and \(\iota(U(g))\) is a scalar
multiple of \(U(g)\), so the two operations commute.
For a local Artin \(E\)-algebra \(A\) with residue field \(E\),
write a deformation as a continuous cocycle \(z\) for this
action with values in \(1+\operatorname{Mat}_4(\mathfrak m_A)\).
Thus \(z(gh)=z(g)(g*z(h))\). We identify cocycles by strict
changes of basis \(z(g)\mapsto Pz(g)(g*P^{-1})\), with
\(P\) reducing to the identity. On \(G_L\) the deformed
representation is \(z(g)U(g)\). This is the usual fixed
polarization condition in the group with outer automorphism
given by inverse transpose. At a real conjugation \(c\), the
pairing for \(U|_{G_L}\) is the symmetric form \(JU(c)\) of
Proposition~\ref{prop:seedadjoints580}.

Fix a finite set of places containing those above \(\ell\),
the archimedean places, and the ramification of \(U,L/F\).
Require unramifiedness outside the set. At finite places
in the set away from \(\ell\), require \(z\) to be a strict
coboundary on inertia. At each chosen \(w\mid\ell\) include
an invariant full flag of direct summands lifting the indicated
one, and transport flags under changes of basis. Let
\(\epsilon_{v,i}\), \(0\le i\le3\), be the ratio of the
deformed graded character to its original value on local
units, read by local reciprocity.

\begin{theorem}[The polarized weight space and the symplectic equations]
\leavevmode\label{thm:polarweights593}
The ratios \(\epsilon_{v,i}\) identify the polarized ordinary
deformation functor with all tuples of continuous characters
of local units into \(1+\mathfrak m_A\), independently in
the four positions and at every \(v\mid\ell\). Its completed
ring over \(E\) is a power series ring in \(4d\) variables.
The deformations obtained by restricting fixed-multiplier
symplectic deformations of \(U\) over \(F\) form the smooth
closed subspace of codimension \(2d\) cut out by
\[
       \epsilon_{v,0}\epsilon_{v,3}=1,\qquad
       \epsilon_{v,1}\epsilon_{v,2}=1 \qquad (v\mid\ell)
\]
on all the local units. On this subspace the symplectic
deformation over \(F\) is unique up to strict symplectic
conjugacy.
\end{theorem}
\begin{proof}
Let \(V\) be \(\operatorname{Mat}_4(E)\) with the derivative
of the displayed action. On \(G_L\) this is conjugation.
At a real conjugation \(c\) its action is
\(X\mapsto-(JU(c))^{-1}X^{\mathsf T}(JU(c))\), since
\(U(c)^2=1\). We first prove
\[
          H^1_f(F,V)=H^1_f(F,V^\vee(1))=0 .
\]
The first group is the full polarized adjoint Selmer group.
Over \(L\) the representation \(U\) is attached to a regular
algebraic polarized cuspidal representation and remains
absolutely irreducible over \(L(\mu_{\ell^\infty})\), by
Theorems~\ref{thm:cuspidalinitial577} and~\ref{thm:regularind579}.
The identification of its polarization is given in
Proposition~\ref{prop:seedadjoints580}. Thorne's adjoint
vanishing theorem therefore applies directly to \(V\).
\href{https://arxiv.org/pdf/2207.04925}{J.~A. Thorne,
\emph{On the vanishing of adjoint Bloch--Kato Selmer groups
of irreducible automorphic Galois representations}, Theorem~6.1}.

For the dual group use the global Selmer dimension formula.
Only scalars commute with \(G_L\), and a real conjugation
acts by negation on scalars. Thus \(H^0(F,V)=0\). The trace
pairing is invariant, also under negative transpose, so
\(V^\vee\simeq V\). An invariant in \(V(1)\) would restrict
to a nonzero intertwiner from \(U\) to \(U(1)\) on \(G_L\).
Absolute irreducibility and determinants would give
\(\kappa_\ell^4=1\) there, which is impossible. The finite
local condition at \(v\nmid\ell\) has dimension \(h^0(F_v,V)\)
by the unramified cohomology formula. Above \(\ell\), \(V\)
is crystalline and has six positive weights per embedding,
six negative ones and four zero weights. The difference
\(h^1_f-h^0\) is \(6[F_v:\QQ_\ell]\) by the crystalline
dimension formula. The fixed Lie algebra at each real place
is skew-adjoint for a nondegenerate symmetric form of rank
four, of dimension six. The total difference of the two
Selmer dimensions is consequently \(6d-6d=0\).
\cite{XII-IX-V-milne98,XII-IX-VI-bk117}

At \(v\mid\ell\) let \(\mathfrak b\) be the flag-preserving
subalgebra, \(\mathfrak n\) its strictly upper triangular
ideal, and \(\mathfrak t=\mathfrak b/\mathfrak n=E^4\) the
trivial diagonal quotient. The six graded characters in
\(\mathfrak n\) have weights at least two at every embedding.
Those in \(V/\mathfrak b\) are their negatives. Hence
\(H^0(F_v,V/\mathfrak b)=0\), and local duality gives
\(H^2(F_v,\mathfrak n)=H^2(F_v,\mathfrak b)=0\), since
the respective twisted duals have no weight-zero constituents.
The ordinary tangent classes inject into \(H^1(F_v,V)\)
from \(H^1(F_v,\mathfrak b)\), and they surject onto the
four additive characters of local units. The latter
assertion uses the exact sequence with quotient \(\mathfrak t\)
and extension of an additive unit character to \(F_v^\times\)
by choosing its value on a uniformizer.

The kernel has dimension
\(h^0(F_v,\mathfrak b)+(10-4)[F_v:\QQ_\ell]\) by the local
Euler characteristic formula. It contains \(H^1_f(F_v,\mathfrak b)\),
since the finite condition on each diagonal is unramified.
The crystalline dimension formula gives this same dimension
for the finite subspace, because \(\mathfrak b\) has six
positive weights. As \(h^0(F_v,\mathfrak b)=h^0(F_v,V)\),
its image is all of \(H^1_f(F_v,V)\) by that formula again.
Thus the ordinary tangent classes at fixed inertial weight
are exactly the finite local condition. For a small extension
with square-zero kernel annihilated by the maximal ideal,
lift the four diagonal characters using the prescribed units
and arbitrary lifts at uniformizers. In a basis of the lifted
flag, the multiplication defect of upper triangular lifts
lies in a two-cocycle with values in \(\mathfrak n\) tensored
with the kernel. The vanishing of its second cohomology
supplies a correction. The choices of cochains can be made
continuous in unipotent coordinates. The vanishing of
\(H^0(F_v,V/\mathfrak b)\) gives uniqueness of the lifted flag
for a fixed representation.

Away from \(\ell\), undo a strict conjugation to make the
inertial cocycle trivial. The permissible Frobenius correction
is then fixed by the inertia action. The original inertia has
finite image in the action on the Lie algebra. In characteristic
zero the finite logarithm identifies invariant strict matrices
with invariant nilpotent Lie algebra elements, which lift by
averaging over that finite image. A Frobenius correction and
the original inertia define a continuous local lift. To see
continuity, represent the semidirect product with the outer
involution linearly, using the defining module and its dual.
The original image is compact, and the new Frobenius differs
from the old one by a strict matrix. Its eigenvalues on the
successive powers of \(\mathfrak m_A\) are the original ones,
hence units. Its integral powers have compact closure by the
binomial formula on Jordan blocks. The same argument on
inertia and the unramified quotient extends the lift from
the Weil group to the decomposition group. The tangent
condition is unramified cohomology. At the real places
positive-degree cohomology vanishes by averaging over order two.

For completeness, the obstruction to a global lift lies in
\(H^2(G_{F,S},V)\) tensored with the small kernel. It is the
multiplication defect of any continuous lifts of the cocycle.
The local lifts just constructed show it is locally zero.
Poitou--Tate duality identifies the kernel of localization in
degree two with the dual of the everywhere locally zero group
in degree one for \(V^\vee(1)\), which lies in the vanishing
dual finite Selmer group. A global lift therefore exists.
Its discrepancies from the chosen local conditions are measured
modulo the finite conditions in degree one. The localization map
onto the sum of those quotients is surjective by the same duality
sequence and the dual Selmer vanishing. Correcting the lift by
a global class imposes all the conditions with the prescribed
weights. Two corrected lifts have difference in \(H^1_f(F,V)\)
tensored with the kernel and hence are strictly equivalent.
Induction over small extensions proves the bijection. Logarithms
of the four character ratios give four free parameters per
\(\ZZ_\ell\)-generator of the local units modulo torsion and
their prime-to-\(\ell\) part. Their total number is \(4d\).

For a symplectic deformation put \(z(g)=U_A(g)U(g)^{-1}\).
It takes values in the symplectic group, on which \(\iota\)
acts trivially, hence gives precisely a polarized cocycle
of the stated kind. Self-duality of its flags gives the two
displayed unit-character equations. Conversely, given a
polarized deformation satisfying these equations, use its
first two unit characters at every place. Twisting by the
fixed \(b\), Theorem~\ref{thm:seedweight580} gives a unique
symplectic deformation with these two characters, and twisting
back gives the required multiplier \(\nu\). Its four polarized
unit characters agree with the four given ones. The bijection
already proved identifies the two polarized deformations.
In logarithmic coordinates the pairwise-sum equations eliminate
two distinct entries at each local generator, proving the
codimension statement.
\end{proof}

\begin{corollary}[Local parameters from rational primes split on the interval]
\leavevmode\label{cor:polarsplit593}
Choose the minimal odd-order fields \(H_1=HC_{w,\ell}\)
in Theorem~\ref{thm:cuspidalinitial577}. For each initial
point there are \(4d\) distinct good rational primes, all
split in every \(E_x\) on its original interval, and choices
of primes of \(L\) above them, such that their Frobenius
traces minus their values at the point form formal coordinates
for the polarized ring. Among these one may select \(2d\)
giving coordinates on the symplectic subspace. The rational
primes may also be required to split in every auxiliary
starting field of the interval and to avoid any fixed finite set.
\end{corollary}
\begin{proof}
Adjoin to the abelian field \(\mathcal F\) of
Proposition~\ref{prop:commonaux573} the fields \(C_{w,\ell}\)
for the finitely many odd-order coefficient primes on the
interval. The resulting finite Galois field \(P/\QQ\) contains
all the indicated \(L\) and \(E_x\). For any chosen \(w\mid\ell\)
the image of \(G_P\cap I_{L_w}\) in the local unit group
under reciprocity has finite index. This follows by taking
the image of a finite-index closed subgroup of inertia.
On a sufficiently small open subgroup of these units the
original four characters are distinct powers of the
cyclotomic character, since their algebraic weight gaps
are positive. Choose \(a\in G_P\cap I_{L_w}\) with four
distinct original values. Such an element exists since the
cyclotomic image of any finite-index inertia subgroup in
\(G_{L_w}\) contains an open subgroup of principal units.

Let \(\mathcal T\) be the completed polarized ring, and
center each trace at its constant value. Modulo the square
of its maximal ideal the eigenvalues \(a_i\) of the
universal matrix at \(a\) are recovered from its first
four power traces. Indeed Newton's identities give the
characteristic polynomial and its four distinct residual
roots lift uniquely. For \(u\in G_P\cap I_{L_w}\) the
four traces at \(a^j u\), \(0\le j\le3\), recover all
four graded values at \(u\) by inversion of the matrix
\((a_i^j)_{j,i}\), since the representation is triangular
in its flag. Hence the linear parts of the traces on \(G_P\)
span those of the four unit ratios on this open subgroup.
Characters into an \(E\)-vector space are determined on
such a subgroup, which spans the same local unit space
over \(E\). Applying this at every \(v\mid\ell\) and using
the theorem proves that the trace linear parts span the
whole \(4d\)-dimensional cotangent space.

Every \(g\in G_P\) can be approximated by Frobenius elements
of primes of \(L\) whose underlying rational primes split
completely in \(P\). In fact in each finite Galois quotient
over \(\QQ\) refining \(P\), Chebotarev gives a good rational
prime with Frobenius conjugate to the image of \(g\). Choosing
a prime above it realizes that image itself. It restricts
trivially to \(P\), so the residue degree in \(L\) is one.
Trace and its linear part are continuous in the finite-dimensional
coefficient quotient. An \(E\)-subspace of the cotangent space
is closed, so the good split prime traces also span it. Excluding
finitely many primes does not alter the argument. Successively
choose distinct primes giving a basis. In a power series ring
a tuple with invertible linear part is a system of coordinates,
by solving degree by degree for its formal inverse. Its images
span the cotangent space of the symplectic quotient, where the
same reasoning applies to any basis selected from them.
\end{proof}

These coordinates describe complete local rings over \(E\)
at the two automorphic starting points. To reach the strict
groups in Section~\ref{sec:missing}, support must also be
propagated to the relevant integral components at the
prescribed rational weights.

\section{Hecke fibers and local Selmer conditions}\label{nav:chapter-ordinary-hecke-fibers}

\subsection{Common Frobenius primes and ramification at finite precision}
\label{sec:tameprimes594}

Use the interval and triples of Section~\ref{sec:sourcelevels563}.
Let \(\mathcal L\) be the finite set of their coefficient primes.
For the odd-order case take \(F=HC_{w,\ell}\) and in the even-order
case take \(F=H\), as in Section~\ref{sec:polarweights593}.
Write \(U=b^{-1}R\) for the cuspidal point over \(F\), \(L=FK\),
and \(O\) for a coefficient valuation ring, with uniformizer
\(\varpi\) and fraction field \(E\). Use the stable lattice with
the perfect alternating form of the construction. Let \(T\) be
its endomorphism lattice with the polarized adjoint action over
\(F\) from Section~\ref{sec:polarweights593}. In particular
\(T\) has rank sixteen. The statements about its local cohomology
below hold also for the fixed-multiplier symplectic Lie lattice.

\begin{proposition}[Ramified local classes to every precision]
\leavevmode\label{prop:tamerigid594}
At every finite \(v\nmid\ell\) of \(F\), let \(q_v\) be the
residue cardinality, \(I_v\) inertia and \(\Phi_v\) a lift of
arithmetic Frobenius. The operator \(q_v-\Phi_v\) on \(T^{I_v}\)
is invertible over \(E\). Write \(\varpi^{a_i}\) for its elementary
divisors over \(O\), up to units. For every \(n\ge1\),
\[
 \operatorname{length}_O
 \frac{H^1(F_v,T/\varpi^nT)}{H^1_{\rm unr}(F_v,T/\varpi^nT)}
       =\sum_i\min(n,a_i).
\]
At either cuspidal point, allowing unrestricted lifts away from
\(\ell\) at any fixed finite set of places gives exactly the same
ordinary deformation functor on local Artin \(E\)-algebras as in
Theorem~\ref{thm:polarweights593}.
\end{proposition}
\begin{proof}
The inertia image on \(T\) has order prime to \(\ell\).
Indeed \(L/F\) is unramified at finite places. In the odd-order
construction \(C_{w,\ell}\) kills the auxiliary character
\(\theta\), and inertia on the underlying plane away from
\(\ell\) has characters \(1,\xi\) and zero monodromy.
In the even-order construction the inducing extension
\(FD_0/F\) is unramified, and inertia on each plane has
characters \(1,\xi\) restricted to that inertia, with zero
monodromy. These are the conductor calculations in
Theorems~\ref{thm:cuspidalinitial577} and~\ref{thm:regularind579}.
The order of \(\xi\) is prime to \(\ell\) in both cases, and
a scalar twist has no effect on conjugation. Thus averaging
over the inertia image identifies \((T/\varpi^nT)^{I_v}\)
with \(T^{I_v}/\varpi^nT^{I_v}\).

Put \(A=T/\varpi^nT\) and let \(I_0\) be the kernel of the
action of \(I_v\) on \(T\). Restriction identifies \(H^1(I_v,A)\)
with \(\operatorname{Hom}_{\rm cts}(I_0,A)^{I_v/I_0}\) by
inflation--restriction and averaging. Wild inertia has residue
characteristic different from \(\ell\). Its intersection with
\(I_0\) has no homomorphism to \(A\), and the image of \(I_0\)
in tame inertia is an open subgroup of prime-to-\(\ell\) index.
Projection to the \(\ell\)-factor of tame inertia consequently
identifies its maximal pro-\(\ell\) quotient with \(\ZZ_\ell(1)\).
Inertia conjugation on that quotient is trivial and arithmetic
Frobenius acts by \(q_v\). The unramified Galois group has
\(\ell\)-cohomological dimension one. Inflation--restriction
now gives
\[
 H^1(F_v,A)/H^1_{\rm unr}(F_v,A)
   \simeq \ker\bigl(\Phi_v-q_v:A^{I_v}\longrightarrow A^{I_v}\bigr),
\]
after a choice of generator of the tame \(\ell\)-factor.

To see that the elementary divisors are finite, restrict first
to a place of \(L\) over \(v\), where the action on endomorphisms
is conjugation. The regular algebraic polarized cuspidal
representation giving \(U|_{G_L}\) has pure local
Weil--Deligne parameter by Barnet-Lamb, Gee, Geraghty and
Taylor, Theorem~2.1.1(2). Its nilpotent monodromy is zero
by the finite inertia just proved. Thus all Frobenius
eigenvalues have the same complex absolute value. Their
ratios have absolute value one. As \(L/F\) is unramified,
this applies to the action of a power of \(\Phi_v\) of
exponent one or two. It follows also for the eigenvalues
of \(\Phi_v\) on \(T^{I_v}\). Since \(q_v>1\), the operator
in the proposition has nonzero determinant. Reduction of its
Smith normal form modulo \(\varpi^n\) gives the length formula.
The same argument applies on the stable symplectic Lie
subspace. \cite{XII-IX-blggt560}

Over \(E\) the same inertia calculation gives
\(H^1(F_v,T\otimes E)=H^1_{\rm unr}(F_v,T\otimes E)\).
Indeed continuous homomorphisms from compact inertia have
bounded image, and after scaling an integral lattice the same
tame projection computes them. For the assertion about lifts
use induction along small extensions of Artin \(E\)-algebras.
Once the representation over the quotient has fixed inertia
up to strict conjugacy, it has a fixed-inertia lift locally,
as shown in Theorem~\ref{thm:polarweights593}. The difference
from any given unrestricted local lift is a class in
\(H^1(F_v,T\otimes E)\) tensored with the small kernel. Its
restriction to inertia is a coboundary and can be undone by
a strict change of basis. This proves the assertion at each
place and hence for the global functor with those places allowed.
\end{proof}

\begin{theorem}[Split primes preserving strict cohomology and detecting weights]
\leavevmode\label{thm:commonprobes594}
There is one reduced rational residue class on the interval
with the following properties. Every prime \(s\) in it splits
completely in every \(E_x\), in every \(L\) just specified and
in \(D_0\) when it is used, and satisfies \(s\equiv4\bmod\ell\) for every
\(\ell\in\mathcal L\). Finitely many primes can be omitted.
\begin{enumerate}
\item For each indicated \((x,\ell,r)\) and each finite set of
these good primes, allowing ramification above the set leaves
\(\mathcal S_{\ell,r}(E_x)\) and its definition with coefficients
\(\ZZ/\ell^n(-r)\) unchanged, for every \(n\ge1\).
At all these places the ramified quotient of the preceding
proposition is likewise zero to every precision. Every Artin
local deformation of the residual point with finite residue
field of characteristic \(\ell\) is unramified there.
\item At either cuspidal point over \(F\), \(4[F:\QQ]\) distinct
good primes in the class, with choices of places above them,
give Frobenius trace coordinates for its full polarized
ordinary ring over \(E\). Among them \(2[F:\QQ]\) give
coordinates on its symplectic subspace. A single finite list
of primes contains such a choice for every indicated triple
on the interval.
\end{enumerate}
\end{theorem}
\begin{proof}
Retain \(M_*,N^\natural,p_1,p,p',p_3,w\) of
Proposition~\ref{prop:commonaux573} and
Lemma~\ref{lem:cuspidalplane577}. Take \(D\) to be the least
common multiple of \(8N^\natural,M_*,p_1,p,p',p_3,w\),
omitting auxiliary factors if their cases are absent.
Every member of \(\mathcal L\) is at least five and is prime
to \(D\) by the exclusions in Section~\ref{sec:sourcelevels563}
and the progressions defining the auxiliary primes. Prescribe
one modulo \(D\) and four modulo every \(\ell\in\mathcal L\).
This is a reduced class modulo
\(D'=D\prod_{\ell\in\mathcal L}\ell\). Write \(M=\QQ(\mu_{D'})\)
and \(\sigma\in\operatorname{Gal}(M/\QQ)\) for this class.
All the constituents of the common field
\(\mathcal F\) in Proposition~\ref{prop:commonaux573} are fixed
by \(\sigma\). The only ones requiring a check at a coefficient
prime are the quadratic subfields of \(\QQ(\mu_\ell)\), and
four is a square modulo \(\ell\). The fields \(C_{w,\ell}\)
are fixed since \(\sigma\) is one modulo \(w\). This proves
the splitting assertions. Dirichlet's theorem gives infinitely
many primes in the class.

At a place of \(E_x\) over \(s\) the residue size is \(s\).
For its unramified module \(\ZZ/\ell^n(-r)\), the kernel in the
local calculation above is that of multiplication by
\(s^{-r}-s\). The congruences \(\ell-1\mid 2r\) and
\(s\equiv 2^2\bmod\ell\) give \(s^{-r}-s\equiv -3\bmod\ell\).
The same multiplier acts invertibly on \(\QQ_\ell/\ZZ_\ell(-r)\).
Consequently every local class is unramified. A global cocycle
is then already zero on the added inertia groups, because
they act trivially on the coefficients. It factors through
the group obtained by killing their closed normal closure.
This proves the equality of global groups with the prescribed
zero localizations at \(\ell\) in both coefficient systems.

For the initial points the residual diagonal is trivial over
\(\mathcal F\), and \(b\) is congruent to one. Thus all eigenvalues
of the residual Frobenius at \(s\) on \(U\) are one. Here Frobenius
acts on the Lie space by conjugation since \(s\) splits in \(L\).
On the reduced endomorphisms, and on the symplectic Lie subspace,
all its eigenvalues are one. The determinant of \(s-\Phi_v\)
is therefore a unit. Moreover consider any lift over an Artin
local \(O\)-algebra with finite residue field. At a good \(s\)
the residual inertia is trivial. If inertia first becomes
nontrivial modulo consecutive powers of the maximal ideal,
its first nonzero term gives a Frobenius-equivariant
homomorphism from tame \(\ZZ_\ell(1)\) to the residual matrices
tensored with the corresponding graded piece. Wild inertia
has zero image there. Such a homomorphism would give an
eigenvector of residual conjugation with eigenvalue \(s\).
The distinct residue values \(1\) and \(4\) exclude it.
Induction gives trivial inertia over the whole Artin ring.

It remains to obtain coordinates with the prescribed Frobenius
class. Fix a lift \(\gamma\in G_\QQ\) of \(\sigma\). The coset
\(\gamma G_M\) belongs to \(G_L\cap G_{D_0}\) when the inducing
field is used. We claim there are finitely many \(s_j\) in this
coset and coefficients \(c_j\in E\) with
\[
                      \sum_j c_j U(s_j)=1_4 .
\]
In the symmetric cube case the image of \(G_M\) acts absolutely
irreducibly by the commutator argument of
Theorem~\ref{thm:cuspidalinitial577}, since \(M/\QQ\) is abelian.
Burnside's theorem makes its linear span the full matrix algebra.
In the induced case \(G_M\) acts on the two inducing summands
by absolutely irreducible inequivalent representations, by
Theorem~\ref{thm:regularind579} and its proof on
\(G_{\QQ^{\rm ab}}\). Its span is the product of the two full
matrix algebras. Indeed the two simple quotients have distinct
kernels by inequivalence, and the Chinese remainder theorem
applies. This argument may first be made over an algebraic
closure, since membership in a linear span descends to \(E\).
Now \(U(\gamma)^{-1}\) lies in the indicated span, also in the
induced case because \(\gamma\) fixes \(D_0\). Multiplying a
finite linear expansion of this matrix by \(U(\gamma)\)
proves the claim.

Use the ordinary universal representation modulo the square
of the maximal ideal of its \(E\)-ring. Let \(t\) denote its
trace on \(G_L\). At a chosen \(w'\mid\ell\), take
\(h\in I_{L_{w'}}\cap G_M\) whose four original graded values
are distinct. Its four lifted eigenvalues \(a_i\) are the
ordered graded values, with rank-one spectral projectors
\(P_i\), by lifting the distinct roots. Put
\[
             m_k=\sum_j c_j t(s_jh^k)=\sum_{i=0}^{3} d_i a_i^k,
       \qquad d_i=\operatorname{tr}\left(\sum_j c_j U_A(s_j)P_i\right),
\]
where \(U_A\) denotes this first-order representation on \(G_L\).
Each \(d_i\) reduces to one. If
\(\prod_i(X-a_i)=X^4+\sum_{j=0}^3 b_jX^j\), then
\(\sum_{j=0}^3 m_{k+j}b_j=-m_{k+4}\) for \(0\le k\le3\).
The residual determinant of the coefficient matrix is the
square of the Vandermonde determinant of the four original
values, hence is nonzero. The eight moments therefore
determine all \(b_j\), and simple-root lifting determines
each \(a_i\). Their linear parts are in the span of the
linear parts of the traces on \(\gamma G_M\), because
differentiating the linear system and the simple-root
equations only uses \(E\)-linear combinations at first order.

The images of \(I_{L_{w'}}\cap G_M\) in local units have
finite index. On an open inertia subgroup the original
characters have distinct integral cyclotomic exponents.
Choose \(a\) there, also in \(G_M\), with cyclotomic value
of infinite order. For every \(u\in I_{L_{w'}}\cap G_M\)
all but finitely many \(a^k u\) have four distinct original
values, since each pair of characters can coincide for at
most one \(k\). The just proved assertion applied to \(a\)
and \(a^k u\) thus determines the linear parts of all four
unit-character ratios at \(u\). Taking all the places,
Theorem~\ref{thm:polarweights593} shows that the linear
parts of traces on this one coset span the cotangent space.

In any finite Galois quotient over \(\QQ\) refining \(M\),
Chebotarev realizes the image of each coset element as
Frobenius at a chosen good prime above a rational prime,
with any prescribed finite exclusions. Such a rational
prime belongs to the residue class because \(M/\QQ\) is
abelian, and has degree one in \(L\) by splitting.
The trace linear parts are continuous in a finite-dimensional
\(\ell\)-adic vector space, and their span is closed.
The traces at the indicated primes therefore span the same
space. Choosing a basis and inverting its linear part in the
formal power series ring proves the coordinate assertion.
Restriction to the smooth symplectic quotient allows a
subbasis as in Corollary~\ref{cor:polarsplit593}. The finite
union of choices over all triples and fields proves the last
assertion.
\end{proof}

These primes can be included in the tame level when constructing
ordinary families through the two starting points. The argument
controls the cost of adding them to the strict local conditions.
The sizes of the original strict groups and the signed integral
period comparison in Section~\ref{sec:missing} still require bounds.


\subsection{The ordinary Hecke algebra at each cuspidal point}
\label{sec:heckeseeds595}

Keep the fields \(F,L=FK\), the coefficient ring \(O\) and
the representation \(U=b^{-1}R\) of Section~\ref{sec:tameprimes594}.
We use the automorphic normalization in which
\(r_\ell(\Pi)=U|_{G_L}\). In this normalization the
polarized multiplier attached to an automorphic polarizing
character \(\eta_F\) is \(\kappa_\ell^{-3}r_\ell(\eta_F)\)
in rank four. Here \(r_\ell\) on a character denotes its
\(\ell\)-adic avatar with the corresponding reciprocity
convention. We may enlarge \(O\) to contain the embeddings
and character values used below. Write \(d=[F:\QQ]\).

\begin{lemma}[A conjugate self-dual twist at the original auxiliary fields]
\leavevmode\label{lem:unitaryseed595}
There is an algebraic Hecke character \(X\) of \(L\) such that
\(\Pi'=\Pi\otimes(X\circ\det)\) is regular algebraic,
cuspidal and conjugate self-dual. It is unramified at the
places above \(\ell\) and at every finite place inert in
\(L/F\). It descends to the definite unitary group in four
variables over \(F\) that is quasi-split at every finite
place. Its descent has an ordinary refinement at each
place above \(\ell\). One rational prime from
Theorem~\ref{thm:commonprobes594} suffices to make all the
tame levels used for these descents sufficiently small.
\end{lemma}
\begin{proof}
The multiplier \(\nu\) of \(U\) is the restriction of a
character of \(G_\QQ\), namely \((\det W)^3\) in the odd-order
case and \((\kappa_\ell^j\xi)\phi\) in the even-order case.
We use these particular extensions. Rational class field
theory gives an algebraic Hecke character \(\eta\) of
\(\QQ\) whose avatar is \(\kappa_\ell^3\nu\) on \(G_\QQ\).
It is a finite character times an integral norm power.
The pairing of Proposition~\ref{prop:seedadjoints580} and
the arithmetic Satake normalization give
\[
       \Pi^c\simeq\Pi^\vee\otimes
                  ((\eta\circ N_{L/\QQ})\circ\det).
\]
Indeed at each good split place this is the equality of
the two lists of Frobenius eigenvalues, with the
rank-four normalization contributing the third power
in the character just chosen. The pairing gives the
same identity at all good places of \(L\) by applying
it to their Frobenius elements. Strong multiplicity one
then gives the displayed isomorphism. This choice keeps
the specified multiplier even when \(\Pi\) has a
quadratic self-twist. \cite{XII-IX-blggt560}

Let \(\delta_K\) be the character of \(K/\QQ\).
Write the real component of \(\eta^{-1}\) as
\(\operatorname{sgn}(t)^e|t|^m\). Choose \(u\in\{0,1\}\)
so that \(\delta=\eta^{-1}\delta_K^u\) has real component
\(t^m\). We construct a Hecke character \(\lambda\) of
\(K\) with restriction \(\delta\) to the rational ideles
and with archimedean component \(z^m\) at a chosen embedding.
At each finite prime of \(K\) take the full local units,
except above the conductor of the finite part of \(\delta\),
where we take principal units deep enough that their
intersection with the rational units lies in its kernel.
Call the resulting product \(V_0\). On the subgroup
\[
          K^\times\mathbb A_\QQ^\times K_\infty^\times V_0
                      \ \subset\ \mathbb A_K^\times
\]
prescribe the trivial character on \(K^\times,V_0\) and the
two indicated characters on the other factors. To check
consistency, a global \(x\in K^\times\) appearing in a
relation must generate the extension of a rational
fractional ideal, by comparing the finite valuations.
Such an ideal is principal over \(\QQ\), and the only units
in \(K=\QQ(\sqrt{-p})\), \(p\equiv7\bmod8\), are \(1,-1\).
Thus \(x\) is rational and can be absorbed in the rational
idele without changing its character value. The remaining
relation holds by the choice of \(V_0\) and by equality
of the two archimedean characters on \(\mathbb R^\times\).
The subgroup is open of finite index by ray class finiteness.
Extend its character by adjoining coset generators and taking
a root of the prescribed value of each first power lying in
the preceding subgroup. The result is continuous and has
the given algebraic infinity type. Its finite-idele values
lie in a number field. Indeed on the initial subgroup these
values use rational norm powers, finite character values and
powers of elements of \(K\), and the extension uses finitely
many further roots.

Now \(X=\lambda\circ N_{L/K}\) satisfies
\[
              XX^c=\delta\circ N_{L/\QQ}
                   =\eta^{-1}\circ N_{L/\QQ}.
\]
The first equality follows from
\(\lambda(y)\lambda(\bar y)=\delta(N_{K/\QQ}y)\);
the second follows since \(\delta_K\) is trivial on norms.
This gives conjugate self-duality. Twisting preserves
cuspidality and the differences between the four weights.
The possible finite ramification of \(X\) lies over the
conductors of \(\xi,\theta,\phi\) when these occur and over
the discriminant prime \(p\). The former primes all split
in \(K\) by Lemmas~\ref{lem:twoprimes571},
\ref{lem:realquadratic572} and~\ref{lem:cuspidalplane577}.
At \(p\), \(L=FK'\) and \(p\) splits in \(K'\).
All these primes are away from \(\ell\). The original
\(\Pi\) is spherical above \(\ell\): the plane used for the
symmetric cube has prime-to-\(\ell\) level, and the two
local inducing planes at the split places of \(D_0\) in
the even-order case also have prime-to-\(\ell\) level.
The respective unramified Satake transfers and base changes
are spherical. Away from \(\ell\) the conductor calculations
of Proposition~\ref{prop:tamerigid594} and local--global
compatibility give spherical components whenever the
underlying rational prime is outside those tame levels.
This proves the asserted unramifiedness.

The extension \(L/F\) is unramified at every finite place
and all places above \(\ell\) split, by
Theorem~\ref{thm:tametower571} and
Lemma~\ref{lem:realquadratic572}, also after adjoining
\(C_{w,\ell}\) to the real base. The parity condition for
the definite unitary group quasi-split at finite places
holds since \(4d\) is divisible by four. Apply Labesse's
descent in Geraghty's Proposition~2.2.7. It gives the
indicated general linear components at split places and
hyperspecial fixed vectors at the inert places.
The crystalline ordinary flag of \(U\) stays ordinary
after twisting by the avatar of \(X\); this avatar is
crystalline at these places since its smooth character is
unramified there. The ordinary refinement follows from
the level-prime-to-\(\ell\) ordinarity implication in
Barnet-Lamb, Gee, Geraghty and Taylor, Section~2.1,
remark~(7) after Theorem~2.1.1.
At each of the ramified split tame places choose a compact
open subgroup admitting the given local representation.
Choose one further good odd rational prime in the progression
of Theorem~\ref{thm:commonprobes594}. At one place above it
use a principal congruence subgroup inside the domain of the
matrix logarithm. It is torsion-free, so any finite rational
stabilizer of a definite double coset is trivial. This gives
the required small levels for all the fields. \cite{XII-IX-geraghty363,XII-IX-blggt560}
\end{proof}

At this level use the ordinary Hecke action on the space of
definite unitary forms, with all four diagonal weights
variable at the chosen place of \(L\) above each \(v\mid\ell\).
Fix the finite residue-unit diamond character of the starting
point. Write \(\Lambda\) for its weight algebra over \(O\).
It is the completed group ring of the product of four
principal-unit groups at every such \(v\). The local
ramification index over \(\QQ_\ell\) is at most two, since
\(C_{w,\ell}\) is unramified at \(\ell\). Thus the principal
units have no \(\ell\)-power torsion, as \(\ell\ge5\).
The ring \(\Lambda\) is a power series ring in \(4d\)
variables over \(O\).

Let \(\mathbb T\) be the \(\Lambda\)-algebra acting through
the spherical Hecke operators at the good split places,
including the inverse of the last spherical operator.
We use the coefficients of the good Frobenius polynomials
after undoing \(X\). They give the same algebra by the
unramified Satake formula, since the residue cardinalities
and the values of the fixed Galois character are units at
these places. Let \(z\) be the characteristic-zero
eigensystem furnished by the lemma, with weight \(a\).
Enlarge \(E\) so that both are \(E\)-rational, and write
\[
  A=\widehat{\Lambda[1/\ell]_a},\qquad
  C=\widehat{\mathbb T[1/\ell]_z}.
\]

\begin{theorem}[The completed ordinary Hecke algebra at both starting points]
\leavevmode\label{thm:heckeseeds595}
The maps from the completed weight algebra to \(C\) and to the
polarized ordinary deformation ring of \(U\) identify all
three rings. Under these identifications the good split
Hecke polynomials, with \(X\) removed, are the Frobenius
polynomials of the universal deformation. The symplectic
locus is cut out by the two unit-character equations of
Theorem~\ref{thm:polarweights593} at each \(v\mid\ell\).
\end{theorem}
\begin{proof}
Geraghty's ordinary control theorem, Proposition~2.5.3,
Section~2.6 and Lemma~2.6.4, gives a finite free dual
ordinary module \(N\) over \(\Lambda\). The small tame level
allows the first principal-unit subgroup in Proposition~2.5.3.
Projecting onto the desired residue-unit characters is an
integral idempotent since their orders are prime to \(\ell\).
The summand is free over the local ring and nonzero at \(z\).
At dominant algebraic weights its characteristic-zero
specialization is the dual of the corresponding classical
ordinary space. The commuting good spherical operators act
semisimply there by the classical automorphic decomposition.
The subalgebra \(\mathbb T\subset\operatorname{End}_\Lambda(N)\)
is finite and torsion-free over \(\Lambda\), by Noetherianity
and freeness. \cite{XII-IX-geraghty363}

Purely algebraic weights strictly dominant at every embedding
and with the indicated residue-unit characters detect every
element of \(\Lambda\). Here are details of density in the
several weight variables. Choose \(4d\) linearly independent
integral exponent tuples in the interior of the dominance
chamber and multiply them by a common multiple of the
residue-unit orders. One can choose such tuples because a
nonempty open real cone contains a basis of rational vectors.
Add their nonnegative integral combinations to the weight of
the starting point, whose consecutive Hodge gaps are at least
two. In logarithmic coordinates the derivative of the resulting
weight parameterization is invertible over \(E\): for each
\(F_v\), the embedding matrix on a \(\QQ_\ell\)-basis is
invertible by separability, and the chosen exponent matrix
has full rank. The parameterization therefore gives local
analytic coordinates at \(a\), by the logarithm and exponential
on a small polydisc. In that polydisc each integral parameter
can independently run through positive powers of \(\ell\)
tending to zero. Successive use of the one-variable identity
theorem makes any analytic function vanishing on this grid
zero as a germ. A weight power series injects into its
completed local ring at \(a\) by localization and the Krull
intersection theorem. This proves density.

For every eigensystem at these weights take its semisimple
Galois representation on \(G_L\) from Geraghty's
Proposition~2.7.2 and untwist by \(r_\ell(X)\). Include
the starting eigensystem. The proposition applies to the
classical constituents of the definite group, with the
degree-four good Hecke polynomial and with multiplier
\(\kappa_\ell^{-3}\) before untwisting. Thus after untwisting
the multiplier on \(G_L\) is exactly
\[
   \kappa_\ell^{-3}r_\ell(XX^c)^{-1}
        =\kappa_\ell^{-3}r_\ell(\eta\circ N_{L/\QQ})
        =\nu|_{G_L}.
\]
Evaluation embeds \(\mathbb T\) in the product of the
coefficient valuation rings of these classical systems:
if all evaluations vanish, the operator acts as zero on
each of the classical fibers by semisimplicity, and every
matrix entry on \(N\) vanishes by weight density. The finite
\(\Lambda\)-module \(\mathbb T\) is compact and its image in
this Hausdorff product is closed. Good Frobenius elements
at places of \(L\) split over \(F\) are dense up to choice
of places and conjugacy. Indeed any element of \(G_L\) is
realized in a prescribed finite quotient by Chebotarev over
\(F\), using a Galois closure containing \(L\). Taking
coordinatewise traces in the compact product therefore
defines a continuous \(t:G_L\to\mathbb T\), agreeing with
each classical trace everywhere. The representations are
unramified outside a fixed finite set made up of the level
and the places above \(\ell\). \cite{XII-IX-geraghty363}

We reconstruct a representation over \(C\). Choose
\(g_1,\ldots,g_{16}\in G_L\) whose matrices at \(z\)
form a basis, by absolute irreducibility and Burnside's
theorem. The trace-pairing matrix \(H=(t(g_i g_j))_{ij}\)
has determinant nonzero at \(z\). For \(h\in G_L\) define
\(e(h)\) in a free module of rank sixteen over
\(\mathbb T[1/\det H]\) by assigning its coordinates in
the basis \(e(g_i)\) using the equations with right sides
\(t(hg_j)\) and coefficient matrix \(H\). Give this
module the multiplication \(e(g_i)e(g_j)=e(g_i g_j)\).
The product is associative, has identity \(e(1)\), and
\(e(h)e(k)=e(hk)\). These identities hold at every
classical specialization with \(\det H\ne0\), where the
sixteen matrices are a basis. They hold in the localization
since the product of the evaluations is faithful: after
clearing denominators one further factor of \(\det H\)
also kills the evaluations with zero determinant.

After extending to \(C\) this algebra has special fiber
\(\operatorname{Mat}_4(E)\). Lift a rank-one idempotent to
an idempotent \(e\) in the algebra. For example Newton
iteration for \(Y^2-Y\) starting at any lift converges in
the complete ring, since \(2Y-1\) reduces to an involution
and the iterates commute. The left ideal cut out by right
multiplication by \(e\) is a free module of rank four,
as a projective module over the local ring \(C\). Left
multiplication identifies the algebra with its endomorphism
algebra, since reduction gives that isomorphism and both
sides are free of rank sixteen. This gives a representation
\(\rho\) lifting \(U|_{G_L}\), and its trace is \(t\).
For the last assertion, the trace of left multiplication
on the rank-sixteen algebra is four times the prescribed
trace at all the same classical points, hence in the
localization. In a matrix algebra it is four times the
ordinary matrix trace. The construction is continuous on
every finite Artin quotient of \(C\): the trace was continuous
in the compact Hecke algebra, and evaluation of a weight
variable there has constant term in the maximal ideal of
\(O\) plus a nilpotent term. It gives a continuous map from
\(\Lambda\) to that finite-dimensional \(E\)-algebra with
its \(\ell\)-adic topology, and hence from a finite
\(\Lambda\)-module by using module generators.

Conjugate duality holds over \(C\). The traces of
\(\rho^c\) and \(\nu\rho^\vee\) agree on all of \(G_L\)
by their classical evaluations. Two such representations
with absolutely irreducible isomorphic special fibers and
equal traces are isomorphic. Indeed the trace-pairing
basis construction identifies the two matrix algebras and
their group elements, and a matrix algebra automorphism
over a local ring is inner. The latter follows by choosing
a generator for the image of the first diagonal matrix
unit and using the other matrix units to carry it to a
basis of the module. Thus an invertible pairing
\(\mathcal J\) satisfies
\(\rho(h)^{\mathsf T}\mathcal J\rho(chc)=\nu(h)\mathcal J\)
for \(h\in G_L\). Scale it to reduce to \(J U(c)\).
Its transpose is another such pairing, so is a scalar
multiple of it since the matrices span the full algebra.
This scalar squares to one and reduces to one, hence is one.
A strict change of basis makes the symmetric pairing equal
to \(JU(c)\): the endomorphism \((JU(c))^{-1}\mathcal J\)
has a square root congruent to the identity given by the
binomial series, self-adjoint for \(JU(c)\).
With this basis take the polarized cocycle to equal one
at \(c\) and \(\rho(h)U(h)^{-1}\) at \(h\in G_L\).
The pairing equation verifies the cocycle relation at
\(chc\), and \(c^2=1\) verifies the remaining relation.

At each chosen \(w\mid\ell\) let \(\chi_i\) be the four
unit characters in the universal weight, in the order of
the initial flag, including the fixed untwist. For clarity,
Geraghty's Corollary~2.7.8 gives for the \(j\)-th
character before untwisting, \(1\le j\le4\), the product
of \(\kappa_\ell^{1-j}\) and the character which on units
\(x\) is \(\prod_{\tau:w}\tau(x)^{-\lambda_{\tau,5-j}}\)
at algebraic weight \(\lambda\). The four varying factors
are precisely the independent slot characters of the
weight algebra in his Section~2.6, up to its fixed
normalizing characters. Ratios to their values at \(z\)
therefore give formal coordinates on \(A\). The ordinary
flag in the corollary applies at the regular weights used
for density. In fact the residue-unit part may be fixed at
its algebraic value throughout our grid, so the classical
space being used is the ordinary space with Iwahori
invariants and this algebraic weight. \cite{XII-IX-geraghty363}

Choose \(h_0\) on open inertia with four distinct graded
values at the starting point. Put \(a_i=\chi_i(h_0)\) and
\(p_i(Y)=\prod_{j\ne i}(Y-a_j)/(a_i-a_j)\). The denominators
are invertible at \(z\). By the classical flags, the
characteristic polynomial at \(h_0\) is \(\prod_i(Y-a_i)\),
and for \(g\in G_{L_w},\,h\in G_L\) the trace of
\[
           p_j(\rho(h_0))\,\rho(g)\,p_i(\rho(h_0))\,\rho(h)
\]
vanishes when \(j>i\). Indeed expanding the polynomials
writes the expression in group traces with weight
coefficients; the identity on the open set of classical
points passes to the localization by clearing denominators
as above. Nondegeneracy of the trace pairing over \(C\)
then gives zero for the product before the last factor.
The ordered partial sums of the four eigenlines of \(h_0\)
are consequently stable. For \(u\in I_{L_w}\), interpolation
of \(t(h_0^k u)=\sum_i a_i^k\chi_i(u)\), \(0\le k\le3\),
and inversion of the Vandermonde matrix identify the
graded characters on inertia. This proves the required
ordinary condition on the entire completion, even where
the four characters have coincident reductions in
characteristic \(\ell\).

At the finite places away from \(\ell\) the unrestricted
local lifts just obtained have fixed inertial type on every
Artin quotient by Proposition~\ref{prop:tamerigid594}.
The polarized ordinary universal property in
Theorem~\ref{thm:polarweights593}, applied to these
compatible Artin quotients, now gives a homomorphism from
its ring to \(C\). In the four unit-character coordinates
this is exactly the weight homomorphism \(A\to C\).
In particular all normalized good spherical coefficients
in \(C\) are images of elements of \(A\), by taking the
characteristic polynomials in the universal representation.
Their last coefficients are units. Since \(\mathbb T\)
is generated over \(\Lambda\) by these operators and their
indicated inverses, \(A\to C\) is surjective. More
explicitly \(\mathbb T\otimes_\Lambda A\) is finite over
the complete local ring \(A\) and is the product of the
completed local factors above \(a\). Projection to \(C\)
is surjective, and all the algebra generators have just
been accounted for. Finally \(\mathbb T\) embeds into a
finite free \(\Lambda\)-module. Flatness of localization
and completion gives a torsion-free finite module after
tensoring with \(A\), and the nonzero factor \(C\) is
torsion-free as well. Thus \(A\to C\) is injective.
The claimed symplectic locus is now the one in
Theorem~\ref{thm:polarweights593}.
\end{proof}

\begin{corollary}[Hecke relations on the components through the initial point]
\leavevmode\label{cor:seedcomponents595}
Let a Noetherian integral-coefficient family of marked
ordinary symplectic representations of \(G_F\) have the
specified fixed multiplier and have \(R\) at a
characteristic-zero point \(z'\). Assume continuity on
the local Artin specializations at \(z'\) and ramification
confined to a fixed finite set. Compute the weights and
the good split traces on \(G_L\) after undoing \(b\).
Every polynomial relation over \(\Lambda\) in the good
spherical coefficients that vanishes in \(\mathbb T\)
vanishes on each irreducible component through \(z'\).
\end{corollary}
\begin{proof}
Complete the local ring of the given family at \(z'\),
extending the coefficient field if necessary. Forgetting
the markings and applying symplectic descent as in
Theorem~\ref{thm:polarweights593} gives a polarized
ordinary deformation at its Artin quotients. Extra tame
places are allowed by Proposition~\ref{prop:tamerigid594}.
The preceding theorem identifies its trace and weight
specialization with a homomorphism from \(C\). Thus the
relation is zero in the completion. Faithful flatness
of completion gives zero in the local ring, so its
value in the original family is annihilated by an element
outside the prime of \(z'\). Every prime contained in
that prime also avoids this annihilator and hence
contains the value. This applies to all the minimal
primes of components passing through the point.
\end{proof}

This proves the trace relations on the components passing
through the two characteristic-zero starting points.
Propagation to the other components encountered at
characteristic \(\ell\), the strict lengths at the original
weights, and the signed period estimate on \(E_x\) in
Section~\ref{sec:missing} are still open.


\subsection{The Hecke fiber at the prescribed rational weights}
\label{sec:heckefiber596}

We specialize the family of Section~\ref{sec:heckeseeds595} at
the weights of the original strict-cohomology problem. Fix
one of the primes of \(K\) above \(\ell\), and in defining the
unitary level choose above each \(v\mid\ell\) of \(F\) the
place of \(L=FK\) lying over that prime. These choices are
available at all \(v\), and the original local representations
over \(F_v\) identify with those at the chosen places. Write
\(\psi_0,\ldots,\psi_3\) for the four ordered inertia
characters of the original rational weight, restricted to
these places. On open inertia their cyclotomic exponents
in order are
\[
                            r+1,\quad r,\quad 1,\quad 0.
\]
Their outer and inner products both equal \(\mu\) on inertia.
The Hecke algebra \(\mathbb T\) uses the untwisted normalization
of the preceding section. Its starting point is \(U=b^{-1}R\),
and the target in its weight space has ordered unit characters
\(b^{-1}\psi_i\).

\begin{proposition}[An integral component reaching the prescribed weight]
\leavevmode\label{prop:horizontal596}
In \(\Lambda\) impose the two symplectic equations on all
local units from Theorem~\ref{thm:polarweights593} and denote
the quotient by \(\Lambda_{\rm sp}\). It is a power series
ring over \(O\) in \(2d\) variables.
There is a unique irreducible component of
\(\operatorname{Spec}(\mathbb T\otimes_\Lambda\Lambda_{\rm sp})\)
through the characteristic-zero cuspidal point \(z\).
Its reduced coordinate ring \(B\) is a finite domain
containing \(\Lambda_{\rm sp}\). Its fiber over the
exact weight \((b^{-1}\psi_i)_i\) is nonempty in
characteristic zero and is finite over \(O\) before
inverting \(\ell\). Throughout \(B\) the trace on \(G_L\)
is invariant under conjugation by \(c\in G_F\) and satisfies
\[
             t(h)=\nu(h)t(h^{-1})\qquad(h\in G_L).
\]
\end{proposition}
\begin{proof}
Use a basis of the principal units modulo torsion in each
slot. For each basis element, the ratios of the two outer
characters to their starting values are required to be
mutual inverses, and likewise for the two inner characters.
These equations eliminate two of the four group-ring
variables by solving for their units \(1+T\). They hold
on the whole unit group by continuity and the fixed
residue-unit characters. This proves the description of
\(\Lambda_{\rm sp}\).

The completion of the local ring of the indicated tensor
product at \(z\) is the quotient of the ring in
Theorem~\ref{thm:heckeseeds595} by the same equations.
It is the completed characteristic-zero weight ring of
\(\Lambda_{\rm sp}\), a power series domain. Since the
local ring injects into this completion, it is a domain.
The primes of the tensor product contained in the prime
of \(z\) correspond to the primes of that local ring.
There is thus one minimal prime there, giving the domain
\(B\). The inclusion of \(\Lambda_{\rm sp}\) in its own
completed local ring at the weight of \(z\) shows that
this minimal prime has zero contraction. Finiteness
follows from finiteness of \(\mathbb T\) over \(\Lambda\).

The characters \(b^{-1}\psi_i\) have the same reductions
as the ordered characters of \(U\) on inertia by the
residual flag assertions in Theorems~\ref{thm:cuspidalinitial577}
and~\ref{thm:regularind579}. On prime-to-\(\ell\) torsion,
equality after reduction implies equality of the original
values. Their two paired products are \(b^{-2}\mu=\nu\).
Thus the exact tuple defines an \(O\)-valued weight of
\(\Lambda_{\rm sp}\). Tensoring \(B\) at this weight
gives a finite \(O\)-algebra. Lying over for the integral
injection \(\Lambda_{\rm sp}\subset B\) gives a prime above
the kernel of the weight. Its characteristic-zero residue
field is finite over \(E\) since the fiber is finite.
Finally \(B\) embeds in its localization at \(z\) and
in the completion of that localization. On the latter
the universal representation on \(G_L\) comes from a
symplectic one over \(F\) with multiplier \(\nu\), by
Theorems~\ref{thm:polarweights593} and~\ref{thm:heckeseeds595}.
The two trace identities there descend along the injection.
\end{proof}

The scalar \(b\) can have half-integral cyclotomic exponent.
Consequently classical control at the indicated point
requires a choice of scalar on the two sets of places of
\(L\) above \(\ell\). We make this choice on \(L\) itself.

\begin{lemma}[An anticyclotomic translation of the ordinary Hecke algebra]
\leavevmode\label{lem:halftwist596}
There is a continuous \(\ell\)-adic character
\(\beta:G_L\to O^\times\), unramified away from \(\ell\),
with \(\beta^c=\beta^{-1}\) and trivial reduction, with
the following properties.
\begin{enumerate}
\item Its restriction to inertia at the chosen places above
\(\ell\) is \(b\). The character \(b|_{G_L}/\beta\) is the
avatar of an algebraic Hecke character of \(L\).
\item Multiplication of the four ordered unit characters
by \(\beta\), and of each good degree-\(j\) Frobenius
coefficient by \(\beta(\operatorname{Frob})^j\), acts by
an automorphism of the ordinary Hecke algebra \(\mathbb T\).
\end{enumerate}
\end{lemma}
\begin{proof}
The character construction of Lemma~\ref{lem:unitaryseed595}
also gives an algebraic character of \(K\) with complex
component \(z/\bar z\), trivial on rational ideles and
on all finite local units. Indeed its restriction at
infinity to \(\mathbb R^\times\) is one, and the same
ideal argument verifies the relations with \(K^\times\)
when every finite unit group is taken. Use this character
or its inverse and pull it back by \(N_{L/K}\). Denote
the resulting Hecke character by \(Y\) and its avatar by
\(\Xi\), with the sign chosen so that \(\Xi\) equals
\(\kappa_\ell\) on inertia at the chosen places.
To see the restriction, at the two split places of \(K\)
the algebraic unit factors of its avatar are the two
opposite degree-one powers of local reciprocity. Pullback
by the local norm gives the rational cyclotomic character
or its inverse on the units of \(L_w\). The smooth
finite component is unramified. The identity
\(YY^c=1\) follows from its trivial rational restriction.

Write \(m=3j\) for the symmetric cube and \(m=j\) for the
induced representation, so \(\nu\) has cyclotomic exponent
\(m\). Put \(a=r+1-m\) and \(s=a/2\in\ZZ_\ell\). For a
unit character \(\gamma\), write \(\langle\gamma\rangle\)
for its quotient by the Teichm\"uller lift of its reduction.
On principal units at every \(v\mid\ell\) the defining
equation \(b^2=\mu/\nu\) gives \(b^2=\kappa_\ell^a\),
since all finite factors there are unramified or
Teichm\"uller characters. On residue-unit torsion \(b\)
is trivial since it reduces to one. Taking the unique
square root on principal units therefore gives
\(b=\langle\kappa_\ell\rangle^s\) on all of inertia
in the local abelianization. Set
\[
                          \beta=\langle\Xi\rangle^s.
\]
The power exists continuously in the abelian pro-\(\ell\)
group of principal coefficient units, by taking limits of
integer powers. This argument also applies in the presence
of \(\ell\)-power torsion. The choices give every local and
conjugation assertion about \(\beta\).

At the chosen places \(b/\beta\) is trivial on inertia.
At the conjugate places it has cyclotomic exponent \(2s=a\)
on principal units. To verify algebraicity explicitly,
choose an algebraic character of \(K\) whose avatar has
exponent zero at the chosen place and \(a\) at its conjugate.
The same idele construction supplies it: take the integral
one-sided infinity type with these exponents, and on the
rational ideles take a norm power with the necessary real
sign, adjusting the sign by \(\delta_K\). Ramification at
the discriminant prime is permitted in this construction.
After pulling back to \(L\), its quotient into \(b/\beta\)
has finite inertia at \(\ell\) and at every other finite
place. Away from \(\ell\) this follows also from the square
equation for \(b\), since its continuous character there
has finite inertial image when its square does. It is
unramified outside a finite set. Thus the quotient factors
through a finite ray class group of \(L\): take open kernels
on units at the finitely many ramified places and full units
at the rest, and use global class field theory. It has finite
order. This proves the assertion about \(b/\beta\).

For the automorphism first use \(Y^k\) with \(k\) an integer
divisible by the order of the reduction of \(\Xi\).
This character is trivial on \(\mathbb A_F^\times\) inside
\(\mathbb A_L^\times\), since the norm of a real-base idele
to \(K\) is its rational norm. Thus \(Y^k\) is pulled back
from the norm-one torus by \(z\mapsto z/\bar z\).
Local Hilbert 90 gives surjectivity of this map on ideles:
at almost all places an integral norm-one unit has a unit
preimage, since the extension is unramified and one can
adjust a preimage by a base uniformizer. Hilbert 90 over
the number fields similarly gives triviality on rational
points of the torus. The resulting automorphic character
of the torus is trivial on finite maximal compact subgroups
at nonsplit places by this unit-preimage argument, and on
finite units at split places by unramifiedness of \(Y\).
Composing with the determinant of the unitary group
therefore twists classical eigensystems at the same tame
level. Its base-change character at a split place is
\(Y^k\) on the chosen factor, as seen on ideles \((x,1)\).
On weights it adds one common exponent to the four slots
at each embedding. This preserves dominance and ordinarity
(the four graded characters receive the same scalar);
our choice of \(k\) preserves the residue-unit summand.
\cite{XII-IX-geraghty363,XII-IX-blggt560}

On the normalized good Frobenius coefficients and the four
unit characters the induced operation is multiplication
by the indicated powers of \(\Xi^k\). This respects all
relations in \(\mathbb T\): evaluate them on regular
algebraic classical eigensystems and use the faithful
product of evaluations in Theorem~\ref{thm:heckeseeds595}.
The inverse is given by \(-k\). Choose such integers \(k\)
tending to \(s\) in \(\ZZ_\ell\), possible since their
required divisibility is prime to \(\ell\). The resulting
scalars \(\Xi(g)^k\) converge to \(\beta(g)\) for each
\(g\). The weight substitutions converge on formal power
series, since they send each group generator to itself
times a unit congruent to one and converge in the
maximal-ideal topology. Passing to the limit in every
polynomial relation among the Hecke generators over
\(\Lambda\) is legitimate in the finite complete
\(\Lambda\)-algebra. This gives the substitution claimed
in the lemma, with inverse obtained using \(-s\).
\end{proof}

\begin{theorem}[Classical realizations in the fiber of the original weight]
\leavevmode\label{thm:classicalfiber596}
At every geometric characteristic-zero point \(y\) of
\(\mathbb T\) above \((b^{-1}\psi_i)_i\), multiplying its
normalized trace by \(b|_{G_L}\) gives the trace of a
semisimple Galois representation arising from a classical
automorphic system over \(L\), with the original weight
differences. If this representation is absolutely irreducible,
it is attached to a regular algebraic cuspidal representation
of \(\mathrm{GL}_4(\mathbb A_L)\) and is pure at finite places
away from \(\ell\). There is at least one such classical
realization, allowing reducible Galois representations,
on the component of Proposition~\ref{prop:horizontal596}.
The number of distinct geometric eigensystems in this
whole characteristic-zero fiber is at most the
\(\Lambda\)-rank of the ordinary module \(N\) of
Theorem~\ref{thm:heckeseeds595}.
\end{theorem}
\begin{proof}
Translate \(y\) by the automorphism of
Lemma~\ref{lem:halftwist596}. In the normalization with
\(X\) undone, its four unit characters are now exactly
\(\psi_i\) at the chosen places. Restoring the fixed
algebraic \(r_\ell(X)\) shifts all four local exponents by
the same integer at each embedding. The exponents have
successive gaps \(1,r-1,1\), all at least one since \(r\ge2\).
Subtracting the consecutive shifts in the ordinary weight
formula of Geraghty's Corollary~2.7.8 thus gives a dominant
algebraic coefficient weight. The remaining unit characters
are finite order. In particular this is an arithmetic
weight even when \(a\) in the preceding lemma is odd.

Ordinary control for a dominant weight times a finite-order
diamond character gives the classical ordinary space at
sufficient depth above \(\ell\). Every character of the
good Hecke algebra over this weight occurs on this space.
Indeed its finite module \(N\) is faithful over \(\mathbb T\),
so its support is all of \(\operatorname{Spec}\mathbb T\).
Localize at the character and reduce modulo its maximal
ideal. Nakayama's lemma gives a nonzero fiber, also after
extending scalars. The good operators are semisimple on
the classical space and its dual. The dimension of the
specialized dual is the constant rank of \(N\), which
bounds the number of different characters. Translation
preserves that number. \cite{XII-IX-geraghty363}

Let \(D\) be the semisimple Galois representation attached
to the translated classical system in its unitary
normalization by Geraghty's Proposition~2.7.2. At good
split primes its trace is \(r_\ell(X)\beta\,t_y\);
Chebotarev gives the identity on \(G_L\).
Twist \(D\) by the algebraic character whose avatar is
\((b/\beta)r_\ell(X)^{-1}\). The resulting representation
has trace \(b\,t_y\). On the automorphic side this is an
algebraic twist of the weak base change of the classical
definite unitary system. In the description of that base
change used in Proposition~2.7.2, several cuspidal
constituents give a direct sum of their Galois
representations, and a discrete block of length greater
than one contributes the various norm twists of its
constituent. These decompositions follow on semisimplifying
from the good Frobenius polynomials. Absolute irreducibility
therefore forces a single cuspidal constituent of rank four
with block length one. Its distinct weights give regularity.
The resulting algebraic twist is still polarized over
\(L/F\), with its polarizing character multiplied by the
product of the twisting character and its conjugate.
Purity follows from Theorem~2.1.1(2) of Barnet-Lamb, Gee,
Geraghty and Taylor. The existence statement on the
specified component follows from the nonempty fiber in
Proposition~\ref{prop:horizontal596}. \cite{XII-IX-geraghty363,XII-IX-blggt560}
\end{proof}

In particular an absolutely irreducible restriction to
\(G_L\) of a marked representation at the original weight
has this cuspidal realization whenever its normalized
traces factor through the indicated Hecke algebra with
that weight. Establishing this factorization for every
relevant component and bounding the full integral strict
lengths are still necessary for the argument in
Section~\ref{sec:missing}. The integral theta comparison
there concerns the periods on the original \(E_x\).


\subsection{Regular Frobenius on ordinary curves in positive characteristic}
\label{sec:regularcurves597}

Return to the minimal fields \(F,L\) in
Section~\ref{sec:tameprimes594}. Write \(D=F\) or \(D=L\),
\(G=G_D\), \(G'=G_{L(\mu_{\ell^\infty})}\). Let \(k\) be a
finite field of characteristic \(\ell\), \(A=k[[t]]\), and
\(K_1=k((t))\). Consider a continuous representation
\(\rho:G\to\mathrm{GL}_4(A)\) unramified away from finitely
many places, whose reduction has the four residual
characters of the indicated case as its semisimplification.
A common scalar twist of those residual characters is
allowed. Suppose that over \(K_1\) it has an ordinary flag
at the chosen places above \(\ell\), with ordered continuous
unit characters \(\theta_{v,i}\) taking values in \(A^\times\).
For each \(v\) let \(I_v^0\) be the subgroup of its inertia
on which the full \(\ell\)-adic cyclotomic character and
the original four residual diagonal characters are trivial.
These inertia groups lie over split places of \(L/F\).
Assume that for each pair \(i\ne j\), the characters
\(\theta_{v,i},\theta_{v,j}\) differ somewhere on \(I_v^0\),
and that their reductions agree there. These are the
separation hypotheses on the curves in
\srcloc{Development/AutomorphicArithmetic/AutomorphicVarieties/ModularityArgument/ThresholdArguments/DenseThresholdCases/ModularDeformationQuotientRegularAtPoint}{26}.
For a matrix, a regular semisimple spectrum below means
that its characteristic polynomial has four distinct roots
over the algebraic closure.

\begin{theorem}[Density of regular spectra on the cyclotomic kernel]
\leavevmode\label{thm:regularcosets597}
If \(\rho\) is absolutely irreducible over \(K_1\), then
its restriction to \(G'\) is absolutely irreducible. In
every nonempty open subset of \(G'\), as well as in every
nonempty open subset of \(G\), there is an element of
regular semisimple spectrum. The local ordered characters
are pairwise different and a character in one position is
never the product of another with the residual cyclotomic
character. At any good prime from the progression of
Theorem~\ref{thm:commonprobes594}, the local representation
has no nonzero intertwiner to its cyclotomic twist, also
over \(\overline{K_1}\).
\end{theorem}
\begin{proof}
On \(I_v^0\) each quotient of two different local characters
has nontrivial image in \(1+tk[[t]]\). The latter group
is torsion-free: it is pro-\(\ell\), and
\((1+x)^{\ell^m}=1+x^{\ell^m}\) in this domain.
Each of the six kernels is closed of infinite index,
hence has empty interior. Their union cannot cover the
compact group \(I_v^0\), by the Baire theorem. Choose
\(a\) outside the union at any one of the places.
Then \(a\in G'\), \(\rho(a)\) has distinct eigenvalues
and all its positive powers have distinct eigenvalues.
The asserted local character inequalities follow by the
same restriction at each place, since the cyclotomic
character is one on \(I_v^0\).

The semisimplified residual characters take distinct values
at one \(g_0\in G'\), and one of their ratios has order
not dividing four. For \(F=H\) this follows from
Lemma~\ref{lem:cyclospaces575} and
Theorems~\ref{thm:tametower571} and~\ref{thm:eventower572}.
For \(F=HC_{w,\ell}\) take the \(\ell\)-th power of the
element for \(H\). It fixes \(C_{w,\ell}\), and raising
residual values to the \(\ell\)-th power preserves
distinctness and their prime-to-\(\ell\) orders. In the
odd-order case the ratio value \(\chi(g_0)\) has order \(2m\ge6\)
with \(m>1\) odd, and in the even-order case
\(\chi(g_0)^4\ne1\). Common scalar twists do not affect
these comparisons.

Work over \(\overline{K_1}\). Restriction to the normal
subgroup \(G'\) is semisimple: the sum of its simple
submodules is nonzero and \(G\)-invariant. In each
isotypic component multiplicity is one, since multiple
copies force repeated eigenvalues of \(a\). The number
of components divides four because \(G\) permutes them
transitively. Its permutation quotient is abelian since
\(L(\mu_{\ell^\infty})/D\) is abelian, and is regular
on the components: a stabilizer in a transitive abelian
permutation group fixes every point. The action is
continuous. Each component is a sum of eigenlines of
\(\rho(a)\), and its stabilizer is closed by continuity
of the matrices and hence open by finite index.
If there were two or four components, this quotient
would admit a nontrivial character \(\eta\) of order two.
Indexing the components by the quotient, multiplication
on each by its value of \(\eta\) would intertwine
\(\rho\) with \(\eta\rho\). Reduction of their characteristic
polynomials and Brauer--Nesbitt would then make \(\eta\)
permute the four residual characters. It is trivial at
\(g_0\), where the values are distinct, so the permutation
fixes each character. This forces \(\eta=1\), a contradiction.
Thus \(G'\) acts irreducibly.

Let \(H\) be the Zariski closure of \(\rho(G')\) and
\(H^\circ\) its identity component. The latter is normal
under \(\rho(G)\). The identity component of the closure
of the powers of \(\rho(a)\) is a torus in \(H^\circ\),
in the eigenbasis of \(a\). Its four weights are distinct:
a quotient character trivial on that torus would have
finite image on the cyclic group, contrary to our choice
of \(a\). A cocharacter of the torus can be chosen outside
the finitely many hyperplanes on which two weights have
equal pairing. Thus \(H^\circ\) contains
\[
             D(T)=\operatorname{diag}(T^{w_1},T^{w_2},
                                      T^{w_3},T^{w_4})
\]
as a one-parameter subgroup, with four distinct integers
\(w_i\). Restriction to \(H^\circ\) is semisimple by the
same simple-submodule argument, using the irreducibility
on \(H\). The simple constituents are pairwise inequivalent
and of equal dimension, since a power of \(a\) in \(H^\circ\)
has simple spectrum and \(H\) permutes the components
transitively.

We exclude four lines for \(H^\circ\). Otherwise \(G\)
permutes the four lines transitively, with quotient
\(P\subset S_4\) and kernel \(N_0\). Over a finite
extension take the direct sum of the projections of
a stable lattice onto the four lines. This is a stable
monomial lattice, since \(G\) permutes its summands.
The four reduced characters on \(N_0\) are the restrictions
of the original residual characters as a multiset, by
the characteristic polynomials and Brauer--Nesbitt.
They are therefore individually invariant under \(G\).
Transitivity on the monomial lines forces them all to be
equal. The reduced image of \(N_0\) thus consists of
scalars of order prime to \(\ell\). As \(|P|\) divides
24 and \(\ell\ge5\), the entire reduced matrix group has
order prime to \(\ell\). By Maschke's theorem its
representation is a direct sum of the four characters,
so its image is abelian. Its permutation action on the
reduced summands still surjects onto \(P\), making \(P\)
abelian and regular, of order four. All ratios of the
original residual characters are trivial on \(N_0\) and
would consequently have order dividing four. This
contradicts the ratio at \(g_0\).

It remains to prove that for each \(g\in G\) the coset
\(\rho(g)H^\circ\) contains a matrix with simple spectrum.
There are two possible decompositions under \(H^\circ\):
one four-dimensional constituent or two planes.
If there is one constituent, the matrices in \(H^\circ\)
span the full matrix algebra by Burnside's theorem. If
there are two and \(g\) preserves both, the restrictions
of \(H^\circ\) on each span the full endomorphism algebra
of its plane. In either case, in the chosen eigenbasis
each diagonal entry is nonzero somewhere on
\(\rho(g)H^\circ\). The same is true for each diagonal
entry of the inverse, by applying this observation to the
inverse coset. Since \(H^\circ\) is an irreducible variety,
we can choose a single \(x\) in the coset with all eight
entries nonzero. In particular all four diagonal cofactors
of \(x\) are nonzero.

Order the four \(w_i\) increasingly. In the elementary
symmetric coefficients \(e_j\) of the eigenvalues of
\(xD(T)\), a principal \(j\)-row minor of \(x\) is
multiplied by the power of \(T\) with exponent equal to
the sum of its \(j\) weights. Thus the \(T\)-adic
valuations of \(e_1,e_3,e_4\) are \(w_1\),
\(w_1+w_2+w_3\), and \(\sum_iw_i\), respectively.
If the leading two-row minor is nonzero, the successive
slopes of the elementary-coefficient Newton polygon
(with points \((j,v_T(e_j))\)) are the four distinct
\(w_i\), giving distinct roots. If that minor vanishes,
\(v_T(e_2)\ge w_1+w_3\). The middle edge then joins
positions one and three and has slope \((w_2+w_3)/2\).
Its residual equation is a binomial of degree two with
both terms nonzero, since the middle coefficient lies
strictly above the edge. Its roots are distinct in
characteristic \(\ell\ne2\). More explicitly, after
putting \(T=u^2\), substituting the middle slope in the
polynomial and dividing by its lowest power of \(u\),
the nonzero residual roots lift by Hensel's lemma.
The outer edges have length one and distinct slopes.
This again gives four distinct roots.

If \(g\) swaps the two planes, the coset has the block
form \(\left(\begin{smallmatrix}0&B_1\\B_2&0\end{smallmatrix}\right)\).
Using Burnside on each plane and irreducibility of the
coset as a variety, choose \(x\) for which every entry
of the two blocks is nonzero. Both blocks are invertible.
Write \(D_1,D_2\) for the two diagonal blocks of \(D(T)\).
The characteristic polynomial of \(xD(T)\) is
\[
       Y^4-\operatorname{tr}(B_1D_2B_2D_1)Y^2
                   +\det(B_1B_2)T^{\sum_i w_i}.
\]
Indeed eliminating one block in \(Y-xD(T)\) gives the
two-by-two determinant in \(Y^2\). The least exponent
in the trace is the sum of the smaller weight in each
plane, attained by a unique pair, with nonzero coefficient.
Twice this exponent is strictly smaller than \(\sum_iw_i\).
The discriminant of the quadratic in \(Y^2\) is therefore
nonzero, its two roots are nonzero, and \(\ell\ne2\)
gives four distinct roots in \(Y\).

In all cases the discriminant is nonzero for the generic
\(T\), so it is nonzero at some element of the coset.
For any open subgroup \(V\) of \(G'\), the Zariski closure
of \(\rho(V)\) contains \(H^\circ\), since finitely many
translates cover \(H\). An open neighborhood in \(G'\) or
in \(G\) contains a translate of such a subgroup of \(G'\).
The closure of its image contains the corresponding coset
of \(H^\circ\), on which the discriminant is not identically
zero. This proves both density statements. Finally the
Frobenius at the auxiliary tame prime in
Theorem~\ref{thm:commonprobes594} acts on residual
endomorphisms with all eigenvalues one, while the
cyclotomic factor there is four in \(k\). Hence the
determinant of the twisted Frobenius minus one is a unit
in \(A\). No local intertwiner can survive even after
passing to the algebraic closure of \(K_1\).
\end{proof}

\begin{corollary}[Simultaneous Frobenius choices at all cyclotomic depths]
\leavevmode\label{cor:commoncurves597}
Take any finite list of curves satisfying the theorem on
the given interval, allowing different fields, uniformizers
\(t_i\), and coefficient primes \(\ell_i\). There exist
fixed integers \(b_i\ge0\) with the following property.
For every positive integer \(N\), infinitely many good
rational primes \(q\) split in all original period fields
and all the auxiliary starting fields, satisfy
\(q\equiv1\bmod\ell_i^N\) for every \(i\), and admit choices
of Frobenius with characteristic-polynomial discriminant
of \(t_i\)-adic valuation exactly \(b_i\) on every curve.
Finite exclusions at each depth are permitted.
\end{corollary}
\begin{proof}
Use one finite abelian field \(P/\QQ\) containing the
period and starting fields as in
Corollary~\ref{cor:polarsplit593}. Let \(P_\infty\)
be its compositum with all the \(\ell_i\)-power cyclotomic
towers in the list and put \(H_* = G_{P_\infty}\).
For curve \(i\), \(H_*\) is normal in
\(G'_i=G_{L_i(\mu_{\ell_i^\infty})}\). In
\(G'_i/H_*\) the \(\ell_i\)-primary part is finite.
Indeed this abelian quotient embeds by restriction into
the product of a finite group and the finitely many
\(\ZZ_p^\times\) at the other coefficient primes \(p\).
Each such factor has finite \(\ell_i\)-primary part,
and its complementary pro-\(p\) part has no nontrivial
finite \(\ell_i\)-quotient. The same assertions hold for
closed subgroups of that product by the primary
decomposition of profinite abelian groups.

The compact image \(\rho_i(G'_i)\) has an open pro-\(\ell_i\)
subgroup, by reduction modulo \(t_i\). Its quotient by
\(\rho_i(H_*)\) is therefore both virtually pro-\(\ell_i\)
and a quotient of the preceding abelian group. It is
finite. Consequently the image of any open subgroup of
\(H_*\) has finite index in \(\rho_i(G'_i)\); its closure
in the algebraic group contains the identity component
used in the theorem. The same coset argument shows that
regular spectra for curve \(i\) form a dense open subset
of \(H_*\). Openness follows from continuity of the
discriminant. Intersect these finitely many dense opens
and fix \(g\) in the intersection. Take \(b_i\) to be
the discriminant valuation at \(g\). It is determined
by the matrices modulo \(t_i^{b_i+1}\).

For a given \(N\) adjoin to \(P\) the required
\(\ell_i^N\)-th roots of unity and the normal closures
over \(\QQ\) of the finite fields cut out by these fixed
matrix reductions over each base field. In this finite
Galois extension Chebotarev supplies infinitely many
good rational primes with Frobenius conjugate to the
image of \(g\). Choose a place above each so as to
realize that image itself. These primes split in \(P\)
and in the indicated cyclotomic fields, and the matrix
congruences at the same place give the specified
valuations on all curves. The precisions \(b_i+1\) were
fixed before \(N\) was chosen.
\end{proof}

This gives the simple spectra and the separate tame place
needed to test the irreducible ordinary curves in characteristic
\(\ell\). Transporting the Hecke relations to all components
through these curves also requires control of global
cohomology in the patching argument. The discriminant
valuations here concern a chosen finite list of curves.


\subsection{Dual Selmer tests on irreducible ordinary curves}
\label{sec:dualtests598}

Use a finite list of the absolutely irreducible curves of
Section~\ref{sec:regularcurves597}. For each one restrict
to \(G_L\), write \(A=k[[t]]\), \(K_1=k((t))\),
\(A_m=A/(t^m)\), and put \(W=\operatorname{End}_A(A^4)\)
with conjugation action through \(\rho\). Set
\(\omega=\bar\kappa_\ell\) and \(W_m=W/t^mW\). Work with
\(G_{L,S}\) for any fixed finite set containing the
coefficient places, the ramification of the representation
and the infinite places. If necessary enlarge \(S\) to contain
the ramifications of the common finite field \(P\) of
Corollary~\ref{cor:commoncurves597} and all coefficient primes
in the list. Inflation embeds the cohomology from the
smaller set, since the coefficient action already factors
through its quotient.

At an element \(g\) with \(\omega(g)=1\), define
\[
             e_g(y)=\operatorname{tr}\bigl(y(g)\rho(g)\bigr)
               \qquad \bigl(y\in Z^1(G_{L,S},W_m(1))\bigr).
\]
The notation \((1)\) multiplies conjugation by \(\omega\).
The evaluation vanishes on coboundaries, by cyclicity of
trace. The same definition makes sense with \(K_1\)
coefficients. For the even-order case let \(\phi\) cut out
\(D_0/\QQ\), so that \(\phi=\tau\) on \(G_L\). Say that
the curve has this quadratic self-twist if
\(\rho\simeq\phi\rho\) over \(\overline{K_1}\) on \(G_L\).

\begin{proposition}[The quadratic exception to tests at completely split primes]
\leavevmode\label{prop:splitexception598}
On \(G_{P_\infty}\) of Corollary~\ref{cor:commoncurves597},
the curve representation remains absolutely irreducible
in the odd-order case and in every even-order case without
the indicated quadratic self-twist. In the exceptional
case it splits there into two inequivalent absolutely
irreducible planes. The subgroup \(G_{LD_0}\) preserves
both and the other coset of \(G_{LD_0}\) in \(G_L\)
interchanges them. For a cocycle valued in the two
off-diagonal matrix blocks, \(e_g\) vanishes identically
on \(G_{P_\infty}\).
\end{proposition}
\begin{proof}
Write \(H_*=G_{P_\infty}\). The proof of
Corollary~\ref{cor:commoncurves597} shows that
\(\rho(H_*)\) has finite index in
\(\rho(G_{L(\mu_{\ell^\infty})})\). It therefore contains a
positive power of the matrix \(\rho(a)\) with simple spectrum
constructed in Theorem~\ref{thm:regularcosets597}.
Restriction of the irreducible representation on \(G_L\)
to \(H_*\) is semisimple, and its constituents have
multiplicity one by this simple spectrum. If it is
reducible, the abelian quotient \(G_L/H_*\) acts regularly
and transitively on two or four constituents. A nontrivial
quadratic character of this permutation group gives a
self-twist of \(\rho\) by the sign change on constituents
used in Theorem~\ref{thm:regularcosets597}. This quadratic
character survives reduction and must permute the four
residual characters by multiplication.

In the odd-order case they are, up to a common twist,
\(1,\chi,\chi^2,\chi^3\), with \(\chi\) of order \(h\)
divisible by \(2m\ge6\), as seen at \(g_0\).
A quadratic permutation character must
belong to \(\langle\chi\rangle\) by looking at the first
entry. Its nontrivial possibility adds \(h/2\) to the
exponents modulo \(h\). This does not preserve the set
\(\{0,1,2,3\}\): for \(h=6\) the image contains four,
and for \(h\ge8\) the image of zero is outside the set.
In the even-order case the four are \(1,\chi,\phi,\chi\phi\),
up to a common twist. Of the three nonidentity candidates
only \(\phi\) has order two since \(\chi^4\ne1\).
As \(\ell\) is odd, the quadratic sign over \(K_1\) is
uniquely determined by its reduction. This proves the
irreducibility assertions.

If there is the self-twist, scale an intertwiner to have
square one, after finite coefficient extension. Its square
is scalar by Schur's lemma. Its eigenspaces are interchanged
by each element with \(\phi=-1\), hence are two planes,
and \(G_{LD_0}\) preserves them. The group \(H_*\) is
in this latter subgroup. Its algebraic image closure
contains the identity component of the cyclotomic-kernel
image. The proof of Theorem~\ref{thm:regularcosets597}
excluded a decomposition of that component into four lines
and proved multiplicity one. Thus the two planes are
irreducible and inequivalent even on \(H_*\). Multiplying
an off-diagonal matrix by the block diagonal \(\rho(g)\)
has trace zero, giving the last assertion.
\end{proof}

\begin{theorem}[A bounded loss for the entire finite-coefficient dual group]
\leavevmode\label{thm:dualcurve598}
For each curve there are a fixed number \(h\) of test places
and a constant \(C\) depending on that curve and \(S\).
For every \(m,N\ge1\) one may realize the tests at \(h\)
distinct good rational primes \(q_j\) splitting completely
in \(L\), with \(q_j\equiv1\bmod\ell^N\), and with chosen
Frobenius elements \(s_j\) at places above them. Then
\[
  \operatorname{length}_A
  \ker\left(H^1(G_{L,S},W_m(1))
        \xrightarrow{(e_{s_j})_j} A_m^{\,h}\right)\le C .
\]
The matrices \(\rho(s_j)\bmod t^m\) agree with the
reductions of fixed matrices with simple spectrum over
\(\overline{K_1}\). For curves without the quadratic
exception all the rational primes can be the same across
the list, split in all period and starting fields and
congruent to one modulo every coefficient prime to the
power \(N\). In the quadratic exception the tests for
the diagonal blocks likewise use these common primes.
The additional tests for the off-diagonal blocks can be
taken inert in \(D_0/\QQ\).
\end{theorem}
\begin{proof}
First use \(K_1\)-coefficients. For a curve without the
exception the span of \(\rho(H_*)\) is
\(\operatorname{Mat}_4(K_1)\) by the preceding proposition
and Burnside's theorem. The quotient \(G_L/H_*\) is abelian,
and \(\omega\) factors through it and is nontrivial.
Indeed at a place above \(\ell\) the ramification index
over \(\QQ_\ell\) is at most two, so the inertia image
in \(\mathbb F_\ell^\times\) has order at least
\((\ell-1)/2>1\). Apply Theorem~\ref{thm:derivation367}
to the matrix algebra, identified with its dual by trace.
It says that vanishing of all the evaluations on \(H_*\)
forces a coboundary. The theorem applies to the image
of \(H_*\) in the restricted-ramification group as well,
by the same span and quotient assertions. By continuity
it suffices to test on any dense subset, in particular
the simultaneous regular subset in
Corollary~\ref{cor:commoncurves597}.

In the exceptional case extend the curve fraction field
finitely to define the two planes. The block diagonal
space and the space of the two off-diagonal blocks are
stable under conjugation by \(G_L\). For the diagonal
space first restrict to \(G_{LD_0}\), an injective map
on first cohomology since restriction followed by
corestriction multiplies by two. On this subgroup the
representation takes values in the units of the product
of two full matrix algebras. Its values on \(H_*\) span
that product by the inequivalence and absolute
irreducibility on \(H_*\). The character \(\omega\)
is still nontrivial over \(LD_0\), since \(D_0\) splits
at \(\ell\). The same derivation theorem applies to
this product with its trace dual. Thus tests at elements
of \(H_*\) detect the diagonal cohomology, and their
values on the off-diagonal part are zero.

For the remaining off-diagonal part apply the derivation
theorem on \(G_L\) with normal subgroup
\(G'=G_{L(\mu_{\ell^\infty})}\). Here its span is the
full matrix algebra by Theorem~\ref{thm:regularcosets597}.
On \(G'\cap G_{LD_0}\) all the off-diagonal evaluations
are already zero. Therefore the evaluations on the open
subset of \(G'\) with \(\phi=-1\) separate its cohomology
classes. Regular elements are dense there. If a class
with both parts passes all the diagonal tests, subtract
the diagonal coboundary and apply these remaining tests.
This proves separation of all the classes over the
finite coefficient extension, hence over \(K_1\):
the equations for a coboundary are linear over \(K_1\),
so solvability descends to \(K_1\).

We pass from this separation to a bound at every
precision, including classes which do not lift from
\(A_m\) to \(A\). We give the details since reduction
of a collection of tests on the rational cohomology
alone would not establish this assertion. The action
on \(W(1)\) and the continuous cocycles on \(W_m(1)\)
for every \(m\) factor together through a fixed
topologically finitely generated profinite quotient
of \(G_{L,S}\). In fact take the open normal kernel
of the action modulo \(t\). Each cocycle combined with
the linear action has pro-\(\ell\) image on this kernel.
The latter's maximal pro-\(\ell\) quotient is finitely
generated, since its \(H^1\) with constant
\(\mathbb F_\ell\)-coefficients is finite-dimensional.
For this last fact, pass to the corresponding number
field, adjoin \(\mu_\ell\) by a prime-to-\(\ell\)
extension, and use the \(S\)-units and the finite
\(S\)-class group in the Kummer sequence. A restricted
ramification quotient only decreases the dimension by
inflation. The kernel of the maximal pro-\(\ell\)
quotient is characteristic in the open normal subgroup,
so dividing it out gives the desired common quotient.

Take \(u\) profinite generators of this quotient. On the free
profinite group on their lifts, any assigned \(u\) cocycle
values extend continuously by the semidirect product law.
A relator has cocycle value an \(A\)-linear expression in
these \(16u\) scalar variables. The rows of all such
expressions are generated by finitely many of them over
the Noetherian ring \(A\). The resulting finite matrix
\(M\) cuts out all cocycles over \(A\) and over every
\(A_m\), since tuples modulo \(t^m\) lift before the
relations are imposed. The matrix \(D\) of coboundaries
is given by applying the generators minus one. Over
\(K_1\) the quotient \(\ker M/\operatorname{im} D\)
after scalar extension is finite-dimensional, and the
separation already proved allows a finite list of regular
test elements. One may take their union on \(H_*\) for
all curves and the additional \(\phi=-1\) list separately
for each exceptional curve. Enlarging the common list
only adds tests, and all its elements can be chosen
regular for every curve.

Adjoin the evaluation rows for these tests to \(M\), using
fixed group elements, and call the new matrix \(M'\).
Then \(\ker(M'\otimes K_1)=\operatorname{im}(D\otimes K_1)\).
Let \(b_1\) be the maximum valuation of a nonzero elementary
divisor of \(M'\), or zero if there are none, and choose
\(b_2\) with \(t^{b_2}\ker M'\subset\operatorname{im}D\)
over \(A\). Such \(b_2\) exists by finite generation.
In Smith coordinates \(M'v=0\bmod t^m\) implies
\(t^{b_1}v\in\ker M'+t^m A^{16u}\). Thus each class
passing the tests modulo \(t^m\) is killed by
\(t^{b_1+b_2}\). Its entire kernel has at most \(16u\)
generators: it is a quotient of a submodule of
\(A_m^{16u}\), and the bound follows from the elementary
divisor theorem. Its length is therefore at most
\(16u(b_1+b_2)\), independently of \(m,N\).

Finally there are only finitely many cocycles at the
given precision, by the generator presentation. Take
the joint finite quotients recording their values via
the semidirect product, the curve matrices modulo
\(t^m\), and, if necessary, the matrices to higher fixed
precision detecting the nonzero discriminants.
For a common test take all these quotients over the
fields in the list and their normal closures over
\(\QQ\); adjoin \(P\) and the \(\ell_i^N\)-th roots
of unity. Chebotarev matches the chosen element of
\(H_*\) in this joint finite quotient with Frobenius at
a good rational prime and a choice of place. For a test
with \(\phi=-1\), do the same over the single field
with \(L,D_0,\mu_{\ell^N}\) included, matching its chosen
element of \(G'\). Then the prime is split in \(L\),
congruent to one modulo \(\ell^N\) and inert in \(D_0\).
Repeating with finite exclusions gives distinct primes.
All evaluation rows modulo the indicated precision are
identical to the fixed tests. This proves the theorem.
\end{proof}

\begin{corollary}[An ordinary Selmer length bound at the auxiliary levels]
\leavevmode\label{cor:ordinarycurve598}
Suppose the curve has an integral polarized adjoint \(V\)
over \(F\) restricting to \(W\) over \(L\), with the real
involutions acting by negative adjunction for nondegenerate
symmetric forms. This includes the fixed-multiplier
symplectic curves over \(F\) with odd multiplier, by using
the rank-four polarized adjoint. Take the old finite set
stable under complex conjugation, including the
ramifications of this action. On \(V_m=V/t^mV\) impose
unramified cohomology at the old finite places away from
\(\ell\) and the image of flag-preserving cohomology
at \(v\mid\ell\). At one place of \(F\) from each of
the \(h\) rational test primes allow full local cohomology.
For the choices of the theorem, the resulting Selmer group
and its dual with annihilating local conditions satisfy
\[
  \operatorname{length}_A\operatorname{Sel}(F,V_m)
       \le 4\bigl([F:\QQ]+h\bigr)m+O(1),\qquad
  \operatorname{length}_A\operatorname{Sel}^*(F,V_m^\vee(1))
       =O(1).
\]
Both constants are independent of \(m,N\).
\end{corollary}
\begin{proof}
The trace form identifies the polarized adjoint with its
dual, also for the negative-adjunction action. Composing
the perfect \(A_m\)-pairing with the trace to \(\mathbb F_\ell\)
of the coefficient of \(t^{m-1}\) gives its finite-group
dual pairing. At a new test place the dual local condition
is zero. Restriction to \(L\) is injective by
corestriction since \([L:F]=2\). The restricted dual
class is zero locally at every new place over the chosen
ones and unramified away from the old set at all other
places. It therefore factors through the original
\(G_{L,S}\): at the new places inertia acts trivially on
the coefficients and the cocycle is zero on inertia.
Its Frobenius evaluations vanish there. The dual estimate
follows from Theorem~\ref{thm:dualcurve598}.

For the first estimate apply the global Selmer dimension
formula and divide logarithms of cardinalities by
\(\log|k|\) to obtain lengths over \(A\). The global
invariants of \(V\otimes K_1\) vanish, since the centralizer
over \(L\) is scalar and real conjugation negates scalars.
Consequently \(H^0(F,V_m)\) has bounded length: finitely
many matrices \(g-1\) already have zero common kernel
over \(K_1\), so their stacked Smith form gives a uniform
bound at all precisions. The subtracted global twisted
dual term can be discarded for an upper bound. At the old
places away from \(\ell\), unramified cohomology has the
same length as \(H^0\), since the kernel and cokernel of
Frobenius minus one on the finite inertia invariants
have the same order.

At \(v\mid\ell\), intersect the flag over the fraction
field with the stable lattice. Its saturated steps give
a full flag of direct summands over \(A\). Its preserving
submodule \(B_v\subset V\) has rank ten and free quotient.
On its dual twist the graded characters are
\(\omega\theta_{v,j}/\theta_{v,i}\) for \(i\le j\).
At different indices they are nontrivial by the
separation on \(I_v^0\); at equal indices \(\omega\) is
nontrivial by the ramification index bound at \(\ell\).
Hence the dual twist has zero invariants over \(K_1\)
and bounded invariant length at finite precision by the
same matrix argument. Local Euler characteristic and Tate
duality give
\[
  h^1(F_v,B_v/t^m)-h^0(F_v,B_v/t^m)
        =10[F_v:\QQ_\ell]m+O(1),
\]
where \(h^i\) denotes length. Its \(H^0\) injects into that
of \(V_m\), so the image condition contributes at most
this amount. At a real place the invariant skew-adjoint
matrices have rank six, giving the term \(-6m\).

At a new place \(q_j\) the unramified adjoint action
is ordinary matrix conjugation and \(q_j=1\) in \(k\).
The full local \(h^1-h^0\) is the length of the twisted
dual invariants by local Euler characteristic. This is
the length of the centralizer of the indicated Frobenius
matrix modulo \(t^m\). It equals \(4m+O(1)\) since the
matrix agrees modulo \(t^m\) with a fixed regular
semisimple matrix. Indeed its commutator map has rank
twelve over \(K_1\). Its twelve nonzero elementary
divisors contribute at most their total valuation in
addition to the four free kernel terms. Summing all
places, the terms over \(\ell\) and at infinity give
\((10-6)[F:\QQ]m\), and the \(h\) new places give at
most \(4hm+O(1)\). This proves the bound.
\cite{XII-IX-V-milne98}
\end{proof}

The constants in these estimates belong to curves over
\(k[[t]]\). Section~\ref{sec:missing} requires a bound
on integral cohomology at the prescribed rational
weights over \(E_x\). Passing to those weights involves
the support of the ordinary Hecke module through the
curve primes and the two-constituent cases.


\subsection{Rank-one local conditions at primes split in the period fields}
\label{sec:projectiveprobes599}

Let \(P\), \(P_\infty\), \(A=k[[t]]\), \(A_m=A/(t^m)\),
\(\omega=\bar\kappa_\ell\), and \(W=\operatorname{End}_A(A^4)\)
be as in Section~\ref{sec:dualtests598}. At a prime split
in the quadratic inducing field, taking the ordinary trace
can discard the whole off-diagonal part of a cocycle.
We can recover its matrix entries by changing the local
condition at the test prime. We do this over the original
interval for all the curves of Corollary~\ref{cor:commoncurves597}.

\begin{lemma}[Restriction to the kernel of the adjoint action]
\leavevmode\label{lem:projectivekernel599}
Write \(G=G_{L,S}\), where \(S\) includes all places
ramified in \(P_\infty/L\). Let \(J\) be the common kernel
of the maps from \(G\) to \(\operatorname{Gal}(P_\infty/L)\)
and to \(\operatorname{Aut}_A(W)\) arising from conjugation
by \(\rho\). For every \(m\ge1\), restriction gives an
injection
\[
 H^1(G,W_m(1))\ \longrightarrow\
          \operatorname{Hom}_{\mathrm{cont}}(J,W_m),
 \qquad W_m=W/t^mW .
\]
\end{lemma}
\begin{proof}
Choose one place of \(L\) above \(\ell\) and let \(I\)
be its inertia image. The four residual characters are
equal on \(I\), by Theorems~\ref{thm:tametower571}
and~\ref{thm:eventower572}. This equality persists under
the common scalar twist. In a composition-series basis
over a finite residue extension the residual inertia
matrices are scalar times upper unitriangular. Their
conjugation action has \(\ell\)-power order. The kernel
of reduction in \(\operatorname{Aut}_A(W)\) is pro-\(\ell\),
as successive congruence quotients are additive matrix
groups over \(k\). Thus the adjoint image of \(I\) is
pro-\(\ell\). On the other hand \(\omega(I)\ne1\), since
the local ramification index divides two and
\((\ell-1)/2>1\). There is an element \(z\in I\) acting
trivially on \(W\) with \(\omega(z)\ne1\). Otherwise
\(\omega|_I\) would factor through a pro-\(\ell\) group
and have both prime-to-\(\ell\) and \(\ell\)-power order.

The group \(G/J\) embeds in the product of the adjoint
image and the abelian group \(\operatorname{Gal}(P_\infty/L)\).
Consequently the image of \(z\) is central in \(G/J\).
It acts on \(W_m(1)\) by \(\lambda=\omega(z)\) with
\(\lambda-1\) a unit. If \(y\) is a cocycle on \(G/J\),
evaluating it at \(zg=gz\) gives
\[
                 (\lambda-1)y(g)=(g-1)y(z).
\]
This makes \(y\) the coboundary with witness
\((\lambda-1)^{-1}y(z)\). On \(J\) both the adjoint
action and \(\omega\) are trivial, the latter because
\(\mu_\ell\subset P_\infty\). A cocycle restricting to
zero on \(J\) descends to \(G/J\), proving the lemma.
\end{proof}

\begin{theorem}[One matrix coefficient per completely split test prime]
\leavevmode\label{thm:splitprojectors599}
Fix a finite list of the curves, including any with a
quadratic self-twist. For curve \(i\) there is a finite
list of rank-one idempotents \(D_{ij}\) in
\(\operatorname{Mat}_4(A_i)\), of length \(h_i\). Given
positive precisions \(m_i,N_i\), each index \((i,j)\)
can be assigned a different good rational prime \(q_{ij}\).
Every such prime splits in all the interval fields \(E_x\)
and in all starting fields, and is one modulo
\(\ell_a^{N_a}\) for every curve \(a\). For the assigned
curve choose a place of \(L_i\) above \(q_{ij}\) with
Frobenius \(s_{ij}\). Its matrix is scalar modulo
\(t_i^{m_i}\), and the following map is injective:
\[
 H^1(G_{L_i,S_i},W_{i,m_i}(1))
   \longrightarrow A_{i,m_i}^{\,h_i},\qquad
 y\longmapsto \bigl(\operatorname{tr}(D_{ij}y(s_{ij}))\bigr)_j .
\]
The number of tests and the idempotents are independent
of both precisions.
\end{theorem}
\begin{proof}
Work first with one curve in an integral basis. The sixteen
idempotents \(E_{aa}\) and \(E_{aa}+E_{ab}\), \(a\ne b\),
are rank one and span the matrix algebra over \(A\):
subtracting \(E_{aa}\) gives \(E_{ab}\). The pairing
\((D,Y)\mapsto\operatorname{tr}(DY)\) is perfect in
this basis, at every precision. Thus the tests
\(\operatorname{tr}(D y(n))\) for these idempotents
and all \(n\in J\) detect all classes by the lemma.

It is possible to choose finitely many pairs \((n,D)\)
once for all precisions. Use the common quotient with
\(u\) profinite generators constructed in the proof of
Theorem~\ref{thm:dualcurve598}. It carries the action and
all cocycles at every precision. For each \(n\in J\) choose
a lift of its image to the free profinite group on these
generators. Evaluation at that lift is \(A\)-linear in
the \(16u\) generator coordinates. Indeed the cocycle law
is linear in those coordinates on words, and on profinite
limits this follows by continuity, or by evaluating the
coordinate basis cocycles in the compact semidirect
product. Reduction gives the same row on tuples over
\(A_m\); the relators act trivially for every actual
cocycle. Taking the trace with \(D\) gives one row in
\(A^{16u}\). The submodule generated by all the rows is
finitely generated, since \(A\) is Noetherian. A finite
subset of the rows generates it, by expanding each of
a finite set of generators in the original rows. Their
reductions therefore suffice at every \(m\). All tests
vanish on coboundaries since \(J\) acts trivially.

Lift each chosen \(n\) to the absolute Galois group.
It fixes \(P_\infty\) because \(P_\infty/L\) is unramified
outside \(S\). Its matrix is scalar, since a matrix commuting
with all matrix units over \(A\) is scalar. For the given
precisions form a finite Galois extension of \(\QQ\)
containing \(P\), the required roots of unity and all
the normal closures of the finite extensions recording
the curve matrices and the cocycles at their respective
precisions. There are only finitely many cocycles, as
each is determined by a finite tuple over a finite ring.
As in Theorem~\ref{thm:dualcurve598}, adjoining the
semidirect-product quotients records actual cocycle
values. Chebotarev gives infinitely many good rational
primes admitting a place with Frobenius equal to the
chosen image of \(n\) in this finite group. Choose
distinct primes successively. They have the asserted
splitting, congruences and Frobenius matrices. At an
assigned place the twisted adjoint action of Frobenius
is trivial modulo \(t^m\), so the matrix entry test
descends to cohomology. Its value agrees with the
chosen row. This proves injectivity.
\end{proof}

\begin{corollary}[Vanishing of the dual Selmer group with one tame direction]
\leavevmode\label{cor:rankonedual599}
For a curve with the integral polarization of
Corollary~\ref{cor:ordinarycurve598}, keep its old local
conditions over \(F\). At the selected place over each
of its own \(h\) assigned primes, allow arbitrary
unramified classes and the single tame direction of
the indicated idempotent \(D\). Use the unramified
condition at all other new places. For these
finite-coefficient Selmer conditions one has
\[
  \operatorname{Sel}^{*}(F,V_m^\vee(1))=0,\qquad
  \operatorname{length}_A\operatorname{Sel}(F,V_m)
       \le \bigl(4[F:\QQ]+h\bigr)m+O(1).
\]
The constant depends on the curve and its old local
conditions, uniformly in \(m,N\) and the primes chosen
by the theorem.
\end{corollary}
\begin{proof}
At its assigned place the identification with \(W_m\)
uses the chosen prime of \(L\) above it. The inertia
action is trivial and Frobenius acts trivially on the
adjoint and on \(\mu_\ell\). Local class field theory
gives two independent scalar additive characters modulo
\(\ell\), one unramified and one tame: a uniformizer
gives the unramified coordinate, the cyclic residue
units give the tame coordinate since \(\ell\mid q-1\),
and the principal units are pro-\(q\). Fix a generator
of each coordinate. With coefficients \(W_m\), the
local condition consists of unrestricted unramified
values and tame values in \(A_mD\). This line is free
of length \(m\) because \(D\) has a unit entry.

In the local Tate pairing two unramified classes pair
to zero, since the procyclic residue group has
\(\ell\)-cohomological dimension one. The pairing of
the scalar unramified and tame generators is nonzero
by perfectness on the two-dimensional local character
space, after identifying \(\mu_\ell\) with
\(\mathbb F_\ell\) at the place. For the matrix
coefficients this pairing multiplies the matrix trace
pairing by the same nonzero scalar and applies the
perfect additive functional on \(A_m\) used in
Corollary~\ref{cor:ordinarycurve598}. Thus annihilating
all unramified values forces a dual class to be
unramified. For its Frobenius value \(Y\), annihilating
all multiples of the tame direction is equivalent to
\(\operatorname{tr}(DY)=0\) in \(A_m\). At the remaining
good places the annihilator of unramified cohomology is
unramified cohomology for the dual. \cite{XII-IX-V-milne98}

A dual Selmer class therefore restricts over \(L\) to
an old global class, with every test of
Theorem~\ref{thm:splitprojectors599} zero. Its restriction
vanishes. Restriction to \(L\) is injective on global
first cohomology by corestriction, since \([L:F]=2\)
is invertible. This proves the first assertion. At
an assigned place the local condition has length
\(16m+m\), while the invariants have length \(16m\);
its contribution to the Selmer dimension formula is
\(m\). All other new places have unramified conditions
and contribute zero. The old finite and real places
and the global invariants satisfy the estimates proved
in Corollary~\ref{cor:ordinarycurve598}, giving
\(4[F:\QQ]m+O(1)\) as their combined upper bound.
This proves the second assertion.
\end{proof}

\begin{proposition}[The tame lifting equation at a scalar Frobenius]
\leavevmode\label{prop:rankonelift599}
At an assigned prime let a tame inertia generator act
as \(T=1+(\delta-1)D\), where \(D^2=D\) and
\(\delta^{q-1}=1\). For a Frobenius lift with matrix
\(\Phi\), the local tame relation is equivalent to
\[
                    (\delta-1)[\Phi,D]=0.
\]
At scalar residual Frobenius, an unrestricted first-order
change in Frobenius and a first-order change in this
inertia character can be obstructed at second order.
\end{proposition}
\begin{proof}
On the image of \(D\), \(T\) is multiplication by
\(\delta\); on its kernel it is the identity. Hence
\(T^q=T\), and \(\Phi T\Phi^{-1}=T^q\) is the displayed
equation. In a basis with \(D_0=E_{11}\) at the closed
point, let the residual Frobenius be \(\lambda I_4\).
Over \(k[s,u]/(s,u)^2\) take \(\delta=1+s\),
\(\Phi=\lambda I_4+uE_{12}\), and \(D=D_0\). These
matrices satisfy the relation. In \(k[s,u]/(s,u)^3\),
the term of degree two in the left side of the displayed
equation is \(su[E_{12},D_0]=-suE_{12}\), which is
nonzero. Corrections of degree two cannot alter it,
since the degree-zero Frobenius is scalar. The order
constraint on \(\delta\) allows the displayed choice
in this ring since \(\ell\mid q-1\) and \(\ell\ge5\).
Thus the two given first-order changes do not lift together.
\end{proof}

The cohomological vanishing applies on the same period
interval, even on a curve induced from the quadratic
field. Its local conditions vary with the coefficient
precision. To deduce the strict cohomology estimate at
the prescribed characteristic-zero weights one still
needs an integral comparison through those weights.
The tame lifting equation gives an explicit constraint
on a deformation construction using these local conditions.



\subsection{Exact cyclotomic depth at the split Selmer primes}
\label{sec:exactselmer614}

Fix a finite list of the ordinary curves in
Corollary~\ref{cor:commoncurves597} on one of the period intervals.
On curve \(i\) the representation has dimension four and coefficient
ring \(k_i[[t_i]]\), where \(k_i\) is a finite field of odd
characteristic \(\ell_i\). Let \(P/\QQ\) be the finite Galois field
used there, containing all the period and starting fields on the
interval. Theorem~\ref{thm:splitprojectors599} assigns to each
curve finitely many rank-one probes. At a precision \(m_i\ge1\),
they detect its full finite-coefficient dual cohomology by evaluating
one matrix entry at each assigned prime. The splitting requirements
in that theorem prescribe lower bounds on the cyclotomic depths.
For the integral tame characters one needs the sizes of the actual
cyclic groups at those primes.

\begin{theorem}[Split Selmer tests with prescribed exact depths]
\leavevmode\label{thm:exactselmer614}
For each distinct coefficient prime \(\ell\) put
\begin{align*}
 d_\ell={}&v_\ell([P:\QQ])+
 \max\Bigl(\{0\}
   \cup\{\lceil\log_\ell(17m_i)\rceil:\ell_i=\ell\}\\[-2pt]
 &\hspace{87pt}{}\cup
      \{v_\ell(|\mathrm{GL}_{17}(k_i)|):\ell_i\ne\ell\}\Bigr).
\end{align*}
Choose one integer \(N_\ell\ge1+d_\ell\) per distinct prime.
All the probes of Theorem~\ref{thm:splitprojectors599} can be
realized by different good rational primes \(q\) such that
\[
              v_\ell(q-1)=N_\ell
                 \quad\hbox{for every coefficient prime }\ell .
\]
The primes split in \(P\). The curve matrices and the dual
cocycle evaluations at their assigned precisions are the same
as in that theorem. Any further finite list of primes can be
excluded.
\end{theorem}
\begin{proof}
Let \(K/\QQ\) contain \(P\) and be generated over it by all the
conjugates of the fields recording the given matrices and the
cocycles at these precisions. We use the smallest such compositum.
There are finitely many cocycles at each precision by the finite
generator argument in Theorem~\ref{thm:splitprojectors599}.
One cocycle \(y\) with values in the twisted endomorphisms of
the four-dimensional representation is recorded by the matrices
\[
        \begin{pmatrix}a(g)&y(g)\\0&1\end{pmatrix}
        \in\mathrm{GL}_{17}(k_i[[t_i]]/(t_i^{m_i})),
\]
where \(a(g)\) is the action on the sixteen endomorphism coordinates.
The cocycle identity is precisely the multiplication law for these
matrices. The original four-dimensional matrices can also be
embedded in this group by adjoining an identity block. Since
\(P\) contains every conjugate of the starting fields,
\(\operatorname{Gal}(K/P)\) embeds by restriction into a product
of these matrix groups, including their conjugate copies.

For \(\ell_i=\ell\), an element \(g\) of \(\ell\)-power order acts
unipotently over \(k_i\) on the underlying vector space of dimension
\(17m_i\). Indeed its minimal polynomial divides a power of
\(X-1\) in characteristic \(\ell\). The nilpotent matrix \(g-1\)
has nilpotence index at most \(17m_i\); taking the power
\(\ell^{\lceil\log_\ell(17m_i)\rceil}\) gives the identity.
For \(\ell_i\ne\ell\), the kernel of reduction of the matrix
group is an \(\ell_i\)-group by its successive congruence
filtration. An \(\ell\)-element therefore has the same order
as its reduction in \(\mathrm{GL}_{17}(k_i)\). The order bound
in a direct product uses the largest of these \(\ell\)-exponents.
For an \(\ell\)-element of \(\operatorname{Gal}(K/\QQ)\),
first raising it to \(\ell^{v_\ell([P:\QQ])}\) puts it in
\(\operatorname{Gal}(K/P)\). Thus its order divides
\(\ell^{d_\ell}\). This bound passes to \(\ell\)-elements of
every quotient, since a Sylow subgroup of the quotient is the
image of a Sylow subgroup of the original group.

Set \(C=\QQ(\mu_{\prod_\ell\ell^{N_\ell+1}})\).
By the Chinese remainder theorem the tuple
\((1+\ell^{N_\ell}\bmod\ell^{N_\ell+1})_\ell\) defines an
automorphism of \(C\). Each individual component fixes \(C\cap K\).
In fact the image of \(1+\ell\ZZ_\ell\) in the Galois group
of this intersection has exponent dividing \(\ell^{d_\ell}\)
by the preceding quotient bound. For odd \(\ell\) the logarithm
maps \(1+\ell\ZZ_\ell\) isomorphically to \(\ell\ZZ_\ell\),
and \(1+\ell^{N_\ell}\) is consequently an
\(\ell^{d_\ell}\)-th power in that group. Here the local factor
acts trivially at all the other factors of the cyclotomic modulus.

For each selected probe take the element \(n\) from the proof
of Theorem~\ref{thm:splitprojectors599}. It fixes \(P\) and the
full cyclotomic towers for all the coefficient primes. Its image
on \(K\) and the indicated automorphism of \(C\) agree on
\(C\cap K\). They extend to an automorphism of \(KC\), by the
fiber-product description for two finite Galois extensions.
Chebotarev's theorem supplies infinitely many rational primes
with a choice of place realizing this automorphism as arithmetic
Frobenius. On \(K\) the required matrices and cocycle values
are those of \(n\). On \(C\) the action is given by the
residue class of \(q\), proving the exact depth. Choosing the
places successively gives distinct primes and allows the stated
exclusions.
\end{proof}

We next compute the local integral intersection for these probes.
Let \(\mathcal O\) be a complete discrete valuation ring finite
over \(\ZZ_\ell\), with uniformizer \(\pi\) and residue field
containing \(k_i\) at a probe of characteristic \(\ell=\ell_i\).
Take the residual scalar Frobenius and trivial residual inertia
there. In a fixed basis lift its rank-one idempotent to
\(D=\operatorname{diag}(1,0,0,0)\). Such a basis is obtained
by lifting bases of the residual image and kernel. Write
\(\Phi\) for the freely variable framed Frobenius matrix lifting the
given scalar, and require the tame generator to act by
\(1+uD\), with its character of order dividing
\(M=\ell^{v_\ell(q-1)}\). The variable \(u\) reduces to zero.
Here these are local integral parameters, independent of the
curve parameter \(t_i\).

\begin{proposition}[The integral intersection at one test prime]
\leavevmode\label{prop:tamebranches614}
Let \(Y\) denote the ten coordinates in the diagonal blocks
of \(\Phi\) for \(D\), translated by a lift of its residual value.
Let \(X=(X_1,\ldots,X_6)\) be its two off-diagonal blocks.
Set \(F_M(u)=((1+u)^M-1)/u\) and \(S=\mathcal O[[Y,X,u]]\).
The complete ring for the specified local data is
\[
 R=S/(uF_M(u),uX_1,\ldots,uX_6).
\]
It is reduced and flat over \(\mathcal O\). Its unramified
branch \(S/(u)\) and its union of ramified branches
\(S/(F_M,X)\) meet along \((\mathcal O/M\mathcal O)[[Y]]\).
Precisely, \(R\) is the fiber product of the two branch rings
over this intersection.
\end{proposition}
\begin{proof}
Invertibility of \(\Phi\) holds over every local coefficient
algebra because its residual determinant is a unit. The order
condition on inertia is \(uF_M(u)=0\). It also implies
\((1+uD)^q=1+uD\), as \(M\mid q-1\) and \(D\) is idempotent.
The tame relation is consequently \(u[\Phi,D]=0\), precisely the
six equations \(uX_a=0\) by Proposition~\ref{prop:rankonelift599}.
Conversely these equations give a homomorphism of the tame
decomposition group with the assigned generator images. This
can be checked first on finite local coefficient algebras, where
the inertia image is cyclic of order dividing \(M\) and the
Frobenius image has finite order. Passing to inverse limits
gives the stated complete ring.

The polynomial \(F_M\) is monic of degree \(M-1\) and reduces
to \(u^{M-1}\) in the residue field. Division by this distinguished
polynomial makes \(S/(F_M,X)\) free of rank \(M-1\) over
\(\mathcal O[[Y]]\). Multiplication by \(u\) on this free
module is injective. Over the fraction field the constant
coefficient \(F_M(0)=M\) makes \(u\) invertible, and freeness
gives the injection before passage to that field. Thus
\((u)\cap(F_M,X)=u(F_M,X)\) in \(S\). The exact sequence for
the intersection of two ideals identifies \(R\) with the asserted
fiber product, since \(S/(u,F_M,X)=(\mathcal O/M\mathcal O)[[Y]]\).
Both branch rings are \(\mathcal O\)-torsion free, so \(R\) embeds
in a torsion-free module and is flat over this discrete valuation
ring. The unramified branch is a power series domain. The other
branch embeds by the same freeness in the algebra obtained over
the fraction field of \(\mathcal O[[Y]]\), where \(F_M\) has the
distinct roots \(\zeta-1\) for the nontrivial \(M\)-th roots of
unity over an algebraic closure. Both branches, and hence their
subring \(R\), are reduced.
\end{proof}

\begin{corollary}[The exact cost of isolating unramified inertia]
\leavevmode\label{cor:selmercongruence614}
Take the \(h\) probes for a single curve and one common
\(\mathcal O\). At each probe fix an integral Frobenius
specialization commuting with its specified projector and lifting
the residual scalar. On this product of local character rings,
consider integral functions vanishing on every characteristic-zero
tuple at which at least one inertia character is nontrivial.
Their values at the all-unramified tuple form the ideal
\[
        \left(\prod_{a=1}^h M_a\right)\mathcal O,\qquad
        M_a=\ell^{v_\ell(q_a-1)} .
\]
The same ideal is obtained by first taking functions on the full
product of the framed local charts that vanish on all such tuples,
and then evaluating at the fixed unramified point. For the exact
depth \(N_\ell=1+d_\ell\) in Theorem~\ref{thm:exactselmer614}
its colength is \(h\,v_\pi(\ell)(1+d_\ell)\).
For the fixed list and interval this is
\(O(\log(2+\max_{\ell_i=\ell} m_i))\) as the precisions vary.
\end{corollary}
\begin{proof}
At the specified Frobenius values all the \(X_a\) are zero.
The remaining ring for a single inertia generator \(z=1+u\)
is \(\mathcal O[z]/(z^M-1)\). Indeed its presentation in the
complete ring by the distinguished polynomial \(uF_M(u)\)
has the same basis of \(M\) powers. The ring for the product
is therefore \(\mathcal O[G]\) for \(G=\prod_a\ZZ/M_a\ZZ\).
Its characteristic-zero points are the characters of \(G\)
after an extension containing their values. If a group-ring
element takes value \(b\) at the trivial character and zero
at all the others, Fourier inversion makes every coefficient
equal to \(b/|G|\). Integral coefficients force
\(b\in |G|\mathcal O\); conversely \(\sum_{g\in G}g\) attains
\(|G|\). Any separating function on the full charts restricts
to such a function on this product. The element
\(\prod_aF_{M_a}(u_a)\) of the full charts vanishes at every
characteristic-zero point with some nontrivial inertia and
attains \(\prod_a M_a\) at the unramified point. This gives both
equalities. The valuation of the displayed integer and the
definition of \(d_\ell\) prove the colength statements.
\end{proof}

Compatibility of the prescribed Frobenius with the projector
is essential in the specialization. In fact fix an arbitrary
integral \(\Phi_0\) with the same residual value and let
\(I\subset\mathcal O\) be the ideal of its six off-diagonal
coordinates. The specialized module is
\[
       \mathcal O\ \oplus\
         u\bigl(\mathcal O[[u]]/(F_M(u),I)\bigr).
\]
The first summand is evaluation at zero. For the second, cancel
\(u\) in the power series ring before taking the quotient by
\(u(F_M,I)\). For \(I=(\pi^b)\ne0\) the second summand has
length \(b(M-1)\) by the same monic division, and disappears
in characteristic zero. The fixed commuting specialization in
Corollary~\ref{cor:selmercongruence614} describes the local
cost at actual split test primes. Relating these local charts
to a global Hecke algebra at the prescribed characteristic-zero
weights, with constants uniform across intervals, is still needed
for the strict length bound of Section~\ref{sec:missing}.


\subsection{Paired tame inertia in the symplectic trace algebra}
\label{sec:pairedinertia651}

The four-dimensional representations used for the colliding
weights have a prescribed symplectic similitude character.
In the source this condition is imposed on every Galois
element in
\srcloc{Development/AutomorphicArithmetic/AutomorphicVarieties/SymplecticGSp4/DecompSigmaOrd}{142}.
At an auxiliary prime where the multiplier is unramified, a
ramified inertia character on one line must be paired with
its inverse on the dual line. We give the corresponding local
calculation and its effect at an unramified system of traces.

Let \(A\) be a local ring in which two is a unit, let \(V\)
be free of rank four with a perfect alternating form
\(\langle\, ,\,\rangle\), and write \(B^\dagger\) for the
adjoint of an endomorphism. Put
\(\mathfrak s=\{Y:Y^\dagger=-Y\}\).

\begin{lemma}[Paired lines give the same tests on the symplectic adjoint]
\leavevmode\label{lem:pairedlines651}
For any rank-one idempotent \(D\) on \(V\), its adjoint is a
rank-one idempotent with \(DD^\dagger=D^\dagger D=0\). The two
images span a hyperbolic plane and \(H=D-D^\dagger\) belongs
to \(\mathfrak s\). For \(Y\in\mathfrak s\),
\[
                 \operatorname{tr}(HY)=2\operatorname{tr}(DY).
\]
In particular all the tests of Theorem~\ref{thm:splitprojectors599}
on a curve preserving such a form up to similitude detect its
symplectic adjoint cohomology at the same primes and precisions
when \(D\) is replaced by \(H\).
\end{lemma}
\begin{proof}
Choose \(e\) generating the image of \(D\). Write
\(D(v)=\langle v,f\rangle e\) using the perfection of the form.
Since \(D(e)=e\), one has \(\langle e,f\rangle=1\).
The adjoint is \(D^\dagger(v)=\langle e,v\rangle f\).
These formulas prove the claims about the two lines. Taking the
trace after adjunction and cyclically permuting the factors
gives \(\operatorname{tr}(D^\dagger Y)=-\operatorname{tr}(DY)\).
Both the formula and the decomposition into the eigenspaces
of the adjoint involution commute with reduction. If the
Galois representation preserves the form up to similitude,
then conjugation commutes with adjunction. The map
\(B\mapsto(B-B^\dagger)/2\) consequently splits off
\(\mathfrak s\) as a Galois module, also after the common
cyclotomic twist. Its first cohomology injects into the
endomorphism cohomology. The displayed factor of two is a unit
at the odd coefficient primes of that theorem. Applying the
existing injectivity proves the last assertion.
\end{proof}

Take a complete discrete valuation ring \(\mathcal O\) finite
over \(\ZZ_\ell\), with \(\ell\) odd. Let \(q\ne\ell\) be a
rational prime with \(M=\ell^{v_\ell(q-1)}>1\). Suppose that the
residual Frobenius is a nonzero scalar and that residual inertia
is trivial. Fix an unramified multiplier lifting the residual
one. Choose a symplectic basis for two paired lines and their
orthogonal complement. Prescribe the matrix of a tame inertia
generator in that basis to be
\[
        J(z)=zD+z^{-1}D^\dagger+(1-D-D^\dagger),
              \qquad z=1+u,\qquad z^M=1.
\]
Here \(D\) projects to the first line, \(D^\dagger\) to the
second, and \(u\) belongs to the maximal ideal. Wild inertia
and the other tame factors have trivial image. Frobenius is
allowed to vary in the framed fixed-multiplier group.

\begin{proposition}[The two local branches with prescribed multiplier]
\leavevmode\label{prop:pairedchart651}
Use four variables \(Y\) and six variables \(X\), and put
\(F_M(u)=((1+u)^M-1)/u\).
The complete ring for these framed data is
\[
 \mathcal O[[Y,X,u]]/(uF_M(u),uX_1,\ldots,uX_6).
\]
It is flat over \(\mathcal O\) and reduced. Its unramified
branch and the closure of the locus \(z\ne1\) meet in the
ring \((\mathcal O/M\mathcal O)[[Y]]\).
\end{proposition}
\begin{proof}
The prescribed Frobenius multiplier has a square root lifting
the residual scalar by Hensel's lemma. Divide Frobenius by
this scalar lift to obtain a symplectic matrix \(g\) congruent
to one. The Cayley transform
\(Z=(g-1)(g+1)^{-1}\) has \(Z^\dagger=-Z\), since
\(g^\dagger=g^{-1}\). Conversely \(g=(1+Z)(1-Z)^{-1}\)
recovers every such matrix from a skew-adjoint \(Z\)
congruent to zero. The denominators in both transformations
are invertible. Commuting with \(J(z)\) is the same condition
for \(Z\) and for \(g\).

Because \(M\mid q-1\), the tame Frobenius relation is just
commutation with \(J(z)\). Its three diagonal scalars on the
line, the orthogonal plane and the dual line are \(z,1,z^{-1}\).
Their pairwise differences are \(u\) times units, including
\(z-z^{-1}=u(1+z^{-1})\). Thus \(u\) times each off-diagonal
block of \(Z\) vanishes. The block-diagonal skew-adjoint
matrices consist of a scalar and its negative on the two
lines and an element of \(\mathfrak{sp}_2\) on the plane.
They have rank four. The off-diagonal blocks form a direct
complement of rank six. Indeed adjunction preserves the
block-diagonal and off-diagonal spaces, and the full
skew-adjoint space has rank ten by identifying it with
symmetric bilinear forms on \(V\). These direct summands are
free over the local coefficient ring. This gives precisely the
stated equations. Conversely the equations realize the tame
presentation on every finite local coefficient algebra, and
passing to inverse limits gives a continuous representation.

In \(S=\mathcal O[[Y,X,u]]\), the ring \(S/(F_M,X)\)
is free of rank \(M-1\) over \(\mathcal O[[Y]]\) by monic
distinguished division. Over its fraction field the roots of
\(F_M\) are distinct and nonzero. Thus this ring is reduced
and multiplication by \(u\) is injective on it. Consequently
\((u)\cap(F_M,X)=u(F_M,X)\). The two branch rings are
\(\mathcal O\)-torsion free and reduced, and the complete ring
embeds in their product. This proves flatness and reducedness.
The quotient by the sum of their defining ideals is
\(S/(u,F_M,X)=(\mathcal O/M\mathcal O)[[Y]]\), as asserted.
\end{proof}

There is a separating element expressible entirely through
the trace of inertia. Define \(T_0(X)=2,\ T_1(X)=X\) and
\(T_{i+1}(X)=XT_i(X)-T_{i-1}(X)\). For an odd integer \(M>1\)
put
\[
                 U_M(X)=1+\sum_{i=1}^{(M-1)/2}T_i(X).
\]
Induction gives \(T_i(z+z^{-1})=z^i+z^{-i}\). Since
\(\operatorname{tr}J(z)=2+z+z^{-1}\), on the ring \(z^M=1\)
one has the integral identity
\[
     U_M(\operatorname{tr}J(z)-2)=\sum_{a=0}^{M-1}z^a .
\]

\begin{theorem}[Integral comparison at an unramified trace system]
\leavevmode\label{thm:pairedtrace651}
Take finitely many such primes \(q_i\), put
\(M_i=\ell^{v_\ell(q_i-1)}\), and let \(d=\prod_i M_i\).
On the product of the local character rings, every integral
function which vanishes at all characteristic-zero tuples with
some \(z_i\ne1\) takes a value in \(d\mathcal O\) at the
all-unramified tuple. The product of the displayed trace
polynomials attains \(d\). This also gives the exact ideal of
values in the subalgebra generated by the inertia traces.

Now let \(B\) be a finite torsion-free \(\mathcal O\)-algebra
realized in a product of characteristic-zero parameter fields,
containing the traces of the specified inertia generators.
Assume inertia at each \(q_i\) factors through the cyclic group
of order \(M_i\) and its generator has the characteristic
polynomial of the prescribed paired type at each parameter.
In characteristic zero a finite-order matrix is semisimple,
so its eigenvalues determine whether it is the identity.
Let \(B_0,B_1\) be the images on the all-unramified tuples
and on the other tuples respectively.
Fix an \(\mathcal O\)-algebra homomorphism
\(\lambda_0:B_0\to\mathcal O\) and let \(\lambda\) be its pullback.
Write \(I=\ker\lambda,\ I_0=\ker\lambda_0\). For
\(\eta_\lambda(B)=\lambda(\operatorname{Ann}_B I)\), one has
\[
 \begin{gathered}
    d(B_0\oplus B_1)\ \subseteq\ B\ \subseteq\ B_0\oplus B_1,\\
    d\,\eta_{\lambda_0}(B_0)\ \subseteq\ \eta_\lambda(B)
                         \ \subseteq\ \eta_{\lambda_0}(B_0).
 \end{gathered}
\]
For every \(k\ge1\) the map
\(I^k/I^{k+1}\longrightarrow I_0^k/I_0^{k+1}\) is surjective
with kernel killed by \(d\). At \(k=1\) this kernel has length
at most
\[
     v_\pi(d)\,\dim_{\mathcal O/(\pi)} I/(I^2+\pi I),
\]
where \(\pi\) is a uniformizer of \(\mathcal O\).
\end{theorem}
\begin{proof}
The independent variables \(z_i\) with \(z_i^{M_i}=1\) give
the group ring of \(\prod_i\ZZ/M_i\ZZ\). Over a field
containing all its character values, Fourier inversion shows
that an element supported at the trivial tuple with value
\(c\) has constant coefficient \(c/d\) on every group element.
Its integral coefficients therefore imply \(c\in d\mathcal O\).
Each sum \(\sum_{a=0}^{M_i-1}z_i^a\) is zero at every
nonidentity root of unity and equals \(M_i\) at one. Their
product attains the bound, using only the traces.

In \(B\) let \(w\) be the same product of trace polynomials.
On the two sets of components its values are \(d\) and zero.
The embedding in parameter fields therefore identifies
\(w=(d,0)\) and \(d-w=(0,d)\). Multiplying lifts from the
two image algebras by these elements proves the first pair
of inclusions. The image of \(\operatorname{Ann}_B I\)
annihilates \(I_0\) because \(I\to I_0\) is surjective.
In the other direction, for
\(c_0\in\operatorname{Ann}_{B_0}I_0\) the element
\((dc_0,0)\in B\) annihilates \(I\). Evaluating both
observations at the distinguished system proves the second pair.

Surjectivity on \(I^k\) follows by lifting products.
If \(v\in I^k\) maps into \(I_0^{k+1}\), subtract from it a
lift in \(I^{k+1}\). The resulting \(v'\) vanishes in \(B_0\),
so \(wv'=0\). Since \(d-w\in I\), we get
\(dv'=(d-w)v'\in I^{k+1}\). This proves the annihilation.
A submodule of a finitely generated module over a discrete
valuation ring has at most as many generators as its ambient
module, by lifting its preimage to a finite free cover.
Nakayama's lemma gives the stated generator count for
\(I/I^2\). Each generator of the kernel contributes at most
\(v_\pi(d)\) to its length.
\end{proof}

For the primes of Theorem~\ref{thm:exactselmer614} the loss
per prime in this calculation is exactly
\(v_\pi(\ell)v_\ell(q_i-1)\) on the independent local character
rings. The inclusions give an upper bound for the loss on the
indicated global trace images. Applying this to the entire
strict cohomology group still requires an integral comparison
between that group and the relevant trace algebra. In
particular the embedding of \(B\) in characteristic-zero
fields used to prove \(wv'=0\) must be justified for whichever
algebra is used in that comparison.


\subsection{An integral annihilator for finite symplectic inertia}
\label{sec:inertiaann654}

At an auxiliary tame prime the quotient imposing unramifiedness
is defined by the entries of the inertia matrix minus the
identity. This description lets us compare trace algebras
before passing to their characteristic-zero images. Throughout,
\(\mathcal O\) is a complete discrete valuation ring of odd
residue characteristic \(\ell\). Let \(M=\ell^a\), \(a\ge1\).
Let \(R\) be a commutative \(\mathcal O\)-algebra, and
\(U\in\operatorname{Sp}_4(R)\) satisfy \(U^M=I_4\).
This models a cyclic tame inertia image of order dividing
\(M\) at a prime where the similitude multiplier is unramified.
The symplectic module has a fixed symplectic basis with
alternating form matrix \(J\).

\begin{theorem}[A trace polynomial killing the entire inertia ideal]
\leavevmode\label{thm:inertiaann654}
Write \(S=\sum_{j=-(M-1)/2}^{(M-1)/2} U^j\). The scalar
\[
 w_M(U)=\frac{(\operatorname{tr}S)^2-2\operatorname{tr}(S^2)}8
 \qquad\text{satisfies}\qquad
 w_M(U)(U-I_4)=0,\qquad w_M(I_4)=M^2.
\]
It belongs to the \(\mathcal O\)-algebra of traces of powers
of \(U\). The ideal \(M^2\mathcal O\) of possible values at
the trivial character is sharp for universal integral scalar
annihilators when both symplectic planes can ramify independently.

If \(U\) is the identity on a nondegenerate direct summand of
rank two, one may use
\[
       w'_M(U)=U_M(\operatorname{tr}U-2),\qquad
       w'_M(U)(U-I_4)=0,\qquad w'_M(I_4)=M .
\]
Here \(U_M\) is the integral polynomial of
Section~\ref{sec:pairedinertia651}. The value ideal in this
case is \(M\mathcal O\). All the identities hold over rings
with nilpotents and with \(\ell\)-power torsion.
\end{theorem}
\begin{proof}
Let \(X^\dagger=J^{-1}X^{\mathsf T}J\). Since
\(U^\dagger=U^{-1}\), the symmetric range of summation gives
\(S^\dagger=S\). The geometric sum and \(U^M=I_4\) give
\(S(U-I_4)=0\). The matrix \(A=JS\) is alternating.
Its Pfaffian is \(af-be+cd\) if the upper entries, in row
order, are \(a,b,c,d,e,f\). Direct multiplication gives
\[
 \begin{pmatrix}
 0&-f&e&-d\\ f&0&-c&b\\-e&c&0&-a\\ d&-b&a&0
 \end{pmatrix} A=\operatorname{Pf}(A) I_4 .
\]
Thus the scalar \(\operatorname{Pf}(JS)\) kills \(U-I_4\).
The same four by four determinant expansion gives
\(\det A=\operatorname{Pf}(A)^2\).
As \(J\) is unimodular, \(\operatorname{Pf}(J)\) is a unit.
Consequently
\[
 \det(TI_4-S)=
       \left(\frac{\operatorname{Pf}(TJ-JS)}
                         {\operatorname{Pf}(J)}\right)^2 .
\]
The polynomial inside the square is monic quadratic. Write
it as \(T^2-c_1T+c_0\). Its constant term is
\(\operatorname{Pf}(JS)/\operatorname{Pf}(J)\). Comparing
the cubic and quadratic coefficients gives
\[
 2c_1=\operatorname{tr}S,\qquad
 c_1^2+2c_0=
       \frac{(\operatorname{tr}S)^2-\operatorname{tr}(S^2)}2.
\]
The second expression on the right is the sum of the principal
two by two minors. Solving for \(c_0\) gives exactly \(w_M(U)\).
This proves the annihilation over \(R\). Both traces in the
formula are sums of traces of powers of \(U\), and evaluating
at \(S=MI_4\) gives \(M^2\).

In the second case the orthogonal complement of the fixed plane
is stable and nondegenerate. The calculation is local on the
coefficient ring, so one may choose a symplectic basis of this
complement. Its matrix \(V\) has determinant one.
Cayley--Hamilton gives \(V+V^{-1}=(\operatorname{tr}V)I_2\).
The recurrence defining \(T_j\) in the preceding section
therefore gives
\(V^j+V^{-j}=T_j(\operatorname{tr}V)I_2\).
Summing shows that \(S\) on this plane equals
\(U_M(\operatorname{tr}V)I_2\), where
\(\operatorname{tr}V=\operatorname{tr}U-2\). Multiplying
by \(V-I_2\) gives zero. On the fixed plane \(U-I_4\)
already vanishes. Evaluation at the identity gives \(M\).

For sharpness take the group ring
\(\mathcal O[(\ZZ/M\ZZ)^2]\) with its two tautological units
\(z_1,z_2\), and use the matrix
\(\operatorname{diag}(z_1,z_2,z_1^{-1},z_2^{-1})\) in a
hyperbolic basis. A scalar annihilating its difference from
the identity is fixed by multiplication by \(z_1\) and
\(z_2\). Translation on the group basis forces every
coefficient to be the same. Its augmentation consequently
belongs to \(M^2\mathcal O\). With a fixed plane the same
argument uses \(\mathcal O[\ZZ/M\ZZ]\) and yields
\(M\mathcal O\). The constructed elements attain these
values. These lower bounds therefore apply in particular
to functions formed from traces.
\end{proof}

\begin{corollary}[Passage to unramified traces over an integral base]
\leavevmode\label{cor:verticaltrace654}
Let a representation over \(R\) have fixed multiplier
unramified at a finite list of auxiliary primes. Suppose
inertia at prime \(i\) is generated topologically in the
image by \(U_i\) with \(U_i^{M_i}=I_4\), where \(M_i\)
is a power of \(\ell\). Let \(K\subset R\) be the ideal
generated by the entries of every \(U_i-I_4\), and let
\(B\subset R\) be an \(\mathcal O\)-subalgebra containing
the traces of their powers. Define \(B_0\) to be the image
of \(B\) in \(R/K\).
Choose \(a_i=1\) at a prime where \(U_i\) fixes a
nondegenerate plane pointwise, and \(a_i=2\) at every
other prime. Put \(D_0=\prod_i M_i^{a_i}\).

For any \(\mathcal O\)-point \(\lambda_0:B_0\to\mathcal O\)
with pullback \(\lambda:B\to\mathcal O\), one has
\[
 D_0\,\eta_{\lambda_0}(B_0)\ \subseteq\ \eta_\lambda(B)
                    \ \subseteq\ \eta_{\lambda_0}(B_0).
\]
For \(I=\ker\lambda,\ I_0=\ker\lambda_0\) and every \(h\ge1\),
the surjection \(I^h/I^{h+1}\to I_0^h/I_0^{h+1}\)
has kernel killed by \(D_0\).
When \(I/I^2\) is finitely generated over \(\mathcal O\),
the kernel at \(h=1\) has length at most
\[
             v_\pi(D_0)\dim_{\mathcal O/(\pi)} I/(I^2+\pi I).
\]
\end{corollary}
\begin{proof}
Let \(w\) be the product of the scalars in the theorem, using
\(w'_{M_i}\) at a fixed plane and \(w_{M_i}\) otherwise.
This element belongs to \(B\), satisfies \(wK=0\) in \(R\),
and maps to \(D_0\) in \(R/K\). In particular it annihilates
\(\ker(B\to B_0)=B\cap K\), even for a trace supported
entirely in residue characteristic \(\ell\).
The map \(I\to I_0\) is surjective. The image of an
annihilator of \(I\) consequently annihilates \(I_0\).
Conversely if \(a_0 I_0=0\), lift \(a_0\) to \(a\in B\).
Then \(a I\subseteq B\cap K\), so \(wa I=0\), and
\(\lambda(wa)=D_0\lambda_0(a_0)\).
These give the assertion about the value ideals
\(\eta_\lambda(B)=\lambda(\operatorname{Ann}_B I)\).

Surjectivity on powers follows by lifting products. Given
\(x\in I^h\) mapping to \(I_0^{h+1}\), subtract a lift from
\(I^{h+1}\) to obtain \(x'\in I^h\cap K\).
Then \(D_0 x'=(D_0-w)x'\in I^{h+1}\), since \(D_0-w\in I\).
This proves the assertion in every degree. The generator
bound over a discrete valuation ring in the proof of
Theorem~\ref{thm:pairedtrace651} gives the last inequality.
\end{proof}

In particular the cost of removing these prescribed finite
tame inertia groups can be bounded inside the original
integral coefficient algebra. One must still identify the
strict classes with the relevant directions in that algebra,
bound the unramified quotient at the required rational
weights and treat any local type with unrestricted monodromy
by a separate argument. These requirements enter the local
estimate of Section~\ref{sec:missing}.


\subsection{Ramified inertia at the character conductor}
\label{sec:conductorinertia655}

Let \(\mathcal O\) be the ring of integers of a finite extension \(K\)
of \(\QQ_\ell\), where \(\ell\) is odd. Write \(\pi\) for a uniformizer,
\(v=v_\pi\), \(e_\ell=v(\ell)\), and \(v(0)=+\infty\).
Fix a finite order character \(\psi\) with values in \(\mathcal O^\times\).
At a prime \(p\ne\ell\), denote its inertia and decomposition groups
by \(I_p\subset D_p\). We consider the primes occurring in the
ramified local condition of the four-dimensional argument, where
\(\ell\mid M=|\psi(I_p)|\). Choose an inertia element \(\tau\) such
that \(\zeta=\psi(\tau)\) generates this cyclic image, and a lift \(F\)
of arithmetic Frobenius. Write \(\delta=\zeta-1\).

The four characters at the split point, in their given order, are
\(\alpha\Psi,\Psi,\alpha,1\). Here \(\Psi=\chi_\ell^r\psi\),
\(\alpha=\chi_\ell\omega^{-1}\xi\), \(\chi_\ell\) is the
cyclotomic character, \(\omega\) its Teichm\"uller lift modulo \(\ell\),
and \(\xi\) is unramified at \(p\). The form pairs the first coordinate
with the fourth and the second with the third. Thus at \(\tau\) the two
Lagrangians at the split point have characters \(\zeta\) and \(1\),
each of multiplicity two. The multiplier at \(\tau\) equals \(\zeta\).
Write \(a=\alpha(F)\) and \(\beta=\Psi(F)\). The Frobenius separation
in this local condition gives
\[
  \bar a\ne1,\qquad \bar a\bar\beta\ne1,\qquad
  \bar a\bar\beta^{-1}\ne1 .
\]
Indeed \(\bar\alpha=\bar\xi\), and the separation asserts
\(\bar\xi(F)\bar\Psi(F)^i\ne1\) for all integers \(i\).

The local component used here has a concrete meaning. A representation
on it is obtained from a representation over an integral domain
\(B\) over \(\mathcal O\) that specializes to the split point over
\(\mathcal O\). The representation over \(B\) has the prescribed
multiplier and is trivial on \(I_p\cap\ker\psi\). A second
specialization of \(B\) can take values in a ring with torsion.
For the definitions in the source, see \texttt{OnComp},
\texttt{condP} and \texttt{IsPt} in \texttt{DecompSigmaOrd.lean}.

\begin{proposition}[The inertia matrix on the component through the split point]
\leavevmode\label{prop:conductorinertia655}
Let \(U\) denote the image of \(\tau\). Over every coefficient ring
obtained by such a specialization one has
\[
   (U-I)(U-\zeta I)=0,\qquad U\rho(g)=\rho(g)U
                  \quad(g\in D_p).
\]
After inverting \(\delta\), the operator \((U-I)/\delta\) projects
onto the \(\zeta\)-eigenspace along the \(1\)-eigenspace.
Both are invariant Lagrangian summands of rank two.
\end{proposition}
\begin{proof}
The character \(\psi\) identifies \(I_p/(I_p\cap\ker\psi)\)
with its image, so conjugating an inertia element by any element
of \(D_p\) preserves its coset in this quotient. This gives the
commutation relation and \(U^M=I\). The map \(\mathcal O\to B\)
is injective because the split point is a retraction to \(\mathcal O\).
Over the fraction field of \(B\), \(U\) has finite order dividing
\(M\) and is diagonalizable with eigenvalues among the distinct
elements \(1,\zeta,\ldots,\zeta^{M-1}\). Its characteristic
polynomial therefore has coefficients in \(\mathcal O\).
Evaluating at the split point preserves those coefficients, so the
polynomial is \((T-\zeta)^2(T-1)^2\). Diagonalizability gives the
quadratic identity over the fraction field, hence over \(B\) and
every specialization. Inverting \(\delta\) now gives the stated
idempotent. Its two summands have rank two by the characteristic
polynomial over each residue field of this localization. On the
first summand the form scales under \(U\) both by \(\zeta^2\) and
by \(\zeta\). On the second it scales both by \(1\) and by
\(\zeta\). Since their differences are units there, both summands
are isotropic, and nondegeneracy pairs them perfectly.
\end{proof}

\begin{theorem}[The inertia restriction of a symplectic deformation]
\leavevmode\label{thm:symplocal655}
Put \(W_h=\mathcal O/\pi^h\), \(h\ge1\),
\[
 c_h=\min(h,v(\delta)),\qquad
 d_h=\min(c_h,v(\beta-1)).
\]
At the split point consider first-order lifts on this component
to \(\mathcal O\oplus W_h\), with \(W_h^2=0\) and fixed multiplier.
Identify inertia restrictions up to symplectic change of basis
equal to the identity modulo \(W_h\). The module generated by
these restrictions has length at most \(2c_h+2d_h\). More
precisely their classes inject into
\[
 W_h[\delta]^2\ \oplus\
       \big((W_h/\delta W_h)[\beta-1]\big)^2 ,
\]
where \(N[u]\) denotes the kernel of multiplication by \(u\) on \(N\).
\end{theorem}
\begin{proof}
Choose paired bases for the two Lagrangians, so the symplectic form
has off-diagonal blocks \(I_2,-I_2\). Write
\(U=(I+Y)U_0\), \(\rho(F)=(I+Z)F_0\), where the subscript denotes
the split point. Preservation of the multiplier makes \(Y,Z\)
lie in the symplectic Lie algebra over \(W_h\). Thus
\[
 U_0=\operatorname{diag}(\zeta I_2,I_2),\qquad
 Y=\begin{pmatrix} A&B'\\ C'&-A^{\mathsf T}\end{pmatrix},
 \qquad (B')^{\mathsf T}=B',\quad (C')^{\mathsf T}=C'.
\]
In the linear part of \((U-I)(U-\zeta I)=0\), the diagonal blocks
give \(\delta A=0\), and the off-diagonal blocks cancel. Conjugation
by \(I+T\) adds \((1-\operatorname{Ad}U_0)T\) to \(Y\).
It leaves \(A\) fixed and changes \(B',C'\) by arbitrary symmetric
matrices with entries in \(\delta W_h\). The class of \(Y\) with
these identifications determines the whole restriction, since
inertia acts through the cyclic group generated by \(U\).

Commutation with Frobenius gives
\[
  (1-\operatorname{Ad}F_0)Y=(1-\operatorname{Ad}U_0)Z .
\]
On the first Lagrangian the Frobenius matrix is
\(\operatorname{diag}(a\beta,\beta)\), and on its paired dual
it is \(\operatorname{diag}(1,a)\). Thus the two off-diagonal
entries of \(A\) vanish, since \(a-1\) is a unit. On the three
symmetric upper coordinates the Frobenius ratios are \(a\beta\),
\(\beta\) and \(\beta/a\). On the lower coordinates they are their
inverses. After passage to \(W_h/\delta W_h\) the right side of the
display vanishes. The separation conditions leave only the
\(\beta\)- and \(\beta^{-1}\)-coordinates, both killed by \(\beta-1\).
These coordinates and the two diagonal entries of \(A\) give the
asserted injection. Their respective cyclic coefficient modules
have lengths \(d_h\) and \(c_h\).
\end{proof}

\begin{theorem}[The exact scalar cost of the strict local condition]
\leavevmode\label{thm:scalarconductor655}
Give \(W_h\) or \(W_\infty=K/\mathcal O\) the scalar action of
\(\Psi^{-1}\). At \(p\), consider the local classes represented by
cocycles vanishing on \(I_p\cap\ker\psi\). Their lengths over
\(\mathcal O\), at finite precision and with divisible coefficients,
respectively, are
\[
             2d_h\qquad\text{and}\qquad
             d_\infty=\min\bigl(v(\delta),v(\beta-1)\bigr).
\]
Consequently, imposing local triviality at any list of these primes
on a global scalar group already satisfying this inertia-kernel
condition decreases its length by at most the sum of the corresponding
local lengths, whenever the global group has finite length.
\end{theorem}
\begin{proof}
The action is trivial on \(I_p\cap\ker\psi\), so the condition
descends the cocycle to \(D_p/(I_p\cap\ker\psi)\). In that group
the inertia quotient is central cyclic of order \(M\). The quotient
by inertia is \(\widehat{\ZZ}\). The continuous map from
\(\widehat{\ZZ}\) sending \(1\) to \(F\) gives a splitting because
it composes to the identity on \(\widehat{\ZZ}\). Thus it suffices to
specify the two values \(u=z(\tau)\), \(w=z(F)\).
Put \(s=\zeta^{-1}-1\) and \(t=\beta^{-1}-1\). The commutation
relation reads \(tu=sw\). The relation \(\tau^M=1\) adds none,
since \(\sum_{j=0}^{M-1}\zeta^{-j}=0\) in \(\mathcal O\).
Conversely all such pairs give continuous cocycles. On finite
coefficients the affine transformations defined by the two values
commute, the first has order dividing \(M\), and the second has
finite order. For \(K/\mathcal O\) any two values lie in a common
finite submodule stable under both characters. The coboundary
pairs are precisely \((sy,ty)\).

Let \(g\) generate \((s,t)\) in \(\mathcal O\). A change of basis in
\(\operatorname{GL}_2(\mathcal O)\) sends the column \((s,t)\) to
\((g,0)\), since \((s/g,t/g)\) is primitive. The relation \(tu=sw\)
says that the determinant of these two columns vanishes, a condition
preserved by this change of basis. On either coefficient module \(W\)
the group of classes is consequently isomorphic to
\[
                         W/gW\ \oplus\ W[g].
\]
Here \(v(g)=\min(v(\delta),v(\beta-1))\). Both summands at finite
precision have length \(d_h\). On \(K/\mathcal O\) multiplication by
the nonzero \(g\) is surjective and its kernel has length \(v(g)\).
Finally the quotient of the original global group by the subgroup
with trivial restrictions injects into the direct sum of the
computed local groups.
\end{proof}

\begin{corollary}[A conductor bound uniform in weight and coefficient prime]
\leavevmode\label{cor:sumconductor655}
Fix a finite order character in characteristic zero, allowing all its
\(\ell\)-adic realizations, and a finite set \(S\) of rational primes.
At one weight, sum over odd coefficient primes the preceding local
length bounds at \(p\in S\), \(p\ne\ell\), \(\ell\mid|\psi(I_p)|\),
counting one realization per coefficient prime. Weight the length by
\((\log\ell)/e_\ell\). For scalar classes the sum is at most \(|S|\)
with divisible coefficients and at most \(2|S|\) at any finite
precision. For the symplectic inertia restrictions under the
Frobenius separation above it is at most \(4|S|\). Over \(t_0\)
weights the same bounds acquire a factor \(t_0\).
\end{corollary}
\begin{proof}
The order \(M\) of the inertia character is independent of the
embedding. If \(M\) has a prime-to-\(\ell\) factor greater than one,
the reduction of \(\zeta\) is nontrivial, since reduction preserves
the order of roots of unity of order prime to \(\ell\). Thus
\(v(\delta)=0\). In the remaining case \(M=\ell^j\), \(j\ge1\),
and \(v(\delta)/e_\ell=1/\varphi(\ell^j)\). Indeed the cyclotomic
polynomial evaluated at \(1\) gives
\(\ell=\prod_{i\in(\ZZ/\ell^j)^\times}(1-\zeta^i)\).
Each factor has the same valuation as \(1-\zeta\), by the finite
geometric sum for a positive representative of \(i\) and its inverse
modulo \(\ell^j\). For each \(p\) at most one coefficient prime
therefore contributes, and its normalized contribution from one
copy of \(\mathcal O/(\delta)\) is at most
\(\log\ell/(\ell-1)\le1\). The length formulas give the bounds.
\end{proof}

These estimates apply to the ramified conductor terms of the
specified character at the original weights. The kernel of all
the calculated restrictions can still contain global strict
classes. Its order enters Section~\ref{sec:missing} even when
every contribution of this corollary is zero.

\clearpage
\section{References}\label{nav:bibliography-XII}
\begingroup\def\section*#1{}\begin{thebibliography}{99}\small
\bibitem{XII-IX-mult560} H. Jacquet and J. A. Shalika, \emph{On Euler products and the classification of automorphic representations}, Parts I and II, Amer. J. Math. \textbf{103} (1981), 499--558 and 777--815. \href{https://www.math.columbia.edu/~hj/On%20Euler%20products%20II.pdf}{Part II}.
\bibitem{XII-IX-blggt560} T. Barnet-Lamb, T. Gee, D. Geraghty and R. Taylor, \emph{Potential automorphy and change of weight}, Ann. of Math. (2) \textbf{179} (2014), 501--609, Section~2.1, Theorem~2.1.1, and Section~5.1. \href{https://annals.math.princeton.edu/2014/179-2/p03}{Published paper}.
\bibitem{XII-IX-sign560} J. Bella\"iche and G. Chenevier, \emph{The sign of Galois representations attached to automorphic forms for unitary groups}, Compos. Math. \textbf{147} (2011), 1337--1352, Theorem~1.2. \href{https://gaetan.chenevier.perso.math.cnrs.fr/articles/sign.pdf}{Authors' paper}.
\bibitem{XII-IX-lgp361} T. Barnet-Lamb, T. Gee, D. Geraghty and R. Taylor, \emph{Local-global compatibility for \(l=p\), II}, Ann. Sci. \'Ec. Norm. Sup\'er. (4) \textbf{47} (2014), 165--179, Theorem A. \href{https://www.numdam.org/articles/10.24033/asens.2212/}{Published paper}.
\bibitem{XII-IX-newvector561} H. Jacquet, I. I. Piatetski-Shapiro and J. A. Shalika, \emph{Conducteur des repr\'esentations du groupe lin\'eaire}, Math. Ann. \textbf{256} (1981), 199--214, Theorem~5. \href{https://lesesvre.perso.math.cnrs.fr/newforms-references/jpss.pdf}{Original paper}.
\bibitem{XII-IX-clozel561} L. Clozel, \emph{Motifs et formes automorphes. Applications du principe de fonctorialit\'e}, in \emph{Automorphic Forms, Shimura Varieties, and \(L\)-functions}, Vol.~I, Perspectives in Mathematics~10, Academic Press, 1990, 77--159, Lemme~3.14.
\bibitem{XII-IX-franke561} J. Franke, \emph{Harmonic analysis in weighted \(L_2\)-spaces}, Ann. Sci. \'Ec. Norm. Sup\'er. (4) \textbf{31} (1998), 181--279. \href{https://www.numdam.org/article/ASENS_1998_4_31_2_181_0.pdf}{Published paper}.
\bibitem{XII-IX-V-bs} A. Borel and J.-P. Serre, \emph{Corners and arithmetic groups}, Comment. Math. Helv. \textbf{48} (1973), 436--491. \href{https://www.e-periodica.ch/cntmng?pid=com-001:1973:48::32}{Original paper}.
\bibitem{XII-IX-V-stick} L. C. Washington, \emph{Introduction to Cyclotomic Fields}, 2nd ed., Springer, 1997. \href{https://link.springer.com/book/10.1007/978-1-4612-1934-7}{Chapter 6, Stickelberger's theorem}.
\bibitem{XII-IX-V-milne98} J. S. Milne, \emph{Arithmetic Duality Theorems}, 2nd ed., BookSurge, 2006. \href{https://www.jmilne.org/math/Books/ADTnot.pdf}{Part I, Sections 2 and 4}.
\bibitem{XII-IX-V-deligne} P. Deligne, \emph{Formes modulaires et repr\'esentations \(\ell\)-adiques}, S\'eminaire Bourbaki, Exp. 355 (1969). \href{https://www.numdam.org/item/SB_1968-1969__11__139_0.pdf}{Paper}.
\bibitem{XII-IX-ordinary361} C. M. Skinner and A. J. Wiles, \emph{Ordinary representations and modular forms}, Proc. Natl. Acad. Sci. USA \textbf{94} (1997), 10520--10527. \href{https://doi.org/10.1073/pnas.94.20.10520}{Published paper}.
\bibitem{XII-IX-VI-saito40} T. Saito, \emph{Modular forms and \(p\)-adic Hodge theory}, Invent. Math. \textbf{129} (1997), 607--620. See also \href{https://www.ms.u-tokyo.ac.jp/~t-saito/pp/h.pdf}{\emph{Hilbert modular forms and \(p\)-adic Hodge theory}}, Compos. Math. \textbf{145} (2009), 1081--1113.
\bibitem{XII-IX-I-ac123} J. Arthur and L. Clozel, \href{https://www.claymath.org/library/cw/arthur/pdf/30.pdf}{\emph{Simple Algebras, Base Change, and the Advanced Theory of the Trace Formula}}, Annals of Mathematics Studies 120, Princeton University Press, 1989, Chapter III.
\bibitem{XII-IX-VI-bk117} S. Bloch and K. Kato, \href{https://virtualmath1.stanford.edu/~conrad/BSDseminar/refs/BKTamagawa.pdf}{\emph{L-functions and Tamagawa numbers of motives}}, in \emph{The Grothendieck Festschrift}, Vol.~I, Progr. Math. \textbf{86} (1990), 333--400, Section~3.
\bibitem{XII-IX-geraghty363} D. Geraghty, \emph{Modularity lifting theorems for ordinary Galois representations}, Math. Ann. \textbf{373} (2019), 1341--1427, Section 2. \href{https://doi.org/10.1007/s00208-018-1742-4}{Published paper}. Proposition and lemma numbers refer to the \href{https://citeseerx.ist.psu.edu/document?doi=8faa2d675c94ae78d0ac8ebb3f9d4f69ede42d44\&repid=rep1\&type=pdf}{original manuscript}.
\end{thebibliography}\endgroup
\end{document}
